% Edición latina — Seitz 1903, II. Disputatio Luther-Eck (impresas 55-218)
% Compilar: xelatex edicion-latina.tex (dos pasadas)
\documentclass[11pt,a4paper,twoside]{article}
\usepackage{fontspec}
\setmainfont{Libertinus Serif}
\usepackage[lmargin=3cm,rmargin=3cm,top=2.8cm,bottom=2.5cm,headheight=14pt]{geometry}
\usepackage{setspace}
\setstretch{1.08}
\setlength{\parskip}{0.4em}
\setlength{\parindent}{1.2em}
\usepackage[hang,flushmargin]{footmisc}
\usepackage{fancyhdr}
\usepackage[hidelinks]{hyperref}
% llamada con la etiqueta ORIGINAL del impreso + nota con etiqueta repetida
\newcommand{\apnota}[2]{\textsuperscript{#1)}\,\footnote{[#1] #2}}
\pagestyle{fancy}
\fancyhf{}
\fancyhead[LE]{II.\ Disputatio Luther--Eck}
\fancyhead[RO]{Seitz 1903}
\fancyfoot[LE,RO]{\thepage}
\begin{document}
\begin{titlepage}
\centering
{\LARGE Der authentische Text\\der Leipziger Disputation\par}
\vspace{1em}
{\large II.\ Disputatio Luther--Eck (4--14 de julio de 1519)\par}
\vspace{2em}
{\normalsize Transcripci\'on diplom\'atica de Seitz 1903 \textperiodcentered{} impresas 55--218\par}
\vfill
{\small Edici\'on de trabajo \textperiodcentered{} \today\par}
\end{titlepage}
\noindent\textit{Fuente: Seitz, Der authentische Text der Leipziger Disputation (Berl\'in 1903), ejemplar Internet Archive. Transcripci\'on p\'agina por p\'agina con aparato cr\'itico \'integro (siglas M\,U\,Uv\,Mc\,L\"o\,P\,R.); u/v, i/j y puntuaci\'on originales; [\dots] = ilegible; [NNN] = referencias del impreso. Anomal\'ias documentadas en \texttt{lotes/RESOLUCIONES.md}.}

\bigskip
II.

In nomine domini Amen\apnota{1}{Diese Vorrede findet sich nur in P, vgl. oben p. 4.}. Anno ab eodem incarnato M. D. XIX. indictione septima, die vero lunae, quarta mensis Julii, pontificatus sanctissimi in Christo patris et domini nostri, domini Leonis divina providentia Papae decimi anno septimo, hora septima antemeridiana, postquam per aliquot dies, nempe a vicesima septima Junii, celebrata esset disputatio inter egregios et praeclaros viros ac dominos Johannem Eccium, canonicum Eystetensem et procancellarium Ingolstatensem, et Andream Rodolphi Bodenstein Carolostadium, canonicum et archidiaconum Wittenbergensem, artium decretorum et theologiae doctores, in materia de libero arbitrio: coepta est disputatio inter magnificos D. doctores D. Johannem Eccium praefatum et dominum Martinum Lutherum Augustinianum Wittenbergensem in nostra notariorum publicorum et testium infrascriptorum praesentia in hunc qui sequitur modum. Sed quia egregius dominus doctor Martinus in suae disputationis exordio ad priores praefatorum doctorum protestationes se retulit, ideo eaedem ex praecedenti disputatione transsumptae ad verbum sequuntur huius disputationis initio praeinsertae\apnota{2}{Hier folgen in P die Protestationen Karlstadts und Ecks, vgl. oben p. 14 f. [W. A. 254]}.

Dominus D. Martinus in hunc modum orsus est\apnota{5}{U Mc haben hier: Incipit disputatio excellentium theologorum, Johannis Eccii et Martini Lutheri Augustiniani, quae coepta fuit quarta die Julii M. D. XIX. hora septima. Protestatio patris M. Lutheri [Mc doctoris Martini].}.

'In nomine domini Amen\apnota{6}{U ἀρχήν.}. Protestationem utriusque egregii domini et\apnota{7}{fehlt in M} Andreae Carolostadii et Johannis Eccii amplector et sequor. Hoc unum addo, quod pro reverentia summi pontificis et Romanae ecclesiae libens\apnota{10}{M videns.} hanc materiam non necessariam et mire invidiosam praetermisissem, nisi per propositionem\apnota{12}{M protestationem.} E. D. Doctoris Johannis Eccii in eam pertractus fuissem.

Doleo etiam eos non adesse, quos maxime oportuit\apnota{14}{M oportet.}, Qui cum et privatim et publice toties me crimine haereseos\apnota{15}{M in haereseos.} prophanarunt, nunc cum instet cognitio causae se subtraxerunt, haereticae pravitatis inquisitores dico, qui fraternam monitionem\apnota{17}{M Lücke, Mc paternam admonitionem.} et doctrinam postposuerunt criminationibus suis.

[255] Dominus doctor Eccius\apnota{e3}{U Scopus. De potestate, immo de primatu Romani pontificis (Uv a. R. Additio). Eccii protestatio, Mc [rot] = U.}.

In nomine tuo, dulcis Jesu\apnota{e4}{fehlt in M.}.

Antequam in harenam descendam protestor coram, vobis illustrissimis, nobilibus, magnificis et excellentissimis dominis omnia per me dicenda et dicta, inprimis esse subiecta iudicio primae sedis et domini in ea sedentis, deinde\apnota{e8}{M dehinc.} quorumlibet aliorum, quorum est errantes corrigere et ad veritatis cognitionem reducere. Et quia reverendus pater in praefatione sua, quasi se excusando, testatur se libenter 'hanc' obmisisse\apnota{e10}{U omisisse.} materiam pro summi pontificis reverentia, nisi per meam pertractus fuisset propositionem: at meminerit reverendus pater: nisi prius ipse in resolutorio ante tempora Sylvestri negasset Romanum\apnota{e13a}{M Romanam.} pontificem aliis\apnota{e13b}{M aliis esse.} superiorem, non fuisset necessarium me tredecimam illam posuisse propositionem. Et quod in actis coram legato\apnota{e14}{M legatis.} sedis apostolicae beatum Pelagium pontificem evangelica scripta\apnota{e15}{M dicta.} torsisse causatur. Qui tamen maxime omnium, ad sanctorum patrum intentionem verba Christi accepit. Frustra ergo reverendus pater in me causam torquere nititur, cui ipse non semel occasionem praestitit: sed iam ambagibus resectis\apnota{e18}{P reiectis.} principale institutum, deo duce, aggrediamur. Reverende pater, vestra conclusio XIII. meae contradicens, fatetur Romanam ecclesiam' esse' aliis superiorem iuxta frigidissima Romanorum pontificum decreta, intra quadringentos annos\apnota{e22}{M quadringenta annis.} nata. Quibus obstare dicitis, textum sacrae scripturae et historias approbatas mille et\apnota{e23}{et fehlt in M.} centum 'annorum'. Contra quam sic oppono.

Monarchia et unus principatus in ecclesia dei est de iure divino et a Christo institutus. Quare textus sacrae scripturae vel historiae approbatae ei non adversantur, quoniam ecclesia illa militans, quae est velut unum corpus iuxta divi Pauli sententiam, est instituta et facta ad imaginem ecclesiae triumphantis: in qua est una monarchia

omnibus per ordinem dispositis, usque ad unum caput, scilicet deum. Quare et talis ordo a Christo in terris est constitutus, cum fateatur Johannis 5 [v. 19] filium non facere quicquam, nisi quod viderit patrem facientem. Quare non de caelo est, qui capiti subesse recusat. Sicuti non de caelo, sed de lucifero est, qui deo non vult subiici. Quae omnia latissime confirmare possum\apnota{34}{U confirmari possunt.} maxime per sacram illam animam beatum Dionysium Areopagitam, libro de ecclesiastica hierarchia, ubi ait: Nostra enim hierarchia a deo traditis ordinibus sancte disposita, sanctis [256] et caelestibus hierarchiis conformis est. Ita Gregorius Nazianzenus in Apologetico, sacrosancta mysteria, ait ad\apnota{2a}{fehlt in M.} instar caelestis imaginis\apnota{2b}{P maginis.} celebrari. Per quae\apnota{2c}{M terque [gestrichen], Mc per quae.} utique in terris caelestibus ordinibus associamur. Nam quod monstrum esset ecclesiam esse acephalam\apnota{4}{M acevolam, U a. R. Id est, sine capite.}: quod omnes ferme haeretici moliti sunt, ut sanctus Cyprianus ad Rogatianum et Puppianum\apnota{5}{M Rogatianum et Cupianum, U rogationum et Puppianum.} innuit, ut sic debilitato capite errores suos et virus impune possent mentibus hominum inculcare. Et haec fuit ratio principalis cum aliis annexis, propter quam Parisiense felix\apnota{7}{M Parisiensis velut.} studium Johannem Tornacensem damnavit, negantem primatum ecclesiae Romanae. Sic et error fuit Vuiclefficus\apnota{9}{M wiewerbiticus, Mc wicleffiticus.} Romanam ecclesiam aliis non esse superiorem de lege evangelica\apnota{10}{U evangelii.}.

D. doctor Martinus.

Quando dominus doctor arguit esse omnino\apnota{12}{M U omnino esse.} caput ecclesiae universale, optime facit. Et si est aliquis, qui privato\apnota{13a}{M a. R. praefato.} pacto convenerit cum\apnota{13b}{fehlt in M.} dominatione sua ad defendendum contrarium: surgat in medium, ad me nihil pertinet.

D. doctor Eccius.

Quia reverendus pater dicit nihil ad se pertinere de contrario illius, quod intendebam probare, de iure divino esse monarchiam in ecclesia militante: sicut et triumphante, in quo ego\apnota{18}{fehlt in U.} eum collaudo, quod et in hoc divo Johanni in Apocalypsi consentit: Vidi civitatem sanctam novam, descendentem etc. At\apnota{20a}{M ad, Mc at.} propius ad rem accedentes: Si ecclesia militans non fuit sine monarchia\apnota{20b}{U monarcha.}, vellem audire, quis esset iste monarcha\apnota{21}{M ista monarchia.} alius aut unquam fuisset, nisi Romanus pontifex, aut quae\apnota{22}{M atque (statt ant quae).} alia prima sedes, nisi sedes Petri et eius successorum, iuxta illud beati Cypriani, qui in epistola tertia\apnota{23}{M P U secunda.} ad Cornelium Romanum pontificem scribens, contra Novatianos subdole Romam accedentes, inquit: Post ista adhuc insuper, pseudoepiscopo\apnota{25a}{P U cedo episcopo, M sedo episcopo.} sibi ab haereticis constituto\apnota{25b}{W. A. setzt fälschlich hinzu Romam und Anm. „Romam fehlt"; corrigiert am Ende des 2. Bandes.} navigare audent et ad Petri cathedram\apnota{26}{U. a. R. non ad Petris cabellum.} atque ecclesiam principalem, unde sacerdotalis unitas exorta est, a schismaticis et prophanis literas ferre, nec\apnota{27}{U ne.} cogitare eos esse Romanos, quorum fides ab Apostolo laudata est, ad quos perfidia habere non possit accessum. Et Hieronymus idem testatur contra Luciferianos, ecclesiae salus, ait, in summi sacerdotis dignitate pendet: cui si non exors\apnota{31a}{M exsors.} quaedam et ab omnibus eminens detur potestas, tot in ecclesiis efficientur\apnota{31b}{M tot in ecclesiis tot efficienter (so!)} schismata, quot sacerdotes. Et quod ille summus sacerdos sit Romanus, liquet apud eundem beatum Hieronymum in duabus epistolis ad Damasum papam, cuius ferme singula verba ad institutum faciunt\apnota{34b}{M a. R. facio.}. Verum brevitatis studio illa signemus: Cum successore piscatoris et discipulo\apnota{35}{U discipuli.} Christi loquor. Ego nullum praemium sequens nisi Christum, beatitudini tuae id est cathedrae [257] Petri consocior\apnota{1b}{M consocio.}, super illam petram fundatam ecclesiam scio; et infra: Quicumque tecum non colligit ille dispergit. Ex quibus quisque bonus christianus facile coniicit et intelligit\apnota{3b}{et intelligit fehlt in U.} unitatem sacerdotalem a Romano fluere pontifice, et illam fuisse semper cathedram principalem et omnibus aliis praelatam\apnota{4b}{M praelatis.}, et quod illa sit petra, sicut Hieronymus dicit se\apnota{5b}{M se dicit.} scire, super quam fundata sit ecclesia. Vel reverendus pater alium monarcham\apnota{6a}{M aliam monarchiam.} ecclesiae priscis temporibus assignet\apnota{6b}{M U assignet (M assignat) alium monarcham etc.}.

Dominus D. Martinus.

Monarchiam ecclesiae militantis prorsus confiteor eiusque caput non hominem sed Christum ipsum idque auctoritate divina prima ad Corinthios 15 [v. 25]: Oportet autem\apnota{11}{fehlt in M U.} illum regnare, donec ponat omnes inimicos sub pedibus eius. Et paulo ante: Deinde finis, cum tradiderit regnum deo et patri et evacuaverit omnem principatum. Quod beatus\apnota{13ab}{fehlt in U.} Augustinus libro I.\apnota{13bb}{libro I fehlt in M.} de trinitate cap. ultimo\apnota{13cb}{fehlt in M, ultimo cap U.} exponit de regno Christi praesentis temporis. Ita videlicet, quod Christus caput ecclesiae per fidem transferet nos, qui regnum eius sumus, ad\apnota{15}{ad P M Uv, per U.} speciem. Sic Matth. ultimo: Ecce ego vobiscum sum usque ad consummationem saeculi. Item actuum 9 [v. 4]: Audivit Paulus de caelo: Saule, Saule,

quid me persequeris? ubi iterum Augustinus dicit, caput loqui pro membris suis\apnota{18}{fehlt in M.}. Quare prorsus audiendi non sunt, qui Christum extra ecclesiam militantem trudunt in triumphantem\apnota{20a}{Mc ecclesiam triumphantem.}, cum sit regnum fidei\apnota{20b}{M filii.}. Hoc est: quod caput nostrum non videmus et tamen\apnota{21a}{M attamen.} habemus. Juxta illud Psalmi 122\apnota{21b}{P CXIX, M U 119}: Illic sederunt sedes in iudicio super domum David. Multae scilicet\apnota{22a}{Uv sunt.} sedes, in quibus unus sedet Christus\apnota{22b}{Mc scilicet Christus.}. Sedes videmus, non sessorem vel regem. Deinde\apnota{24a}{M U proinde.} ad auctoritates domini\apnota{24b}{M U egregii domini.} D. veniendo\apnota{24c}{M veniendo ad auctoritates etc.} quando asserit in ecclesia militante\apnota{25a}{M triumphante, Mc militante.} esse unum principatum divino iure et a Christo institutum\apnota{25b}{M instituto.}, sua dicit, sed nihil probat. Nam prima eius auctoritas Pauli, nimirum ad\apnota{26}{fehlt in P U.} Ephes. 4 [v. 15], quae Christum caput ecclesiae dicit, pro me, contra eum facit\apnota{27}{U valet.}, cum certissime loquatur ibi\apnota{28a}{M ibi loquatur.} de ecclesia\apnota{28b}{M ecclesiae.} militante et Christum appellet\apnota{28c}{M U appellat.} eius caput. Est etiam\apnota{29a}{M Est est.} eadem sententia, huic contraria\apnota{29b}{M U contraria huic.} I. ad Corinth. 3 [v. 5; 1 v. 13]: Quid est Apollo? quid Cephe? quid Paulus? divisus est Christus?\apnota{30}{U add. etc.} Ubi manifeste prohibet aliud caput quam Christum.

Secunda eius\apnota{32a}{fehlt in M.} auctoritas Johannis 5.: Non potest\apnota{32b}{M poterit.} filius facere quicquam\apnota{32c}{M quicquam facere.}, nisi quod\apnota{33}{fehlt in M.} viderit patrem facientem, neque de ecclesia militante neque triumphante loquitur, sed omnium doctorum sententia de aequalitate sua cum patre, scilicet\apnota{35a}{U sed, Uv scilicet.} quod nihil possit facere pater, quin idem possit filius\apnota{35b}{M quod nihil faciat nec facere possit, quod non idem facit filius, U quod nihil facit aut possit facere pater, quin idem possit et filius.}. Omitto\apnota{36a}{M 'Omitto' omitto.} quod dixit\apnota{36b}{fehlt in M.} non esse de caelo, qui capiti subesse recuset\apnota{36c}{U recusat.}, et esse de lucifero, 'qui\apnota{37a}{M est, qui.} deo nolit\apnota{37b}{U non vult.} subiici'; quia sicut praecedentia sunt male\apnota{37c}{U male sunt.} assumpta, ita\apnota{38a}{M sic, Mc ita.} haec\apnota{38b}{U et haec.} male illata.

Tertia, 'quae et Dionysii', nihil contra nos. Non enim negamus hierarchiam ecclesiasticam\apnota{40a}{M coelestem, Mc ecclesiae.}, sed de capite disputamus monarchiae, non hierarchiae\apnota{40b}{M monarchia, non hierarchia.}.

[258] Quarta\apnota{1a}{U Quarto.}, quae est\apnota{1b}{fehlt in M U.} Gregorii Nazianzeni inducitur\apnota{1c}{M dicitur Mc additur (adducitur U), al. inducitur.}, quod per sacrosancta mysteria sociamur\apnota{2a}{M sociamur.} ordinibus caelestibus\apnota{2b}{Mc coelestibus ordinibus, U sanctis ordinibus coelestibus.} etc., 'a quovis grammatico intelligitur' nihil sonare neque de monarchia neque 'de capite'\apnota{3}{M capite hierarchia, letztes Wort gestrichen.}. Jam quod addidit\apnota{4a}{M nam quod adducit} esse monstrum, si\apnota{4b}{M quod, Mc si.} ecclesia acephala\apnota{4c}{M A scevola.} esset, fateor\apnota{4d}{M (Mc) 'fateor' fateor.}; sed hoc caput nec ipse\apnota{5a}{fehlt in M.} D. doctor aliud dare\apnota{5b}{M dari.} possit quam Christum, quod probo evidenter: quia si suum caput, quod Romanum pontificem appellant\apnota{7a}{M vocat romanum pontificem, U Romanum pontificem appellat.}, moritur, ut\apnota{7b}{Mc quia.} est homo, iam ecclesia est acephala\apnota{7c}{M acephola.}. Si autem Christus interim caput est ecclesiae, donec alius eligatur, non minus monstrum\apnota{8b}{M U est.}, Christum 'cedere' vivo pontifici et succedere mortuo.

Quinta beati Cypriani, qui haereticos insectatus\apnota{10b}{M insectatur, Mc insectatus.}, quod debilitato capite amoliti\apnota{11}{U moliti.} sunt errores suos impune hominibus inculcare: prorsus nihil facit ad propositum. Loquitur\apnota{12b}{M U de quolibet.} enim non de Romano capite, sed quolibet\apnota{12c}{M U de quolibet.} capite cuiuslibet\apnota{13a2}{U cuiusque.} 'episcopatus'. Et\apnota{13b2}{M Additur et, Mc a. R. Atamen (so!) et.} si egregius\apnota{13c}{fehlt in M.} D. doctor voluerit stare auctoritati 'Cypriani', hac hora finiemus disputationem\apnota{14a2}{fehlt in U.}. Ipse enim Cornelium Romanum\apnota{14b}{U a. R. Eccius dixit non fiet.} pontificem nunquam aliter salutavit quam\apnota{15b}{U salutat nisi.} charissimum fratrem. Deinde describens episcoporum electiones et confirmationes\apnota{16a}{M electionem et confirmationem.} per multas epistolas ex divinis literis efficacissime\apnota{16b}{M longissime, Mc efficacissime.} probat, ad plebem pertinere et vicinos duos episcopos\apnota{17b}{M suos episcopos duos.} vel tres, quemadmodum et sancitum est in sacratissimo concilio Niceno\apnota{18b}{U Niceno concilio; M add. ut beatus Augustinus allegat, gestrichen, weil vom Disputator widerrufen.}. Quin idem beatus martyr, ut beatus Augustus allegat libro 2 ca. 2 de baptismo\apnota{19a}{Mc add. parvulorum.}, sic dicit\apnota{19b}{M dicens.}: Neque enim quisquam nostrum episcopum se essee piscoporum constituit\apnota{20c}{M constituat.}, aut tyrannico terrore\apnota{21a2}{U errore.}, ad 'obsequendi\apnota{21b2}{Mc obsequendam statt ad obsequendi.} necessitatem' collegas suos adigit\apnota{21c}{M addidit.}, quando habet omnis episcopus pro licentia libertatis et potestatis suae arbitrium proprium, tanquam ab alio iudicari non possit\apnota{23b}{M posse.}, quomodo nec 'ipse potest alterum iudicare'. Sed expectemus universi\apnota{24a2}{M omnes.} iudicium domini nostri Jesu Christi\apnota{24b2}{iudicium domini nostri Jesu Christi fehlt in M, dafür eine Zeile frei.}.

Quod autem dicit de Romana et Petri cathedra ortam esse\apnota{25b2}{fehlt in M.} sacerdotalem unitatem\apnota{26b}{M dignitatem.}, libentissime admitto, quantum pertinet ad ecclesiam occidentalem. Nam revera etiam Romana ecclesia orta est ex Hierosolymitana: haec est proprie matrix\apnota{28d}{U matrix proprie, M nur matrix.} omnium ecclesiarum.

Sed nec valet consequentia: Ex Romana ecclesia orta est sacerdotalis unitas; ergo ipsa est caput et domina omnium\apnota{30}{U add prima.}. Alioquin insuperabiliter concluderet Hierosolymitanam esse caput et dominam\apnota{31}{et dominam fehlt in M.} omnium.

Ultima auctoritas divi Hieronymi non recte inducitur ab egregio domino doctore, etiam si per omnia esset\apnota{33a}{M esse, Mc esset.} vera auctoritas Hieronymi\apnota{33b}{M beati Hieronymi.}, quia d. doctor intendit probare ecclesiae Romanae potestatem monarchicam iure divino et\apnota{34}{fehlt in U.} a Christo institutam. Hoc verba Hieronymi non habent, quia dicit: Cui\apnota{35a}{fehlt in M} si non exors quaedam et\apnota{35b}{fehlt in U.} ab omnibus eminens\apnota{36a}{M eminus, Mc eminuens.} detur potestas, tot exorientur\apnota{36b}{M efficenter (so!), Mc efficerentur, U in ecclesiis efficientur.} in ecclesia schismata quot sacerdotes\apnota{37a}{M add sunt.}. Detur (inquit)\apnota{37b}{M add non habet.}, hoc est iure humano possit\apnota{37c}{M U posset.} fieri, consentientibus ceteris omnibus fidelibus. Nam nec ego hoc nego\apnota{38}{M et ego non nego.}, si consenserint totius orbis fideles\apnota{39a}{M omnes fideles, a. R. omnes totius orbis.} in Romanum\apnota{39b}{M Romanam.} vel Parisiensem\apnota{39c}{U pariensem.} vel Magdeburgensem vel quemcumque [259], ut esset primus pontifex 'et summus': hunc propter reverentiam totius ecclesiae fidelium sic consentientis\apnota{2a}{M consentientes.} habendum esse summum monarcham\apnota{2b}{M unam monarchiam, a. R. summum monarchum (so!)}. Hoc autem nec factum est unquam neque fit neque fiet\apnota{3a}{M factum est nunquam neque fit neque fiet, a. R. = P U.}, cum ad nostra usque\apnota{3b}{M U usque ad nostra.} tempora Graeca ecclesia non consenserit nec tamen habita sit\apnota{4a}{M U sit habita.} haeretica. Et hanc esse\apnota{4b}{M video esse.} divi Hieronymi sententiam\apnota{5a}{U sententiam divi [fehlt in M] Hieronymi.} probo ex epistola ad Euagrium\apnota{5b}{M ephasum.}, ubi dicit: Ubicumque episcopus fuerit sive Romae sive Eugubii seu Constantinopoli sive Regii sive Alexandriae sive Tanis, eiusdem meriti et eiusdem sacerdotii est. Potentia divitiarum et humilitas paupertatis vel sublimiorem vel inferiorem facit. Ceterum omnes apostolorum successores sunt\apnota{9}{U a. R. Martinus: Das sie den Hiero: nicht lang verbrant haben.}. Recitatur haec epistola in decretis non frigidis\apnota{10a}{M vidi.} dist. XCIII. c. legimus. Idem super epistolam ad Titum\apnota{10b}{M deditum, Mc Titum.}: Idem\apnota{11a}{M Sicut ergo est vor Idem, gestrichen.} est ergo\apnota{11b}{M ergo est.} presbyter qui episcopus. Et antequam diaboli instinctu studia in religione fierent et diceretur in populis\apnota{12a}{M populo.}: Ego sum Pauli, ego 'Apollo, ego Cephae'\apnota{12b}{U Ego sum Pauli, ego sum Cephae.}, communi presbyterorum consilio ecclesiae gubernabantur\apnota{13}{M ecclesia gubernabatur.}. Postquam vero unusquisque 'eos', quos baptizabat, suos esse putabat, in toto orbe decretum est, ut unus de presbyteris electus supraponeretur\apnota{15}{M unus vir super presbyteros superponeretur, U = P, nur superponeretur.}. Et inductis auctoritatibus scripturae dicit in fine: Sicut ergo presbyteri sciunt se ex ecclesiae consuetudine ei, qui sibi praepositus fuerit, esse subiectos, ita episcopi noverint se magis consuetudine quam dispensationis divinae\apnota{18a}{U dispositionis dominicae, W. A. Anm. 2 vermutet die richtige L. A.} veritate presbyteris\apnota{18b}{M presbyteres.} esse maiores. Ideo quod dominus doctor addidit, a Hieronymo intellectum esse 'summum' pontificem Romanum, dum dicit: Cum successore piscatoris et discipulo\apnota{20d}{U discipuli.} Christi loquor et beatitudini tuae id est cathedrae tuae\apnota{21a}{U cathedrae Petri, M capite Petro, a. R. catedrae.} consocior\apnota{21b}{fehlt in M, U consotius.}, super\apnota{21c}{M et super.} illam petram fundatam ecclesiam scio: nihil facit. Non sequitur, consocior huic ecclesiae, ergo haec est sola prima. Nec sequitur, est fundata supra\apnota{23}{M super.} petram, ergo sola est fundata.

Accedit ad hoc decretum concilii Africani dist. XCIX. c. Primae: Primae, inquit, sedis episcopus non appelletur princeps sacerdotum aut summus sacerdos aut aliquid huiusmodi, sed tantum primae sedis episcopus. Universalis autem pontifex nec Romanus appelletur. Quare si divino iure staret monarchia Romani pontificis\apnota{28}{M Romana ierarchia summorum pontificum.}, haec omnia essent haeretica: quod est temerarium asserere. Et in fine audiamus ipsum dominum, qui Luc. 22 [v. 24-26] dicit: Facta est autem contentio\apnota{62}{U contentio discipulorum.}, quis eorum videretur esse maior. Dixit autem eis: Reges gentium dominantur eorum\apnota{32}{U earum.}, et qui potestatem habent super eos, benefici\apnota{33}{M heretici.} vocantur: vos autem non sic, sed qui maior est vestrum, fiat sicut minor\apnota{62}{M add. Ecce hoc (haec?) est mea sententia.}.

D. D. Eccius\apnota{35}{In M beginnt mit diesem Abschnitt eine neue Seite Bl. E 3a. Auf der vorangehenden findet sich ein andrer (von Eck widerrufener?) Anfang: „Reverendus pater debens;" dann folgt ein leerer Raum von etwa 4 Zeilen.}.

Reverendus D. pater satis instructus descendit in harenam, materiam suam habens bono ordine contextam in libello impresso et per eum composito [260]. Quare illustres\apnota{1a}{U illustrissimae.} dominationes vestrae, magnificentiae et dominationes\apnota{1b}{M et magnificentiae statt magnificentiae et dominationes.} veniam dabunt Eccio, multo iam tempore aliis negotiis occupato, si non tam rotunde et accurate e vestigio\apnota{62}{PU evestigio} tanta cumulare poterit, quanta\apnota{62}{P quantam.} reverendus pater iam congessit. Venit enim disputaturus, non librum editurus. Verum quid reverendus pater dixerit, per ordinem dispiciamus.

Primo omnium caput esse ecclesiae Christum probare intendit: quod tamen fuerat supervacaneum, cum nullus hoc negare praesumat, nisi qui Antichristus sit\apnota{8a}{fehlt in M.}. Demiror\apnota{8b}{MU miror.} tamen vehementer, quod hoc

non perpendit, sicut pollicetur in actis coram legato sedis apostolicae posse se\apnota{9a}{M se posse.} praestare quendam iuristotheologum\apnota{9b}{M iuristam theologum, U iuristatheologum.}, plura posse esse\apnota{10}{U esse posse, M om. posse.} capita subordinata. In quo mystici capitis ratio aut etiam symbolici\apnota{11}{M et symbolici, U aut symbolici.} a ratione capitis naturalis deficit. Quare praeter caput Christum et aliud oportere quaeri caput in ecclesia statim probabitur. Neque ei suffragatur, quod apostolum inducit I. ad Cor. 3 [v 5; 1 v. 13]: Divisus est Christus etc.\apnota{13}{fehlt in U.} Nam etsi Paulus ibi Petri meminerit, non tamen falsum dixit beatus Hieronymus libro 1. contra Jovinianum, mihi columna 18: Unus eligitur (loquitur de Petro) ut capite constituto schismatis\apnota{16}{M ut schismatis.} tollatur occasio. Clare appellat Petrum caput in ecclesia constitutum\apnota{17a}{U esse constitutum.}. Sed haec missa faciamus\apnota{17b}{M U facimus.}, solutiones ad nostra inducta repulsuri\apnota{18}{M refusuri.}.

Primo cum respondit\apnota{19a}{M U respondet.} ad illud Johannis 5: Neque enim potest filius\apnota{19b}{M U filius potest.} facere quicquam, nisi quae viderit patrem facientem, dixit secundum sanctos patres exprimi hic\apnota{21}{fehlt in M.} aequalitatem patris et filii. Verum legat attentius quaeso reverendus pater beatum et inadulabilem\apnota{22}{M inadulabilem, Mc adulabilem.} patrem Bernhardum libro 3. de consideratione ad Eugenium, ubi de forma ecclesiae loquens et eam esse de iure divino probans, nostram formavit rationem\apnota{24a}{M formavit probationem, U rationem formavit.}, mihi columna 7.\apnota{24b}{mihi col. VII fehlt in M.} Nec vilem reputes\apnota{25}{U reputans, Uv reputes.} formam hanc, quia in terra est: exemplar habet in coelo. Neque enim filius potest facere quicquam, nisi quae viderit patrem facientem. Praesertim cum ei sub Moysi nomine dictum sit: Vide, omnia\apnota{27}{M ne omnia.} facias secundum exemplar, quod tibi in monte monstratum est. Viderat hoc, qui dicebat\apnota{28}{U dicebar.}: Vidi civitatem sanctam\apnota{29}{U add. etc.}. Ego enim propter similitudinem dictum reor: quod sicut illic Seraphim et Cherubim\apnota{30a}{M cherubim et seraphim.} et 'ceteri' quique\apnota{30b}{U ceteri qui; M ceteri, constitutique, gestrichen constituti.} usque ad angelos et archangelos ordinantur sub uno capite deo, ita hic quoque sub uno summo pontifice primates vel patriarchae, episcopi, presbyteri\apnota{32}{M Archipresbyteri statt episcopi, presbyteri; U Archiepiscopi, episcopi, presbyteri.} vel abbates et reliqui in hunc modum. Tunc\apnota{33}{fehlt in M.} 'subdit Bernhardus': Non est parvi pendendum, quod et deum habet\apnota{34}{U habeat.} auctorem, et de caelo ducit originem. Quis iam non intelligit\apnota{35}{U intelligat.}, hanc ecclesiasticam hierarchiam secundum Bernhardum a Christo institutam? Et velut Christus\apnota{36a}{M U deus.} caput est in caelo, ita summus pontifex caput est\apnota{36b}{M U sit.} in militante ecclesia, neutiquam\apnota{37}{M nec dicam.} Christum excludendo, cuius se fatetur vicarium. Sed ad ratiunculam plebeiam\apnota{38a}{U pleiam, Uv plebeiam.}, quam inducit\apnota{38b}{U induxit.}, ecclesiam manere acephalam\apnota{38c}{M ad Scepolam.}, [261] papa mortuo, nisi ridicule diceremus Christum cedere papae vivo et succedere papae mortuo, 'ridicula prorsus ratiuncula', quae\apnota{2}{M et quae.} in re tam seria, inter tot praestantes viros, vix est digna renarrari\apnota{3}{M U enarrari.}, cum ab initio dixerim, caput illud symbolicum a veri capitis 'naturalis'\apnota{4}{M naturalis 2 mal.} ratione in aliquibus convenientiis deficere. Neque Christus, cuius regnum manet in aeternum et sempiternum est eius\apnota{5}{U et cuius sempiternum est.} sacerdotium, cedit papae vel succedit, cum ei sit data omnis potestas in caelo et in terra (Matthaei ultimo) et defuncto papa modo coetus cardinalium, sicut defuncto episcopo capitulum, iura illa teneant\apnota{8}{M U teneat, P Uv teneant.}, usquequo novus pontifex deligatur\apnota{9}{M delegatur.}.

Quod vero secundo loco reverendus pater dicit, Cyprianum loqui de quolibet episcopo, non de Romano pontifice, plurimum miror, cum intelligentia dictorum ex causis dicendi\apnota{12a}{M dicendis.} sit accipienda et Cyprianus in locis per me\apnota{12b}{M a me.} adductis eos increpet, qui a Cornelio defecerunt, qui certe Romanus fuit pontifex. Ob id sciat reverendus pater me non nudis\apnota{14a}{M nullis, Mc nudis.} verbis satiari\apnota{14b}{M satiare, Mc satiari.}, ut sophistas pascere solemus. Quae ex Cypriano pro se affert, opponendo arbitror cumulatius efficiet. Nam quod Cyprianus Cornelium fratrem appellat et apostolos fratres fuisse, nemo nescit. Attamen Petrus sicut et eius successor Cornelius caput fuit apostolorum, apex et vertex iuxta beati Dionysii ca. 3. vel. 7. de divinis nominibus sententiam\apnota{64}{fehlt in M (Lücke).}. De electione et a\apnota{64}{fehlt in U (ist in Uv zugefügt).} concilio Niceno et a Cypriano conscripta nec promovent nec impediunt negotium. Et multo minus Augustinus lib. 2. cap. 2. de baptismo parvulorum debebat adduci post Cyprianum. Nam Augustinus arrogantiam et temeritatem eorum, qui per ambitionem et superbiam ad ecclesiasticas praelaturas se ingerunt, reprehendit: quod non debeant se constituere nec alios cogere ad se constituendum, cum quisque praelatus expectare debeat, ut vocetur, sicut Aaron.

Tertio aliam Cypriani sententiam, ex tertia epistola\apnota{26}{P M U ex secunda epistola.} ad Cornelium, diluturus\apnota{64}{M dilaturus.} respondit unitatem quidem sacerdotalem in occidentali ecclesia a Romana exortam, non in orientali. Obticuit hic reverendus pater, quod Cyprianus praecedenter Romanam ecclesiam appellavit cathedram Petri\apnota{64}{M U petri cathedram.} atque ecclesiam principalem. Sed quid sua solutio possit, medullam verborum inspicienti est manifestarium. Nam reverendus pater more grammatico ortum unitatis sacerdotalis intellexit\apnota{32a}{M U intelligit.} quo ad auspicia\apnota{32b}{U auspicium.} et initium\apnota{32c}{M initia.}, cum profecto Cyprianus voluerit explicare ortum commissionis,

subordinationis seu influxus: ut ab uno Petro, velut a capite, in ceteros omnes iurisdictio sit derivata. Alioquin non dabit unum sacerdotem nec Hierosolymis. Taceo, quod eum non relevat glossula adiecta de occidentali ecclesia, cum beatus Hieronymus ab initio suae epistolae ob id orientalem ecclesiam, ex oriente scribens, schismaticam dicat, quae indiscissam domini tunicam et desuper contextam minutatim per frusta\apnota{39a}{W. A. notiert: „per frustra alle Gesammtausgaben und Löscher"} discerpit et\apnota{39b}{fehlt in U.} Christi vineam (inquit) vulpes exterminant, denotans scilicet Hieronymus, quod sponsa in canticis [2 v. 15] conquerebatur [262]: Capite nobis vulpeculas, quae demoliuntur\apnota{1a}{M demoliantur.} vineam. Taceat ergo quaeso\apnota{1b}{M quaero.} reverendus pater et nobis non insultet cum Graecis et orientalibus, qui a Romana deficientes ecclesia a fide quoque christiana facti sunt exules. Consectarium est, quo pacto accipienda sit illatio\apnota{4}{M accipiendo sit dilatio.}: Est radix, ergo domina. Non enim de radice loquimur initii vel temporis, sed de radice influxus et principalitatis.

Quarto loco nititur reverendus pater se evolvere ex verbis Hieronymi et illa vitare: nam summo sacerdoti tribui\apnota{8}{M U tribuit.} quidem et dari summam dignitatem, sed hoc fiat iure humano. Sed cur tunc b. Hieronymus Damasum piscatoris appellat successorem et cathedrae\apnota{10}{M cathedra.} Petri vult associari, illud divinum citans Matth. 16: Super illam petram fundatam ecclesiam scio? Quod ita reliquis ecclesiis appropriari non potest, ut Bernhardus ratiocinatur. Et proh dolor, cum summa christianorum\apnota{13}{M christiana.} iniuria sumus experti, portas inferorum praevaluisse ecclesiae Hierosolymitanae, Antiochenae, Alexandriae, addo quoque Bohemicae. Quod tamen de ecclesia, quae fundata est supra\apnota{15}{M super.} petram, veritas incorrupta Christi non patitur. Sed verissimum est in his\apnota{16}{U iis.}, quae sunt fidei, quod Hieronymus in eadem epistola affirmat: ubicumque fuerit corpus, ibi congregabuntur et aquilae. Profligato a mala sobole patrimonio apud vos solos incorrupta patrum omnino servatur auctoritas. At hoc in principali adhuc evidenter probabimus\apnota{20a}{U probavimus, Uv probabimus.}: non iure humano primatum obtinere ecclesiam Romanam\apnota{20b}{M Romanam ecclesiam.}, sed divino. Praestat tamen diluere, quae venerandus pater pro se ex Hieronymo adduxit. Primo ad Euagrium, ubi episcopum Romae et Eugubii\apnota{23}{M U egubii.}, Constantinopoli et Regii eiusdem meriti dicit et sacerdotii. Hoc scivimus, antequam Theognis\apnota{25a}{M diagonos.} nasceretur. Neque\apnota{25b}{Uv a. R. (von moderner Hand?) atqui, namque, sed.} papatus est ordo ultra episcopatum. Nam sic et alio loco aequales dicit fuisse apostolos, ubi tamen primatum sancto Petro non abstulit. Sed quod reverendus pater tantopere a me efflagitabat, ne

diverticula quaererem, quod nunquam facere sum solitus, ipse\apnota{28}{U ipsum.} cano. Legimus\apnota{29a}{M legem.} adducens 93. dist. me ad hoc ducit Brocardum\apnota{29b}{M U procardum; W. A. erinnert an das französ. brocard.}, quod domini canonistae et theologi controvertunt\apnota{30a}{fehlt in U, W. A. ergänzt aus Löscher = M commentantur, a. R. alii antevertunt.}: an episcopatus ordo sit et\apnota{30b}{fehlt in M U.} proprio charactere insignis sacerdotio additus; de quo in praesentia, velut impertinente, nihil decerno. Hoc tamen dico apparere mihi, salvo\apnota{32}{M solvo.} semper meliori iudicio, non talem fuisse in primitiva\apnota{33}{M primativa.} ecclesia confusionem, ut episcopus a sacerdote non separaretur, cum duodecim apostoli LXXII excedant discipulos, in cuius rei testimonium beatum affero Dionysium, Hieronymo antiquiorem et primitivae ecclesiae hierarcham, qui lib. de ecclesiastica\apnota{36}{U ecclesiae.} hierarchia inter sacros ordines episcopatum ponit et supremum\apnota{37}{M supremam.} hierarcham, et quo pacto is debeat consecrari: cui assentiens\apnota{38a}{M assentient.} episcopos ab ipso ecclesiae initio\apnota{38b}{M indicio.} commune sacerdotium exceluisse opinor.

[263] Quinto\apnota{1a}{P M U Tertio, also kommt dieser Fehler auf Ecks Rechnung.} adduxit ca. concilii\apnota{1b}{M consilii.} Africani XCIX. distinct.\apnota{1c}{U add. can. 1., in Uv gestrichen.}, ubi concilium\apnota{2}{M consilium.} prohibet, nec Romanum pontificem debere dici universalem, et quod Christus hoc prohibuerit, Lucae 22 [v. 24ff.]: »Reges gentium\apnota{3}{U add. etc.}.« Respondeo verum esse, fastuosum universalis episcopi\apnota{4}{M episcopi universali (so!).} nomen fuisse prohibitum: non quod Romanus pontifex unquam a vero Christiano habitus sit non primus et summus pontifex, sed quod episcopus peculiariterRomanae ecclesiae non sit proprius cuiusque ecclesiae episcopus, sed primus, quod alioquin\apnota{7}{M alioqui.} debitus honor episcopis inferioribus non tribueretur. Verum nihil criminis inest, si quispiam Romanum pontificem universalem dixerit, pro primo. Rectius vero dicetur, eum non episcopum universalem, sed universalis ecclesiae episcopum\apnota{10}{U dicetur non episcopus universalis, sed universalis ecclesiae episcopus.}, veluti est Christi vicarius. quod dominus ambitiosam redarguit apostolorum contentionem, qualis inter saeculares quaeritur, hoc primatum ecclesiae Romanae non aufert. Sed quod beatus Gregorius primus fecit et agnovit se ideo ecclesiae catholicae praelatum, ut servum servorum se esse\apnota{14}{fehlt in M.} agnosceret, in qua re sequentes pontifices non convitiis\apnota{15}{M cum vicii.} lacessere, sed orationibus ut tales fiant, a deo impetrare studeamus.

Hora secunda pomeridiana continuabitur disputatio famulo publico universitatis proclamante\apnota{17a}{Die Zeitangabe fehlt in M U.}.

Eadem hora secunda pomeridiana\apnota{17b}{U Hora secunda pomeridiana quarta die Julii 1519.}.

D. D. Martinus.

Primam responsionem meam, quam ex auctoritate Pauli 1. ad Corinth. 3 ostendi\apnota{21a}{M U ostendi ab apostolo.} prohibitum, ne fideles sibi arrogarent vel Cephe\apnota{21b}{M Sceve.} vel Paulum\apnota{22a}{M Pauli.} vel Apollo, ad istum modum confutavit E. D. doctor\apnota{22b}{M Egregius doctor, U dominus doctor egregius.}, quod etsi\apnota{22c}{M si.} Paulus ibidem Petri\apnota{23}{M Petri ibidem.} meminerit, non tamen falso Hieronymus contra Jovinianum dicit: unus eligitur, ut\apnota{24}{U quod.} capite constituto occasio schismatis tollatur. Clare\apnota{25a}{M quare.} appellat Petrum caput\apnota{25b}{M caput Petrum.} in ecclesia constitutum. Adiecit: sed haec missa faciamus.

Respondeo. Non patior propter minorem auctoritatem inductam\apnota{26}{M adductam.} me divelli\apnota{27}{M afelli.} a maiore. Nec tantus est Hieronymus, ut propter eum Paulum deseramus\apnota{28}{M U deseram.}. Non ergo Paulus solum meminit ibi Petri, sicut conatur extenuare ista confutatio; sed cum plena auctoritate docet et prohibet, ne quis dicat se\apnota{29}{M se dicat.} esse Petri. Unde et idem capitulum\apnota{30}{M capitulo.} concludit: Omnia vestra, sive Paulus, sive Cephe\apnota{31}{M scevela (?).}, sive mors, sive vita: vos autem Christi, Christus autem dei. Stat ergo ista responsio\apnota{32}{responsio ista.} mea adhuc invicta. Et nisi robustius confutata fuerit, oppono eam omnibus praeteritis et futuris argumentis D. doctoris. Verbum enim dei super omnia verba hominum est. — Ad Hieronymum autem dico, quod et ego missa facio\apnota{35}{M missefacia.} ea, cum sit ambiguissimus locus, ut bene sensit D. doctor.

Secundam responsionem meam ad auctoritatem Johannis 5\apnota{37}{M 15.}, ubi dixi Christum loqui de aequalitate potentiae paternae: iussit me D. doctor attentius [264] legere beatum Bernhardum, qui hanc auctoritatem inducit pro ecclesia militante, sicut auditum est. Respondeo, divum Bernhardum veneror et eius sententiam non contemno. Sed in contentione accipiendus est sensus genuinus et proprius scripturae, qui stare in acie possit. A quo sancti patres nonnunquam locupletandae orationis gratia digrediuntur et sine culpa\apnota{5a}{et sine culpa fehlt in M (Lücke).}. Nunc\apnota{5b}{M Nunc, Mc setzt dafür Nunquam.} autem ex praecedenti et sequenti textu clarum fit, Christum loqui de aequali sua potentia\apnota{7a}{M potentia Christi, U omnipotentia sua.} cum patre; (textus est)\apnota{7b}{textus est fehlt in M U.}: Propterea persequebantur Judaei\apnota{7c}{M Judaei persequebantur.} Jesum, quia haec faciebat\apnota{8a}{M U faciebat haec.} in sabbato\apnota{8b}{M sabbata.}. Propterea\apnota{8c}{fehlt in U.} ergo magis quaerebant eum Judaei interficere\apnota{8d}{M Judaei interficere eum.}, quia non solum solvebat sabbatum, sed et patrem suum dicebat deum, aequalem se faciens deo. Respondit itaque

Jesus et dixit eis: Amen, amen, dico vobis, non potest filius a se facere quicquam, nisi quod viderit patrem facientem. Et sic clare patet, quod Bernhardus alio sensu tractat hoc verbum Christi\apnota{12}{fehlt in M.}.

Ad tertiam videlicet plebeiam\apnota{13}{M blebeiam.}, 'ut dixit', et ridiculam ratiunculam, qua dixi etiam sine papa ecclesiam habere caput suum, opposuit, quod non esset digna renarrari\apnota{15a}{M narrari, U enarrari.} coram tantis viris et in re tam seria, respondeo: Sit plebeia\apnota{15b}{M blebeia.} et ridicula, modo sit invicta. Necdum video eam confutatam. Non enim intelligo, si ecclesia ad tres vel quatuor menses non est acephala\apnota{17}{M a scevola statt non est acephala.} mortuo papa, modo sint alii episcopi: quomodo etiam sit acephala, si nullus sit papa\apnota{18}{M modo quomodo non sit a sceevola si nullus est papa, U quomodo etiam non [in Uv gestrichen] sit acephala, cum nullus sit papa.}. Nam quod adduxit de cardinalibus, penes quos sit ius\apnota{19}{U vis.} eligendi etc., responsionem meam confortat, cum hinc sequatur eo tempore, quo nondum erant cardinales, 'ut tempore Hieronymi', non potuisse esse papam\apnota{21}{M caput.}.

Ad quartam Cypriani auctoritatem, ubi dixi eum loqui de quolibet pontifice, obiecit, quod ex textu pervincatur eum locutum esse de Romano Cornelio\apnota{24a}{U Cornelio Romano.} pontifice\apnota{24b}{fehlt in M.} contra Novatianum\apnota{24c}{M U Novatianos.}, respondeo: Mea non refert, nec eam epistolam in memoria habeo. Hoc autem scio, quod divus Cyprianus per multas epistolas in hoc totus est, ut asserat cuilibet ecclesiae\apnota{26}{M episcopo (gestrichen) ecclesiae.} suum caput et suum episcopum per suffragium plebis et iudicium vicinorum episcoporum ordinari. Ideoque si est, ut D. doctor allegat, contra\apnota{28}{M de (gestrichen) contra.} Novatianos de Cornelio dictum\apnota{29a}{U dico.}, certum est eum locutum de 'capite' Romanae ecclesiae\apnota{29b}{M de Romana ecclesia, Mc = P.}, non universalis ecclesiae. Item et illud confutavit, quod Cyprianus Cornelium semper compellat\apnota{31}{M appellat.} fratrem, nunquam autem dominum, ut nunc episcopi faciunt, ponentes relativum sine correlativo. Obiecit, quod et Petrus apostolos\apnota{32}{P apostolus.} habuerit fratres et tamen caput et vertex apostolorum fuerit, iuxta Dionysium de divinis nominibus\apnota{33}{U a. R. de divinis nominibus.}. Respondeo: Si poterit egregius dominus doctor probare, quod Petrus unquam unum apostolorum ordinaverit, immo unum ex septuaginta discipulis aut unum\apnota{35}{M illum illorum unquam, U ullum illorum unquam.} miserit, cedo\apnota{36a}{M concedo.} omnia et volo\apnota{36b}{M volo volo st. et volo.} esse victus. Si autem ego probavero, quod ne omnes quidem apostoli potuerunt\apnota{37a}{M U potuerint.} unum\apnota{37b}{nec unum.} apostolorum mittere vel ordinare, rogo concedat mihi Petro nihil fuisse potestatis super ceteros apostolos. Quo sequitur multo minus potestatis esse episcopo, successori Petri\apnota{39}{M U Petro.}, super [265]

episcopos, successores\apnota{1}{fehlt in M.} ceterorum apostolorum. Nunc autem clarissimus textus est actuum 1 [v. 23ff.], quod Matthias apostolus ab universo coetu apostolorum et discipulorum\apnota{3a}{et discipulorum fehlt in M.} ordinari non potuit\apnota{3b}{M poterat.}, sed de caelo sicut et ceteri omnes electus et ordinatus est\apnota{4}{M ordinatus et electus [est fehlt].} a Christo. Ita et capitulo 13. Saulus et Barnabas segregante spiritu sancto assumpti sunt\apnota{5}{U sint.} in opus. Error ergo est apertissimus, quod Petrus habuerit potestatem super apostolos. Hoc sane fateor, apostolum Petrum fuisse primum in numero apostolorum et ei deberi honoris praerogativam, sed non potestatis. Aequaliter electi sunt et aequalem potestatem acceperunt. Ita et de Romano pontifice sentio, quod honoris praerogativa ceteris debeat anteferri, salva cuiusque aequali potestate, et non sicut Pelagius in frigidissimo suo decreto infert: Ubi maior auctoritas, ibi maior potestas, et manet ceteros\apnota{12a}{M ceteris.} obsequendi necessitas\apnota{12b}{M potestas.}.

Ad quintam\apnota{13}{M quintum.}, ubi electionem episcopi ex Cypriano et Niceno concilio attuli, E. D. doctor satis rhetorice contempsit, dicens, quod nec promoveant nec impediant negotium. Sed per hoc non dissolvitur mea responsio. Stat ergo adhuc decretum Nicenum, aut si non stat et contra ius divinum statuerunt, non catholicum concilium, sed diabolicum conciliabulum erit\apnota{17}{U appellandum erit.}. Sic cum putaret\apnota{18a}{M confutaret statt cum putaret.} Augustinum non fuisse inducendum\apnota{18b}{M dicendum.} et pulcherrima glossa sua Cyprianum, per Augustinum adductum, ita interpretatus est, quod Cyprianus ambitionem et superbiam eorum reprehenderet, qui se ipsos ingererent, antequam vocarentur sicut Aaron: hoc mera fiducia dixit. Ceterum\apnota{21}{M veterum.} textus satis clarus est, quod nullus episcopus, qui iam sit episcopus, se constituere debeat aliorum episcoporum episcopum\apnota{23a}{M se aliorum se constituere episcoporum (so!)}. Stat ergo mea adhuc\apnota{23b}{M U adhuc mea.} responsio.

Ad sextam graviter\apnota{25a}{M Sextam graviter, Mc ad sextam ubi graviter.} me reprehendit egregius dominus doctor\apnota{25b}{M U idem egregius dominus (fehlt in M) doctor.}, quod in altera auctoritate Cypriani\apnota{26}{fehlt in M.} obticuerim hoc vocabulum (principalem ecclesiam); deinde velut grammaticum me naso\apnota{27a}{P sub naso, M me nasum.} suspendit, quod dixerim\apnota{27b}{M dixerit, Mc dixerim.} ortam esse unitatem sacerdotalem\apnota{28a}{M unitatem sacerdotalem ortam esse.} ex Petri sede. Ideo novus dialecticus\apnota{28b}{M dialectius.} vel philosophus potius ortum hunc interpretatur 'commissionis, subordinationis'\apnota{29}{M ordinationis.} seu influxus. Alioquin, inquit, nec ex Hierosolymis dabit unum sacerdotem. Respondeo: Sive tacuerim\apnota{31a}{U docuerim, W. A. obticuerim.} sive locutus fuerim hoc vocabulum (principalem)\apnota{31b}{M U principale.}, idem

est. Nec enim potest principalis vocari respectu orientalis ecclesiae, ut satis dictum est.

Deinde commentum suum\apnota{33a}{fehlt in M.} de ortu influxus eadem facilitate\apnota{33b}{fehlt in M (Lücke), Mc füllt diese Lücke durch contemptu aus.} contemno, qua\apnota{33c}{M quae.} ipse finxit. Et non est difficile mihi dare\apnota{34}{mihi dare fehlt in M.} unum sacerdotem ex Hierosolymis, scilicet ipsum Christum, a quo incepit\apnota{35}{M U add et hinc orta, in Uv gestrichen.} et profecta est ecclesia, iuxta illud Esaiae 2 [v. 3]: De Sion exibit\apnota{36}{M ex Sion exivit.} lex 'et verbum domini de Hierusalem'. Jam quod addidit auctoritate Hieronymi orientalem ecclesiam fuisse schismaticam et indiscissam\apnota{38a}{fehlt in M.} domini tunicam minutatim\apnota{38b}{M minutatim domini tunicam.} discerpentem, nescio sane, quid velit. Nec enim potest dicere, quod tota orientalis ecclesia et semper fuerit [266] schismatica, sicut\apnota{1a}{U hinc, Uv sicut.} nec negare\apnota{1b}{M negari.} potest, quod et latina ecclesia sua habuit aliquando schismata et tamen mansit ecclesia. Proinde nihil est, quod iubeat me tacere et non insultare per Graecam ecclesiam, quia deficientes a Romana ecclesia simul facti sunt et a fide Christi exules. Ego potius rogo dominum doctorem Eccium, ut pro Ecciana modestia\apnota{5}{M modestiam.}, quam iactat, parcat tot milibus sanctorum, cum usque ad nostra tempora durarit ecclesia Graeca et sine dubio usque hodie durat et durabit. Non enim Christus accepit medium terrae Romanae, sed omnes fines terrae in possessionem\apnota{8a}{M possionem.} et hereditatem a patre Psal. 2 [v. 8]\apnota{8b}{fehlt in M (Lücke).}.

Ad septimum\apnota{9}{M Septimum, U Ad septimam.}, quod ex Hieronymo de summo sacerdote opposuit, meam responsionem dixit esse evasionem. Ideo confirmaturus dictum suum adiecit, cur beatus Hieronymus Damasum piscatoris successorem dicat\apnota{11}{successorem dicat fehlt in M, Mc hat appellat successorem.} et cathedrae Petri vult associari, illud divinum citans Mat. 16: Super illam petram fundatam ecclesiam scio, quod ita\apnota{13}{fehlt in M.} reliquis ecclesiis appropriari non potest. Et deinde complorat\apnota{14a}{M U comploravit.} Hierosolymitanae, Antiochenae\apnota{14b}{fehlt in M.}, Alexandrinae, tandem et Bohemicae 'casum', et quod eodem Hieronymo teste apud solos Romanos incorrupta patrum servaretur auctoritas\apnota{15}{M et apud Romanos Pauli servaretur auctoritas (z. t. von Mc aufgefüllt); in Uv ist patrum gestrichen.}. Respondeo et rogo dominum doctorem egregium, velit dicta patrum allegare secundum conscientiam, ne pro theologis 'videamur' esse sophistae\apnota{18}{M sophisticare.}. Nam summum sacerdotem eo loci Hieronymus appellat episcopum quemlibet, etiam ut qui de numero reliquorum sacerdotum sit elevatus\apnota{19}{M assumptus.}. Ideo nihil pertinet proprie ad

Romanum pontificem\apnota{20a}{proprie ad Romanum pontificem fehlt in M (Lücke); Mc setzt proprie ad reverendum patrem.}. Deinde illa\apnota{20b}{fehlt in M.} auctoritas Matth. 16 non appropriatur ecclesiae Romanae\apnota{21}{M U Romanae ecclesiae.} tantum, quod expresse indicant verba Christi, dum\apnota{22a}{U cum.} dicit: »ecclesiam meam«. Quaecunque ergo est ecclesia Christi\apnota{22b}{fehlt in U, Uv fügt es bei.}, est aedificata supra petram, non sola\apnota{23}{M et non sola, U et non solum.} Romana. Aut si non convenit ceteris ecclesiis hoc verbum, Romana ecclesia\apnota{24}{M ecclesia Romana.} erit sola et per consequens non prima. Quare unitas ecclesiae non ab unitate primatus Romani, sed longe melius iuxta Apostolum Ephes. 4 [v. 5] ab unitate fidei, baptismatis, domini pendet, sicut et frequenter Cyprianus in epistolis suis sentit. Nec mansit apud solos Romanos\apnota{28a}{M sola apud 'Romanos' (das letzte Wort steht über der Zeile), U apud Romanos solos.} incorrupta patrum auctoritas\apnota{28b}{M integritas, Mc auctoritas.}, nisi forte eo tempore, quo scripsit Hieronymus; immo tradunt historiae Liberium Romanum pontificem concessisse Arianis, et illud Hieronymus in viris illustribus scribit, quod Achatius Caesariensis, episcopus Arianus, Eusebii Ariani discipulus\apnota{31}{M episcopus discipulos (so!)}, auctoritate Constantii imperatoris\apnota{32}{Constantii imperatoris fehlt in M.} ordinavit Felicem in Romanum pontificem.

Octavo confutans auctoritatem Hieronymi a me inductam ad Euagrium dicit se scivisse\apnota{34}{M U omnes eiusdem.}, eiusdem meriti eiusdemque sacerdotii omnes episcopos fuisse\apnota{35}{M fuisse vel esse, U fuisse et esse, Uv fuisse.}, sed papatum esse ordinem ultra episcopatum. Sed non evicit responsionem meam, quia Hieronymus sublimitatem et inferioritatem episcoporum non iuri divino, sed consuetudini et potentiae divitiarum tribuit; ideo adhuc maneo cum Hieronymo.

Ad nonum: Super canonem Legimus 93. dis.\apnota{39}{P M U can. 99. dist., W. A. verbessert Ad nonam: Super etc.} dicit sibi non apparere [267] in ecclesia primitiva fuisse talem confusionem, ut episcopus a sacerdote non separaretur. Respondeo: Quid hoc ad me? Pugnet cum Hieronymo et canonibus. Quod autem Dionysium inducit inter sacros ordines episcopatum recensere, miror, quod non etiam inducit ex eodem auctore Romanae ecclesiae monarchiam, cum habeat tantum momenti ad hierarchiam, ut sine hac similitudo triumphantis ecclesiae consistere non posset. Decuit autem professorem\apnota{6}{M scriptorem.} scribendae hierarchiae argumento suscepto praesertim in praestantiore eius parte satisfacere. At Dionysius usque ad episcopum solum deducit.

Decimo\apnota{9a}{U ad decimam.} ad ca. »Primae«\apnota{9b}{fehlt in allen Quellen, ist nach 263, 1 ergänzt; cf. W. A.} 99. dis., ubi respondi prohibitum fuisse, ne Romanus pontifex universalis\apnota{10}{U universalis ecclesiae, Uv = P.} episcopus appellaretur, confutavit in hunc modum: Non quod Romanus pontifex non sit

primus et summus, sed quod episcopus peculiariter Romanae ecclesiae non sit proprius cuiuslibet ecclesiae episcopus. Respondeo: Quasi ulli mortalium possit tam stultus incidere affectus, ut una persona velit omnibus ecclesiis et singulis praeesse, ut tantam insaniam\apnota{15a}{M pro tanta insania.} necesse fuerit prohibere\apnota{15b}{M providere.}. Tamen\apnota{15c}{W. A. tum gegen alle Quellen.} relinquens hanc confutationem rectiorem dedit, scilicet non episcopum universalem, sed universalis ecclesiae episcopum Romanum\apnota{17a}{M U eum Romanum, W. A. esse Romanum.} pontificem dicendo\apnota{17b}{U dicendum.}. Nisi parcerem ei ego\apnota{18a}{so W. A., et ego alle Quellen.}, suspenderem et hanc solutionem. Sed iudicent iudices\apnota{18b}{fehlt in M.} et auditores.

Ultimo, ad auctoritatem Christi, ubi dixit: vos autem non sic, dixit ambitionem fuisse redargutam, sed non primatum. Respondeo: Haec est petitio principii, quasi iam obtinuerit esse scilicet principatum\apnota{22}{M iam principatum.}. Deinde clarus est textus, quod non tantum ambitionem, sed ipsam maioritatem super ceteros prohibuit.

D. D. Eccius.

Ad defensiones solutionum a reverendo patre datarum dico de primo, non increpasse Paulum eos, qui dixerant\apnota{27a}{M dixerint.} se esse Petri, si\apnota{27b}{U quia, Uv si.} habuissent intentionem primatus\apnota{28}{M U principatus.} apostolorum, sed quia singularitatem aspiciebant personarum: quod verba ipsa clare indicant divisionis et schismatis. Et\apnota{29}{fehlt in M.} quamquam merito Paulum praeferat Hieronymo, pie tamen credendum est Hieronymum eo loco sententiam Pauli bene intellexisse. Nec locus est ambiguus, ob id caput in ecclesia constitutum, ut schismatis tolleretur occasio, quod 'grammatico\apnota{32}{M quod grammaticus, Mc aliquod quo grammaticus.} satisfacit\apnota{33}{U satisfecit.} lectori', quam disciplinam praecipuam prae aliis partibus philosophiae ad theologiam subservientem\apnota{34}{M theologiae sufficientem, Mc add alii subvenientem, U theologiae subservientem.} dixit in disputatione quadam reverendus pater\apnota{35}{von dixit an fehlen die Worte in M, Mc dixit d. doctor in quadam disputatione.}.

Secundo. Nemo negavit Christum Johannis 5. aequalitatem cum patre affirmasse\apnota{37a}{M confirmasse.}, nisi Arianus\apnota{37b}{U Arrianos, W. A. Arriani.}. Neque alio sensu inducit Bernhardus, sed quod reverendus pater arbitratur sanctos patres sanctam induxisse scripturam [268] locupletandae orationis gratia\apnota{1a}{M gratia arbitratur.}, respuimus: quoniam talem\apnota{1b}{M pauli.} ostentationem de eis suspicari non licet.

De tertio de papa mortuo, tunc ecclesiam esse acephalam\apnota{3}{M ascevolam.}, dico: Nunquam negatum Christum esse caput ecclesiae, iuxta etiam glossam Canticorum 5 [v. 11]: Caput eius aurum optimum, glossa: Caput, id est

Christus, papa vero est eius vicarius. In consistorio autem episcopus et vicarius computantur pro eadem\apnota{7a}{U una.} persona; unde appellare a vicario ad episcopum interdicitur\apnota{7b}{M U non licet.}. De cardinalibus vero dixi, quod nunc hoc est iuxta ecclesiam institutam\apnota{8}{M U iam institutam.}, electio ad cardinales venerit\apnota{9}{M venit.}, Nicolao 'quodam pontifice ordinante', tamen tempore Hieronymi fuisse cardinales credo, nisi Hieronymus non fuerit presbyter\apnota{10}{fehlt in M.} cardinalis\apnota{11}{U a. R. Martinus respondit, Hieronymum non fuisse cardinalem.}.

Quarto de Cypriano dicit\apnota{12}{fehlt in U, M dicit quod iam non viderit, zitiert bei Löscher p. 346 [eam statt iam].}: certe fieri non potest, ut sancti martyris verba in Romanae dioecesis angustias coarctet\apnota{13a}{M coarctat.}, quia ad episcopum maximum\apnota{13b}{ad fehlt in P M U, W. A. vermutet episcopi Novatiani.} ex Numidia, quae regio est Africae ultra Atlantem\apnota{14}{U a. R. D. Martinus dixit citra Atlantem.}, ut Ptolemaeus et Strabo auctores sunt, ad Romam adventarunt\apnota{15}{M adventavit.}. Quod vero Cyprianus Cornelium fratrem appellavit, arbitror collectorum\apnota{16}{M U collectoris.} fuisse, non Cypriani mentem. Nam si sanctorum episcoporum epistolas legerimus, inveniemus plane\apnota{17}{U clare.} ea tempestate usitatiores fuisse conscriptiones illas laudabiles et magnificas, quam modo fit\apnota{18}{M U fiat.} Romano pontifici: ut de Ambrosio, Augustino, Hilario, ceteris\apnota{19}{M U ceterisque.} liquet patribus. Nam mutuo se vocant beatissimos, sanctissimos, deo amabilissimos\apnota{20}{U beatissimus, sanctissimus, deo amabilissimus.} etc. Ad adiunctum causatur, quod quaeram diverticula et ad rem non pertinentia. Cum bona venia\apnota{21}{U vinea, Uv venia.}, libenter dicerem: Turpe\apnota{22}{M turpi.} est doctori etc. Petit a me, ut probem Petrum ordinasse aliquem apostolum. At\apnota{23a}{U sed.} hoc instituto nostro\apnota{23b}{U ad institutum nostrum.} est impertinens. Non enim quaerimus, quis unum vel alterum ordinaverit\apnota{24}{U aut alterum ordinavit.}, sed quis primatum super alios a domino Jesu accepit\apnota{25}{M susceperit.}. Subsequens omnino respuo, cum huiusmodi fecit\apnota{26}{W. A. gegen alle Quellen facit.} illationem: Petrus nullum potuit ordinare apostolum; quare nec successor Petri successorem alterius potuit\apnota{27}{M poterit.} ordinare vel super eo habere potestatem\apnota{28a}{P patestatem.}. Nam assumptum est verum, illatum est manifestissime\apnota{28b}{M manifestarie, U manifeste.} falsum, cum summus pontifex modo potestatem habeat et ordinet alios episcopos. Verum\apnota{30a}{M vero.} ratio enodandi nodum erit illa\apnota{30b}{M erat ista, U erit ista.}, quod apostolatus officium fundamentum ecclesiae plus dicit quam esse episcopum. Quare Leo Decimus successor est Petri apostoli\apnota{32}{fehlt in M.}, non apostolus. Sed quod fatetur Petrum bene fuisse primum apostolum\apnota{33}{fehlt in M.} numero et honoris praerogativa et non potestatis, illud non satisfacit instituto. Primo quia evangelistae numerum apostolorum non aequaliter inceperunt, sicut

Chrysostomus super Matthaeum annotavit\apnota{35}{MU super Matthaeo annotat.}. Et glossa videatur Matth. X. Dein\apnota{36}{M U Deinde.} quod de potestate et honoris praerogativa variat [269], expresse est contra\apnota{1}{est contra fehlt in P.} sanctum martyrem Cyprianum, qui in tractatu De simplicitate praelatorum contra Novatianum, loquens de diaboli astutia, increpat eos, qui velut iustitiae ministri asserunt noctem pro die, interitum pro salute, desperationem 'sub obtentu\apnota{4a}{U obtutu (falsch, cf. W. A. Anm. 2), Uv obtentu.} spei', perfidiam sub praetextu fidei. Et\apnota{4b}{fehlt in M.} infra: Et quamvis apostolis\apnota{5}{M apostolis, Mc Christus.} omnibus post resurrectionem suam parem tribuat potestatem et dicat: sicut me misit pater etc., tamen ut unitatem manifestaret, unitatis eiusdem originem ab uno incipientem auctoritate sua disposuit. Hoc erant utique et\apnota{8a}{utique et fehlt in M.} ceteri apostoli\apnota{8b}{U discipuli.} quod\apnota{8c}{M quo.} Petrus: pari 'consorcio praediti'\apnota{8d}{Mc praediti doppelt.} et honoris et potestatis (hoc ponderandum\apnota{9}{M ponderando.}), sed exordium ab unitate proficiscitur, ut ecclesia una esse monstretur\apnota{10}{M demonstretur.}. Et infra: Hanc unitatem qui non tenet, legem dei non tenet, non tenet patris et filii fidem, vitam non tenet et salutem. Haec notatu\apnota{12}{P M notata.} dignissima Cypriani verba, qui inter praerogativam honoris et potestatis apostolos non variat.

Ad quintum. De electione dico, ut prius, nihil nos disputare de modo eligendi, sed de qualitate aut potius quantitate ipsius electi. Quare concilium Nicenum\apnota{16}{M Niceni.} non fuit conciliabulum. Sed quae iactum et conversationem respiciunt, pro temporis, personarum et locorum qualitate mutari posse\apnota{17}{M U possunt.}, ut in multis canonibus videmus.

Ad sextum. Quod observandus d. pater\apnota{19}{U doctor.} uti dialecticum opponit me finxisse distinctionem de duplici ortu, iam audivimus Cyprianum, et loquitur\apnota{20}{M U et loquatur, W. A. quod loquatur.} pro Eccio, cui non tantum suppetit ingenium, ut nova fingat, sed vetera\apnota{21}{M cetera.} sanctorum dicta pro virili explanet\apnota{22a}{M U explanat.}. Dum vero unum\apnota{22b}{U omnium, Uv unum.} mihi dat\apnota{22c}{M U dat mihi.} sacerdotem Christum, primo Cypriani intentionem\apnota{23}{U attentionem, Lö Anm. richtig.} non assequitur nec Hieronymi, quoniam Petrum primum apostolorum constitutum volunt, unde reliquorum sacerdotum potestas defluat, non quidem interius conferendo, quod Christi est velut capitis, sed potestatem ecclesiasticam communicando\apnota{26}{M communicatum.}.

Septimo, quod desiderat Eccianam modestiam, Graecos et orientales pro damnandis iudicantem, respondeo: Graecos longo tempore non solum fuisse schismaticos, sed\apnota{29a}{M et.} haereticissimos, ut tot errorum cumulus\apnota{29b}{M cumulos.} et pervicax obstinacia testatur\apnota{30}{U add Clementina de summa trinitate, in Uv gestrichen; M add de summa trinitate.}, ut de spiritu sancto, de confessione,

de trium evangelistarum falsitate et innumeris aliis, quamvis saepe fictam fecerint ecclesiae Romanae\apnota{32a}{U Romanae ecclesiae.} obedientiam, ut et\apnota{32b}{fehlt in M U.} in concilio 'Florentino Eugenii IV.' temporibus factum est\apnota{33}{M factam statt factum est.}. Quare si pauci Christiani apud nos multorum sententia\apnota{34a}{M sententias.} salvari debeant\apnota{34b}{M debeat, U debent.}, quanto magis paucissimi aut nulli in Turcia salvantur, nisi religiosi aliqui et eis adhaerentes, qui Romanam tenent 'obedientiam'.

De octavo, rogat reverendus pater, ut auctores allegem secundum conscientiam, in quo nihil haesitet\apnota{37}{M haesito.}. Vellem, quod eos adducere possem, etiam [270] secundum scientiam. At\apnota{1}{M ad.} Hieronymum agnoscere Damasum velut summum pontificem\apnota{2a}{U pontifice, Uv pontificem.}, nulli potest esse dubium\apnota{2b}{P dubimu (so!)}. Sic quoque ecclesiam universalem\apnota{2c}{fehlt in P.} esse fundatam supra petram nemo ambigit\apnota{3}{U ambigat, Uv ambigit.}. Sed petram illam esse Petrum et eius successorem\apnota{4a}{U successores.}, alio tempore probabo\apnota{4b}{M probamus.}.

Sugillat tamen nonnihil Hieronymum\apnota{5}{M Hieronymo.}, cum dixisset: Apud vos solos incorrupta patrum\apnota{6}{U prima, Uv primum statt patrum.} omnino servatur auctoritas; quasi et Romani pontifices non fuerint sine naevo. Si ad tempus Hieronymi refert, Liberius enim\apnota{7}{fehlt in M U.} et Anastasius eum praecesserunt. Hoc dico, quod merito in mentibus fidelium est admirandum, nullum Romanum pontificem, quantumcunque\apnota{9}{U quantum cuique, Uv quantumcumque.} malum vel errantem, definivisse\apnota{10a}{M diffinivisse.} iudicialiter et ex constituto, quod unquam meminerim, quod esset\apnota{10b}{U est.} contra fidei Christianae mandata, quamvis facto ipso saepe erraverint, imo talia iudicia erronea facere attentarunt\apnota{12a}{M attemptaverant.}, ubi\apnota{12b}{W. A. setzt ubi vor talia.} divino quoque correpti sunt iudicio, ut de Leone Ariano adversante Hilario, ut de Anastasio can. Anastasius dist. 19 forte vel 20\apnota{14}{U 20. dist., Uv add forte vel 19.; can. Anastasius etc. fehlt in M.}.

Ad nonum, de canone Legimus 93. distin. forte reverendus pater me non intellexit. Nunquam venit mihi in mentem, papatum esse ordinem ultra episcopatum, sed dignitatem. Quod vero dicit, mihi bellum esse cum Hieronymo et canonibus, dixi, quid sentiam\apnota{18}{M sentiunt.}, praefero hoc loco Dionysii auctoritatem velut senioris. At cum rhetoricetur reverendus pater de Dionysio, cur non monarcham\apnota{20a}{M monarchem.} ecclesiae descripserit\apnota{20b}{P U descripsit.}, imo ultra episcopatum non progrediatur, facile est quod dicam. Dionysius enim sacramenta ecclesiae perstringit\apnota{21}{P M perstringens.}, papatum autem cum ordinem esse negaverim, episcopatus\apnota{22}{U episcopatum.} inter ordines omnium consensu supremum obtinebit locum.

Decimo, quod neminem tam insanum existimat, qui crediderit aliquem esse episcopum proprium singulis ecclesiis\apnota{25a}{M per ecclesiis, Mc per ecclesias.}: ecquid\apnota{25b}{M U quid.}, est quod ego conqueror, nisi quod stultorum et qui singularitati student infinitus est numerus? Legat dominus\apnota{27a}{U reverendus.} pater Alvarum de planctu ecclesiae, Iohannem de Turre cremata\apnota{27b}{M clamato.} in summa ecclesiae, Vuilhelmum Occam\apnota{28}{U de Ocken.} in dialogo, et inveniet aliquando homines\apnota{29}{U homines quandoque.} illius fuisse insaniae.

Quod vero suspendere vult solutionem nostram de eo, quod dixi, papam universalis ecclesiae pontificem melius, quam universalem episcopum, feci hoc divi Bernhardi auctoritate et summorum pontificum consuetudine. Ait enim Bernhardus lib. 2. de considera. ad Euge.\apnota{33a}{M Eusegenium, Mc Eugegenium (so!)} col. 7: Nempe signum singularis pontificii Petri. Et infra\apnota{33b}{M U etc. statt et; infra fehlt in M.} ita: Cum quisque ceterorum suam ecclesiam habet\apnota{34}{U habet ecclesiam, M habet (ecclesiam fehlt).}, tibi una commissa est, grandissima navis facta ex omnibus, ipsa universalis ecclesia toto orbe diffusa.

Ad undecimum\apnota{36}{U Ad decimum.}, verbo Christi non satisfactum Lucae. 22., quasi sit commentum meum. At ne sine auctoritate loquar\apnota{37}{U non sine auctore loquor.}, adduco ei Richardum Armacanum, lib. 7. cap. 3. de quaestionib. Armeniorum, locum istum ita intelligentem [271] b. Leonis auctoritate. Et quod hoc sit verum, verba indicant: Et qui maior est inter vos etc. Ergo praesupponebat Christus aliquem esse maiorem. Quis autem esset maior, non tunc indicavit, sed posterius, cum Petro diceret de daemone eos expetente\apnota{4}{M expotente.}, et quomodo rogaret pro eo, ut conversus confirmaret fratres suos, ubi maioritatem eius declaravit.

Crastina hora septima antemeridiana continuabitur disputatio iuxta edictum publicum semper fieri solitum.

Die Martis quinta mensis Julii hora septima antemeridiana.

D. D. Martinus orsus est ut sequitur\apnota{7}{Die Zeitangabe fehlt in M; U Martinus Lutherus mane hora 7. die Martis, quae erat 5. Julii.}.

Confutaturus egregius D. D. responsionem\apnota{8a}{U responsionem egre. d. d. (meam fehlt, in Uv nachgetragen), M ad responsionem meam.} meam\apnota{8b}{U add in causa, in Uv gestrichen.}, ad illud\apnota{8c}{fehlt in M.} 1. Corinth. 3: Quid est Paulus? quid est Cephe etc.\apnota{9}{fehlt in U.}? dixit ibi non esse intentam principalitatem, sed singularitatem personae, idque indicari per verba schismatis; respondeo: Me plus movet ipse textus Pauli, quam

tam violenta et extorta distinctio, quae nulli prorsus innititur\apnota{12}{M U nititur.} auctoritati, cum tamen velit e divino iure arguere. Clarum est autem contentionem eorum fuisse de praestantia\apnota{14}{U add vel principalitate, in Uv gestrichen.} personarum, quod indicat ipsa confutatio Pauli, per tapinoses et comparationes urgens et dicens: Quid est Cephe? quid Paulus? Ministri, per quos credidistis. Itaque qui plantat, qui\apnota{16}{U et qui.} rigat, unum sunt. Neque qui plantat\apnota{17a}{U add est aliquid, in Uv gestrichen.}, neque qui rigat, est aliquid\apnota{17b}{M aliquid est.}, sed qui incrementum dat, deus. In quo manifeste detrahit personam, id est, occasionem schismatis. Quomodo et ad Gal. 2 [v. 6] cum titulo praestantiae Petri et aliorum apostolorum Galatae essent seducti, audet et dicit: Qui videbantur, quales aliquando fuerint\apnota{21a}{M fuerit.}, id est, quam magni\apnota{21b}{id est quam magni in M a. R.}, nihil mea refert, deus enim personam hominis non accipit. Quasi diceret: Sive Petrus\apnota{22a}{M add sive Paulus.} sive\apnota{22b}{U vel.} quicunque apostolus primus vel ultimus sit, 'nihil ad rem'. Vellem ergo egregius D. D. scripturas rectius legeret, antequam\apnota{24}{M quam.} 'sic' allegaret. Stat ergo, quod primatus iste seu persona nihil pertineat ad ecclesiam de iure divino saltem.

Transeo illud, quod dicit Hieronymum\apnota{26a}{U Hieronymum dicit, M dicit Hieronymus (so!)} Pauli\apnota{26b}{fehlt in M.} sententiam intellexisse\apnota{26c}{M concessisse.}. Item\apnota{27}{Der Satz Item et illud etc. fehlt in M} et illud, quod concessit Iohan. 5. Item et illud de plebeia ratiuncula transeo, quia non est confutatum\apnota{28a}{M U confutata, Uv. a. R. confutatum habent alia.}. Quod enim\apnota{28b}{fehlt in U.} de vicario et episcopo\apnota{28c}{M U episcopo et vicario.} in consistorio adduxit, nihil ad rem est. Mea ratio hoc voluit: Si ecclesia non est acephala\apnota{29}{M a cephola.} mortuo papa, nec acephala erit\apnota{30}{fehlt in U, in Uv nachgetragen.} nullo papa. Transeo illud de cardinalibus, quia omnibus notum, quando coeperint\apnota{31a}{quando coeperint fehlt in M.}. Transeo et\apnota{31b}{M transeat statt transeo et.} illud de Cypriano, quia ex propria saliva D. D.\apnota{32a}{doctor fehlt in M.} dixit nescio\apnota{32b}{M nescit.} quae de Numidia ultra Atlantem, quae non est in rerum natura. Et illud, quod Cornelius frater 'appellatur' a Cypriano\apnota{33}{U Damaso, Uv wie oben, M a. R. de Cypriano.}, transeo\apnota{34a}{fehlt in P.}, quia non confutatum est, licet dictum sit, quod honestissimis\apnota{34b}{M dignissimis, Mc honestissimis.} titulis se honoraverint, ut 'beatissimos', sanctissimos\apnota{35a}{M add etc.} appellando, nam et hoc ego\apnota{35b}{M ego hoc.} dico. Non autem soli Rom. pontifici, sicut hodie, sic scribebatur. Hoc enim probandum erat.

Transeo et illud, quod dixerit\apnota{37}{M quod impertinens esse dixerat.} impertinens esse ad institutum, an Petrus apostolus ordinarit\apnota{38}{U an Petrus apostolum aliquem ordinaret.}, quia hic unicus nervus [272] est et invictus contra D. Eccium, ideo non inconsulte transit\apnota{1}{Mc transeo.}, ne irrecuperabiliter

impingat. Item admitto ea\apnota{2}{fehlt in U, in Uv nachgetragen.}, quae ex martyre Cypriano adduxit de honore et potestate aequali apostolorum, et accipio tanquam gratiam mihi factam\apnota{3}{M factam mihi.} a D. D. Simili cura, quod transiit\apnota{4a}{M add electionem.} a Niceno concilio et Cypriano\apnota{4b}{M U a Cypriano.} descriptam electionem\apnota{5}{U ecclesiam.} episcoporum, placet, quia non potest confutari. Quod autem repetiit 'de distinctione' ortus sacerdotalis, videlicet\apnota{6}{U unde patet.}, quod intentio Cypriani sit de defluxu\apnota{7}{M de fluxu, Mc de defluxu.} potestatis sacerdotalis, quia nihil dixit nisi sua verba, apud me nihil concludit\apnota{8a}{M quae apud me nihil concludunt.}.

Illud\apnota{8b}{U illud quoque.} magis attendendum est\apnota{8c}{fehlt in U.}, quod Graecos haereticissimos ausus est appellare satis modeste, cum in universa ecclesia nulla pars dederit\apnota{10a}{P dedit.} plures\apnota{10b}{fehlt in M.} excellentiones scriptores, quam Graeca. Nam quod toties mihi submordens Bohemiam opproperat aliasque\apnota{11}{M alias quam, Mc alias que.} calumnias intertexit, ad sophistas\apnota{12}{M zophistos.} relego, indigniora sunt, quam ut in gravi et theologica disputatione memorentur, nedum obiiciantur\apnota{13}{M obicerentur.}.

Quare illud Matth. 16: Super hanc petram etc. videamus. Ubi dico, aut petra significat hoc loco potestatem vel fidem (confido enim dominum D. nunquam satisfacturum suae promissioni, quod petra significet papam seu successorem Petri); si potestatem, superflue dicitur\apnota{17}{P dicitut (so!)} infra: Tibi dabo claves regni caelorum\apnota{18}{regni caelorum fehlt in M.}, id est, potestatem, nisi dixeris, potestatem potestati traditam. Deinde si omnino potestatem significat\apnota{19}{P significet.}, eadem erit omnium ecclesiarum, quia dicit: Super hanc petram aedificabo ecclesiam meam, non Romanam\apnota{20}{U add tantum.}. Aut significat fidem (quod verum est), iterum eadem est fides omnium ecclesiarum. Ita patet, quod\apnota{22}{M Jam hoc facit, Mc ita patet quod.} hoc solum pronomen »meam« communem facit petram, quidquid significetur per petram.

Ideo frigidissime\apnota{23}{M frigidissimam.} inducunt hanc 'auctoritatem' decreta pontificum pro singularitate principatus, quae tam\apnota{24}{U tamen.} constanter defendit\apnota{25a}{U defendunt.} communitatem petrae\apnota{25b}{M U petri.}, et sic concordat cum apostolo ad Ephesios: Una fides, unum baptisma\apnota{26}{M unum baptisma, una fides.}, unus dominus. Non enim, ut b. Hieronymus ait\apnota{27a}{fehlt in M.} ad Euagrium, alia fides Rom. ecclesiae, alia\apnota{27b}{M et alia.} Britanniae et totius orbis, ita nec alius Christus, nec alia petra. Quodsi eadem fides et idem dominus et idem baptisma omnium ecclesiarum, sequitur, quod etiam eadem omnia reliqua, quae ad fidem et ad Christum et ad baptisma\apnota{30}{U et ad baptisma et ad Christum.} sequuntur.

Ultimo\apnota{31a}{M Ultimo quod erat auctoritas Matthaei.} ex Richardo Armacano\apnota{31b}{P U Amarcano, M a marcano.} interpretatus est verbum Christi de ambitione, non de maioritate, quia inquit, praesupponebat Christus, esse unum inter eos maiorem, licet non indicaret\apnota{33a}{M U indicarit.}, quis\apnota{33b}{M quid.} esset. Respondeo: Quidquid sit de Armacano\apnota{34}{P M Amarcano.}, praevalet clarissimus textus obscuro auctori, quia scribitur, quod contentio facta sit 'inter' eos, quis eorum videretur\apnota{35}{M videtur.} esse maior. Et Christus: Qui vult inter vos esse maior, satis indicat, nullum suppositum fuisse esse\apnota{36}{fehlt in M.} maiorem, et ex consequentia textus patet\apnota{37a}{M patuit.}, quod noluit ullum 'esse maiorem etc.'\apnota{37b}{Mc ceteris statt etc.}. Quia\apnota{38}{U quod.} autem pro confirmatione sua induxit illud Lucae: Rogavi pro te, Petre, et tu conversus confirma etc.\apnota{39}{conversus fehlt in M U, Uv add a. R. confirma; et tu confirma fehlt in M.}, pro me valet, quia fratres Petri Petro confirmandos praecipit, non maioritatem.

[273] D. D. Eccius.

De apostolis quod attulit reverendus pater, ut semel dicam, ex Cypriano audivimus et Hieronymo, quod et s. papa et martyr\apnota{3}{M U martyr et papa.} Anacletus testatur cap.\apnota{4a}{M U canone.} In novo\apnota{4b}{Mc nova.}, dist. 21: Aequales fuisse eos in apostolatu, quod nemo negat\apnota{4c}{M negarvit.}. At aequalitas illa apostolatus non praeiudicat primatui et principatui, nisi impudenter dicatur, s. martyrem Anacletum in eodem decreto sibi contradixisse. Quod vero tantum ponderis et praesidii locat in ordinatione apostolorum, scio, quod Paulus ingenue scripsit: Postquam ascenderit Hierosolymam, nihil se recepisse ab his\apnota{9a}{U iis.}, qui videbantur\apnota{9b}{Mc sibi videbantur.} aliquid esse. Verum si nervus iste tam invincibiliter Eccium ligare debet, utatur eo, dum\apnota{10}{U cum, Uv dum.} copiam opponendi habuerit. Sed quod gratiam sibi factam ait ex Cypriano, apostolos pari consortio praeditos honoris et potestatis; si istam appellat gratiam, solutiones suas interimere, paratus sum saepe ei talem gratiam exhibere. Nam distinxerat reverendus pater, Petrum fuisse priorem praerogativa honoris, non potestatis, Cyprianus in utroque eos aequiparat.

De Numidia et Mauritania\apnota{16a}{M U De Numidia, quam ait esse inter Mauritaniam et; Mc Uv wie oben.}, 'quod attulit', negotio\apnota{16b}{U negotium, Uv negotio.} est impertinens, et diu cosmographicum negotium non agens, non omnium etiam memini. Hoc scio Mauritaniam Tingitanicam\apnota{18}{U Turganicam, M Tinganicam (so!)} et Caesariensem ab Atlante vergere usque ad mare, neque Atlantem poni\apnota{19}{M ponit.} terminum Africae, sed Africam pro magna sui parte dividere\apnota{20}{M dividerit.} ad plures gradus; sufficiat venientes ad Cornelium ex Numidia non fuisse de dioecesi Romana, sed

primam adiisse ecclesiam. Ob id invicta stat adhuc Cypriani illa sententia nondum soluta. De influxu miror, quod affert verba mea fuisse, cum Cyprianum contra Novatianum adduxerim\apnota{24a}{M induxerim, Mc adduxerim.}, ad quem\apnota{24b}{M quam.} reverendus pater nihil respondit.

De Graecis, fateor eos olim fuisse Christianissimos et doctissimos, cum Rom. ecclesiam primam appellabant sedem. Sed superbia elati et invidia infecti ab obedientia Rom. sedis se subtrahentes in pessimos inciderunt errores, et simul fidem cum imperio perdiderunt. Quod vero divinatur, me non exhibiturum pollicita, et argumentum fecit ex Matth. 16, miror venerandum patrem sophistis tam infestum, ut tamen sophistarum cautelis pulchre utatur\apnota{30}{U et . . . utitur statt ut . . . utatur.}. Et cum sit respondens, se ad partem opponentis transfert. Quare nihil iam respondeo, sed quod petit per me probari\apnota{32}{M probare.}, probabo, ne semper tempus deliberandi ad partes suas veniat.

Postremo obscurum doctorem Richardum Armacanum reiicit, illius oblitus, quod dixi Richardum\apnota{35}{M U Richardum dixi.} id facere b. Leonis auctoritate. Porro id ex litera non evincit, quod petit, contentionem inter discipulos factam Christus merito reprehendit. Quare et verba Christi sic sunt accipienda, ut contentionem reprimant, primatum non auferant. Et textum posterius inductum credit [274] reverendus pater facere pro se, quod\apnota{1}{M quo.} Christus apostolos vocaverit fratres et ita eum non fecerit maiorem. At si verba textus ponderaret, uti alias facere solet, iam videret Petrum aliis apostolis superiorem, cum maior sit confirmans confirmato. Hoc de solutionibus dictum\apnota{4}{Hoc de solutionibus dictum fehlt in M.}.

Venio ergo ad principale, quod petit, probaturus primatum ecclesiae Romanae esse de iure divino et constitutione Christi, ita, quod Petrus fuerit monarcha ecclesiae a Christo institutus cum suis successoribus. Pro quo repeto rationem Bernhardi inductam et nondum solutam. Repeto dictum Cypriani pariformiter, et tertio probo per illa verba Christi, Matth. 16: Tu es Petrus, et super hanc petram aedificabo ecclesiam meam. Ubi glossa ordinaria: specialiter illam potestatem Petro concessit, ut ad unitatem nos invitaret, ideo enim eum principem apostolorum constituit, ut ecclesia\apnota{12}{M ecclesiam.} unum principalem Christi haberet vicarium, ad quem diversa membra recurrerent\apnota{13}{M concurrerent.}, si forte inter se dissentirent\apnota{14}{U dissentirent (so!).}, quod, si diversa capita essent, vinculum unitatis rumperetur.

Ita b. Augustinus in epistola contra Donatum\apnota{15}{W. A. Donatistas.} exposuit: Tu es Petrus, et super hanc petram, id est, Petrum, aedificabo ecclesiam meam, et quamvis idem Augustinus alicubi exposuerit super hanc petram, id

est, Christum, ad mentem apostoli: Petra autem erat Christus, tamen in lib. Retracta. expositionis suae primae eum non poenituit, ita et b. Hieronymus libro primo contra Pelagia. mihi col. 5. ait: Quid Platoni et Petro? ut\apnota{20}{fehlt in M.} ille enim princeps philosophorum, ita et hic apostolorum fuit, super quem ecclesia Domini stabili mole fundata est, quae nec impetu fluminis\apnota{22}{U fluminum.} nec ulla tempestate concutitur. Ita b. Ambrosius Sermo. 47. Petrum dicit fuisse petram

Huic quoque sententiae Chrysostomus adstipulatur ab initio eiusdem capitis. Quid igitur Petrus, inquit, omnium apostolorum os, vertex totius consortii. Et infra: Sublimiora ipsum sapere fecit, et ecclesiae futurae pastorem constituit, et post\apnota{27a}{M subdit statt post.}: hunc universo orbi terrarum Christus praeposuit\apnota{27b}{M proposuit.}. Fatetur hoc s. martyr Cyprianus ad Cornelium papam dicens: Petrus, super quem aedificata fuerat\apnota{29a}{M fuerit.} ab eodem domino ecclesia\apnota{29b}{M ecclesiae}, unus pro omnibus loquens et ecclesiae voce respondens, ait: Domine, ad quem ibimus\apnota{30a}{M add tu habes verba vitae, U statt dessen etc.}? Propterea s. Leo 24.\apnota{30b}{M 34.} q. 1. cap.\apnota{31}{fehlt in U.} Cum beatissimus, ingenue fatetur, Petrum apostolum accepisse primatum ecclesiae a Domino. Ita omnes consentiunt sancti Petrum a Christo primatum totius ecclesiae obtinuisse, iuniores praetereo, Bedam, Bernhar. et similes. Loquatur s. papa et martyr Anacletus, qui non a quadringentis annis frigidissima decreta composuit, sed a mille et quadringentis annis\apnota{35}{fehlt in P.} ita intonat: Sacrosancta Romana et apostolica ecclesia non ab apostolis, sed ab ipso Domino et salvatore nostro primatum obtinuit, sicut b. Petro apostolo dixit\apnota{38}{M dixerit.}: Tu es Petrus, et super hanc petram etc. Et infra: Haec autem apostolica sedes caput et cardo, ut praefatum est, a Domino et non ab aliis constituta est. Transcribuntur\apnota{40a}{U transferibuntur.} haec verba can.\apnota{40b}{M in canone.} Sacrosancta 22. dist.

[275] Ita s. Marcellus quartus ante Silvestrum can. Rogamus, 24. q. 1. scribit episcopis Antiochenis: licet prima sedes fuerit apud Antiochiam, posterius Domino iubente Romam translata est. Sic sanctus\apnota{3a}{U beatus.} Julius 3. q. 6. can. Dudum\apnota{3b}{P M U can. b. Petrus.}, a Christo testatur ecclesiam Romanam habere primatum et esse caput aliarum ecclesiarum.

Sic Pelagius papa non a quadringentis annis, sed nongentis et 28 s. patrum secutus sententiam, verba Christi: Et super hanc petram etc.\apnota{6}{fehlt in P.} eodem modo accepit. Hinc inter damnatos et pestiferos errores Johannis Wicleff damnatus est et ille: Non est de necessitate salutis credere Romanam ecclesiam esse supremam inter alias. Sic inter pestilentes Johannis Huss errores ille quoque connummeratur\apnota{11a}{M enumerantur.}: Petrus non est nec

fuit caput ecclesiae\apnota{11b}{U Romanae ecclesiae.} s. catholicae. Et alius: Non est scintilla apparentiae, quod oporteat esse unum caput in 'spiritualibus regens ecclesiam'\apnota{12}{Mc rei ecclesiae statt regens ecclesiam.}, quod semper cum militante\apnota{13}{M conmilitante, U in militante.} ecclesia conversetur.

Et iste\apnota{14a}{U ista.}: papalis dignitas a Caesare inolevit. Et: Papae praefectio\apnota{14b}{P M perfectio.} et institutio a Caesare emanavit. Ita Bonifacius VIII. haeresim condemnans Lugdunensium contra errorem eorum decernit, esse de necessitate salutis, omnem creaturam humanam\apnota{17}{M U humanam creaturam.} subesse Romano pontifici, ut super hoc habetur decretalis: Unam sanctam. Johan. quoque XXII. haeresim Marsilii Paduani\apnota{18}{fehlt in M (Lücke).} condemnans, hunc quoque reprobavit, b. Petrum apostolum non plus fuisse caput ecclesiae, quam fuerit alius apostolus\apnota{20a}{M U quam alios apostolos.}. Ita longa serie incipiendo a primitiva\apnota{20b}{M primativa.} ecclesia semper in confesso fuit apud bonos Christianos ecclesiam Romanam obtinuisse primatum\apnota{22}{M U primatum obtinuisse} a Christo, non iure humano et populi consensu.

Fateor, quod Bohemi in suorum errorum\apnota{23a}{M erroro (so!)} pertinaci defensione illa commemorant et his\apnota{23b}{M illis.} armis virulentis se defendunt, ut est videre in eorum oblatione facta in concilio Basiliensi, et in alia disputatione facta coram rege et optimatibus\apnota{25}{M aptimatibus.} regni inter Ragusium\apnota{26a}{M in terra — statt inter Ragusium.} et illum iniquitatis ministrum Johannem Rockenzcenam\apnota{26b}{P Rockezanum}. Unde veniam precor a venerando patre, si Bohemis (non loquor de Christianis, sed de\apnota{28a}{fehlt in U.} schismaticis) sum infestus, tanquam ecclesiae inimicis, et quod eorum in\apnota{28b}{fehlt in M U.} praesenti memor sum disputatione, quoniam et conclusio ipsa et ea, quae heri allata\apnota{29}{M illata.} sunt, primatum ecclesiae iure humano constitutum\apnota{30}{M constituto.}, meo parvo et exili iudicio, plurimum favent erroribus eorum, et, ut fama est, de hoc plurimum gratulantur. Haec modo in praesentia afferre libet, audituro\apnota{32}{M U auditurus.} reverendi patris et sententiam et allegatorum solutionem.

D. D. Martinus.

Quod\apnota{35}{U Primum diluam contumeliam quod.} me egregius D. doctor insimulat Bohemicae factionis studiosum et plane patronum, parcat ei Dominus, praesertim in tanta\apnota{37a}{fehlt in M.} corona tantorum virorum. Nunquam mihi placuit neque in aeternum\apnota{37b}{in aeternum fehlt in M.} placebit quodcunque schisma\apnota{38}{M scissa.}; inique faciunt Bohemi, quod se autoritate propria separant a nostra unitate, etiamsi ius divinum pro eis staret\apnota{39}{M stat.}, cum supremum [276] ius 'divinum' sit 'caritas' et unitas spiritus\apnota{1}{fehlt in M.}.

Hoc ego quaevisi solum, et rogo quemlibet Christianum\apnota{2}{M U bonum Christianum.}, dignetur Christiana caritate perpendere, an non longe sit\apnota{3}{M U sit longe.} impudentissimae iniquitatis, tot milia martyrum et sanctorum per annos mille et quadringentos in Graeca ecclesia habitos extra ecclesiam eiicere et nunc demum etiam\apnota{5a}{fehlt in M.} regnantes in caelo velle deturbare\apnota{5b}{M reprobare.}? Nam si etiam\apnota{6a}{M etiam si.} insaniant omnes adulatores Romani pontificis\apnota{6b}{M Romani adulatores pontificis.}, negare non possunt, ecclesiam Christi fuisse 20 annos fundatam, coronatam\apnota{7}{M et coronatam.} per multam orbis terrarum\apnota{8a}{M terrarum orbis.} partem, 'antequam ecclesia Romana'\apnota{8b}{M U Romana ecclesia.} fieret ex Petro\apnota{8c}{M a. R. fieret petra.}, ut clarissime patet ex epistola ad Gal. [1. v. 18 ff.], ubi scribit Paulus\apnota{9}{fehlt in M.}, se post tres annos venisse ad Petrum, deinde post quatuordecim\apnota{10}{M a. R. quadraginta.} annos iterum ascendisse ad Petrum. Qui si conferantur, invenientur\apnota{11}{M invenienter.} ferme decem et octo anni post ascensionem Christi, quando Petrus adhuc erat Hierosolymis, ut taceam annos, quibus sedit Antiochiae, ut non possit dici, Romanam ecclesiam esse primam et caput iure divino. 'Sed\apnota{14a}{M U Jam.} illud\apnota{14b}{M illud quod.} magis urget, quod Graeca ecclesia' usque ad nostra\apnota{14c}{M ista.} tempora nunquam accepit episcopos suos confirmatos ex Roma\apnota{15}{U, W. A. Romana.}. Ideo si fuisset ius divinum, per tantum tempus omnes episcopi Alexandriae, Constantinopolis aliquot sanctissimi, ut Gregor. Nazian. et ceteri quam plurimi\apnota{17}{M. plures.} essent damnati, haeretici et Bohemici, qua blasphemia nihil potest detestabilius dici.

Ad argumenta\apnota{20}{U Ad argumenta solutionum.}.

Quando egregius D. D. dicit: Aequalitas apostolatus non praeiudicat primatui\apnota{22}{M primatum.}, et adducit s. mart. Anacletum distin. 21. ca. in Novo test.\apnota{23}{fehlt in M.}, dico breviter: hoc est unum de frigidissimis decretis, quod impugno, nec ullus mihi persuadebit, hoc decretum esse sanctissimi 'pontificis\apnota{24}{M illius sanctissimi pontificis (Lücke), U huius sancti pontificis et martyris.}.

Ad secundum'\apnota{26}{M ad secundum dico.}: Non posui vim in ordinatione apostolorum, de quo in 'opponendo videbimus'.

Tertio, quod Cyprianum induxit aequiparentem apostolos tam in honore quam in potestate, cum ego dixissem\apnota{29}{M dixisset.} Petrum honore fuisse priorem, libentissime admitto, et si opus est, libenter\apnota{30a}{fehlt in M.} errare volo\apnota{30b}{M volo errare.}, modo D. D. teneat\apnota{30c}{U obtineat, in M fehlen die Worte modo d. d. bis principatus.} hoc pro veritate et dicat mihi, ubi postea\apnota{31}{U postea (statt mihi), ubi postea.} maneat principatus. Dico tamen, quod aequalis est honor apostolorum

ad alios, verum inter 'se merito' Petro primum locum dederunt. Unusquisque enim apostolorum in sua sorte, sicut\apnota{33}{U et sicut.} unusquisque episcopus in sua dioecesi aequalem habet honorem.

Quod exprobrat, me nihil respondisse ad Cyprianum, quaesivi epistolam et non inveni. Tamen satis responsum est ex praecedentibus, quod Cyprianus cuilibet ecclesiae tribuit suum caput per multas epistolas. Hi\apnota{37}{M hic, Mc hii.} qui fuerunt ex Numidia non probant ius divinum esse, sed factum ostendunt, sicut in simili Achatius de facto constituit Felicem Romanum pontificem\apnota{39}{M Romanum pontificem felicem.}, 'quod dictum est heri'. Imo Epiphanius episcopus Cypri deposuit maiorem episcopum Constantinopolitanum, scilicet Johan. Chrysostom., ut est in tripar. histor. Non tamen ex facto isto factum est ius, et multa alia exempla.

[277] Satis miror, egregium\apnota{1}{U d. egregium d. (im Druck als 2) gezählt; 1) ohne eigene Definition)} D. D. instituisse probare ius divinum, et usque hodie ne unam syllabam quidem\apnota{2}{M ne quidem unam sillabam, U ne unam quidem syllabam.} scripturae inducit, sed tantum dicta et facta patrum eademque sibi ipsis\apnota{3}{M U sibi ipsi, Uv sibi ipsis.} repugnantia.

Quod Graecos olim confitetur\apnota{4a}{U fatetur.} Christianissimos\apnota{4b}{M fuisse Christianissimos, U Christianissimos fuisse.}, postea recedentes ab Romanae ecclesiae obedientia\apnota{5}{U obedientia Romanae ecclesiae, in M fehlt Romanae ecclesiae (Lücke).} fidem cum imperio perdidisse, idem est, quod heri dixit eadem auctoritate Ecciana, praevaluisse\apnota{6a}{M praefalnuisse.} adversus eos portas inferi\apnota{6b}{M eam adversus inferos.}, sic 'tractans scripturas', ut intelligat per praevalentiam inferorum adversus ecclesiam\apnota{7}{Mc inferiorem ecclesiam; in U fehlt adversus ecclesiam, in Uv nachgetragen.} terrenae vitae et rerum temporalium amissionem\apnota{8a}{M U vel rerum amissionem.}. 'Bellissima glossa certe', quasi non potuerit fides permanere\apnota{8b}{M U manere.} amisso imperio et ita liceat divinare, quod nulli in 'Graecia sint Christiani\apnota{9}{M U sint Christiani in Graecia.}, quia nullum imperium'. Eadem ratione 'diceret\apnota{10}{U dicet.}, et martyres ab inferis esse superatos'.

Exprobrat etiam\apnota{12a}{U etiam dicens.}, quo dex respondente factus sim\apnota{12b}{U sum, Uv sim.} opponens\apnota{12c}{M quod fecerim me ex respondente opponentem.}, gratus accipio monitionem, nec ideo feci, quod\apnota{13}{U at hoc ideo fecerim ut.} tempus haberem 'deliberandi'. Non sunt tanta\apnota{14a}{M non tanta sunt gestrichen, a. R. non sunt tanta.} Ecciana fulmina ('ut et ego quid glorier'\apnota{14b}{Mc 'ut ego mihi quid glorier'.}), ut necessaria mihi sit\apnota{15a}{U sit mihi, in M fehlt mihi.} 'deliberatio'.

Illud etiam, quod Richardus Armacanus\apnota{15b}{P M U Amareanus.} non sua, sed Leonis auctoritate verbum Christi tractavit\apnota{16}{M U tractaverit.}, parum movet. Ostendat ipse ex textu ipso\apnota{17}{M U ipso textu.} (hoc est, iure divino) propositum, et ero contentus.

Ad ultimum, ubi sic ratiocinatur: Confirmans est maior confirmato,

ergo Petrus est\apnota{19a}{fehlt in M U, in Uv nachgetragen.} maior apostolis\apnota{19b}{U apostolus, Uv apostolis.}, forte quod ex Aristotele putet\apnota{19c}{M putes.} agens esse praestantius suo passo\apnota{20a}{M praestantius esse suo passo, U suo passo esse praestantius.}. Sed argumentabor\apnota{20b}{P U argutabor, M argumentabar.} et ego plane Aristotelice: Confirmans est bene maior per se, sed per accidens bene\apnota{21a}{fehlt in M (Lücke).} minor, nisi fortassis\apnota{21b}{M confortasset.} confirmationem hic intelligat sacramentum confirmationis, quod non credo, alioquin non rarum est superiorem per\apnota{23}{M vel. statt per.} inferiorem adhortari, consolari, confirmari. Hoc ad confutationes eius\apnota{24}{hoc ad confutationes eius fehlt in M.}.

Ad principale\apnota{25}{Ad principale fehlt in P Uv.}.

Ad principale accedens probavit, Petrum esse monarcham iure divino, hoc est\apnota{27}{U et hoc ex, Uv hoc est.} auctoritate Bernhardi superius inducta, Cypria., Hier., Ambro., Chryso., Augustini, deinde omnium sanctorum consensu, item multis decretis et decretalibus diversorum pontificum. Respondeo: D. D. Eccius voluit iure divino probare, et mox sui oblitus incidit in auctoritates patrum, quas pro maiore parte iam tractavimus, et vidimus eos in diversis locis diversa aliquando sensisse, et multo\apnota{32}{M multum.} plures et saepius pro me quam pro D. D. Eccio.

Secundo videamus tamen per ordinem.

Prima Bernhardi superius satis est dissoluta, quod videlicet non probat, sed suadet tantummodo, eo quod nitatur in alieno sensu scripturae, ut ipsemet D. doctor concessit\apnota{36}{doctor. fehlt in M; U concessit doctor.}. Similiter et\apnota{37a}{fehlt in M.} Cypriani auctoritas de principalitate et de ortu sacerdotalis unitatis\apnota{37b}{M dignitatis.} satis visa est.

Tertio probavit per illa verba: Tu es Petrus, et super hanc petram aedificabo\apnota{39}{fehlt in M U.} etc., quae Augustinus sic exposuerit: Super hanc petram, id est Petrum, et hanc non retractaverit. Respondeo: Quid ad me? Si vult pugnare contra me [278], conciliet ipse primum dicta contraria. Certum est enim, Augustinum saepius exposuisse petram Christum et fortassis vix semel Petrum, ideo plus pro me, quam contra me facit.

Quodsi etiam Augustinus et omnes patres Petrum intellexerint\apnota{4a}{U intellexerunt.} per petram, resistam ego eis\apnota{4b}{M U eis ego.} unus auctoritate apostoli Pauli\apnota{4c}{fehlt in M U, in Uv nachgetragen.}, id est divino iure, qui scribit 1. Cor. 3 [v. 11]: Fundamentum aliud nemo ponere potest, praeterquam quod positum est, quod est Jesus Christus, et auctoritate Petri, primae\apnota{7}{U Petri 1. 2.} 2 [v. 4 ff.], ubi Christum lapidem vivum et angularem appellat, docens, ut superaedificemur in domum spiritualem. Alioquin si Petrus esset fundamentum ecclesiae, lapsa fuisset ecclesia ad unius ancillae ostiariae vocem, quam tamen nec portae inferorum expugnare poterunt.

Sequitur ergo, quod sancti patres, quando Petrum appellant petram hoc loco, vel humana patiuntur vel aliquem alium sensum habent, de quo non pronuntio. Ideo illud Ambrosii dicentis, Petrum esse petram, facile admitto, cum et quilibet Christianus sit petra propter Christum, in cuius soliditate firmatur et unum cum eo efficitur.

Quando autem Chrysostomus Petrum appellat\apnota{15}{M vocat.} pastorem futurae ecclesiae et constitutum praepositum universo orbi, sum contentus, modo intelligatur ista pastura non totius ecclesiae, ne Paulum apostolum excommunicemus, qui multo plures ecclesias pavit quam Petrus, et Petrum esse primum in honore in toto orbe confiteor, et hoc etiam Chrysostomus tangit, quando dicit, eum esse verticem totius sacri\apnota{20a}{fehlt in M U, in Uv nachgetragen.} consortii apostolici. Vertex non est 'caput verticis', 'sed est\apnota{20b}{fehlt in U.} pars capitis'. Imo apertius appellat eum apostolorum os, quod et Hieronymus et Cyprianus asserunt, quia non in sua persona solum, sed omnium apostolorum et totius ecclesiae audivit: Tibi dabo claves etc.\apnota{23}{fehlt in M.}.

Ad decreta nihil dico, quae dixi frigidissima, et praesertim istius\apnota{24}{fehlt in M.} Anacleti multum iactati hac hora, quod bonus Christianus non credat\apnota{25a}{Mc tenet credat.} esse Anacleti\apnota{25b}{U Anacleti esse.} martyris, qui Cephas interpretatur caput et 'Romanam ecclesiam vocat cardinem'.

In fine, quandoquidem adeo displicent Bohemi D. D. egregio, ostendat memoriam et ingenium suum, scribat contra eos. Satis ego admiror\apnota{28}{U miror.}, tam multos inveniri Bohemorum criminatores et hostes, nullum tamen esse, qui fraterna caritate dignetur eorum errorem confutare in gloriam Romanae ecclesiae.

[D. Martinus petiit Eccium, ne velit impingere tantam contumeliam, ut eum Bohemum faceret, quia sibi semper invisi fuissent, ideo quod ab unitate dissentiant]\apnota{34}{Die eingeklammerten Worte fehlen in P M, sind in Uv als „Additio" bezeichnet. Da sie nur den Inhalt des Gesagten zusammenfassen, sind sie wohl irrtümlich in den Text von U aufgenommen.}.

[279] Hora secunda continuabitur disputatio post meridiem iuxta publicum edictum famuli universitatis.

Eadem hora secunda pomeridiana

dominus D. Martinus prosecutus est hodie coepta in hunc modum\apnota{3}{Die Zeitangabe fehlt in M; U hat Hora secunda continuata est disputatio eadem 5. die Julii Martinus Lutther.}.

Adiecit\apnota{4a}{M U obiecit, Uv = P.} egregius D. D. in fine articulos Vicleff et\apnota{4b}{Vicleff et fehlt in M.} Johannis Huss damnatos et Bonifacium octavum\apnota{5}{fehlt in U.} damnatorem 'eorundem'.

Respondeo, sicut prius, me non velle nec posse tueri Bohemorum\apnota{6}{M me non velle Bohemorum defendere, U me non velle nec posse defendere Bohemorum.} schisma, sed Graecam ecclesiam mille et quadringentorum annorum, sive cum ea senserint Bohemi sive non\apnota{7}{sive non fehlt in U, in Uv nachgetragen.}, nihil ad me, certum habeo, quod nec Romanus pontifex nec omnes eius adulatores possint\apnota{8}{M possunt.} tantum numerum sanctorum sub potestate Romani\apnota{9}{M Romanae.} pontificis nunquam agentium de\apnota{10a}{M non possint de (so!).} caelo deturbare\apnota{10b}{M probari, Mc deturbari.}.

Secundo, et hoc certum est, inter articulos Johannis Huss vel Bohemorum multos esse plane Christianissimos et evangelicos, quos non possit universalis ecclesia damnare, velut est ille\apnota{13}{M illa.} et similis, quod tantum est una ecclesia universalis. Haec enim agentibus impiissimis adulatoribus inique est damnata, cum oret universitas ecclesiae: Credo in spiritum sanctum, sanctam ecclesiam catholicam, sanctorum communionem. Hunc nobilissimum articulum fidei inter articulos Johannis\apnota{17a}{fehlt in M.} Huss numerant.

Deinde ille\apnota{17b}{M illam.}: Non est de necessitate salutis, credere Rom. ecclesiam 'aliis esse\apnota{18}{U esse aliis.} superiorem', sive sit Vicleff, sive Huss, nihil\apnota{19}{U non.} curo. Scio, quod salvati sunt Basilius magnus, Gregorius Nazianzenus\apnota{20a}{M U Greg. Naz., Basil. Magn.}, Epiphanius Cyprius\apnota{20b}{Mc Cypriae.} et innumerabiles alii episcopi Graeci\apnota{20c}{M episcopi, U Graeciae episcopi.}, et tamen hunc articulum non tenuerunt. Nec est 'in potestate Romani pontificis' aut inquisitoris\apnota{22}{M inquisitorem, U inquisitorum.} haereticae pravitatis, 'novos condere articulos fidei', sed secundum conditos iudicare. Nec potest fidelis Christianus\apnota{23}{M Christianus fidelis.} cogi ultra sacram scripturam, quae est proprie ius divinum, nisi accesserit nova et probata revelatio. Imo ex iure divino prohibemur credere, nisi quod sit probatum, vel per scripturam divinam vel per manifestam revelationem, ut Gerson etiam\apnota{26}{M enim.} 'etsi recentior' in multis locis asserit, et divus Augustinus antiquior pro singulari canone observat, dicens ad divum Hieronymum: Ego solis eis libris didici hunc honorem deferre, qui canonici appellantur, ceteros autem ita lego, ut, quantalibet\apnota{29}{U quantolibet, Uv quantalibet.} doctrina sanctitateque praepolleant, non ideo verum existimem, quia illi sic senserunt, sed si ex libris canonicis vel probabili ratione mihi persuadere potuerunt.

Quin etiam ipsi iuristae, de quibus minus videretur, in can. Significasti, de electione et elect. pote.\apnota{33a}{et elect. pote. fehlt in M U.} statuunt\apnota{33b}{U statuerint, W. A. statuerunt.}, praevalere unius privati hominis sententiam\apnota{33c}{M sententiae.} tam pontifici Romano quam concilio et ecclesiae, si meliori auctoritate nixus fuerit vel ratione. Ideo nihil est,

quod egregius dominus doctor, volens ex iure divino contra me arguere, dimisso iure divino arguit contra me ex collectaneis haereticae pravitatis inquisitorum\apnota{37a}{M inquisitoribus.}.

Proinde et\apnota{37b}{fehlt in U.} ista propositio Johannis Huss\apnota{37c}{M add vel Wickleff.}: 'Papalis dignitas' a Caesare inolevit\apnota{38}{a Caesare inolevit fehlt in M.}, si est falsa, eradatur Platyna in vita Benedicti [280] secundi, ubi scribit Constantinum IV. Graecorum imperatorem\apnota{1}{M imperatorum (so!) Graecorum, U imperatorem Graecorum.} sanxisse, Romanum pontificem\apnota{2}{U pontificem Romanum.} esse vicarium Christi generalem, quanquam nec sic sit observatum\apnota{3a}{M servatum.} a Graecis\apnota{3b}{U Graeciae.} episcopis.

Quare, quantum me urget egregius D. doctor per Bohemos nondum centum annorum, tantum ego urgeo eum per orientalem ecclesiam, meliorem partem universalis ecclesiae et mille quadringentorum annorum. Si illi\apnota{6}{M ille.} sunt haeretici, quia Romanum pontificem non agnoverunt. haereticum accusabo adversarium, qui tot sanctos per universalem ecclesiam celebratos audet asserere damnatos. Per eadem dico ad Bonifacium VIII., qui qualis pontifex fuerit, et qua fide eius gesta recipienda, satis probant historiae.

Proinde concludo et rogo D. doctorem velit Romanos pontifices concedere fuisse homines et non constituere deos, praesertim quoties iudicaverunt in causa propria: deinde non per se ipsos, sed per indoctorissimos adulatores, quando divus Gregorius multis epistolis, etsi Romanus pontifex, reiecit\apnota{13}{M reiicit.} a se primatum totius orbis, allegans ad hoc praedecessorem suum\apnota{14}{praedecessorem suum fehlt in M, die Lücke von Mc ausgefüllt mit dicens (falsch: doppelt).} Pelagium, dicens inter cetera, quod veneranda synodus Chalcedonensis obtulit hunc primatus\apnota{15a}{fehlt in M.} honorem Romano\apnota{15b}{M Romani.} pontifici, et nullus tamen\apnota{16}{U tam.} ausus est acceptare. Si ergo ego erro, errat mecum Gregorius primus cum suis praedecessoribus\apnota{17}{M praecessoribus.} et damnabiliter peccaverunt, quod oblatum primatum non assumpserunt.

Per haec volo probatum, quod ex decretis, damnationibus, approbationibus recentioribus\apnota{19}{M recensoribus.} ecclesiae Romanae nihil contra me agitur, cum sint suspectissima omnia et antiquae veritati et consuetudini per omnia contraria, nihilominus tamen pro reverentia et vitando schismate libentissime tolero et toleranda persuadeo, modo non tanquam\apnota{22}{U tantum.} iure divino tot sanctos praecedentes damnemus. Haec habui, quae dicerem de articulis\apnota{24}{Haec habui etc. fehlt in M.}.

D. D. Eccius.

Quod reverendus pater, honorem suum excusaturus, negat se Bohemorum patronum, si facta verbis responderent, magnificarent

eum\apnota{27}{M magnificarem deum, U magnificarem eum.}, at ultima primis non concordant, cum pestilentissimos 'Hussitarum errores' non Christiane dicat\apnota{28}{M dico, U dicit, Uv dicat.} Christianissimos. At\apnota{29}{M ad.} de 'his' posterius.

Conditionem tamen illam odio\apnota{30}{M odeo.}, quod schismatici Bohemi et Picardi tanquam datum\apnota{31a}{U deum, Uv datum.} acceptare possent\apnota{31b}{M accipere possunt.}: »si pro eis staret ius divinum.« Graecos laudatissimos et sanctos martyres semper laudavi. At revend. pater, artis coquinariae minus instructus, commiscet sanctos\apnota{33}{M sanctis, Mc sanctos.} Graecos cum schismaticis et haereticis, ut fuco sanctitatis patrum haereticorum tueatur 'perfidiam'\apnota{34}{Mc perfidiam tueatur.}; quoniam mille et quadringentos\apnota{35}{M quagrigenta annos (so!).} simul inculcat annos,

D. D. Martinus\apnota{36a}{Die Zwischenbemerkung in M a. R.; U Überschrift Carolostadius, gestrichen und Martinus überschrieben; U D. Martinus protestabatur dicens.}. Protestor (inquit)\apnota{36b}{fehlt in M U.} coram vobis omnibus\apnota{37a}{[M U add et publice.}, quod egregius D. D. haec\apnota{37b}{U hoc.} impudenter et mendaciter de me loquitur\apnota{37c}{M falso et mendaciter dixerit, U mendaciter et impudenter de me loquitur.}. Et egreg. D. D. Eccius protestatur rursus, quod velit haec probare\apnota{38a}{rursus fehlt in M U; M se probare, U se velle probare.} scriptis et dictis\apnota{38b}{scriptis et dictis fehlt in M; hier ist der Einschub zu Ende, nicht wie W. A. nach U angibt hinter confessores.},

in [281] quibus longo tempore maior Graecorum pars et ecclesia\apnota{1}{et ecclesia fehlt in M.} fuit haeretica et schismatica, sed non est conventio lucis ad Belial et schismaticorum ad sanctos martyres et confessores.

De ecclesia ante Romanam viginti annos existente dicam, cum\apnota{4}{U quam, Uv cum.} reverendus pater mihi opponet\apnota{5}{U componit, Uv opponet.}, non moveor, quod episcopi Graeci a Romano pontifice non fuerint confirmati, nam et plebani seu parochiani\apnota{6}{seu parochiani fehlt in M (Lücke, von Mc ausgefüllt mit plebei).} sacerdotes a papa iam non confirmantur. Sed extremae esset dementiae, dicere ob hoc summum pontificem non habere primatum\apnota{8a}{M potestatem.} super plebanos\apnota{8b}{M plebeiis, Mc. a. R plebanos, U plebeianos.} sacerdotes.

De frigidissimo Anacleti decreto dicam inferius, alia quoque decreta defensurus. De Numidia reverendus pater dixit esse factum non ius, quod est praeter Cypriani sententiam, qui eos incusat velut schismaticos, qui audeant ad cathedram Petri et principalem ecclesiam, unde unitas orta est, accedere. Quod Epiphanius Cyprius Johannem Chrysostomum episcopum superiorem ab episcopatu expulerit, viderit ipse, hoc, inquam, factum non ius\apnota{14}{U iustum.} fuisse\apnota{15}{M esse.} censeo.

Praeterea\apnota{16}{M praetereo.}, quod in me cavillos torquet, quasi non idonee sacras literas tractantem, qui ob amissionem\apnota{17}{M omissionem.} rerum temporalium et divitiarum dixerim\apnota{18a}{M dixerim Graecos (so!).} adversus Graecos portas inferi praevaluisse, nunquam cogitatum mihi\apnota{18b}{M mihi cogitatum.} imponit, nam haereses, schismata, errores contra Graecorum ecclesiam invaluisse dixi. Verum est, quod et imperii amissio cum magno Christianitatis pudore subsecuta est.

Dein\apnota{22}{M U deinde.} quod gloriatur reverendus pater, se iure divino loqui Luc. 22; ego obscurum virum afferam Richardum Armacanum, etiam Leonis auctoritate fretum, dico\apnota{24a}{U dicat; Uv = P.} et me eandem\apnota{24b}{U eundem, Uv = P.} credere evangelii scripturam et ius divinum. At reverendus pater suo nixus intellectu me respuit antiquorum sequentem\apnota{25}{M me sequentem (so!).} intelligentiam. Indubitatum est et Arium haereticum et Athanasium habuisse evangelium (Pater maior me est)\apnota{26}{in U steht der Zwischensatz Pater maior me est hinter Arium, in M fehlt me; Uv verbessert nach P.}, verum Arius erronee, Athanasius, ut spiritus sanctus efflagitabat, intellexit. De adiecta confirmatione\apnota{28}{M de obiecta confutatione.} non refert, cum quisque prudens intelligat eum, qui auctoritate superioris sibi concessa alios confirmat, praestare in hoc confirmatis\apnota{30}{U confirmatos.}. De per se et per accidens nescio, quo pertineant\apnota{31a}{U a. R. nisi forte ad currum Carolostadii.}. Hoc pro primo\apnota{31b}{Hoc pro primo fehlt in M.}.

Assumpturus reverendus pater argumenta mea contra principalem intentionem exprobrat mihi immodice, institutum meum fuisse, probare iure divino ecclesiam Romanam aliis 'esse'\apnota{34}{M U esse aliis.} praelatam, et tamen solum in dicta patrum et sanctorum inciderim, quasi amphoram facturus urceolum accipiam\apnota{35}{M urceolam accipiam (W. A. incipiam).}. Parcat mihi reverendus pater, si intentionem meam noluerit vel non potuerit assequi. Nam hoc fere\apnota{37}{alle Quellen fore, W. A. fere.} de iure divino satis fit, dum tot s. patres de iure divino esse dixerunt, licet ius divinum non obticuerim\apnota{38}{U obticuerimus.} Matth. 16 [v. 18]: Supra quam\apnota{39a}{U hanc.} petram, auctoritates\apnota{39b}{U ubi auctoritates.} s. patrum induxi.

[282] De Bernhardo, cuius auctoritas et simul ratio est invincibilis, nescio quomodo reverendus pater sui oblitus dixerit, me admisisse Bernhardum alio sensu locutum, quod nullis armis posset efficere\apnota{3}{M possit efficere, U possit devincere, Uv = P.}; servato enim vero et genuino scripturae sensu, aequalitatis filii ad patrem, optime procedit s. Bernhardi ratio.

De Augustino et aliis, qui Petrum dixit esse petram, tanquam contraria dicentem insinuat se non recipere. Contra quem ego dico: Quomodo audeat tam sanctum, tam doctum patrem credere sensisse

contraria in eodem libro, in eodem capite, in utriusque sententiae collatione, lib. 1. retract. cap. 21.? Sed quam modeste et humiliter reverendus pater Augustinianus responderit, aliorum sit iudicium, cum unus se promiserit tot sanctis patribus se oppositurum. Hoc est verum\apnota{12}{W. A. verbessert virus.} Bohemicum, plus velle intelligere sacram scripturam, quam summi pontifices\apnota{13}{U pontificis.}, concilia, doctores et universitates in magno vigore existentes, cum tamen spiritus sanctus ecclesiam suam non deseruerit. Et mirum esset, si illam veritatem deus tot sanctis et martyribus occultasset usque ad adventum reverendi patris. Neque evincit\apnota{16}{M evinci.} quippiam b. Pauli inductio, quod eius dicta s. patribus et doctoribus minime adversentur. Est Christus petra, est fundamentum, est lapis angularis, est caput ecclesiae indubitatum\apnota{19}{M indubitato.}. Porro illa non debere attribui eius vicario, contra s. patrum et nostram\apnota{20}{M U martyrum.} attestationem, non est acceptandum.

Ratiunculam adiecit etiam nuper in vulgari sermone\apnota{21}{M sermone vulgari.} disseminatam: Si super Petrum\apnota{22a}{U petram.}, quomodo super Petro\apnota{22b}{M U Petrum.} ecclesia aedificari potuerit\apnota{22c}{U potuit, Uv = P.}, qui ad unius ancillulae\apnota{23a}{M aucilluae.} vocem Christum et fidem Christianam\apnota{23b}{fehlt in M.} negavit\apnota{23c}{U abnegaret.}. Parcat mihi reverendus pater, cum philosophiam Aristotelicam contemnat\apnota{24}{M contemnit.} et grammaticam tanti faciat, cur hic lynceis suis oculis verbum »Aedificabo« futuri temporis non respexit\apnota{26}{M U perspexit.}? Non enim sponso praesente, ubi filii laetabantur, opus erat vicario, neque tunc Christus ei potestatem dedit clavium, sed potius promisit. Aute ergo clavium et potestatis dationem ab ostiaria interpellatus fuit s. Petrus\apnota{29a}{M pater.}, quod et Ambrosius et dein\apnota{29b}{M U dehinc.} Gregorius testantur\apnota{29c}{M attestantur.}.

De Chrysostomo\apnota{30}{M U Ad Chrysostomum.} respondit, Petrum superiorem fuisse honore, nam et Paulus plures paverit ecclesias. Quasi reverendus pater non plus praedicaverit\apnota{31}{M iam praedicaverit.}, quam iam summus pontifex, et ob hoc non\apnota{32}{W. A. verbessert ob hoc sit dicendus.} sit dicendus maior pastor quam papa\apnota{33}{M pater.}. Sic nimis subtiliter se torquet in dictione »vertex«, quasi Chrysostomus non primatum voluerit\apnota{34}{U voluit, Uv = P.} intelligere. Se dquis hoc feret, cum symbolicum sit, et vertex pro summitate 'in tali translatione' et metaphora usurpetur? Sic et os apostolorum dicitur, quod saepe pro omnibus apostolis fuerit locutus, quod et b. Chrysostomus ponderat\apnota{37}{M perpenderat.}.

De decretis\apnota{38}{M decreto.} dicit, illa esse frigidissima decreta, de quo plurimum miror, cum in disputatione sua posuerit\apnota{39}{M posuerat.}: Romanam ecclesiam esse

omnibus aliis [283] superiorem, probatur ex frigidissimis Rom. pontificum decretis intra quadringentos annos natis, et ego ei multum\apnota{2}{M multa.} vetustiora adduxerim, antequam incepit refrigescere caritas multorum\apnota{3a}{M sanctorum.}.

Subterfugium suum non accipio, quod\apnota{3b}{M U quo.} negat constitutionem illam esse Anacleti, nam sic omnium conciliorum\apnota{4}{M consiliorum.} et summorum pontificum decreta nisi sigillata\apnota{5}{U mihi sigillatim, Uv = P.} eluderentur, et sic tota facultas iuridica falsitatis insimulatur\apnota{6a}{W. A. verbessert insimularetur.}, quasi falsis fundamentis sit innixa in lectionibus\apnota{6b}{M illectionibus (statt in lectionibus).}, decisionibus, iudiciis et similibus. Et sic iura eorum pontificia erunt de nigra faba, nisi consensu\apnota{8}{M cum consensu.} totius christianitatis fuerint approbata. Unde nullo pacto recipienda est reverendi patris sententia, cum omnium summorum pontificum decreta et conciliorum in pluribus locis reperiantur, et ut est audax hominum genus, dudum fuissent reperti, qui simiam illam Anacleti leonina pelle incedentem diripuissent. Nisi ergo ostenderit in originalibus illa non haberi, fidem ei non adhibeo.

Postremo de Bohemis dixit, certe non sine christianorum doctorum contumelia, plures esse criminatores Bohemorum, ubi sint qui ex caritate et bono zelo contra eos scribant et eos exhortentur; ibi exerceant ingenium\apnota{16}{M et ingenium.} et memoriam. Quid surdo narrarem fabulam? At Christianus negare non debet plurimos optimo zelo motos contra Bohemos scripsisse, quales fuerunt patres in Constantiensi concilio deputati, qualis fuit doctor respondens Bohemis\apnota{20a}{M prohemis.} in concilio Basiliensi\apnota{20b}{U Herbipolensi, Uv = P.}, qualis Ragusius, qualis optimae probitatis et doctrinae Johannes Capistranus, divi Francisci\apnota{21}{M divi s. francisci.} sacerdos, omnibus fidelibus notissimus, Nicolaum Cusanum Germanorum doctissimum\apnota{22}{U Nicolaus Cusanus Germanorum doctissimus.} cum pluribus aliis, ut obiicit mihi reverendus pater, inquisitoribus, praetereo\apnota{23}{W. A. verbessert quos praetereo; P M haben eine andere Konstruktion, s. o.}. Quare non defuerunt Bohemis, qui bona scriberent, sed ipsi defuerunt, qui in haeresi obstinati bona non sequerentur.\apnota{25}{W. A. verbessert bona sequerentur.}

De Graecis sanctis diximus saepius, sed hoc horrendum omnibus Christi fidelibus arbitror\apnota{27}{U esse arbitror.}, quod reverendus pater contra tam sanctum et laudabile Constantiniense concilium tanto consensu totius Christianitatis congregatum non veretur dicere, articulos aliquos Hussiticos et Viclefficos fuisse Christianissimos et evangelicos\apnota{30}{M evangelica.}

Protestatus est hoc loco D. D. Martinus, non esse verum, quod contra Constantiense concilium dixerit. Eccius contra offert se hoc probaturum ex scriptis et dictis\apnota{32}{M. a. R. 'protestor' Martinus (ursprünglich Carolostadius) non est verum quod contra Constantiense concilium dixerim. Eccius offert se probaturum ex dictis et scriptis [Mc rot fügt bei: ‚per notarios'], U a. R. Protestatur Martinus, im Text Non est verum, quod contra Constantiense concilium locutus sim. Eccius contra offert se probaturum ex dictis et scriptis.}.

quos non possit universalis ecclesia damnare; sicut pessime sonat, inique esse damnatum articulum illum\apnota{33}{M istum.} Hussiticum, de necessitate salutis non\apnota{34a}{fehlt in P.} esse Romanam ecclesiam credere superiorem aliis\apnota{34b}{M U de necessitate salutis Rom. ecclesiam esse superiorem aliis ecclesiis [U aliis superiorem], Mc (rot) 'non est' de . . . Rom. 'credere' ecc.}, de quo quidem dicto Bohemi non immerito exsultant\apnota{35}{M gaudent.} et hoc orationibus suis apud\apnota{36}{fehlt in M.} deum postulaverunt, verum magna ecclesiae iactura. Quod (sicut b. Augustinus alibi\apnota{37}{M alibi Augustinus.} ratiocinatur) si ullum mendacium ad sacras scripturas fuerit admissum, 'tota erit' de veritate suspecta. Ita et damnati Hussitae reverendi Patris fulti\apnota{39a}{M vulti, Mc fulti.} patrocinio

D. D. Martinus intercipiens dixit impudens esse mendacium cum protestatione\apnota{39b}{Die Zwischenbemerkung fehlt in M; U a. R. Martinus Luther protestatur, im Text Mendacium est impudentissimum.}

procul dubio [284] dicturi sunt: Si concilium\apnota{1}{M U consilium.} erravit in his duobus articulis Christianissimis, ita eius apud nos in aliis articulis vacillabit\apnota{2}{M vacillabit.} auctoritas.

Unde in re prius condemnata nolo plura verba prodigere, ad quid cogendus sit Christianus aut admittendus, hoc dico auctoritate concilii vel Romani pontificis fieri, ut sententia sine suspicione haeresis defendi non possit, quae alioquin citra fidei laesionem impune posset\apnota{6}{M possit.} defendi. In promptu est exemplum: An essentia in divinis generet. Quam sententiam Richardus amplexus lib. de 'trinitate' minime inculpatur\apnota{8a}{M inculpavit.}. At decisione per concilium\apnota{8b}{M consilium.} facta nullus iam sine haeresis suspicione essentiam diceret generare, forte in\apnota{9}{U credo in (statt forte in) fehlt in M.} capitulo: Damnamus vel firmiter\apnota{10a}{vel firmiter fehlt in M U.} de summa 'trinitate et fide catholica'\apnota{10b}{et fide catholica fehlt in U.}.

Ad Augustinum de legendis\apnota{11a}{P legdis (so!).} canonicis scripturis inductum nihil moror\apnota{11b}{M modo, Mc movet, U moveor, Uv = P.}, quoniam conciliorum et summorum pontificum decreta non excludit\apnota{12}{M excludo.}. Quid domini mei iureconsulti in c. Significasti, de electione ponderent, iam non memini. At suae professionis memores velut iusticiae cultores non dabunt, ut arbitror, leges pontificias pessundari. Quod vero Platynam adducit, quasi ille plus sit, quam summus pontifex vel Augustinus vel Cyprianus, qui in vita Benedicti II. renarret Constantinum IV. sanxisse, Rom. ecclesiae deberi primatum; legat\apnota{18}{M lege.}, quaeso,

reverendus pater, quae domini iureconsulti annotant in canone Ecclesia s. Mariae\apnota{19}{M in canone S. Mariae, U in canone ecclesiae S. Mariae.} de constitut., quid, quantum et quomodo valeant saecularium constituta in ecclesiasticis rebus, et videbit, quo sensu verba Platynae (ne dicam Romani imperatoris) de ecclesiis et haereticis\apnota{22a}{M) ne dicam Romanae ecclesiis [Mc rot ecclesiae] et de haereticis.} debeant\apnota{22b}{U debent.} accipi.

Ultimo loco proponit, me urgere Bohemos, urgere\apnota{23}{fehlt in U.} Graecos, fateor schismaticos. Nam s. patres Graecos canonisatos non damno, sicuti gigantes volebant in fabulis Jovem e caelo pellere. Sed tunc vere\apnota{25a}{M vero.} occurret\apnota{25b}{U occurit.} nobis\apnota{25c}{fehlt in M.} reverendus pater, quando dabit Graecum inobedientem et rebellem ecclesiae Romanae, qui sit canonisatus. Ecquis\apnota{27}{M U et quis.} est hic, et non laudabimus eum? Frustra ergo assumit se tolerare monarchiam, quasi non esset de iure divino, sed quadam tolerantia populi et consensu fidelium inducta.

Quod permaxime liquet ex Gregorio inducto, qui primatum inter fideles sibi oblatum reiecit\apnota{30}{M reiicedit, Mc (rot) reiecit.}. At sensum sequatur reverendus pater, non verborum corticem, et reperiet in epistolis summam optimi patris humilitatem, ut plus mititate et humilitate sua superbiam patriarchae Constantinopolitani\apnota{33}{M Constantinopolitanae.} infringeret, quam quod suspicio elationis aut superbiae de eo haberetur. Unde etiam servum servorum se primus\apnota{34}{M primum.} scripsit\apnota{35}{fehlt in M.}. Et quod hoc sit verum, si tam diligens est reverend. pater Gregorii lector, potuit in eiusdem\apnota{36a}{M eiusmodi.} epistola legere (sicut transponitur\apnota{36b}{M transsumitur.} in decret. canonicis 2. q. 6. Decreto\apnota{37}{P M U de cetero statt decreto.}), ubi potestatis suae memor humillimus pater, divitiarum et gloriae contemptor, scribit: Alios episcopos, praeterquam Romanum, in partem vocatos sollicitudinis, non in plenitudinem potestatis. Nemo est, qui hic dubitare possit, quin reliquis episcopis partem, Rom. vero pontifici plenitudinem [285] ex grammatico sensu dare velit. Ob id rogo, reverend. pater ea, quae dicta sunt vel per alios vel per me, non semper velit reiicere in adulationis studium, quippe adulari non didici, adulari etiam nescio.

Crastina luce hora septima continuabitur disputatio.

Actum ut supra praesentibus ibidem honorabilibus viris Michaele Gabler Vueysenburgense et Nicolao Meyer Leyphanense, laicis Eystettensis et Augustensis dioecesis ad praemissa testibus requisitis.

Johannes Graumann Notarius ad praemissa subsignavit\apnota{4a}{Die Worte von Crastina luce hora an fehlen in M U.}.

Hora septima antemeridiana die Mercurii VI. mensis Julii.

Ante exordium strenuus D. Caesarius Pflugk praefatus est, ut in fine post producta D. D. Martini signabitur\apnota{4b}{Zeitbestimmung und Zwischenbemerkung fehlen in M, letztere auch in U, die Zeitbestimmung lautet in U: Sexta die Julii hora septima mane}.

D. D. Martinus\apnota{5}{U Martinus exorsus est.}.

Postquam heri egregius D. D. non partis, sed iudicis officium usurpans contra conventionis pactum et voluntatem illustrissimi principis Georgii\apnota{7}{M Gerogii.}, patroni nostri, me toties 'definivit et' proclamavit haereticum, cum sui officii fuisset solum adductis rationibus et auctoritatibus iudicibus relinquere, essem haereticus necne; in quo, si non est violata\apnota{10a}{M vigolata, Mc violata.} publica fides, viderint hi\apnota{10b}{U ii.}, quorum interest.

Ad causam primum obiecit, quod errores Huss pestilentissimos vocaverim Christianissimos, de quo testor meam innocentiam. Nec hoc poterit unquam probare, et expostulo, ut assignet eosdem articulos, quos pestilentissimos vocaverim Christianissimos, aut vocem suam revocet.

Secundo criminatur, quod miscuerim\apnota{16}{M U commiscuerim.} sanctos Graecos cum schismaticis. Quid aliud faceret, qui non haberet, quod diceret, cum certissimum sit tempore novissime vastatae Constantinopolis\apnota{18}{M fastatae Constantinopolitanae.} fuisse integerrimos Christianos in Graecia, postea in Italia\apnota{19a}{M alia, Mc Italia.} receptos\apnota{19b}{M receptus.}. Atque si etiam\apnota{19c}{M U etiam si.} hoc non urgeret, adhuc stat, quod\apnota{20a}{adhuc stat fehlt in M, quo (statt quod) M.} usque ad concilium Nicenum\apnota{20b}{M Niceni.} saltem fuerint integri Christiani per totum\apnota{21}{M totam.} orientem, non subiecti Rom. pontifici, ut clarissime testatur decretum eiusdem Niceni concilii\apnota{22a}{M U concilii Niceni.} dicentis lib. 10. ca. 3\apnota{22b}{P can. (Druckfehler für can. iii), can. 3 fehlt in U, in Uv nachgetragen, M can 6; dicentis und historiae ecclesiasticae fehlen in M.} hist. Eccl.: Et ut apud Alexandriam vel in 'Urbe Roma\apnota{23a}{M Romano (so!)} vetusta' consuetudo servetur\apnota{23b}{U conservetur.}, ut ille Aegypti, hic\apnota{24}{P M U hoc, Uv hic.} suburbanarum 'ecclesiarum' sollicitudinem gerat. Idem concilium non Romano, sed Hierosolymitano tribuit honoris\apnota{25a}{U honorum.} primatum, dicens ibidem\apnota{25b}{fehlt in M.}: Et ut\apnota{26a}{M quod, Mc ut.} episcopo Hierosolymorum antiquitus tradita honoris\apnota{26b}{U honorum.} praerogativa servetur.

Quodsi nec ista satis sunt, hoc\apnota{27}{U quod.} tamen adhuc concludit, quod egregius doctor satis callide tacuit, quod ecclesia Christi viginti annos\apnota{28}{M annus, Mc annos,} ad minus fuit, antequam Rom. ecclesia nasceretur\apnota{29a}{M fuerit.}: ideo\apnota{29b}{M igitur.} ridicula

est eius confutatio, quod fingit me miscere sanctos cum schismaticis, cum suos schismaticos assignare non possit.

Tertio, cum confutasset rationem\apnota{31a}{M confutationem.} meam, quia videlicet\apnota{31b}{M U scilicet.} non confirmarentur episcopi Graeciae a Rom. pontifice, dixit hoc nihil esse, cum nec plebani passim confirmentur\apnota{33}{M confirmarentur.} a Rom. pontifice. Quis non videt haec perdendi temporis gratia dicta? cum tamen plebani\apnota{34}{U plebani tamen.} ab episcopis ordinentur. Simul haec omnia remitto ad prius dicta de Hieronym. ad Euag. Eadem causa 'et illud' toties repetit de Numidia, cum illi scribantur venisse ad Romam non vocati et accusentur, quod venerint. Et adhuc D. D. ius 'ex hoc\apnota{37}{fehlt in P.} facto' Rom. pontifici tribuit divinum\apnota{38}{U divinum tribuit, M stellt divinum hinter ius.}. Sed et illud, quod Epiphanius Cyprius Johannem Chrysostomum deposuit, non pro iure, sed pro\apnota{39a}{fehlt in M U.} facto accipit\apnota{39b}{U accepit.}, et sic vult confutasse, [286] cum ille egerit secundum statutum Niceni concilii\apnota{1}{M consilii.} et auctoritatem Cypriani de ordinandis episcopis citra auctoritatem Rom. pontificis.

Stat ergo adhuc invicte, quod Rom. primatus aut non est iure divino firmatus\apnota{4a}{U formatus.}, aut totius orientalis ecclesiae sanctos esse in aeternum\apnota{4b}{M in aeterum esse.} damnatos. Illud praetereo, quod praevalentiam portarum inferi interpretatur invalentiam haeresium, quod non curo, nihilominus non ostendit Graecos fuisse haereticos, cum vulgatum sit etiam per iura, Graecos non haberi pro haereticis. Transeo et illud Leonis\apnota{8a}{M leoni, Mc leonis.} et Armacani\apnota{8b}{P U Amarcani, M domini marcani.}, ubi dicit se sequi intelligentiam antiquorum, accipiendo numerum pro numero, scilicet unius Leonis singularis. Item de confirmatione fratrum Petro commissa nihil dico\apnota{10a}{U dicam, Uv dico.}, satis dictum est\apnota{10b}{fehlt in P M.}, quomodo\apnota{11a}{U cum.} nihil probet, nec adhuc confutatum\apnota{11b}{U confutatum sit.}.

Ad rem principalem dicit, se non tacuisse\apnota{12}{U non tenuisse, W. A. tenuisse.} ius divinum, quia illud Matth. 16 [v. 18]: Tu es Petrus et super etc.\apnota{13a}{et super etc. fehlt in M U.}, cum sententiis patrum sit\apnota{13b}{M sic.} secutus, praesertim Ambrosii et Augustini, qui dixerunt\apnota{14}{M U dixerint.}, Petrum esse petram, ausus etiam hoc addere, quod Augustinus non retractaverit\apnota{15a}{M retractavit.}, ego post\apnota{15b}{M U postea.} consulens librum Retractat. invenio contrarium. Vere enim retractat et dicit, Petrum non esse petram, sed confessum esse petram. Idem hoc dicit in homilia, quam omnes sacerdotes orant in die Petri et\apnota{18}{fehlt in M.} Pauli, ubi dicit: Super petram, non super te, sed super petram, quam confessus es. Idem invenio in Ambrosio, quanquam et ipse variet. Ad hoc est simul aurea illa glossa, ut gloriantur, super capitulum: 'Ita

Dominus'\apnota{21a}{M innovo, a. R. ita dominus.}, dicens: Et super hanc petram\apnota{21b}{U petram etc.} per hanc dictionem non credo Dominum aliud demonstrasse quam haec verba, quae Petrus respondit Domino, cum dixit\apnota{23a}{domino cum dixit fehlt in M.}: Tu es Christus, filius dei vivi, quia super illum articulum\apnota{23b}{M U super illo articulo.} fidei fundata est ecclesia, ergo super se ipsum\apnota{24}{M U ipso.} fundavit Christus ecclesiam.

Nihil est ergo\apnota{25}{M ergo est.}, quod gloriatur egregius D. D. dicta patrum a se stare, cum multo fortius inveniantur a me stare. Poterat ergo parcere linguae suae et auribus nostris, quando rhetorico boatu exclamavit: Me velle\apnota{27}{M vellem.} unum esse doctiorem omnibus, plus velle intelligere scripturam, quam concilia, universitates, doctores, quam Rom. pontifices\apnota{29}{M quam consilium (so!), doctores, pontifices, universitates etc.. U quam doctores, universitates, concilia et Romanum pontificem.}, et mirum esset, si mihi soli tam diu occulta veritas revelata esset. Hoc enim non fuit disputare, sed invidiam movere.

Quod vero ridet meam ratiunculam, ubi dixi, lapsam esse ecclesiam negante Petro, si super eum esset\apnota{32}{M super eam esse.} aedificata, adducens grammatice verbum futuri temporis »aedificabo«, quasi post mortem\apnota{33}{W. A. ergänzt Christi.} aedificata sit ecclesia super Petrum, transeo, quod quisque facile videat, quid valeat. Nihilominus tamen etiam post spiritum sanctum missum gravissimo scandalo fidei cecidit Petrus, ad Gal. 2, cum redargueretur a Paulo, in quo permansit integra fides et confessio, in Petro quidem fides, sed simulatio adversus veritatem evangelii.

Admitto tamen, quod egregius D. D. hac auctoritate praecipue nixus nunc se ipsum excludit ab ea, ut sic nihil iuris divini relinquatur ei. Nam si promissio tantum\apnota{40a}{M tanta.} facta est, eundum\apnota{40b}{M eundem, a. R. eundum.} est ad locum, ubi exhibita promissio est [287], et invenietur illud Iohan. ult. [20 v. 22 f.], ubi non dicit: Petre, accipe spiritum sanctum, sed aequaliter omnibus: Accipite spiritum sanctum, et\apnota{2}{fehlt in M U.} quorum remiseritis etc. Transeo et illud de Chrysostomo, ubi Petrum honore superiorem fuisse\apnota{4}{fehlt in M.} dixi, D. D. confutaturus opposuit: cum Paulus plus praedicavit quam Petrus, plus etiam habebit honoris, quasi ego de multitudine laboris ac non\apnota{6}{M U non potius.} de primitate ordinis dixerim. Transeo de vertice et symbolico suo\apnota{7}{fehlt in M.} capite, quae sunt mera verba.

De decretis, miratur me dicere frigidissima decreta et intra quadringentos annos nata. De quadringentis annis supersedeo, postea dicturus. Nam non fuit opus, ut docerer ante mille et ducentos annos\apnota{11}{fehlt in M U.}, imo quadringentos annos, fuisse decreta, imo fuisse et contentiones de primatu. Miretur interim ipse D. D., quod auctoritatem Matthaei intelligit promissionis verbum esse, non exhibitionis, cum tamen decreta ipsa super

idem verbum tanquam exhibitionis penitus nitantur. Necesse\apnota{15a}{U necessarie.} ergo est, aut decreta improprie tractare, etiam auctore\apnota{15b}{M doctore, Mc auctore.} D. D., aut ipsum errare.

Quod vero non accipit negatum a me decretum Anacleti, et quae de sigillatis decretis 'ibi' ludit, transeo, nondum enim confutavit, quod idem decretum Cepham interpretatur caput, quae inscitia tanto pontifici non est tribuenda, praesertim eo tempore, ubi floruerunt linguae et Judaeorum habebatur copia. Sed et\apnota{21}{fehlt in U.} hoc constat, librum decretorum nondum esse approbatum.

Illud etiam graviter insectatus est, quod dixi, plures esse criminatores Bohemorum quam instructores. Utinam mendacium fuissem locutus. Video quidem multa dici et scribi contra eos, sed satis infeliciter, quod non nominentur fraterno nomine, quo\apnota{25a}{M U quod.} tamen Paulus Galatas\apnota{25b}{M Paulos Galatos.} lapsos in perfidiam Judaeorum non dedignatur. Credo ego Bohemos\apnota{26}{U Bohemus.} esse homines et blandis sermonibus et concessionibus posse attrahi, qui criminationibus et opprobriis haeretici nominis magis indurantur. Non est ergo\apnota{28}{M ergo est.}, quod excusemur non esse narrandam fabulam surdo, cum praecipiente Paulo instandum sit opportune, importune etc.\apnota{29}{U importune, opportune; etc. fehlt in MU.} [2. Ti. 4 v. 2].

De Graecis sanctis saepius\apnota{30a}{M saepe.} diximus. Illud nunc tractandum\apnota{30b}{U autem [statt nunc] tractandum [Uv tractandum].} est, quod ad concitandam invidiam multum elevavit\apnota{31}{M me elevavit, U olenavit, W. A. clamavit.}, horrendum esse omnibus Christi fidelibus, quod contra tam sanctum\apnota{32a}{M sanctam.} et laudabile Constantiense concilium\apnota{32b}{M consilium.} non vererer\apnota{33}{U veretur.} dicere, aliquot articulos Huss fuisse Christianissimos et evangelicos, quos non possit universalis ecclesia damnare. Respondeo: Inter articulos Huss est\apnota{35a}{M ut est.} et ille: Unica\apnota{35b}{U unita, W. A. una} est sancta universalis ecclesia, quae est praedestinatorum universitas. Item alius: Universalis sancta ecclesia tantum est una, sicut tantum unus est numerus praedestinatorum\apnota{37}{M praedestinatorum sanctorum, U omnium praedestinatorum.}. Hi duo non sunt Huss, sed Augustini super Johannem 6\apnota{38a}{„6“. fehlt in [M] U, in Uv nachgetragen.} ad verbum proprie\apnota{38b}{M sed beati Augustini ad verbum proprie [P prope].}, et repetuntur per Magistrum sententiarum in 4.\apnota{39a}{M U magistrum 4. sententiarum.} de sacramento Euchar. Tertius est\apnota{39b}{de sacram. eucharistiae und est fehlen in M.}: Duae naturae, divinitas et humanitas, sunt unus Christus. Hos articulos credo 'confitetur' mecum D. [288] Johannes\apnota{1a}{fehlt in M.} Eccius. Quartus: Divisio immediata humanorum\apnota{1b}{M humanorum immediata, Mc notiert die Umstellung.} operum\apnota{1c}{fehlt in M.} est, quod sunt vel virtuosa vel vitiosa, quia si homo est vitiosus et agit\apnota{2}{M U agat.} quicquam, tunc agit vitiose, et si est virtuosus et agit quicquam, tunc agit virtuose\apnota{3}{U add etc.}.

Iste articulus, ut audivimus praecedenti hebdomada\apnota{4}{P hebdomoda.}, triumphatus est per egregium D. D. Carolostadium, ita ut egregius D. D. Johannes\apnota{5}{fehlt in M.} Eccius coactus sit Scotum et Scotistas, Capreolum et Thomistas cum universis suae factionis doctoribus celebrioribus repudiare et ei consentire.

Proinde quae heri in me evomuit, quod essem patronus Bohemorum\apnota{8}{M haereticorum Bohemorum.}, haereticus pestilentissimus, et si qua similia\apnota{9}{Die Worte von haereticus an fehlen in M, a. R. steht 'haereticus pestilentissimus'.}, refero in sinum eius. Hic sentit cum Johanne Huss, quidquid pro se dixerit, pro me contra eius criminationes\apnota{10}{M criminationem.} dictum est. Illud addo, quod recte dixi Christianissimos et evangelicos, praesertim illum\apnota{12a}{M ille.}: Duae naturae, divinitas et humanitas, sunt unus Christus\apnota{12b}{U etc. statt der Worte von divinitas an.}. Quare debuit mihi admittere 'hoc' D.\apnota{12c}{U hoc admittere egreg. dominus.} D., quod pro reverentia concilii Constant. crederem hos et similes articulos non fuisse ibidem\apnota{14}{U ibi, M ibi non fuisse.} damnatos, sed ab aliquo impostore intersertos. Verum ultra hoc\apnota{15}{M ultra haec, U contra haec.} cum ipsummet concilium dicat, aliquos esse haereticos, aliquos erroneos, aliquos blasphemos, aliquos temerarios, aliquos seditiosos, aliquos piarum aurium offensivos, debuit prius egregius D. D.\apnota{17a}{doctor fehlt in M.} pro sua prudentia\apnota{17b}{M reverentia.} quemlibet in ordinem suum\apnota{18a}{U suum ordinem.} redigere, et non contra determinationem concilii\apnota{18b}{M consilii.} omnes in universum\apnota{19a}{M omnes in universo omnes (so!)} haereticos damnare, quos concilium vix temerarios forsitan\apnota{19b}{fehlt in M U.} esse voluit. Nam potest et ipsa purissima veritas temeritatis, scandali, seditionis, offensionis aurium accusari, sicut Christo contigit\apnota{21a}{M Christum [contigit fehlt].}. Nec ideo est falsus\apnota{21b}{U falsus est.} articulus, multo minus haereticus, quia temerarius vel offensivus. Et sic patet, quod nimis praecipitanter et longe citra modestiam Eccianam haereticus accusatus sum, quem vix fortassis\apnota{24}{M U fortassis vix.} offensivum poterit probare, imo cum ad eum non pertineat prorsus de articulis istis iudicare, qui erronei, qui haeretici, qui temerarii, clare\apnota{26}{M U clarum.} patet, quam inique et temere me haereticum et pessimo nomine accusavit\apnota{27}{U accusat.}.

Proinde non esse de necessitate salutis, Romanam ecclesiam esse superiorem aliis, nondum convincitur esse articulus haereticus, etiamsi inter haereticos numeretur\apnota{29}{M enumeretur.}.

Quod autem ratiocinatur Augustini exemplo: Si ullum mendacium in concilio admittatur, vacillabit\apnota{31}{M facillabit.} universa auctoritas concilii, infelix similitudo est, Augustinus de scripturis divinis ratiocinatur, quae est verbum dei infallibile, concilium vero creatura istius verbi. Ideo iniuria fit verbo dei\apnota{33}{U verbo dei fit.} per hanc comparationem, cum concessum sit,

concilium posse errare, ut\apnota{34}{M quod.} notat Panormitanus cap.\apnota{35a}{M U in cap.} Significasti\apnota{35b}{fehlt in M.}. Transeo illud de essentiae\apnota{36}{M substantiae.} divinae generatione, quia nihil ad propositum.

Ad Augustinum, qui excepta scriptura omnium scripta cum iudicio iubet legere, dicit egregius\apnota{38}{fehlt in M.} D. D., non esse exclusa per eum decreta Rom. pontificis et conciliorum. Hoc dicitur, sed non probatur; quin solutionem [289] meam confirmo auctoritate Pauli ad Thess.: [I 5 v. 21] Omnia probate, quod bonum est tenete. Rom. pontifex et concilia sunt homines, ergo probandi sunt et sic tenendi, nec eximendi ab hac regula apostolica.

Quod per me adductum cap. Significasti dicit se non meminisse, et dominos iuristas monet, ut non pessundent\apnota{5}{M pessumdant, Mc pessumdarent.} leges pontificias, quod et ipsum valet ad concitandam\apnota{6}{M concitandum.} invidiam mihi. Nec per hoc pessundantur leges pontificiae, si eis\apnota{7a}{U eius.} praeferantur leges divinae, quod cum faciant\apnota{7b}{M U faciunt.} domini iuristae, pessime faciunt theologi, qui in hac re etiam resistunt optimis iuristis.

Platynam per me allegatum dicit non esse plus, quam\apnota{9}{P plusquam.} Rom. pontificem, Augustinum et\apnota{10a}{fehlt in U.} Cyprianum, quasi per Cyprianum, Augustinum\apnota{10b}{M U Augustinum et Cyprianum.} aut summum pontificem\apnota{11}{aut summum pontificem fehlt in M.} iam probaverit primatum, qui toties iure divino probaturum se\apnota{12}{M se divino iure probaturum, U iure divino se probaturum.} promisit nec nisi decreta frigidissima 'exhibuit' et aliquot patrum false 'intellectas auctoritates'\apnota{13}{M falsas [Mc false] 'auctoritates intellectas.'}. Platynae nihil tribuo, sed historiae, quae est mater veritatis, quam scribit Platyna.

Quod de constitutionibus dicit, transeo\apnota{15}{de constitutionibus dicit transeo fehlt in M.}.

Tandem fatetur, Graecos sanctos non esse damnatos, sed solum schismaticos. Hoc nihil est ad propositum, neque enim ego\apnota{17}{U ego enim} unquam respondi de schismaticis Graecis, sed de sanctis, praesertim Niceni concilii\apnota{18}{M consilii.}, sicut nec credo ipsum sentire de schismaticis Latinis, quando de Rom. ecclesia loquitur.

Gregorium a me inductum dicit esse sequendum secundum sensum, non\apnota{20}{M et non.} secundum corticem verborum, et sic satis est confutata\apnota{21}{M confortata est, Mc confutata und confundata.} 'mea responsio, sufficit enim solum 'verbum Eccianum'\apnota{22}{M 'Sed' sufficiunt 'verba Eccii', U sufficiunt enim sola verba Ecciana.}.

Quod autem induxit 2. q. 6 cap. Decreto\apnota{23}{2 q. 6 fehlt in P, M 22 (statt 2), P M U de cetero.}, ubi scribit idem Gregorius, alios episcopos praeterquam Romanum in partem vocatos\apnota{24}{M vocatus.} sollicitudinis, non plenitudinem potestatis, debuit D. D. ostendere, quod

omnes episcopos intellexerit per orbem et non solum occidentalis ecclesiae episcopos\apnota{26}{M episcopum.}. Quod si quam maxime faceret, non ideo probatur ex iure divino. Gregorius homo est, eoque minus probatur, quo in multo\apnota{28a}{U multo in.} pluribus locis diversum sentit, ut maius\apnota{28b}{U magis.} sit mihi ius dicendi D. doctori, ut\apnota{29}{M quod.} ipse sensum et non corticem in suis 'adductis' sequatur, aut primum eundem sibi conciliet.

Ultimo\apnota{31a}{fehlt in M, U ultimum.}, quando ego contra egregium D. D. induxi\apnota{31b}{M induxi contra doctorem.} Paulum, 1. ad Corinth. 3 [v. 11]: Quod aliud fundamentum\apnota{32}{U fundamentum aliud.} nemo ponere potest, praeter id quod positum est, quod est Jesus Christus, dixit etiam tribuendum vicario Christi, ergo oportet sic dicere: Nemo aliud fundamentum ponere potest, praeter id quod positum est, quod est Rom. pontifex, et sic resistendum erit evidentissimae auctoritati Pauli, aut revocandum, quod haec non sunt\apnota{36}{U sint.} tribuenda vicario Christi.

Similiter et\apnota{37}{fehlt in M U.} illud ad Gal. 2, [v. 6] ubi dixit Paulus, qui videbantur (loquitur de Petro, Jacobo et Johanne)\apnota{38a}{M Johanne et Jacobo.} quales aliquando fuerint, nihil mea\apnota{38b}{M ad me.} refert, deus enim personam hominis non accipit, ibi Paulus evidentissime dicit, Petri magnitudinem et qualitatem nihil referre, nec eam accipi a deo. Si autem [290] est de iure divino, certissimum est, Paulum hic\apnota{1}{M hic Paulum.} mentiri, nam quod ius divinum est, omnino accipitur a deo, et multum refert. Ideoque ista auctoritas Pauli cogit primitatem ullius\apnota{3a}{M U unius.} hominis episcopi non esse de iure divino\apnota{3b}{M est iuris divini.}, et ideo revocandum, quod D. D. dixit esse\apnota{4a}{fehlt in U.} de iure divino\apnota{4b}{M divino iure.}, cum patenter\apnota{4c}{U apparenter.} contradicat apostolo. Haec\apnota{5}{U hic [zum vorigen Satze gezogen!]} sic dico, quod nisi D. D. egregius suam sententiam mutarit, tanquam evidentissime contrariam Paulo, non sum\apnota{6}{M U sim.} eum habiturus pro catholico. Haec satis sint de mea responsione ad obiecta D. D. Johannis Eccii super propositione tertia decima\apnota{8}{M super propositionem 13.} de potestate et primatu summi pontificis, quae relinquo iudicio iudicum et omnium lectorum.

Hora prima pomeridiana mandato dominorum continuabitur disputatio\apnota{9a}{Die Zeitbestimmung fehlt in M U.}.

Sequitur praefatio strenui domini Caesarii Pflugs, de qua supra fit mentio\apnota{9b}{vgl. oben S. 95. Dieser ganze Zwischenfall [Rede Pflugks] fehlt in M U.}.

Ob criminationes hesterna disputatione inter egregios dominos disputatores, quemadmodum ex factis protestationibus liquet, exortas et

in dominorum praesidentium consilium relatas, die sequenti videlicet Mercurii, sexta mensis Julii, in principio horae septimae ante exordium disputationis inito primum consilio, strenuus dominus Caesarius Pflugk nomine illustrissimi principis et dominorum de consilio, in quos super emergentibus huiusmodi defectibus partes compromiserint, vigore huiusmodi compromissi hortatus est et voluit, dominos disputatores posthac a mutuis criminationibus et famae laesione abstinere; itemque ne sanctam ecclesiam et eius concilia temere attractarent et ut huiusmodi materiae disputationem hodierna die finiant, ne disputatio plus satis prorogata externis sit sumptuosior et nostrates sua studia diutius posthabere cogat; illud adiciens, ne partes cum suis adhaerentibus aliquid contra pacem et salvum, ut vocant, conductum tentarent.

Johannes Graumann Notarius subsignavit.

Hora prima pomeridiana eiusdem diei\apnota{9c}{Unterschrift und Zeitbestimmung fehlen in M U.}.

D. D. Eccius.

Cum reverendus pater mihi improperet, quod iudicis usurpaverim\apnota{11}{U occupaverim, Uv = P.} officium, nihil ad rem, sum disputator. Neque haereticum dixi\apnota{12}{M dicta dixi (so!)}, sed dicta sua haereticis\apnota{13}{M haeretica.} et maxime Bohemis favere et patrocinari, potissimum cum illud horrendum dicere auderet, aliquos articulos Johannis Huss per concilium sacrum Constantiense reprobatos, esse Christianissimos et evangelicos. Ad rem potius descendamus. Duos heri protulit\apnota{16}{M heri protulit heri, Mc streicht das erste heri.} huiusmodi articulos, hodie aliquos addidit, de quibus\apnota{17}{M quo.} infra.

Quod vero ad Graecos attinet, mirum\apnota{18}{M miror, Mc mirum.} quo studio reverendus pater eos defendit, quos fuisse bonos\apnota{19a}{M bonos esse.} asserit etiam in excidio\apnota{19b}{M concilio.} Constantinopolitano', qui dein\apnota{20}{M dehinc, U deinde.} in Italiam venerint, oportere me ostendere et nominare, qui fuerint Graeci schismatici et haeretici. Miror, quod reverendus pater in sua epistola obiicit mihi historiarum ignorantiam, et ipse hoc loco nolit scire, quod tota novit ecclesia. Nonne schismaticus fuit\apnota{23a}{M fuit haeretícus.} Macedonius, Nestorius\apnota{23b}{M Nestorius nonne Macedonius, U Nestorius non Macedonius.}, Eutyches\apnota{23c}{M nonne Eutyches}, Achatius, Johannes Constantinopolitanus, et longo tempore in schismate perseverantes sub Eugenio IV. in concilio Florentino obedientiam fecerunt ecclesiae, quamvis pro pertinaci perfidia statim ad vomitum\apnota{26}{M fomitum.} redierint? Alioquin si Graeci non fuissent schismatici et haeretici\apnota{27}{M haeretici et schismatici.}, non haberemus 'decisionem de summa trinitate et\apnota{28a}{trinitate et fehlt in U} fide' catholica contra Graecos\apnota{28b}{contra Graecos fehlt in M.}. Vel forte

reverendus pater adhuc\apnota{28c}{fehlt in M U.} ignorat 'Thomam edidisse librum' de erroribus Graecorum.

Addo, quia induxit iam saepius ecclesiam orientalem non fuisse subiectam Rom. pontifici, et Epiphanium Cyprium\apnota{31}{fehlt in P M.} repulisse s. Chrysostomum ab episcopatu, quod bene factum dicit reverendus pater auctoritate Niceni concilii\apnota{32}{M consilii.} et Cypriani. Sed cecidit in foveam\apnota{33}{M faveam.}, quam fecit, quoniam s. Chrysostomus odio imperatricis fuit deiectus, et\apnota{34}{fehlt in M.} haereticus Arianus substitutus, quem summus pontifex Julius mandavit recedere, et sanctum Chrysostomum\apnota{35}{M sanctum Johannem Chrys., U Johannem Chrys.} restituit. Quod dum factum non esset, Julius papa Johannem Chrysostomum restitui praecepit et imperatorem excommunicavit. Ita et contra alios episcopos haereticos processit, ut scilicet\apnota{38}{fehlt in M U.} non semel, sed decies Graecos ab obedientia Rom. ecclesiae recessisse [291] historiae tradant. Eat\apnota{1a}{M Ea.} nunc reverendus pater et\apnota{1b}{fehlt in M.} dicat Rom. pontificem de ecclesiis orientalibus non disposuisse, vel iustificet tam s. patris\apnota{2}{U patres.} Chrysostomi repulsam\apnota{3}{M repulsum.}.

Tertio, de Nicena\apnota{4a}{M Niceno.} synodo quam ex Historia scholastica\apnota{4b}{W. A. verbessert ecclesiastica.} lib. 10. adduxit, si ista est constitutio, quam in sua voluit conclusione, bona est, inquam, sed ad propositum frigida. Nam\apnota{6a}{U iam, Uv nam.} synodus nulla\apnota{6b}{nullo.} etiam eo tempore legitima arbitrabatur, quae non foret auctoritate Rom. pontificis congregata. Ita Leo, ita Marcellus, ita Julius, sancti, non frigidi papae sanxerunt, circa idem tempus viventes. Constitutiones eorum leguntur dist. 17.\apnota{9}{P M U dist. XIX.} cap. Synodum, et sequentibus, unde quod patriarchae Hierosolymitano praerogativam honoris dederunt, sed non primatum totius ecclesiae, hoc debebat probare dominus pater. Sexta autem synodus expresse primam sedem dat Romanae ecclesiae, quae transumitur XXII dist. Quod et Leo papa contra imperatorem Michaelem et b. Gregorius contra Johannem Constantinopolitanum et Mauricium imperatorem\apnota{15}{fehlt in M.} obtinuit.

Unde hoc loco id quoque excutiam de b. Gregorio, quem ita inducit, ut primatum reiecerit, quod constat esse falsissimum, cum Platynae tantum tribuat, qui sic de Gregorio sancto scribit: Praeterea vero cum Johannes Constantinop. episcopus, habita Graecorum synodo, se ipsum icumenicum\apnota{19}{M economicum, U icomenicum.}, id est universalem patriarcham creasset monuissetque\apnota{20}{que fehlt in U.} Gregorium Mauritius, ut Johanni obtemperaret, respondit homo constantis ingenii et fidei, potestatem ligandi atque solvendi Petro traditam eiusque successoribus, non episcopis Constantinopolitanis, proinde desineret iram dei in se concitare.

Ex quo constat Gregorium, sine dubio sanctum et humilem patrem, primatum non solum non\apnota{25a}{fehlt in M.} reiecisse, sed etiam asseruisse contra imperatorem Constantinop.\apnota{25b}{M contra Const., U et Const., W. A. et episcopum Const.} Unde ego nolo tam sancto patri imponere, quod reverendus pater, ut sui minime constans contraria\apnota{27a}{W. A. quod contraria.} et repugnantia\apnota{27b}{M U. pugnantia.} dixerit Gregorius, sed universalem se esse episcopum negavit, ad sensum datum in priore disputatione, et tamen primatum obtinuit.

Quod autem semper iste fuerit mos, etiam tempore concilii Niceni, observatus, liquet ex epist. 92. Augustini et coepiscoporum ad Innocentium I., ubi inquiunt: Quia te Dominus gratiae suae praecipuo munere in sede apostolica collocavit talemque nostris temporibus praestitit, magnis periculis infirmorum membrorum Christi pastoralem diligentiam quaesumus adhibere digneris. Respondet papa\apnota{35}{U add in epistola sequenti.}: Diligenter ergo et congrue apostolici honoris consulitis fastigium, honoris, inquam, illius, quem\apnota{36}{M U quoniam.} praeter illa, quae sunt extrinseca, sollicitudo manet omnium ecclesiarum super anxiis rebus, quae sit tenenda sententia, antiquae scilicet regulae formam secuti. Quid\apnota{38a}{M quis.} expressius dici posset\apnota{38b}{M U possit.}, [292] papam tempore Augustini\apnota{1}{M U tempore Augustini papam.} habuisse sollicitudinem de omnibus ecclesiis, et quod tunc non incepit, sed sit forma antiquae regulae?

Dein venit\apnota{3}{M Dein tenet; U Deinde venit.} reverendus pater cum suo Hercule, hoc est ratione illa viginti annorum, ubi ecclesia fuerit ante Romanam. Respondeo: Verum esse Christum in evangelio non expressisse Romanam ecclesiam, sed Petrum principem apostolorum constitutum. Unde quia Petrus, cuius erat principatus, iubente Domino transtulit sedem ab Antiochia ad Romam, quare tunc non iure humano, sed iussu dei facta est prima, iure prius a b. Petro habito, ut\apnota{9}{fehlt in P M.} testatur Marcellus sanctus martyr et papa ca. Roga. 24. q. 1: Licet prima sedes in Antiochia fuerit, posterius iubente Domino Romam translata est. Quis enim ita desipit, ut primatum velit loco affigere. Sed haec est s. patrum sententia, quia iure divino Petrus constitutus sit Christi vicarius, eius\apnota{13a}{M omnes [statt eius], U et omnes.} successores, quales sunt Romani pontifices, ubicunque\apnota{13b}{M U ubi.} locorum fuerint, sunt vicarii Christi. Romanorum\apnota{14a}{U Roma: W. A. Romanus.} rex vel imperator non est ob\apnota{14b}{U vel ob.} hoc Romae. Ita videtur sensisse b. Augustinus contra Manichaeos. Palam est inquit, quod in re dubia ad certitudinem fidei valeat auctoritas ecclesiae catholicae, quae ab ipsis fundatissimis apostolorum sedibus usque ad hodiernum\apnota{17}{M U hodiernam.} diem, succedentium sibimet et episcoporum serie et populi consensu firmatur, quare in epistolis catalogum refert b. Augustinus Romanorum pontificum. Haec de praeambulis\apnota{20}{Haec de praeambulis fehlt in M.}.

Nunc\apnota{21}{fehlt in M} ad principale.

Responsurus reverendus pater inductis per me, dixi enim Petrum iure divino ceteris apostolis praelatum, Matth. 16. Allegavi Hieronymum, Bernhardum, Leonem et Cyprianum, ad quos nihil respondit, quamvis manu sua in suo libro fateatur b. Cyprianum sensisse 'ecclesiam' fundatam supra Petrum\apnota{26}{M U petram.}, sed ausus fuit addere, quod b. Cyprianus 'ibi' fallatur. Ad Augustinum nisus est respondere, nam de Cypriano est indubitatum, qui alia epistola ad Cornelium\apnota{28}{U Cornelium 8, Mc hat 8 vor epistola eingeschoben.} Rom. ecclesiam vocat matricem et radicem aliarum.

Impingit mihi, quod dictum citaverim ab Augustino retractatum. Lectoris iudicium obscurare nequit\apnota{30}{M nequid.} reverendus pater, nam adeo eum suae prioris sententiae non poenitet, ut etiam pro firmamento b. Ambrosium introducat, sed quod ab initio induxi\apnota{32}{M induxit, Mc induxi.}, testatur Augustinus, se per petram etiam exposuisse Christum, neutram tamen\apnota{33}{tamen fehlt in U, M nec tamen.} retractans sententiam, neutram etiam praeeligendo\apnota{34a}{in M Lücke, Mc 'neutram sententiam praeeligendam' (so!).} subdit: Harum autem duarum\apnota{34b}{fehlt in U, autem in M gestrichen.} sententiarum quae sit probabilior, eligat lector. Augustinus definire non audet\apnota{35}{M U diffinire; U audit.}, et reverendus pater unam sententiam vult reiicere tot s. patribus 'attestatam' et aliam pro\apnota{36}{U acceptatam et cum [statt attestatam et aliam pro], Uv = P.} suo arbitrio acceptare. Unde Augustinum non admitto in uno\apnota{37}{M U novo.} capite dixisse contraria vel adversa, sed diversa. Quoniam ambas amplector\apnota{38a}{U amplectatur.} sententias, et Christus est petra\apnota{38b}{U quoniam petra est Christus; M fuit, Mc est.}, et [293] Petrus fuit petra. Quare maneo cum Ambrosio, Hieronymo, Cypriano, Bernhardo et aliis sacris conciliis\apnota{2}{M alii [Mc aliis] consiliis.} ac decretis.

Secundo, cum ab ostiaria sumpsisset argumentum reverendus pater, communem\apnota{4}{U meam.} sententiam debellaturus, iussi eum, ut grammatica ratione verba Christi melius ponderaret, omnium enim consensu Matth. 16: super hanc petram aedificabo et tibi dabo etc. Christus promisit Petro, non tunc dedit. Sed dimittens grammaticam suam, quam tamen dixit\apnota{7a}{M tamen dicit, U dixit tamen.} plus valere ad theologiam\apnota{7b}{ad theologiam fehlt U. M valere, Mc pertinere.} reliquis partibus philosophiae\apnota{8a}{Mc theologiae.}, opponit, quare tunc decreta fundent\apnota{8b}{M fundant.} se in illo\apnota{8c}{M hoc.} loco Matth. 16. et ipse me fundarim\apnota{9a}{M ipse met fundaret.}, a 'quo iam recedam'\apnota{9b}{M 'quo tamen iam recedo', U qua iam recedo.}. Respondeo: Quia Christus est via, veritas et vita, ideo indubitate credendum est eum

praestitisse Petro, quod fuerat pollicitus. Quare decreta recte argumentantur ab illo loco, ubi Christus promisit, sed primo post resurrectionem praestitit.

Sed tacite obiicit, Christum non solum Petro, sed omnibus apostolis hoc praestitisse dicendo: Accipite spiritum sanctum etc.\apnota{14}{M U quorum remiseritis etc.}. Noluit reverendus pater exprimere, ubi promissum principatum Petro dederit, nam in hoc conveniunt doctores, in coena Christum discipulos fecisse sacerdotes, dando eis potestatem super\apnota{17a}{U supra.} corpus Christi verum, dicendo\apnota{17b}{M U dicens.}: Hoc facite in meam commemorationem, 'Luc. 22' [v. 19]\apnota{18}{U in meam Luc. 20 (so!) Uv 22.}, et deinde in die resurrectionis dedisse eis potestatem supra corpus Christi mysticum dicendo\apnota{19}{fehlt in M, eis . . . Christi und dicendo fehlen in U.}: Accipite spiritum sanctum. Sed primatum et praelaturam totius ecclesiae commisit\apnota{20}{U promisit.} Petro: Pasce oves meas, Johan. ult. [21 v. 15 ff.] Hoc\apnota{21a}{U hinc, Uv = P.} Gregorius, Chrysostomus, Hieron.\apnota{21b}{fehlt in U.} et alii s. patres. testantur\apnota{21c}{M attestantur.}. Ait enim Chrysostomus: Eximius apostolorum erat Petrus et os discipulorum, vertex collegii. Unde et negatione deleta committit ei\apnota{23a}{U promittit [ei fehlt].} praelationem fratrum. Et homilia 80.\apnota{23b}{U 8.} de poenitentia: Sed eadem nocte Petrus lapsus est et resurrexit, post illum tam\apnota{25a}{fehlt in M.} gravem casum rursus eum ad priorem gradum duxit poenitentia et ei\apnota{25b}{M egit totum orbem terrarum ei, U ei per totum orbem terrarum.} totum orbem terrarum, ecclesiae praesidentiam tradidit\apnota{26}{M dedit.}. Sic glossa 1. Pet. 5 [v. 2]: Pascite qui in vobis est gregem\apnota{27}{fehlt in U.}, ait: sicut Dominus soli Petro totius gregis curam habere commisit etc.\apnota{28}{fehlt in M.}. Quare prius promissum Christus tunc praestit. Sicut et b\apnota{29}{fehlt in P.}. Gregorius consentit: Cura totius ecclesiae et principatus Petro committitur, scilicet: Pasce oves meas. Quodsi etiam\apnota{30}{U coram, Uv = P.} non monstrassem, ubi Christus dedisset, veritas tamen mentiri non potuit\apnota{31}{U poterit.}. Et multa sunt facta, quae non sunt scripta, ut docet eleganter\apnota{32a}{M U eleganter docet.} Alexander 3.\apnota{32b}{3. fehlt in M.} ca. Cum Marthae de celebratione missarum\apnota{32c}{U etc. [statt de celebratione missarum].}.

Sed iterum opponit reverendus pater, quia post spiritus sancti missionem\apnota{33}{M missionem spiritus sancti.} adhuc Petrus peccavit\apnota{34}{U peccaverit.}. Et satis magnificavit peccatum Petri, sicut scribitur ad Gala. 2. Quamvis dissensio fuerit primo inter apostolos Petrum et Paulum, deinde\apnota{36}{M dein.} inter ecclesiae doctores, Hieronymum et Augustinum, tamen non possum in hanc adduci sententiam, ut aliquis apostolorum post missionem spiritus sancti ceciderit in aliquod peccatum mortale, sicut illud de sanctificatis in [294] utero etiam

non admittitur. At illud praeter institutum. Dato tamen et minime concesso, quod sanctus Petrus in simulatione sua peccasset mortaliter, adhuc mansisset petra et caput ecclesiae\apnota{3a}{petra et caput ecclesiae fehlt in M, Mc füllt die Lücke aus petra et ecclesia.}, nisi 'reverendus pater velit' et hunc\apnota{3b}{Mc vellet [et hunc fehlt].} articulum Hussiticum\apnota{4}{M hussiticum articulum.} defendere, quod non credo: Nullus est dominus civilis, nullus praelatus nullus episcopus in peccato mortali; quod summam\apnota{5}{U summum.} faceret Christianae religionis incertitudinem.

Quarto, de decretis apud eum frigidissimis hoc dico: Toties in ecclesia summorum pontificum constitutiones in materia fidei receptas, et quod plus est, olim Rom. pontifex episcopum non confirmabat, nisi mitteret schedam\apnota{9}{M cedam.} protestationis fidei, in qua profitebatur se credere evangelia, quatuor concilia, synodos legitimas\apnota{11}{M legitimos.} et decreta summorum pontificum. Ob eam rem Johannes papa Gilbrechtum Coloniensem\apnota{12}{M gibrechtum, Mc Gilibrechtum [Coloniensem fehlt].} episcopum noluit confirmare, quod hanc fidei schedam\apnota{13a}{M cedam.} non perfecte 'obtulisset', ut legimus can. Optatum 100. dist\apnota{13b}{U 20. dist.}.

Quod vero Anacleti decretum reiicit, diximus heri, qua auctoritate id faciat vel qua probatione. Nam cum originalia sint prae manibus, nullus adhuc dixit, istud decretum s. martyris Anacleti non esse. Sed Nicol. Cusanus Germanorum doctissimus, qui fatetur se concordantiam fidei scripsisse, ex originalibus illud Anacleti allegat. Fucum addidit reverendus pater, quia in tam s. martyre non sit talis inscitia, ut exponat Cepham, id est caput, quasi ad sanctitatem vitae tanta requiratur scientia. Dicam\apnota{20}{U dicamus.} tamen unum, quia Cephe\apnota{21}{U Cephe aut Cephas.} debet esse Syriacum et hebraeis 'quoque familiare', valens tantum sicut Petrus vel solidum\apnota{22}{M solidum vel validum.}, ut Erasmus post Hieronymum refert. Addamus\apnota{23a}{M addamus, gestrichen, Mc. a. R. ad damasum} hoc et singulare\apnota{23b}{et singulare fehlt in M.}, quod Cusanus chaldaice et hebraice doctus in libro\apnota{24a}{U libris.} excitationum testatur, Petrus valere etiam\apnota{24b}{M U etiam valere.} tantum\apnota{24c}{fehlt in P.}, sicut caput domus. Quare non tanta inscitia laboravit Anacletus, si Cepham pro capite interpretatus est\apnota{26}{U interpretatur.}. Sed utcunque sit non video, quo fuco possit negare alia decreta s. martyrum et confessorum, ut Clementis\apnota{27}{M add 22. dist., U 20. dist. W. A. vermutet 80. dist.}, Marcelli, Julii, Pelagii, Nicolai, Agathonis, Symmachi\apnota{28a}{U Simonis.}, Benedicti\apnota{28b}{fehlt in P. U Vigilli Benedicti.} etc. Quare tantis patribus et sacro concilio malo condescendere.

At\apnota{30}{U als Überschrift ad scrupos.} omnium invitissime audivi, quod reverendus pater contra illustrissimi principis iussum, quo intacta manere voluit illustrissimus princeps a sacris conciliis decisa, reverend. tamen pater in hesterna

obduruit sententia, et quatuor proposuit articulos inter Hussiticos condemmatos, quos reputat ipse\apnota{34}{M ipse reputat.} catholicos et evangelicos, pulchrum certe pro Bohemis patrocinium\apnota{35a}{U pulchre certe Bohemis patrocinatur.}. Pro concilii\apnota{35b}{M conciliis} tamen reverentia voluit ab impostore aliquo\apnota{36}{M aliquot.} additos. Primo omnium egregius doctor et nobilis Hieronymus de Croaria\apnota{37a}{fehlt in M, Mc U Croatia.}, qui acta concilii\apnota{37b}{M consilii.} curavit imprimi, transumi fecit sub forma authentica, et cum res ita manifeste\apnota{38}{M manifesta.} acta sit Constantiae\apnota{39a}{M Constantia.}, Hussitae tam 'diu non tacuissent'\apnota{39b}{Mc reticnisse, U reticuissent.} articulos aliquos esse subditicios [295]. Neque 'poterit' suspicari reverendus pater\apnota{1a}{reverendus pater fehlt in M.}, quod heri assumpserat, adulatorum\apnota{1b}{U adulatoris.} pernicie in Rom. pontificem factum, cum Johannes Huss anno XV.\apnota{2}{U a. R. in Julio.} fuerit combustus, Hieronymus collega eius\apnota{3a}{M U suus.} anno XVI.\apnota{3b}{M 16 Julii, U a. R. in Maio.} Martino V. primum\apnota{3c}{M primo.} electo anno XVII. ipso die Martini. Quare s. patres et viri inadulabiles ex omnibus nationibus principalibus ad hoc deputati articulos illos discusserunt, et sacrosancta synodus eos damnavit, reprobavit et auctorem combussit. Ideoque a quolibet bono Christiano pro condemmatis et reprobatis 'habendi sunt'. Nec imponat mihi reverendus pater, quod velim de illis articulis iudicare, quia iam iudicati sunt.

Non relevat, quod synodus\apnota{9}{M U quod cum synodus.} meminit aliquos esse haereticos, alios temerarios, seditiosos et piarum aurium\apnota{10}{M animarum, Mc a. R. aurium.} offensivos; quoniam in quemcunque ordinem redegerit illos articulos, non possunt dici Christianissimi et evangelici, et quamvis non assumpserim hunc 'laborem defendendi' totam\apnota{12}{M totum.} synodum in articulorum condemnatione\apnota{13a}{M condemnationem.}, dicamus tamen aliqua\apnota{13b}{U doceamus, (Uv dicamus) tamen aliquid.}, et brevibus.

Articulum primum iudicat catholicum et esse Augustini cap. 6. super Johannem\apnota{15}{U Johanne.}. Dico: Forte reverendus pater benigniter interpretatur articulum, sed quia res non sermoni, sed sermo rei est subiectus, non fuerunt stupidi concilii\apnota{17}{M consilii.} definitores. Verum est, unam sanctam et universalem esse ecclesiam, sed quod sit tantum una, sicut est unus numerus praedestinatorum ad Hussiticam intelligentiam\apnota{19a}{M Hussiticum intellectum.}, est haereticissimum\apnota{19b}{Mc (haereticissimum) velut amissa fide, die drei letzten Worte gestrichen.}, quo ipse voluit, existentes in peccato mortali, velut amissa\apnota{20}{M a missa.} fide, non esse in ecclesia, cum tamen regnum caelorum comparetur\apnota{21}{M U a Christo comparetur.} decem virginibus, ubi quinque erant prudentes, quinque vero fatuae, cum reliquis parabolis eidem proposito deservientibus.

Neque Augustinus tract. 26. super Johannem facit iota\apnota{23}{M U iotam.} ad Hussiticum propositum, communionem 'eucharistiae'\apnota{24}{M ecclesiae gestrichen, Mc eucharistiae.} commendans, de quo lectoris desidero iudicium.

Ita de alio dicit articulo: Duae naturae, divinitas et humanitas\apnota{25}{U divina et humana.}, unus est Christus. Nihil pro hoc adduxit, nisi quod sit de fide. In Athanasio\apnota{26}{U Athana.; W. A. Athanasiano.} aliter legimus: Deus et homo unus est Christus, non deitas et humanitas.

In alio articulo: divisio immediata humanorum\apnota{28}{fehlt in U, in Uv nachgetragen.} operum etc., quam ipse mire in me retorquet, et quasi a me iudex sit datus\apnota{29a}{U citatus.}, suum\apnota{29b}{M suam.} collegam exclamavit de me\apnota{29c}{de me fehlt in U.} triumphasse. De quo vos omnes, qui interfuistis, potestis dare testimonium, quam verissime 'dixerit\apnota{30}{Mc 'dixit'.} reverendus pater' et encomium prius 'cecinit' ante victoriam\apnota{31}{ante victoriam fehlt in M.} Wittembergae, iam aliud etiam cecinit\apnota{32}{M alium; cecinit fehlt in M U.}, quasi perdita victoria.

Ad rem.

Articulum istum nunquam reputavi\apnota{34}{M U illum; nunquam reputavi fehlt in U}, nunquam reputabo christianum, et ob hoc Gregor. 'Arim.' 28. dist. 2. repuli. Neque actum fuit\apnota{35a}{M fuerit.} inter nos, an 'possit'\apnota{35b}{M fuerit, gestrichen, Mc possit.} dari opus indifferens, vel an omnis vita infidelium\apnota{36a}{P fidelium.} sit peccatum, et\apnota{36b}{M aut, U vel, Uv = P.} similia, quae hunc respiciunt articulum. Unde quod palliando errores Hussiticos innuit aliquando sinistre\apnota{38a}{M sinistra.} verba intelligi, sicut Domino Christo evenerit, quae\apnota{38b}{Mc et quae.}, rogo, comparatio lucis ad Belial? quod\apnota{39}{fehlt in U.} Johannem Huss in hoc comparat Christo, non passus, quoad veritatis sugillationem in parte admissam me comparare\apnota{40}{M admissa ne compares.} sacra [296] concilia\apnota{1}{U consilia.} sacris scripturis, in quibus utrisque est indubitata et infallibilis veritas. Nullum enim Christianum movere\apnota{2a}{M moveri.} debet, quod\apnota{2b}{U et quod.} dicit, concilia\apnota{2c}{M U consilia.} esse homines et ita creaturas et sic peccare posse, nam si errant, ut fuit concilium\apnota{3}{M consilium.} Ephesinum a Leone, credo, papa damnatum, ut Ariminense, 'Aquisgranense', tunc non sunt concilia, sed conciliabula. Unde potius hoc constantissima fide tenere debemus, quidquid concilia legitime congregata in iis\apnota{6}{M U his.}, quae sunt fidei, determinaverint et definiverint\apnota{7a}{U et diffiniverint; et fehlt in P M.}, esse certissimum\apnota{7b}{P certissima, M certissimam.}. Sic enim Christus manet nobiscum usque ad consummationem saculi, et si duo congregati fuerint in nomine meo etc.\apnota{9a}{M Christi (statt meo); etc. fehlt in M U.}. Praerancidum est, quia homines sint\apnota{9b}{M U sunt.} in

\noindent\textit{[test-p03: impresas 59-60 (test pp. 5-6). Claves: serie p.59 simples; serie p.60 con sufijo 2/b donde la numeración del original colisiona (18, 24a, 25a, 25b, 26).]}
\noindent\textit{[test-p04: impresas 61-62 (test pp. 7-8). Serie p.62 con sufijo -62 donde colisiona (31, 34, 3a, 3b). El fragmento cruza el salto 61->62 a mitad de frase ("scripturae / dicit in fine"), sin guion. Termina a mitad de frase ("quod hoc" -> sigue p.63).]}
\noindent\textit{[test-p05: impresas 63-64 (test pp. 9-10). Serie p.64 con sufijo -64 donde colisiona (19a, 19b, 27, 29). Cierra las 10 pp. de prueba (impresas 55-64). Termina a mitad de frase (sigue p.65, fuera del test).]}
\noindent\textit{[nota 32: venía partida entre pp. 92-93; texto completo fusionado en [\textasciicircum{}092-32] de pp-092.md]}

\noindent\textit{Notas del aparato sin llamada en el cuerpo:} [14] P U Mc baptizabat, M baptizabas, W. A. baptizarat.; [18] M sancto Gregorio.

concilio, et ita peccare vel errare posse\apnota{9c}{M et ita peccare (Mc add errare) posse, U et ita errare possunt, Uv = P.}; nam etsi ut homines defectibiles sint, verum concilium legitime congregatum\apnota{10}{M in concilio legitime congregato, U in consilium legitime congregatum.} non humano sensu, sed spiritu divino regi credendum est, unde\apnota{11}{regi credendum est fehlt in U; U ut (statt unde); Uv = P.} sat probata\apnota{12a}{M probata sunt, U probata sint (Uv ? = M).}, quae concilium probat\apnota{12b}{U probavit.}, et non sunt\apnota{12c}{fehlt in U} per cuiuslibet\apnota{12d}{M U cuiusque.} singularitatem et capitositatem\apnota{13}{M singularitate et capitositate, Mc a. R. capciositate, U singularitatem vel capitositatem.} ulterius discutienda, sed captivare debemus intellectum in obsequium fidei. Quare miror\apnota{14a}{M minor, Mc miror.}, quod textum negaturus iurium\apnota{14b}{U Siculam (statt iurium), Uv trägt iurium nach.} commentatorem admittit\apnota{15}{M admittit.} Panormitanum.

Sic quod postremo adduxit ex verbis apostoli Pauli\apnota{16}{U patet} de fundamento, semper dixi, hoc esse de omnium catholicorum sententia, Christum esse fundamentum ecclesiae\apnota{17}{Christum esse fundamentum ecclesiae fehlt in U, ecclesiae fehlt in M.}, Christum esse caput ecclesiae, ille est enim\apnota{18}{MU enim est.} principalis dominus. Sed quia Petrum constituit vicarium, et potestatem ei super tota ecclesia contulit, ut\apnota{19}{U ubi.} s. Leo can. Ita Dominus 19. dist.\apnota{20a}{alle Quellen 21. dist.} testatur\apnota{20b}{M testavit.}. Et miror plurimum, quod relicto textu s. Leonis reverendus pater arripit ibi unam\apnota{21}{fehlt in M.} glossam textui contrariam. Omnino enim convenit Christo ascendente\apnota{22a}{U Immo tantum confert Christo ascendenti (W. A. Christus ascendens).} ad caelos, ut loco sui\apnota{22b}{M suo.} constitueret hic vicarium caput, ad quod esset recursus in dubiis et aliis casibus emergentibus, ne tota Christi ecclesia esset 'unum' confusum chaos\apnota{24}{M caus, Mc chaos.} Anaxagoricum. Istam sententiam tota tenet ecclesia, sicut est in responsorio\apnota{25}{M P responso.}, quo ecclesia de s. Petro utitur: Tu es pastor ovium, princeps apostolorum. Tibi tradidit deus omnia regna mundi. Ita oramus in ecclesia. Imo divus\apnota{27}{M U divinus.} Franciscus, sigillis stigmatum approbatus, fratres suos et totum orbem terrarum docuit obedire summo pontifici, et in regula hoc praecepit. Ideo de eo canit ecclesia: Franciscus vir catholicus et totus apostolicus tenere fidem Rom. ecclesiae docuit. Ita plurimi summi\apnota{31a}{M romani.} pontifices, ita sacra concilia definiverunt\apnota{31b}{M U diffiniverunt.}, quae brevitatis studio non adduco. Plurimae quoque s. patrum sententiae idem sonant et approbant. Quae omnia suo tempore, cum videro reverendi patris tractatum super hoc negotio, adducere non negligam\apnota{34}{M negligabo, Mc negligam.}.

Si tamen pri-[297]matus Rom. pontificis tantum iuri innititur humano et elevatur consensu populi, unde tunc reverendo patri\apnota{2}{U reverendi patris.} mendicitatis privilegium? Unde iste religionis habitus? Unde potestas\apnota{3}{M potestas praedicandi (also doppelt!)}

confessiones audiendi, praedicandi, cum aliis innumeris\apnota{4a}{M innumerandis.} privilegiis, 'quibus ab apostolica sede\apnota{4b}{M U a sede apostolica.} obruuntur', archiepiscopis, episcopis\apnota{5}{U episcopis archiepiscopis.} et curatis frequenter contradicentibus? Qui tamen ut filii obedientiae' Rom. pontificem audiunt et in eo Christum, ut sic per patientiae bonum vitam adipiscantur aeternam.

Haec brevia volui adiicere\apnota{8}{M adducere.} in praesenti disputationis negotio, quantum tempus dedit, plura alioquin allaturus, nisi fastidium me absterreret et illustrissimi principis iussum\apnota{10}{Mc U iussus.}.

D. D. Martinus.

Duo mihi incumbunt agenda, primo\apnota{12}{M U primum.} respondendum confutationibus egregii D. D., secundo pro conventionis pacto opponendum eidem. Et quia iam triduo respondi tempusque absumptum\apnota{14}{M U consumptum.} est, ut nec unum possim explere, dico breviter:

Quod omnia inducta per egregium D. D. copiosissime quidem sunt dicta, sed prorsus nihil\apnota{16}{M nihil prorsus.} ad scopum, cum iam per tres dies promiserit, divino iure se acturum\apnota{17a}{U acturum se.} contra me, nec adhuc audita est una\apnota{17b}{fehlt in M.} auctoritas divini iuris, praeter illam Matth. 16, quam omnes audivimus, quam varie tractetur a s. patribus, et quod maior pars eorumque sanior sententia pro me stet. Unum ipse locum Augustini et alterum Ambrosii pro se adduxit, cetera omnia fuerunt vocabula Anacleti, Marcelli, Pelagii etc.\apnota{21}{M U Pelagii, Anacleti etc. [etc. fehlt in U].}, eorum, qui iura humana conscripserunt. Cyprianus vero fere\apnota{22a}{M fero.} totus mecum sentit\apnota{22b}{M sit.}, quod committo iudicio prudentis lectoris. Similiter 'et clarissimus\apnota{23}{M clarus.} textus' Hieronymi ad Euagrium, et super Titum. Deinde Gregorius in registro per senas ferme epistolas, post haec historiae et usus orientalis ecclesiae. Neque enim quicquam facit ad rem, quod D. D. schismaticos Graecos allegavit, 'Nestorium, Macedonium, Achatium, Eutychem et ceteros'\apnota{26}{Macedonium, Achatium, Eutychem fehlt in M U, Uv = P.}, non erant illi ecclesiae orientalis. Possem 'et' ego sic latinam ecclesiam describere per 'Pelagianos, Manichaeos, Jovinianos', Vigilantianos\apnota{28}{M U possum et ego [M ergo] latinam ecclesiam sic describere per Pelagios [M Pelagianos], Manichaeos [fehlt in M], Jovinianos, Vigilantianos [U add Helpidianos].} et similia monstra. Sed absit, ut propter paucos malos et schismaticos totam aliquam ecclesiam schismaticam appellem\apnota{30a}{M appellam.}.

Cetera 'non potero prosequi'. Et\apnota{30b}{fehlt in M.} reservo mihi praecipue articulos Johannis Huss mihi obiectos 'et inter eos novissimum de humanis actibus' ad calamum et papyrum, et cetera omnia, quae habebam\apnota{32}{M haberem.}

opponere, ut quae hic facere non possum\apnota{33}{U possim.} exclusus spatiis, praestem per scripturam, voloque in hoc cedere et obedivisse illustrissimi principis\apnota{34}{U principi.} studiorum patroni clementissimi\apnota{35}{M patronum [clementissimi fehlt]; P clementissi.} iussis, de quo vos dominos notarios requiro et auditores in testimonium voco.

Actum ut supra, praesentibus ibidem honorabilibus Olavo Roters de Osvaberga et Coelestino Goricio Gorlicensi, clericis Coloniensis et Misnensis dioecesis, testibus specialiter vocatis et legitime requisitis.

Johannes Graumann Notarius subsignavit\apnota{36}{Die Worte von Actum an fehlen in M U.}.

[298] Die Jovis septima mensis Julii, hora septima mane\apnota{2}{Die Zeitbestimmung fehlt in M, in U lautet sie: Die septima Julii anno 1519. mane hora septima.}.

D. D. Eccius.

Quia reverendus pater heri horam habuisset respondendi ad nostra, ut sic illustrissimi principis iussa fuissemus exsecuti, ubi respondendo plurimum per verbum\apnota{6}{per verbum fehlt in M.} »transeat« appellavit ad calamum, tamen deliberatione habita proposuit, se iam plenius responsurum. Cui ego non dissensi pro veritatis amore et elucidatione. Admiror, quomodo reverendus pater magno boatu audet coram tot viris\apnota{9}{M U doctis viris.} intonare, quia dixerim me probaturum de iure divino monarchiam et primatum in ecclesia, cum tamen solum allegaverim dicta s. patrum et\apnota{11a}{fehlt in M U.} decreta pontificum\apnota{11b}{M pontificum et sic iura humana.}, cum tamen ego\apnota{11c}{fehlt in U.} adduxerim ei ius divinum Matthaei 16. de\apnota{12a}{M pro.} promissione aedificationis super\apnota{12b}{U supra.} petram, quod\apnota{12c}{M quo.} Christus specialiter rogavit pro fide Petri, quod specialiter ei commisit\apnota{13}{U promisit.}, ut confirmaret fratres, quod specialiter ei commisit in praesentia Johannis et Jacobi, ut pasceret oves, tribus vicibus. Et hoc ius divinum ita esse intelligendum, tunc adduxi s. patres, sanctos\apnota{16}{fehlt in U.} summos pontifices et martyres, consensum conciliorum et omnium universitatum. Quare ex iure divino probatum est, esse primatum in ecclesia\apnota{18a}{U ecclesia dei.}, nisi quod reverendus pater intelligentiam suam praeponit\apnota{18b}{M U proponit.} intellectui tot s. patrum, summorum pontificum et sacri concilii\apnota{19}{U consilii.}, cum laudabile Constantiense concilium\apnota{20a}{U consilium.} talem articulum inter pestilentes errores\apnota{20b}{M articulos; errores fehlt in U.} Johannis Huss damnaverit\apnota{21a}{M damnavit.}. Ad quod reverendus pater silet. Sed nihilominus\apnota{21b}{Sed nihilominus fehlt in U.} si vult offerre\apnota{21c}{M U affere.} meliora, pro veritatis elucidatione, quod heri potuisset

fecisse iuxta illustrissimi principis iussum\apnota{23a}{M iussu.}, sum paratus audire, eo pacto, quod dum\apnota{23b}{M tum.} fungitur officio respondentis, iam in fine actus non transferat se callide ad personam opponentis, sed quaecunque oppositurus sit contra istam ineluibilem veritatem, paratus sum his adductis respondere et ea diluere, de quo protestor.

D. D. Martinus.

Cum heri statutum esset\apnota{28}{M esse, Mc esset.} nomine illustrissimi principis, patroni nostri, ut materia ista consumaretur\apnota{29}{M consumeretur; ista fehlt in M.} eo die, et egregius D. D. maiorem partem temporis sine necessitate et odiosius quam decebat hoc splendidum auditorium consumpsisset, mihi pro triduana responsione unica fuit\apnota{31}{M mihi.} reliqua horula, qua tantae farragini\apnota{32}{M varragini.} suae responderem simul et opponerem. Placuit postea gratuito affectu eidem illustrissimo\apnota{33a}{M illustrissimi.} principi nostro, quod\apnota{33b}{U ut.} copia daretur mihi\apnota{33c}{M mihi daretur.} respondendi et opponendi\apnota{34a}{M et opponendi et respondendi.}, de quo gratias ago illustrissimae\apnota{34b}{M illustrissimi.} gratiae suae quam maximas. Ideoque ad propositum veniendo, antequam coeptam hesternam responsionem continuem, primo hanc matutinam nebulam egregii D. D.\apnota{36}{U d. d Eccii.} dispellam, ubi more suo satis odiose et semper obiicit Hussiticas pravitates, non ob-[299]stante, quod toties me expurgaverim. Ideoque paulum praefabor Germanice, quod intelligam me pessime audire apud vulgum.

Dicebat\apnota{3}{U Declarabat.} autem se non impugnare, nec Christiane a quoquam improbari\apnota{4}{U impugnari.} posse primatum et obedientiam Romanae ecclesiae. Verum eo pertractum esse, ut declarare se oporteat, eum primatum non esse de iure divino, sicut neque imperatoriam apud Germanos potestatem\apnota{6}{U potestatem apud Germanos.} improbari posse, quamvis in sacra scriptura non fundatam.

D. D. Eccius\apnota{7}{U quamvis d. Eccius.}

dixit, verum esse, eam esse controversiam, an primatus ille\apnota{8}{U primatus ecclesiae Romanae.} sit de iure divino, quod improbare censeat\apnota{9}{U constat.} esse Hussiticum et articulum a Constantiensi concilio damnatum\apnota{10}{U condemnatur.}, esseque dolendum, si Christianorum corda tam frigida essent, ut nemo hoc defenderet.

Deinde D. D. Martinus Latine prosecutus est dicens\apnota{11}{U post hoc}.

Quod ergo inducit hanc auctoritatem Matth. 16, contendens ibi primatum promissum esse, satis responsum est in praecedentibus, quod

D. egregius D. vix duos auctores habet et eosdem ambiguos in suam sententiam, cum tamen maior pars statuat Petrum ibi\apnota{15}{fehlt in M.} fuisse personam omnium apostolorum et fidelium, quod et ipse textus consequentia\apnota{16}{W A. continentia.} sua urget, dum Christus interrogat omnes, et non solum Petrum, dicens: Vos autem quem me esse dicitis? Et ibi Petrus, os apostolorum, ut Chrysostomus ait, et persona apostolorum, ut Hieronymus et Origenes dicunt et Beda et Augustinus et Ambrosius, respondet pro omnibus. Ideo dixi nihil probari ex hac auctoritate de primatu\apnota{21a}{M primatui.}. Item, quod textus\apnota{21b}{M U textus ipse.} non dicit sic: Tu es Petrus, et super te, sed super hanc petram, manifeste indicans petram aliam, a qua Petrus dictus sit. Proinde cum sint vulgata iura et ipsi Rom. pontifices testentur, interpretationem scripturae magis valere apud doctores quam apud se ipsos, in causis vero decidendis\apnota{25}{U descendendis (W. A. verbessert nach Löscher.)} secus.

De articulo illo Hussitico dixi satis heri, quod nondum probatum est, eum esse haereticum. Hoc pro responsione matutina.

Redeundo ad hesternam, quando egregius D. D. secundo loco opposuit depositionem Chrysostomi ideo factam illegitime, quia odio imperatricis facta sit, nego consequentiam; non ideo potestas ordinandi et deponendi per sacrum Nicenum concilium\apnota{31}{M concilium tragens.} statuta fuit illegitima, quia casu forte habuit malum usum. Nam in eadem historia\apnota{32}{M U etiam historia.} scribitur, quod Rom. pontifex Victor\apnota{33}{fehlt in M.} martyr voluit Asiae episcopos excommunicare. At illi rursus tanquam superiores mandaverunt silentium, et ut ecclesiam dei\apnota{34a}{fehlt in U, in Uv nachge-Nicenum.} non perturbarent iusserunt\apnota{34b}{M perturbarent iussarent [Mc iusserunt].}, lib. 5. ecclesiasticae historiae\apnota{35}{libro 5. eccles. hist. fehlt in M.}. Sed et Irenaeus Lugdunensis Galliae\apnota{36a}{M Tirenaeus Luctunens Galli (so!)} episcopus eundem Rom. pontificem\apnota{36b}{U episcopum.} coercuit, et erant tunc tem-[300]poris in Graecia Epiphanius celeberrime laudatus a Hieronymo, Gregorius Nazianzenus, Basilius\apnota{2}{M Basilius etc.} Christianissimi viri, non tamen sub Rom. pontifice unquam fuerunt, sed iuxta Niceni\apnota{3a}{M Nicelmi.} concilii statuta\apnota{3b}{M constituta.} ab episcopis provinciarum ordinabantur.

Quod vero egregius D. doctor inducit Julium I. excommunicasse Arcadium imperatorem, nihil facit ad rem. Nam et Bonifacius VIII. ausus est et regem Francorum velle e sede eiicere. Non probat\apnota{6}{U probatur.} ius esse, quidquid Rom. pontifices pro humana fragilitate quandoque tentaverunt. Nec mirum est, sanctos illos\apnota{8}{M U istos.} viros ambitionis tentatione

vexatos fuisse, quando plus quam semel sanctissimi apostoli etiam praesente Christo laboraverunt eadem.

Quod dicit episcopo 'Hierosolymitano' datam praerogativam honoris, non primatum\apnota{11}{M praerogativam datam (Mc a. R. gratiam) non primatum, U datam honoris praerogativam [non primatum fehlt].}, recte dicit, sed non probat primatum esse\apnota{12}{U esse hinter iure divino, M hat esse an beiden Stellen.} Rom. pontificis, saltem iure divino.

Etiam ibi optime dicit egregius D. D., quod sexta synodus dedit Romanae ecclesiae primatum, non autem praecedentes synodi. Hoc est enim quod volui, quod synodicis statutis et humano iure, consensu\apnota{15}{M consensum, Mc consensu.} fidelium Christi, cui non licet resistere, datus sit iste primatus. Quod et divus Gregorius in registro in haec verba dicit\apnota{17}{U dixit.}: Scitis, quod veneranda synodus Chalcedonensis huic apostolicae sedi obtulit primatum et tamen nullus praedecessorum meorum unquam praesumpsit\apnota{19}{M U praesumpserit.} acceptare. Proinde nihil facit, quod egregius D. D. dicit, Gregorium non reiecisse\apnota{20}{M reiecisse.} primatum, sed asseruisse. Legat pius lector epistolas\apnota{21}{M [epistolam.} Gregorii, et inveniet, Gregorium prohibuisse primatum et oecumenicum\apnota{22a}{M icumenicum [et fehlt], Mc universalem.} pontificatum non vendicasse, ut ex supra relatis dist. 99.\apnota{22b}{P M U dist. 95.} satis probavimus.

Ad auctoritatem Augustini, ad Innocentium Rom. pontificem et huius ad illum respondeo: Legat lector utriusque verba et iudicet, anne egregius D. D. novas significationes vocabulis imposuit.

De viginti\apnota{27}{U „22".} annis ecclesiae ante Romanam dixit, cum Petro fuisse translatum primatum ad Romam, forte volens hoc proverbium: Ubi' papa, ibi Roma. Respondeo, probandus erat primatus, qui transferretur, non translatio, quam libens concedo.

Induxit etiam auctoritatem Augustini contra Manich., quam suspicor esse 2. libro de\apnota{32a}{M beati Augustini 2. lib.} doctrina Christiana\apnota{32b}{U doctrinae christianae.} c. 8., quod in re dubia ad certitudinem fidei valeat auctoritas catholicae ecclesiae\apnota{33}{U ecclesiae catholicae [catholicae fehlt in M].}, quae ab ipsis apostolorum sedibus succedentium sibimet et episcoporum serie et populi consensu firmatur\apnota{34}{M firmati.}. Nihil de Rom. ecclesia loquitur\apnota{35a}{U loquimur.}, sed universali, et pro me facit contra D.\apnota{35b}{fehlt in M.} D., quia successiones diversarum sedium apostolicarum sibimet et consensum\apnota{36}{U consensu.} populi describit Augustinus, non autem confirmationem aliarum sedium per unam sedem.

Ad\apnota{39}{U nunc ad.} rem ipsam.

[301] De auctoritate Matth. 16: Tu es Petrus, et super hanc petram etc.\apnota{1}{et super hanc petram etc. fehlt in M U.} et ad inductos auctores dicit me non respondisse, quod committo

vobis auditoribus et lectoribus, an verum sit. Augustinum in retractationibus\apnota{3a}{M retractatione.} relinquentem electionem lectori\apnota{3b}{fehlt in P.} super duabus sententiis eiusdem auctoritis opposuit dicens\apnota{4}{M dicit.}, se ambas amplecti sententias, quod mihi placet\apnota{5a}{M Uv et ego facio [statt mihi placet].}, et sic nihil contra me probant\apnota{5b}{M U probat.}.

Secunda auctoritas divini iuris a dignatione\apnota{6}{M U dominatione.} sua inducta est haec: Pasce oves meas, Johan. ult. [21 v. 15 ff.], et\apnota{7}{fehlt in M.} gaudeo tandem post tres dies aliquando audiri unam auctoritatem divini iuris, quo firmissime nitatur eius sententia, sed videamus et\apnota{9}{fehlt in M.} ipsam.

Primum ista auctoritas videtur intelligi dupliciter, uno modo\apnota{10}{fehlt in U.} per verbum diligere, id quod est, sibi ipsi fidere\apnota{11}{M videre, P a. R. lectum fuit videre media longa.}, et quidvis licere, et per verbum pascere, idem quod primum esse et regnare. Et sic erit sensus: Petre, si diligis me, id est, si tua quaesieris et omnia feceris, quae adulatoribus tuis placent\apnota{13}{M placet.}, pasce oves meas, id est\apnota{14}{id est fehlt in M.}, esto primus et dominus omnium. Hunc sensum non habeo in meo codice.

Alio modo: Si diligis me, id est, si\apnota{15a}{fehlt in U.} abnegas temetipsum\apnota{15b}{U te ipsum.}, si ponis animam tuam pro me, si etiam omnem dignitatem respueris et non\apnota{16}{M U nihil.} nisi me dilexeris, ut egregie hoc loco exponit Augustinus\apnota{17}{hoc loco fehlt in M; U exponit Augustinus in eodem loco.}: Pasce oves meas, id est, doce, praedica verbum\apnota{18}{M verbo.}, exhortare, ora, bono exemplo praecede. Nam Graecum verbum\apnota{19}{M U verbum Graecum.} hoc loco non simpliciter regere et pascere significat, sed suaviter et molliter curare et omnia facere, ne quid desit ovibus\apnota{20}{M omnibus.}, et hunc sensum credo esse evangelicum.

Proinde rogo D. D.\apnota{21}{M et proinde rogo doctorem.} Eccium, ut mecum oret\apnota{22}{M U roget.} Dominum Jesum Christum, quo velit inspirare non solum summo, sed omnibus\apnota{23a}{U etiam omnibus.} pontificibus\apnota{23b}{non solum summo, sed omnibus pontificibus fehlt in M.}, ut haec verba credant ad se pertinere. Non est dubium, quin totus orbis obviis manibus et profusis lacrimis sit excepturus\apnota{24}{M accepturus.} talem virum, qui secundum haec verba se gerere vellet. Proinde si D. egregius\apnota{25}{U egregius dominus.} D. textum diligenter inspiciat, non ius ibi conceditur et privilegium, ut putatur, sed praeceptum imponitur, et labor hodie intolerabilis iniungitur, et officium demandatur ei, qui de facto est vel in primo vel in medio loco constitutus et non primum constituitur. Quod si urgeat, officium non posse administrari, nisi primum locus potestatis habeatur, placet, sed tunc alio verbo\apnota{31}{U alio modo.} struendus est iste locus, hic hoc verbum non nisi officium commendat.

Proinde recte fecerunt patres sancti, qui hunc locum ad omnes

pertinere dicunt, nam hoc nemo potest negare, si Petro hic sunt\apnota{33}{M sunt hic.} commissae omnes oves, quod ipse non omnes pavit, ut b. Augustinus dicit et Actuum liber evidenter demonstrat\apnota{35}{M evidentissime ostendit.}, et sic Christo inobediens fuit. Nec potest dici, quod licet per se non paverit, tamen per alios subordinatos. Primum, ut donemus, quod possimus reddere, non esse dictum Petro: Pasce per alium, sed ipsimet in persona propria: Pasce tu\apnota{38}{W. A. verbessert nach Walch non esse dictum Petro ipsimet in persona propria „pasce tu", sed „pasce per alium".}, tamen hoc convincit, quod nec apostolus ullus nec Paulus saltem a Petro sunt subordinati, qui multo\apnota{39}{M multos, Mc multo.} plures oves [302] paverunt quam Petrus. Ideo verbi intelligentia ex tota scriptura et circumstantia rerum gestarum petenda est\apnota{2}{M esset.}, quod haec auctoritas non ad solum Petrum pertineat aut non ad omnes\apnota{3a}{M U omnes oves.}.

Tertio dico\apnota{3b}{M tertium; dico fehlt in M U.}, si per ista verba probabitur primatus, incertus et nullus erit primatus, quia non ponitur primatus nisi sub conditione dilectionis, nec enim oportet scripturam dirumpere in diversas partes, sed magis concinnare. Ideoque cum sit incertum, quis diligat Christum, incertum quoque nobis erit\apnota{7}{M U erit nobis}, quis sit pastor noster. Et si ei non est necessaria dilectio, nec nobis necessarium erit agnoscere pastorem. Rectius ergo dicitur, quod hoc verbum sit praeceptum datum universis praelatis ecclesiae, ut contemptis divitiis, dignitate, etiam ipso primatu, tandem vita et morte sese impendant pro ovibus Christi. Quis est hic, et laudabimus eum?

Transeo ea, quae ex Chrysostomo attulit de vertice collegii, et quomodo Petro tradita sit praesidentia orbis\apnota{14}{M tradita praesentia (so!) totius orbis.} terrarum, quia omnia concedo, quod sit primus in honore, sed non solus in administratione, saltem\apnota{15}{in fehlt in P, saltem in M.} iure divino.

De casu Petri, Galat. 2, ubi iterum egregius D. D. odiosius mihi articulum Hussiticum obiexit, nullum esse dominum spiritualem vel civilem, qui\apnota{17}{M qui esset sine, die beiden letzten Worte gestrichen.} sit in mortali peccato\apnota{18a}{M U peccato mortali.}, si dicerem, Petrum ideo\apnota{18b}{M ideo Petrum.} non esse praelatum, quia peccasset mortaliter, respondeo: Hoc non agebatur, et ad rhombum\apnota{19}{M romgum.} est ista disputatio. Scio\apnota{20}{M U scio et optime scio, Uv = P.}, quod praelatus malus non sit reiiciendus. Ideo et ego damno articulum Hussiticum illum\apnota{21a}{P artieu- (so!) Hussiticum illum, U Hussiticum articulum; illum fehlt in M U.}. Sed hoc volui\apnota{21b}{M vult, U voluit.}, quod Petrus, quando in causa fidei praebuit scandalum, si non fuisset emendatus per Paulum, merito debuit amoveri a praelatura, nam haereticus pastor aut is, qui simulat\apnota{24}{M simulum, Mc simulat.} in periculum

fidei, graviter peccat. Nam hac simulatione Petri funditus peribat fides Christi, ut Paulus dicit, ideo potius fuisset faciendum, nisi emendatus fuisset Petrus\apnota{26}{fehlt in M.}, ut Petrus pro nullo haberetur, nedum pro summo pontifice. Peccatum mortale nocet personae propriae, sed haeresis nocet personae communi et\apnota{28}{personae communi et fehlt in M.} communitati. Ideo gratias ago domino D., quod saltem hoc didici ex hac disputatione, quod aedificare supra petram significet\apnota{29}{M U super petram significat.} subesse praelato, sive bono, sive malo. Nam hanc significationem aedificationis et petrae fateor prius ignoratam\apnota{31}{M me ignorasse.}, quia dixit: etiam si cecidisset Petrus mortaliter, adhuc mansisset petra, id est primus\apnota{32}{M primus praelatus, U primus et praelatus.}.

Quod dicit, olim episcopos non fuisse confirmatos, nisi praestita scheda\apnota{33}{M ceda.} fidei, admitto, sum contentus, sed\apnota{34}{M U verum.} nihil ad ius divinum, nihil enim derogo pontifici, quidquid tribuitur ei.

Post haec adduxit Nicolaum Cusanum et excusavit Anacletum s. virum, quod non fuerit ei necessaria scientia tanta\apnota{37}{U tanta scientia.}, tam sancto viro. Fateor. Verumtamen summo pontifici, sive sancto, sive non sancto, turpe est evangelia ignorare, cum sit pastor ovium et doctor evangelii. Cum\apnota{39}{M dum.} autem Cephe interpretetur per Johannem evangelistam ca. 1. [v. 42] Petrus\apnota{40}{M Petrum, fehlt in U.} dicentem: Tu es Simon, tu [303] vocabere Cephe\apnota{1}{M cephas.}, quod interpretatur Petrus, quando ergo filius tonitrui sic tonat, qui novit\apnota{2}{P novat.} omnium linguas, merito debet tacere vanitas terreni funii, sive sit Nicolaus Cusanus, Chaldaeus vel Hebraeus. Non quod reprehendam egregium D. D., qui dixerit Petrum idem valere, sicut caput domus Chaldaice, auctore suo Cusano, sed quod laudem imperterritum disputatorem, qui audeat se committere in harenam theologicam\apnota{6a}{U theologorum.} nec tantum prius in evangeliis\apnota{6b}{U evangelio.} studuerit, ut\apnota{7a}{M quod.} primum caput evangelii Johannis legerit\apnota{7b}{Johannis fehlt in M; U in Johanne legit Evangelista.}. Similiter non habet\apnota{7c}{M habeo.} in libro suo, quid est Cephe? quod ignosco, legat tamen caput primum I. ad Corinth. [v. 12]\apnota{8}{M U cap. 1 ad Cor. 1 (Die letzte Zahl fehlt in M).} et inveniet.

Item ultimo de articulis Bohemicis\apnota{10}{U [item fehlt] Ultimum de articulis Bohe: [W. A. Bohemorum]} non vult iudicare, licet nunquam desistat\apnota{11a}{M dicit [Mc a. R. licet] nunquam desistere.} me sugillare, et ego de eis transeo\apnota{11b}{M ei conscio [statt de eis transeo].}, nihil dicens, nisi quod Gregorium Ariminensem dist. XXVIII. reprobatum ab egregio D. D. ego approbo. Est enim totus aliud nihil quam Augustinus et divina scriptura, resistens quidem omnibus doctoribus scholasticis, tum\apnota{14}{M eum, Mc tum.} maxime Aristoteli, sed nondum ab ullo confutatus\apnota{15}{U confutatus est.}.

Consentio cum domino D., quod conciliorum statuta in his\apnota{16}{U iis.}, quae sunt fidei, sunt omni modo amplectenda, hoc solum mihi reservo, quod etre servandum est, concilium aliquando errasse et posse\apnota{18}{U aliquando posse.} errare, praesertim in his\apnota{19a}{U iis.}, quae non sunt fidei, nec habet concilium auctoritatem\apnota{19b}{fehlt in M.} novorum articulorum condendorum in fide, alioquin tot tandem habebimus articulos, quot\apnota{20}{M quod.} hominum opiniones.

Quod in fine s. Franciscum induxit, debere nos summo pontifici obedire, miror contra quem dixerit\apnota{23a}{M U dixit.}. Item transeo istas criminationes\apnota{23b}{M [item fehlt] taceo istam criminationem.} de fratribus\apnota{24a}{fehlt in M.} mendicantibus, miser iste habitus etiam cruciat D. doctorem\apnota{24b}{U miser [M nisi] iste etiam habitus cruciat egregium d. doct.}. Ego pro mea sententia pronuncio, velle me, nullum esse\apnota{25}{M U vellem me esse nullum.} ordinem mendicantium. Haec\apnota{26}{U Ecce [M etc.?] haec.} habui, quae responderem ad obiecta egregii D. D. super propositione 13. Quare nunc finita responsione restat, ut et ego moveam adversus egregium D. D., quae me movent\apnota{28}{quae me movent fehlt in M.}.

Hic dom. doctor Martinus ob partes opponentis ad se ex pacto, ut dicebat, devolutas ad opponendum aggredi voluit. Sed dom. doctor Eccius praetendebat aliqua inter solvendum iam opponentis more dicta, quae primum vellet diluere iuxta protestationem super hoc factam. Nihilominus permisit, ut procedat iam oppositio dom. doctoris Martini. Coepit ergo\apnota{34}{in U lautet der Bericht über den Zwischenfall: Hic intonuit Eccius, se obiectis velle respondere, quasi et ipse respondens fuisset. Respondit Martinus, hoc pacto nunquam finiendam disputationem, accusans Eccianam petulantiam, quae solet singula etiam verba ancupari. Contigit prorsus virgula divina, ut clementissimus Dux Georgius interesset, cuius nutu et pio sceptro concessum est Martino, impetere Eccianam conclusionem, quae contraria est huic iam discussae. In M fehlt der Bericht ganz.}

D. D. Martinus\apnota{35}{M Martinus opponens Egregie domine doctor.}.

Contra sententiam egregii D. D. videtur id militare, quod satis fortiter transiliit in confutando, et scribitur 1. ad Corinth. 3 [v. 5], ubi Paulus detrahens [304] personam omnibus apostolis, dicit: Quid est Apollo? quid Paulus? quid Cephe\apnota{2}{M Scephe.}? ministri, per quos credidistis. Ubi Paulus prorsus vult tollere occasionem contentionis\apnota{3}{M contentionem, Mc contentionis.} et schismatis; negat ullum ecclesiam ad ullum apostolorum pertinere, sed esse omnia omnibus communia. Unde in fine tertii capitis [3. v. 21 ff.] dicit: Omnia vestra, sive Cephe, sive Apollo, sive Paulus\apnota{5}{M sive Pauli (so!) sive Apollo.}, vos autem Christi.

Nec valet evasio, qua dicitur, eos non de principalitate, sed singularitate personarum contendisse\apnota{7}{M eos contendisse.}, cum manifestus textus sit, quod de dignitatibus personarum contenderunt, quod alii praeferebant omnibus Petrum, alii Paulum, alii Apollo, alii vero solum Christum. Confirmatur per illud Gal. 1. et\apnota{9}{M Gal. 1. 2., U Gala. 2.} 2., ubi simili contentione seducti Galatae propter primitatem Petri commendatam Paulum et doctrinam eius velut indigniorem reliquerunt. Contra quae longo textu probat, nihil ad rem pertinere Petri maiestatem vel aliorum apostolorum dicens, se neque ab homine neque per hominem esse missum ad eos\apnota{13}{M a deo (statt ad eos).}, nec vidisse quidem Petrum nec didicisse ab eo\apnota{14}{M a Petro.}, sed omnia sine Petro habuisse et tradidisse. Si ergo auctoritas Petri fuit\apnota{15a}{M fuerit, Mc fuit.} necessaria ut\apnota{15b}{U et.} ius divinum, erit Paulus in hoc loco manifeste impius et blasphemus, ut qui nolit etiam a deo per hominem mitti videri, et prorsus\apnota{17}{M prorsum.} reiicit auctoritatem Petri. Tertio et\apnota{18}{fehlt in M.} infra clarius, ubi dicit: Mihi qui videbantur, nihil contulerunt, et quales aliquando fuerint, nihil mea refert, deus enim personam hominis non accipit. Ecce hic clare dicit, quod qualitas Petri et aliorum apostolorum nihil sua referat, quod esset impiissimum dicere, si qualitas Petri iure divino fuisset servanda. Eadem impietate diceret: Deus personam hominis non accipit, cum ius divinum et ea, quae sunt iuris divini, etiam sub aeternae maledictionis poena praecipiat. Quare videtur primatus iste et maiestas, vel quocunque nomine censetur Petri persona seu qualitas, non statui iure divino.

D. D. Eccius respondere voluit iam ad aerem repetitis obiectis. Tum iussu dominorum per famulum publicum intercepta disputatio fuit et in horam secundam pomeridianam suspensa.

Hora secunda pomeridiana eiusdem diei\apnota{26}{Zwischenbemerkung und Zeitangabe fehlen in M U.}.

D. D. Eccius.

Dum reverendus pater contra me inducit tres b. Pauli auctoritates, dudum ab adversariis sedis apostolicae ventilatas et per apostolicae sedis defensores enodatas\apnota{29}{M ab adversariis sedis [gestrichen, U sedes] apostolicae sedis defensores enodatos [offenbar in Unordnung].}, facile ei respondeo, annectendo malos fucos, quibus dogma suum nititur palliare et hominibus os et oculos oblinire. Dico primum\apnota{30}{M primo.}, me non transiliisse, sicut imponit, locum apostoli 1. Cor. 3 [v. 5.], quamvis ipse false adduxerit cum Paulo et Apollo Cepham, quia textus ab initio non habet, sed in fine meminit Cephae, 'sed

statim' glossa interlinearis fatetur eum maiorem omnibus. Non ergo suo capite, sed sanctorum accipiat sacram scripturam reverendus pater.

Et dico solutionem meam datam\apnota{35}{meam fehlt in M U, datam fehlt in M.} fuisse optimam et infringibilem\apnota{36}{M infallibilem.}, et fucus reverendi patris, quo simplicibus illudit, nihil facit, testis est b. Hieronymus lib. 1 contra Jovinianum\apnota{37}{P Jovinianum 1.}, col. 18. mihi, ubi expresse dicit, supra Petrum aedificatam ecclesiam. Ibi fortiter transiliit\apnota{38}{M resiliit.} reverendus pater et nihil dixit\apnota{39a}{U dicit.}. Et infra\apnota{39b}{et infra fehlt in M.}: ut capite constituto schismatis tolleretur [305] occasio. Ergo quod apostolus Paulus voluit cavere\apnota{1a}{M vitare.} schismata et\apnota{1b}{fehlt in M U.} contentiones, per hoc primatum capitis non abstulit\apnota{2}{M tollit.}, quod schismata 'melius tolli non possunt'\apnota{3a}{Mc possint.}, nisi per caput. Unde nec verbum facit\apnota{3b}{M dicit.} nec apex ad propositum de negatione\apnota{4a}{de negatione fehlt in M.} primatus. Quare false auspicatus est hodie, me vix\apnota{4b}{M vix in me.} duos adduxisse s.\apnota{5a}{fehlt in M.} patres et eos ambiguos, cum certos et indubitatos adduxerim\apnota{5b}{M induxerim.}, Augustinum, Hieronymum\apnota{6a}{M U Ambrosium, Hieronymum.}, Chrysostomum, s.\apnota{6b}{fehlt in U.} Leonem, Bernhardum etc.\apnota{6c}{M et Bernhardum (statt Bernh. etc.)}. Quod autem laudat me amplecti ambas Augustini sententias libro 1. retract. cap. 21, et affirmat\apnota{8a}{W. A. verbessert affirmat et, wahrscheinlicher ist, dass mit facere der Satz abbricht.} se ipsum facere\apnota{8b}{U facere idem, Uv = P.}. . . . .

Negat D. D. Martinus se id fecisse etc., nisi lapsu linguae et minus praemeditate etc.\apnota{8c}{Die Zwischenbemerkung fehlt in M; U. a. R. Negavit d. Martinus id se fecisse, ut supra protestatus est.}.

Quod vero secundo loco inducit Paulum ad Galatas, quasi ille scripserit\apnota{9}{M U scripsit, Uv = P.} contra Petrum, volens Galatas confirmare in fide, et toto ibi processu ostendat se aequalem fuisse Petro et aliis, quia non ab homine neque per hominem acceperit etc.\apnota{12}{M acceperat; etc. fehlt in M U.}, dico, quod, si reverendus pater tam diligens esset veritatis inquisitor, sicut promittit, didicisset vera fundamenta et intelligentiam scripturarum, ut non opus esset ita iactanter eum in tanto consessu dicere, se unum velle resistere mille\apnota{15a}{M millis.}. Ita se\apnota{15b}{fehlt in U.} habet veritas et negotium. In apostolatu et ordinis dignitate apostoli omnes fuerunt aequales, hoc sacra scriptura, hoc s. patres voluerunt. Paulum\apnota{17a}{U palam Uv. = P.} dicit Hieronymus non reprehendisse Petrum, si sciret se imparem\apnota{17b}{M esse imparem.}. Ita Anacletus cap. In novo dixit\apnota{18a}{U dicit.} apostolos aequali\apnota{18b}{U inaequali.} potestate dotatos. Ita Cyprianus, ita et alii. Cum ergo

Paulus fuerit\apnota{19}{U affuerit.} designatus apostolus a Christo, ingenue dixit, se non ab homine neque per hominem accepisse, et quod Jacobus et Petrus sibi nihil contulerunt, omnia enim ista ad apostolatum referenda sunt.

Aliud vero est de potestate regiminis et administrationis, ubi concors est s. patrum sententia, sacrum concilium\apnota{23}{M sacrum consilium, U et concilium.} contra Johannem Huss, ad ipsum Petrum huiusmodi primatum pertinuisse. Non ergo recte utitur sacrae scripturae testimoniis, sed eas aliorsum\apnota{25a}{W. A. verbessert ea; M aliorsim.} trahit pro suo arbitrio, contra quam spiritus sanctus\apnota{25b}{M U sanctus spiritus.} efflagitat. Quod enim hodie ponderabat Paulum dixisse, nihil referre ad eum, quales fuerint, non hoc quod volebat reverendus pater, sed hoc significabat\apnota{27}{U significat, Uv = P.}, Petrum et Jacobum et ceteros\apnota{28}{et ceteros fehlt in M.} fuisse idiotas et minus peritos, antequam cum Domino ambulassent, et sic deus personam non accipit.

Sed inquiet\apnota{29}{M hinc (statt inquiet).}, ut est mordaculus\apnota{30}{U cordaculus, W. A. audaculus.}, hoc est commentum Eccianum. Legat Ambrosium eodem loco, qui sententiam Pauli non corrumpens ingenue testatur\apnota{31}{M fatetur.}, ecclesiam fundatam super Petrum\apnota{32}{M super petram, U supra petram.}. Quare nisi firmiora adduxerit, rogo, ab instituto desistat. Unde hodie, dum dixissem articulum esse Hussiticum, obedientiam ecclesiasticam non esse expressam in\apnota{34}{fehlt in M U.} auctoritate scripturae sacrae, respondit brevibus, articulum non esse haereticum. Mallem punctim agere 'cum reverendo patre' et quaerere\apnota{36a}{U dicere.}: Iste est articulus damnatus, et ipse tenet eum, itaque articulum damnatum et reprobatum\apnota{36b}{M U tenet articulum. P reprobum.}.

Item quaero ab eo\apnota{37}{ab eo fehlt in U, in Uv nachgetragen.}, si articulus non est haereticus, vel ergo\apnota{38}{fehlt in M U.} est temerarius vel seditiosus vel offensivus piarum aurium; utcunque [306] dederit ipse pater, erit vel temerarius vel 'seditiosus' vel in haeresi suspectus\apnota{1}{U de haeresi suspectus [vel fehlt]; M vel haeresis suspectus.}, aut pias offendit aures\apnota{2}{Mc vel piarum aurium offensivus [Lücke in M ausgefüllt], U [Uv aut] pias offen: (W. A. offendens) aures.} contra theologicum 'institutum'.

Mirum, quo colore pavonico ornet\apnota{3}{M Miror, quo [Mc colore] pavonico colore ornat.} verba Christi Matth. 16. Quia Christus omnes interrogaverit\apnota{4a}{U interrogaverat.} discipulos et ita Petrus\apnota{4b}{M Christus, Mc Petrus.} pro omnibus responderit, quis hoc negat? Sed de sua farina addidit\apnota{5a}{M varina hoc addit.}, propterea Petro\apnota{5b}{U Petrum, W. A. ad Petrum.} 'pro se\apnota{6a}{pro se fehlt; in M.} et' aliis apostolis dictum: Tu es Petrus, et super hanc petram aedificabo\apnota{6b}{fehlt in M U.} etc., quod tanquam commentum Lutheranum\apnota{7a}{P Luterianum} non recipio\apnota{7b}{M accipio.}, cum sanctis patribus adversetur\apnota{7c}{M est adversum adversetur fehlt in U, in Uv nachgetragen.}.

De Chrysostomo, mirum, quam faciat\apnota{8a}{M U facit.} iniuriam s. patri\apnota{8b}{M patrem, Uv streicht quam bis patri, a. R. cum sanctis patribus adversetur (vgl. 7), nachher gestrichen.}, quasi legitime et secundum decreta concilii\apnota{9}{M consilii.} Niceni fuerit destitutus, cum exploratum sit ex historiis\apnota{10a}{M ab historia, U ab historiis.} post Eusebium 'Antiochenum', schismaticum\apnota{10b}{U credo schismaticum.}, cum aliis Rom. pontifici rebellibus Christianos episcopos fuisse expulsos, qui a prima sede refugium quaerentes ab impiis schismaticis repulsam\apnota{12}{U repulsum.} passi sunt, usque ad 'tempora Theodosii Christiani\apnota{13a}{U christianissimi.} imperatoris. Ubi dum\apnota{13b}{U tunc.} schismatici sedem triginta annis\apnota{14}{U 30. annis sedem.} occupassent, s. Gregorius Nazianzenus Constantinopoli 'episcopatum' suscepit\apnota{15a}{U accepit.}, non 'ambitione', sed ne alius schismaticus succederet\apnota{15b}{U alii schismatici succederent.}.

'Sanctus vero' Chrysostomus iniuria per Epiphanium pulsus, Arsatio schismatico substituto Rom. pontificis auctoritate regressus est. Nec valet, quod reverendus pater tunc\apnota{18}{M hunc.} Rom. pontificem et sanctos sibi adhaerentes calumniatur, quasi de facto, non de iure processerint contra illos schismaticos, quod nudis verbis suis\apnota{20a}{M U suis verbis.} ait, et quod\apnota{20b}{fehlt in M U.} tot sancti, etiam a d. patre commendati, Romano pontifici non consensissent\apnota{21}{U tot sanctos . . . commendatos . . . consensisse, Uv = P.}, unde sine causa s. illos patres ante mille annos viventes reverendus pater iam incipit ambitionis insimulare, a quo vitio a\apnota{23}{fehlt in M.} tot seculis ab omni ecclesia habiti sunt immunes.

Ita\apnota{24}{M Ideo.} captiosum est et simplicibus illudit, quod b. Gregorius 'noluerit se oecumenicum' et universalem pastorem appellare\apnota{25}{M appellari.}, nam Platynam tanti fecit pridie, historiae, 'ut dicebat', secutus veritatem. Jam\apnota{26}{M nam.} Platynam pro suo arbitrio reiicit, quod ubique existimat sibi licere.

Decretum allegat pro se; si allegetur, respuit, et dum textus apertissimos\apnota{28}{U apertissimus.} etiam Leonis cap.: Ita dominus 19. dist. cum similibus habeat sibi contrarios, glossam acceptat\apnota{30a}{M habet sibi contrarias glossas [Mc contrariam glossam] accipit [Mc acceptat].}, sicut 'proci Penelopes', ut per iocum\apnota{30b}{M iocam.} dicam, cum domina potiri\apnota{31}{M potire.} non poterant, ancillas sequebantur. Certum est, sanctos quoque pontifices scripsisse se universales episcopos, ut Sixtus, ut Victor, et ut in Chalcedonensi concilio\apnota{33a}{M consilio.} acclamatum fuit Leoni\apnota{33b}{fehlt in M.}: Vivat leo sanctissimus oecumenicus patriarcha, quamvis humilitatis et Christi formam sequentes pontifices ab eo nomine sibi temperarint\apnota{35a}{M U temperaverunt.}, melius se servos servorum fatentes\apnota{35b}{M servus servorum 'praestare'.}, et ne crederentur velle se tollere episcoporum auctoritatem\apnota{36}{P auctoritate.}, ut dudum ad illa 'respondit' s. Thomas lib. de impugnatione fidei.

De epistola Augustini facit vobis dominis auditoribus negotium, ut verba legatis, verba audita fuere: Innocentius ante centum et mille annos dixit esse veterem morem consulendi summum pontificem in his\apnota{40}{U veterum more consulendum pontificem in iis.}, quae sunt fidei [307]. Ita dictum contra Manichaeos voluit reverendus pater transplantare in hortum\apnota{2a}{M ortum.} suum. At hoc non faciet\apnota{2b}{M faciat.}, si epistolam Augustini legat, in qua non seriem universitatis episcoporum, sed Romanorum pontificum enumerat ab ipso Petro usque ad tempora sua.

Solutionem per me\apnota{5}{Solutionem per me in M doppelt, das erste Mal gestrichen.} datam de viginti annis dicit per me non probatam, acutus est disputator 'reverendus pater, hoc nondum scivi debere\apnota{6}{M probare, Mc debere.} respondentem probare solutionem suam, semper putavi, respondentem debere\apnota{7}{Mc U suam solutionem semper arbitratus sum respondentem [debere fehlt].} offerre solutionem opponenti improbandam. Debilis fuit Hercules suus', si alieno Marte nititur\apnota{9}{M utitur.}.

Quod vero tanta laborat impudentia, ut audeat illustrissimis dominationibus et reverendis paternitatibus vestris suggerere, quasi primum Johannis caput\apnota{12}{M U caput Johannis.} non legerim, insignis profecto temeritas, cum puer nondum decennis, demptis prophetis, bibliam totam\apnota{13}{M totam bibliam.} legerim, sed hoc nihil ad rem, quantum quisque legerit.

Dixi\apnota{14a}{M dixit.} Anacletum defendendo non\apnota{14b}{M defendendum [non fehlt].} 'usque adeo inscitam fuisse interpretationem'\apnota{15a}{U interpretationem eius.} Cephas, id est caput, non ignorans, quod κεφαλη Graece caput Latine\apnota{16a}{fehlt in U.} significet\apnota{16b}{M U significat.}, sed quia apud\apnota{16c}{M et apud.} Johannem, quem heri vidi, apud Augustinum et glossam ordin.\apnota{17}{et glossam ordin. fehlt in M.} Cephe interpretatur Petrus, Petrus autem, Cusano auctore, etiam significat caput domus, quae omnia in risum trudit\apnota{18}{M tradit.} reverendus pater, de quo iudicent ordinandi iudices.

At illud praestantius\apnota{20}{M quaestantius.}, quod auctoritatem, »Pasce oves« tantis ambagibus\apnota{21a}{M imbalgibus, Mc ambalgibus (so!)} involvit, solum ut homines videntes faciat non videre, mitto\apnota{21b}{M mitto, Mc omitto.}, quod tam ridiculam attulit expositionem a se fictam. Sed ad alium eius sensum veniam\apnota{22}{M venio.}, ubi multa dixit, at ego non singularitati studens, quae est mater errorum, s. patrum secutus sententiam et s. matris ecclesiae accepi verbum »Pasce«, quod\apnota{25}{M pascere quia.} eum pastorem ibi constituerit et primatum dederit universalis ecclesiae. Ita Gregorius, ita Chrysostomus intellexerunt, ita Ambrosius Lucae ulti. heri per me non allegatus\apnota{27a}{M heri non allegatis.}. Ad quos pulchre 'praetergrediens'\apnota{27b}{P praegrediens, M Lücke, Mc progrediens.} nihil dixit reverendus 'pater, verum veram attulit\apnota{28a}{Mc varam attulit, U [veram fehlt] attuli, Uv attulit.} Augustini sententiam, cuiusmodi debeat esse praelatus vel\apnota{28b}{praelatus vel fehlt in U, in Uv nachgetragen.} pastor, quem et ego admitto, et credo quemlibet fidelem\apnota{29a}{U fidelium.} orare, ut et\apnota{29b}{M et et.} summus pontifex, 'quilibet princeps', quilibet praelatus quaereret honorem dei et subditorum salutem, quod haec\apnota{31}{U hoc, Uv = P, aber gestrichen.} sit verissima pastura, per quam, sicut hic principantur in terris, etiam elevantur, ut regnent\apnota{32}{M regnant.} in caelis. Quare Augustinus non adversa, non pugnantia, sed talia dicit, ex quibus praelatus intelligere debet, qualis sua sit institutio.

Nec valet, quod reverendus pater adduxit sub conditione ei commissas oves, cum nulla ibi adhibeatur conditio, nisi sua paternitas aliud habeat\apnota{36a}{M habet.} evangelium quam ego. At qualis esse debeat\apnota{36b}{M U debeat esse.} pastor, significans Christus dilectionem et amorem dei a Petro quaerit, et sic eum pastorem totius ecclesiae fecit, sicut canit s. mater\apnota{38}{M mater sancta.} ecclesia, et reverendus [308] pater orat\apnota{1a}{U orat vel legit.}: 'Tu' caput ecclesiae, pastor ovium\apnota{1b}{M U ovium etc.}. Quare iure divino habemus hoc loco, debere praelatum diligere deum, habemus Petrum constitutum pastorem ovium, sicuti s. Gregorius, s. Ambrosius et Chrysostomus etiam Graecus\apnota{4a}{M Graecus, Mc Graecos.} intellexerunt.

Non refert praetenuis\apnota{4b}{M praetenuit. Mc quae tenuit.} ratiuncula, in diversum a reverendo patre inducta\apnota{5a}{M inductum, Mc inducta.}, Petrum non omnes pavisse; quod\apnota{5b}{M quo, Mc quod.} nostro nomine agitur, et nos fecisse dicimur, sicut Clemens\apnota{6}{U Clemens 21. ult. dist.} testatur, Petrum ordinasse episcopos, presbyteros\apnota{7}{U presbyteros, episcopos.} et diaconos, qui per universum orbem terrarum evangelium seminarent, ut legitur can. In illis, 80. vel 81. dist.\apnota{8}{M can. in illis 81. dist. vel 80. dist., a. R. in c. in illis 81. dist, U in can. illis 80. vel 81. dist; richtig: 80. dist.} Unde missus a principali, sicut fuit s. Paulus, non cogitur mitti ab ipsius vicario, tamen Paulus\apnota{10a}{M U Petrus, Mc Paulus.} ascendit cum Tito, ad Gala. 2., et contulit cum Petro evangelium\apnota{10b}{M evangelia.}, qui dederunt ei dexteram societatis.

Porro, quod cudit in me glossam nunquam a me cogitatam, quasi dixerim: Aedificare supra\apnota{13}{M super.} petram sit subesse Rom. pontifici, o somnia, o verborum portenta! nunquam hoc fuit mentis meae. Christus enim aedificat ecclesiam, non hi qui subsunt, subditi enim aedificatae ecclesiae subesse\apnota{15}{U esse, Uv subesse, vor aedificatae.} debent.

De Gregorio Ariminensi, cuius sententia mihi displicet, ei placet, nihil iam ad propositum. Nolo ergo in ista diverticula descendere. Sed

ut accepimus\apnota{18a}{P occepimus, M suscepimus.}, defendamus articulum Constantiensis concilii\apnota{18b}{U consilii.} esse verum, et quicunque huic contradixerit, errare et veritati et honori concilii\apnota{19}{U consilii.} contradicere.

D. D. Martinus.

Antequam ad auctoritates a me oppositas veniam, breviter discurram, quae egregius D. D.\apnota{22a}{M dom. egr. doctor.} de hesterna et hodierna responsione mea\apnota{22b}{M meus, Mc mea.} commiscuit, ubi inter cetera, tanquam consumptis omnibus armis\apnota{23}{M armis omnibus.} desperans, in unum articulum Constantiensis concilii extremum fiduciae praesidium collocat, et odiosissime repetit et inculcat, quasi non abunde satis supra ad ipsum sit responsum. Dicam tamen adhuc semel.

Quando articulus de operibus hominum inter damnatos recensetur et a\apnota{28}{fehlt in M.} s. Paulo et Augustino asseritur, deinde per Gregorium Ariminensem, per omnes universitates usque hodie defenditur, non movebor odiosissimis inculcationibus huius articuli, donec egregius D. D. probaverit, concilium non posse errare, non errasse aut etiam non errare, cum concilium facere ius divinum non possit, ex eo, quod natura sua non est ius divinum, ideo neque haereticum est, nisi quod contra ius divinum est. Haec pro isto articulo\apnota{33}{haec pro isto articulo fehlt in M.}.

Quod Chrysostomum legitime\apnota{34}{M a me legitime.} destitutum carpit, non intellexit me, cum clare dixerim hodie, potestatem ordinandi et destituendi episcopi fuisse legitimam in Epiphanio Cypr.\apnota{36a}{P M U in beato Cypriano.} et Theophilo, quamquam an bene tunc usi sint vel male\apnota{36b}{M an tunc bene vel male usi sint, U an bene tum usi sint an male.} hac potestate legitima, nihil disputaverim\apnota{37}{P potestare; M disputavi, U disputaverit.}. Sufficit mihi, quod citra [309] auctoritatem Romani pontificis tot sancti et celeberrimi viri Graeciae\apnota{1}{M Graeci.} catholicissimi, non schismatici, potestatem habuerunt ordinandi et destituendi episcopos.

Elevat etiam egregie, quod antiquos sanctos insimulaverim ambitionis, nec attendit, quod dixi etiam in apostolis hanc fragilitatem fuisse praesente Christo\apnota{6}{U Christiano, Uv = P.}.

Quod Platynam nunc ita\apnota{7}{M U ita nunc.} magnifacit adversus sanctum Gregorium, mea auctoritate satis claret cuilibet vel mediocriter ingenioso, quid concludat. Ego enim ex Platyna historias accipio et non amplius. Quod autem idem Platyna rationes quasdam adducit, respuo propter auctoritatem Gregorii contrariam.

Jam illud transeo, quod Sixtus, Victor\apnota{12}{W. A. et Victor.} scripsere se universales pontifices et Leoni acclamatum sit in concilio Chalcedonensi, de quo

hodie satis dictum est. Illud observandum, quod sequentes pontifices dicit egregius D. D. humilitatis formam in Christo secutos\apnota{15}{M Christi formam secutus (so!)} sibi temperasse ab hoc nomine, quod ego aliter intelligere non possum, quam quod causa humilitatis iuri\apnota{16}{U iure.} divino detraxerint, mendaces et inobedientes facti fuerint, si ius divinum est, hoc est praeceptum, fuisse eos universales pontifices.

Innocentii epistolam responsivam\apnota{19a}{U responsivam epistolam.} ad Augustinum repetens dicit\apnota{19b}{M dicit dicit.}, veterem morem fuisse consulendi Romanum pontificem\apnota{20}{M ro: pontificis, U Romanos pontifices.} in his, quae sunt fidei. Optime dicit, et mihi\apnota{21}{fehlt in M U.} placet, praesertim in Latina ecclesia, sed quid fidei et primatui? hoc est, ad quod\apnota{22}{U quid ad.} propositum?

Causatus est etiam, quod exegerim ab eo probationem\apnota{23}{fehlt in M} solutionis de viginti annis\apnota{24a}{M 20 annorum solutionem.}, dupliciter deficit egregius D. D., primum\apnota{24b}{U primo}, quod ipse fuit opponens semiquatuor dies et hac hora solum coepit esse respondens. Deinde non exegi, ut probaret solutionem, sed translatum primatum cum Petro.

Excutiens etiam se\apnota{27a}{U se etiam.} de Cephe, quod Petrus interpretatur apud\apnota{27b}{M ad.} Johannem, adhuc persistit in eo, quod Petrus etiam significet\apnota{28}{M significat etiam.} caput domus, quasi videlicet Johannes Graecis interpretaturus evangelium non Graece, sed Chaldaice sit 'locutus'\apnota{29}{sit in M gestrichen, Mc locutus locutus', das zweite gestrichen.}, ne forte intelligeretur.

Verbum »pasce« sic accipit, ut constitutionem universalis pastoris significari\apnota{32}{M U significare.} velit, et ad hoc Gregorium, Ambrosium, Chrysostomum heri non allegatos adduxit\apnota{33a}{fehlt in P M U, aus der ed. Viteb. zu ergänzen, vgl. W. A.}, ad quos ego nihil dixerim. Dicam ergo nunc, si fui\apnota{33b}{M fuero.} oblitus, et simul oppono regulam Augustini\apnota{34}{M U divi Augustini.}: Quod omnium scriptorum dicta iudicanda sunt per divinam scripturam, cuius est maior auctoritas, quam totius generis humani capacitas; non quod damnem illustrissimorum patrum\apnota{36}{fehlt in P.} sententiam, sed quod\apnota{37a}{M qui.} eos imiter, qui proxime scripturae\apnota{37b}{M ad scripturas.} accedunt, et prae omnibus, si aperta fuerit scriptura, eam ipsam amplector. Cum ergo clarum sit, pascere vi vocabuli et usu scripturae non significare nisi docere, curare [310] oves sicut Petrus accipit ult.: Pascite eum\apnota{1}{U Sic Petrus accipit ult. Pascite [eum fehlt].}, qui in vobis est, gregem [1. Pe. 5 v. 2]. Non oportet vocabulum aequivocare et simplicem sensum sine necessitate relinquere.

Dissolvit etiam, quod dixi esse conditionem in verbo Christi, Si diligis me, quando ego arbitror\apnota{4}{M U arbitrer.} non modo conditionem, sed et pactum

esse inter Christum et Petrum, dum ille dicit: Diligis me? et hic respondet: Amo te.

Omitto dialecticam, in qua exercitatissimus est D. D., et coniunctionem\apnota{6}{M conditionem.} conditionalem »si«. Vellem audire, quid diceret ad decretalem de elect. c.: Significasti, ubi Paschalis pontifex manifeste dicit hic conditionem esse, si saltem adeo valet\apnota{9}{M valet, verbessert valeat.} auctoritas Rom. pontificis, sicut hucusque contendit, dicit enim\apnota{10a}{U enim ad sensum}. [Citaturus est ex libro hanc decretalem.]\apnota{10b}{Die eingeklammerten Worte in P a. R., M statt dessen Si conscientiarum cognitor et factor [aus d. Text d. Decr.].}

Ad rationem meam, quam\apnota{11a}{Ad fehlt in M, U at; quam fehlt in M U.} sic confutavit, quod Petrus non omnes pavit\apnota{11b}{M U paverit.}, quia, inquit, quod nostro nomine agitur, et nos fecisse dicimur. Quanquam non credo, per hoc dominum D. voluisse satisfactum rationi meae, cum dicere non possit, Paulum nomine Petri praedicasse, qui\apnota{14}{M si, Mc qui.} contra ad Rom. 1. [v. 1ff.] apostolum se dicit pro nomine Christi, quod ex eo coniicio, quia dicit Petrum, auctore Clemente, ordinasse episcopos, presbyteros\apnota{16}{M presbyteros episcopos.}, diaconos, qui per totum orbem evangelium seminarent, inter quos Paulum non numerat\apnota{17}{M enumerat.}.

Sed hoc videtur ad rem facere, quod missus a principali non cogitur mitti ab ipsius vicario, hoc admitto, si fuerit primum probatum, Petrum esse vicarium. Nunc vero non minus Paulus est vicarius Christi\apnota{20}{fehlt in U.} quam Petrus, ut iam dictum est, quod nomine Christi sit apostolus, et sic tot vicarii, quot apostoli.

Ad\apnota{22}{M U Nunc ad.} propositum.

Ad auctoritatem Pauli 1. Cor. 3. [v. 5] primum dicit, Cephe non esse de textu ab initio, sed in fine memorari, ubi glossa fateatur, Petrum maiorem omnibus, leviculum hoc est, transeat\apnota{25}{U transeo.}, cum in principio primi capituli [v. 12] utique ponatur\apnota{26}{P ponitur.} Cephe.

Secundo, dicit aliud nihil, quam quod nec verbum nec apex\apnota{27}{M add de sit, gestrichen.} de negatione primatus sit in textu, et sic vult solutam istam auctoritatem. Adiecit tamen auctoritatem Hieronymi, Augustini, Ambrosii, Chrysostomi, Bernhardi, Leonis\apnota{30}{Ambrosii und Bernhardi fehlt in M; M U Leonis etc.} dicentium, Petrum esse petram. Contra quae 'dico et' replico, hoc secundum dictum nihil est ad propositum et supra satis disputatum, an Petrus sit petra. Similiter et illud non solvit, quod dicit, nullum verbum, nullum apicem esse de negatione primatus\apnota{33}{de negatione primatus fehlt in M, Mc füllt die Lücke durch 'primatum'.}, quod si etiam Hieronymus statuat\apnota{34a}{U constituat.} caput, ut occasio schismatis tollatur,

satis robuste 'contra pugnat'\apnota{34b}{contra pugnat fehlt in M, Mc 'oppugnat'} textus apertus, qui tollit caput, ut occasio schismatis tollatur. Dicit enim: Quid Paulus? Quid Apollo? volens solum Christum esse\apnota{36}{M est.} caput, quo vere tollitur omne schisma. Et sic\apnota{37a}{Et sic fehlt in U.} non est satisfactum hac responsione\apnota{37b}{U sic ad responsionem.}.

[311] Ad secundam auctoritatem Gala. 2. dicit quidem apostolos omnes fuisse\apnota{1}{secundam und fuisse fehlt in U; letzteres in Uv nachgetragen, M omnes apostolos [U apostolus].} aequales, quod et Hieronymus et can. In novo testantur\apnota{2}{M ut Hieronymus . . . attestantur.}. Sed aliud esse potestatem regiminis et administrationis, his tribus verbis solvere tantam auctoritatem, est mihi satis mirabile. Sed et non\apnota{4}{M sed non, U sed et nec.} intelligo, quid velit, quomodo distinguat regimen\apnota{5}{P rem.} et administrationem.

Quod autem dixit, Paulum 'cum Tito' ascendisse ad Petrum, ut conferret\apnota{6a}{U ascendisse cum Tito . . . ut conferat.} evangelium cum eis, clarum est ex\apnota{6b}{fehlt in U.} textu praecedenti\apnota{6c}{M excedenti, Mc 'praecedenti'.} et sequenti, ut et Hieronymus ipse exponit\apnota{7}{ut et H. i. exponit fehlt in M. [Lücke].}, non hoc fecisse Paulum, ut consuleret Petrum super evangelio\apnota{8a}{U evangelium.}, quod iam plus 14. annis\apnota{8b}{M annos} praedicarat et non 'ab homine se\apnota{9a}{fehlt in U.} accepisse iactaverat, sed ut pseudoapostolis\apnota{9b}{pseudapostolorum.} os obstrueret, Galatas confirmaret, eadem esse, quae ipse et ceteri apostoli praedicassent.

D. D. Eccius.

Quod reverendus pater articulum Johannis Huss de obedientia ecclesiastica commiscet alteri articulo de operibus humanorum\apnota{13}{M bonorum.} actuum, auctoritate Gregorii Arimin., qui per omnes universitates defendatur, quem adhuc in eo loco in nulla universitate, quamvis in plurimis fuerim, audivi defensum\apnota{15}{M fui . . . defensam.}; petit tamen a me, ut probem, concilium\apnota{16}{M U ei concilium (M consilium).} non posse errare, nescio quid sibi velit\apnota{17a}{M velit sibi.} ista petitio, an tacite laudabile et gloriosum\apnota{17b}{M an tacite 'an' laudabile et nobile, a. R. 'an nobile', U an tacite laudabile et [gloriosum fehlt].} Constantiense concilium velit habere suspectum. Hoc dico vobis, reverende pater, si creditis, concilium legitime congregatum errasse aut errare\apnota{19a}{M esse [gestrichen] errasse et errare, U errare et errasse.}, estis mihi sicut ethnicus et publicanus\apnota{19b}{M Lücke, Mc 'haereticus et ethnicus', a. R. 'et publicanus'.}. Quid sit haereticus, in praesentia non discutiamus\apnota{20}{M U discutiam.}.

De Chrysostomo, non vult eius expulsionem iustificasse\apnota{21}{M U se iustificasse.}, sed ordinationem episcoporum a Niceno\apnota{22}{M Ciceno, Mc 'Niceno'.} concilio institutam laudasse. Placet mihi pro s. Chrysostomi reverentia, et semper Rom. pontifices in diversis

nationibus diversos permiserunt modos episcoporum eligendorum, tum\apnota{24}{M episcopos [Mc episcoporum] eligendos tamen, U episcopos eligendi tamen.} ex Gregorio Nazianzeno et Chrysostomo restituto satis liquet, intervenisse auctoritatem Rom. pontificis, quod s. Athanasius Alexandrinus, s. Paulus et alii orientalis ecclesiae episcopi testantur, qui se ad Rom. pontificem pro restitutione contulerunt\apnota{28}{M se contulerunt (so!)}. Quod in apostolis quoque quid humanae fragilitatis fuisse assignat\apnota{29}{M signat.}, novimus. Sed dum\apnota{30a}{U cum.} publico fungerentur officio, absque indicio\apnota{30b}{U iudicio. M officio fungerentur absque initio.} et causa tam s. viris non est detrahendum.

De Platyna dicit, se recipere eius historiae auctoritatem\apnota{32}{eius historiae auctoritatem fehlt in M, Mc eius historiam.}, et ego neque ex Platyna rationes adduxi\apnota{33}{U auctoritatem ex Platina adduxi.}, sed historiam.

Ad Innocentii epistolam respondet, eam nihil facere ad propositum, quod de 'Latina probet ecclesia, non de alia'\apnota{35}{Mc probat . . . . aliis, U de latina illum probat ecclesia et non de alia.}. Sed quia haec est glossa Lutherana nulla auctoritate fulcita vel persuasione, ea facilitate\apnota{36}{M facilitatem, Mc facilitate.} contemno, qua adducitur.

De viginti annis respondit, me fuisse opponentem iam triduo\apnota{38}{M per tribuum [Mc triduum], iam fehlt.}, deberem utique fuisse\apnota{39}{M debuisse.}, nisi reverendus pater suis cautelis aliquando fecisset ex oppo[312]nente respondentem\apnota{1a}{M ex respondente opponentem, gestrichen, Mc a. R. wie oben.}. Non enim hoc posui\apnota{1b}{M posui hoc, U hoc opposui.} argumentum, sed ipse posuit\apnota{1c}{fehlt in M.}, 'et certe cathedram illam' Petrus non secum assumpsit ad Romam\apnota{2}{secum und ad Romam fehlt in M, U sumpsit, Uv = P.}, sed primatum transtulit.

De verbo »pasce«\apnota{4a}{U = P.} vult dimittere Gregorium, 'Ambrosium', Chrysostomum, quia\apnota{4b}{M remittere [Mc dimittere] Hieronymum, [a. R. Ambrosium] quia.} simplex verbum non debeat aequivocari. Dico\apnota{5}{M debet aequivocari 'dico' dico.}, quis hoc ignorat, verbum pastoris esse nomen officii, honoris et etiam oneris? Unde et praelationis accepit officium et simul pascendi onus\apnota{7}{M oves.}.

De conditione non recipio allegata per reverendum patrem, quia\apnota{8}{M accipio allegata quia.} nec Christus, sicut ipse semper\apnota{9a}{fehlt in M.} inhaeret literae, conditionem coniunxit\apnota{9b}{U iunxit.} pasturae. Nam\apnota{9c}{M non.} sic error confirmaretur Hussiticus, quod non existens in charitate non esset\apnota{10}{fehlt in U.} episcopus, praelatus, pastor\apnota{11}{M praelatus etc., U praelatus pastor etc.}.

De capitulo Significasti, de elect. scio Gersonem et alios doctores ponderare in materia, 'an concilium\apnota{13}{Mc consilium.} sit supra papam'? Sed quia verba decretalis non sunt adducta, nihil possum dicere, sed coniicio

'nolens me huc alligare'\apnota{14}{M Lücke, Mc a. R. volens (so!) me huc (U hoc) allegare.}, Paschalem conditionem accepisse pro qualitate\apnota{15}{U pastalem conditionem accepisse pro qualitatem (Uv qualitate).}.

De Paulo, qui primatum Petri non agnoverit, sed fuerit apostolus Christi: omnes sumus Christi, arbitror, tamen 'primatum'\apnota{17}{M U ut arbitror tamen Petri (Mc U Petro) 'primatum'.} Petro nullibi abstulit regiminis, sed et\apnota{18}{fehlt in M.} moriturus, ut in epistola de 'transitu apostolorum scribitur', Petrum valedixit ut caput ecclesiae, quae epistola an sit b. Dionysii vel minus, incertum habeo.

Paulum dixit vicarium Christi, sed non generalem, nam etiam\apnota{21}{nam etiam fehlt in M, Mc 'ita etiam'.} in apocalypsi [21 v. 14] ponuntur duodecim fundamenta ecclesiae, et tamen fundamentum aliud\apnota{22}{M alium, Mc aliud.}, scilicet principale, nemo ponere\apnota{23}{M aliud ponere.} potest, quam quod positum est, scilicet Christum.

De secundo dixit, non esse ad propositum, an Petrus sit petra, et est totum negotium nostrum. Tamen solutionem 'pridie per me\apnota{26}{Mc a me.} datam', ad quam me retuli, reticuit\apnota{27}{M U ipse reticuit.}, cum expressum sit 'de litera' Paulum 'improbare contentiones'\apnota{28}{Mc increpare contentionem, Ui mprobare contentionem.}. Sed nescio, quo ingenio 'reverendus pater arbitretur' propterea Paulum negasse primatum, ut tolleretur\apnota{29}{M U tolleret.} schisma. Forte non fuit in 'provincia, quae caruit principe'\apnota{30a}{M fuit in principe [Lücke], Mc fuit in provincia quae caruit principatu.}, fuisset enim contrarium expertus, quot\apnota{30b}{P quod.} dissensiones, contentiones et dissidia emergunt, ubi 'non est princeps et\apnota{31}{Mc vel.} caput', inter causas diiudicans et cognoscens\apnota{32}{U agnoscens.}.

Et mirabile videtur ei, si tam paucis verbis, scilicet apostolatus et regiminis tantam\apnota{33}{U tantum, Uv tantam.} velim solvere auctoritatem. Principia in quantitate sunt minima, in virtute maxima. Si reverendus pater illa duo minuta vocabula, apostolatus et regiminis, capitis principalis\apnota{35}{M Primarii, Mc principalis.} et secundarii 'attendisset', dudum pugnantes, ut sibi videtur, sacrae scripturae et s. patrum sententias reconciliasset, quibus neglectis necesse est exotica\apnota{37}{M est ex hoc scil.} et aliena sequi dogmata.

Crastina hora septima antemeridiana continuabitur disputatio edicto publico.

[313] Die Veneris octava mensis Julii hora septima antemeridiana\apnota{1}{Die Zeitbestimmung von Crastina an fehlt in M, U Die Julii octva (Uv octava).}.

D. D. Martinus.

Ad auctoritates a me inductas, Pauli ad Corinth. et Galatas, egregius D. D. more suo nihil respondet, nisi quod ingeminat concilium Constantiense et Hussitas, per quae non solvuntur\apnota{5}{U solvantur, Uv solvuntur.} mihi auctoritates. Nam sicut dixi, concilium Nicenum alia statuit longe sacratius et celebrius. Ideo quantum ille nititur suo, tantum ego meo et dico, sicut heri dixi, quantumlibet allegetur concilium, nondum habetur ius divinum, de quo solo quaestio est, ideo concilium nihil ad propositum.

Secundo, ut eluderet auctoritates potius\apnota{10}{M alluderet potius auctoritates.} quam solveret, advexit distinctionem apostolatus et administrationis, quae quia humano arbitrio est reperta neque in sacris literis neque ecclesiasticis patribus\apnota{12}{M neque ecclesiasticis, letzteres gestrichen; M U in sanctis patribus.} eas interpretantibus habetur, non sum contentus, eam mihi esse redditam\apnota{13}{M ea . . . . reddita.}, praesertim quando apostolatus non, uti pse cum suis 'sentit', dignitatis nomen sit, sed officii, ut\apnota{14}{M quod.} clarus est apostolus ad Rom. [1 v. 5] et ubique, ubi se dicit\apnota{15}{M dicit se.} accepisse gratiam et apostolatum ad obediendum fidei. Est enim apostolus nuncius verbi, et apostolatus seu ἀποστολή Graece\apnota{17a}{M apostoli graece officium vergit, U apostole, (a. R. apostolē) officium verbi, Uv add graece.} officium verbi, et sic in sacris literis apostolatus et\apnota{17b}{U sive.} administratio fere idem erunt\apnota{18}{M vere idem erit.}; nisi administrationem intelligat contentiones forenses de rebus temporalibus ecclesiasticis, quae non est proprie administratio apostolica. Fateor quidem, si ista permittatur libertas fingendi, fingam\apnota{20}{U fingens, Uv fingam, Lö fingerem.} et ego s. Johannem esse cancellarium\apnota{21}{W. A. Anm. „nach S. 316 Z. 9 f. hat Luther beispielsweise auch noch Marcum sigilliferum hinzugefügt."} et alium apostolum aliud esse. Sed haec omnia non sunt iuris divini. Evanescit ergo principium illud, quod minimum in quantitate iactavit\apnota{23}{U iactitatur, W. A. fügt nach S. 312 Z. 34 zu in virtute maximum.}. Et satis admiror, quod dominus egregius sacras literas ita exhorret et faciem earum\apnota{24}{U eorum.} fugit, divertens in istos cuniculos humanitus inventos.

Proinde adhuc stat auctoritas ad Corinth., quod pro schismate tollendo Paulus tollit personam Petri, Pauli\apnota{26}{M Pauli Petri.}, Apollo etc. et solum Christum relinquit, sibi relicto ministerio dumtaxat verbi, quod est vere apostolatus et administratio ecclesiae. Nec valet, quod hanc robustissimam auctoritatem voluit dissolvere frigida illa\apnota{29}{M fridiga (so!) ista, U frigida ista.} similitudine, quod provincia sine capite exposita sit schismatibus et periculis multis, nam non dabitur ecclesia sine capite, quando dabuntur multi episcopi sine papa.

Quemadmodum videmus fieri in foederibus, ubi eo pulchrius administrantur respublicae, quo minus habent commune caput, nam hoc est ad oculum, quod regnum Franciae, Hispaniae, Ungariae et infinitorum\apnota{33}{U Franciae, Angliae, Britaniae, Galliae et infinitorum.} in mundo sunt sine communi capite et tamen non colliduntur\apnota{34}{M colliguntur, Mc colliduntur.}, quolibet regno suo rege\apnota{35a}{M regi.} contento. Quod si in iis temporalibus rebus\apnota{35b}{M his temporalibus rebus, U iis temporalibus.} fluctuantibus tot diversis ingeniis et moribus hominum, non\apnota{36}{M neque; et moribus fehlt in P.} habentibus neque in caelo neque in terra commune caput, stat pax et concordia terrarum, quanto\apnota{37}{M tanto.} [314] magis id fieri potest in spiritualibus rebus constantissimis, hoc est, una fide, eadem charitate, eisdem moribus praeditis hominibus, deinde commune caput in caelo habentibus, sicut apostolus dicit: Una fides unus dominus, ad Eph. [4 v. 5]\apnota{4a}{ad Ephes. fehlt in M.}. Quare si decem essent\apnota{4b}{P esset.} papae, si mille essent papae, non ideo schismata fierent.

Quare peto D. D. velit adhuc\apnota{6}{M U adhuc d. d. velit.} auctoritatem Pauli melius solvere, similiter et illud\apnota{7a}{M similiter illud, U similiter et illam.} ad Gal. [2 v. 6]: Deus personam hominis non accipit\apnota{7b}{etc. statt non accipit.}, ubi Paulus reiicit, quidquid in Petro et apostolis personatum esse potest seu qualificatum, ut dicunt. Nondum est tacta\apnota{9}{M nondum est facta, U utraque nondum est tacta.}, nedum soluta.

Quibus, ut finiam disputationem, addo et illud, quod Paulus 1 ad Corinth. 12. [v. 28] describens corpus ecclesiasticum, primo loco apostolos, secundo prophetas, tertio doctores et fere postremo loco gubernationes\apnota{12}{M iubernationes.}, quas D. D. administrationes vocat, recenset, si ergo primatus Petri tanta res et tam\apnota{13}{fehlt in M.} necessaria esset de\apnota{14}{fehlt in U.} iure divino, insufficienter descripsisset Paulus ecclesiam, ut qui principalissimum, quod in ea est, omisit.

Insuper et hoc adhuc stat, quod nec Petrus nec universi apostoli Matthiam ordinare\apnota{17}{M ordinari.} potuerunt, Act. 1, nec ei administrationem dare. Ideo Rom. pontifex plus sibi arrogat administrationis conferendae per totum mundum\apnota{19}{M U orbem.}, quam in suo praecessore apostolo non invenit.

Similiter sentio de ordinatione Pauli et Barnabae, Act. 13., quos\apnota{20a}{fehlt in U, in Uv nachgetragen.} spiritus sanctus, non Petrus aut\apnota{20b}{U autem, Uv aut.} ecclesia segregavit. Ideo gloriatur non immerito, se nihil accepisse ab apostolis magnis, imo ad Corinth. scribens dicit, nihil minus se fecisse ceteris apostolis. [II. 12 v. 11.]

Sed et illa auctoritas adhuc Petrum artius stringit, qua scribit ad Galatas: Qui operatus est Petro in apostolatum circumcisionis, operatus

est et mihi inter gentes, et infra: Dederunt mihi dextras et Barnabae societatis, ut\apnota{27a}{M ipsi, gestrichen.} nos in gentes, ipsi autem in circumcisionem, tantum ut pauperum\apnota{27b}{M pauperes.} memores essemus\apnota{28}{M erimus, U simus.} [2 v. 8ff.]. Qua auctoritate si vellem niti, non posset D. D. cogere aut ullus hominum, apostolatum Petri latiorem esse quam in Judaeos, iure\apnota{29}{M de iure.} divino, et per hoc omnes auctoritates, quibus nititur: Pasce oves meas\apnota{30}{U meas etc. et.}, Tu es Petrus etc. referendas\apnota{31}{M reverendas.} esse ad terminos sui apostolatus, quemadmodum Christus clare Joh. 10. [v. 16] oves gentium secernit\apnota{32}{P feceruit.} ab ovibus Israel, dicens: Habeo adhuc alias oves, quae non sunt ex hoc ovili.

Et in fine, quod et D. D. heri tetigit, Apoc. 21. [v. 14] describitur Hierusalem nova cum duodecim fundamentis, praeter illud fundamentum capitale Christum\apnota{36a}{U Christus, Uv Christum.}, et tamen nihil differentiae ibi\apnota{36b}{U ibi nihil differentiae.} ponitur inter fundamenta duodecim, quod omnino faciendum erat, si tantum erat de primatu Petri sentiendum. Ad hoc idem valent plurimae\apnota{38}{M plurimum.} figurae in scripturis, ut duodecim boves maris aenei, 3. Reg. 7.\apnota{39a}{P 3. Reg. 4. M 1. (verbessert 3) regum 4 (a. R. 2 regum 4) duodecim boves mari sereni (so!), U duodecim boves 3. regum 4. maris enei; W. A. verbessert 3. reg. 7.}, item duodecim leunculi\apnota{39b}{P M lenunculi.} throni Salomonis, duodecim lapides Jordanis et similia, in quibus omnibus prorsus omnimoda aequalitas omnium [315] apostolorum descripta est, nec potest ulla inaequalitas probari in contrarium iure divino\apnota{2}{M probari de iure divino; U in contrarium probari iure divino.}.

Haec sint dicta pro ista conclusione, quam ego usque hodie verissimam\apnota{3}{M firmissimam.} credo, imo certo\apnota{4a}{U certe.} scio, nec spero, posse eam unquam convelli\apnota{4b}{M efelli.}. Quare rogo D. D. egregium\apnota{5a}{U d. d. egre:; W. A. d. d. egregius.}, ex quo nullam adhuc opposuit auctoritatem\apnota{5b}{M adhuc nullam auctoritatem opposuit, U nullam adhuc posuit auctoritatem.} divini iuris pro sua sententia\apnota{6}{M sententia sua.}, saltem dissolvat aliquas clare et sine diverticulis, alioqui non habebo eum pro recte sentiente de scripturis theologo et relinquo, imo refero ad iudicium non invidiae neque vulgi, sed ordinandorum iudicum et cuiuslibet prudentis lectoris candidique auditoris.

D. D. Eccius.

Nescio, quid reverendus pater velit, quod ab initio sermonis dicit, me nihil respondisse, et tamen in eius procursu solutiones meas heri datas invictas et robustissimas nititur infringere; an illa sibi constent, ipse viderit.

Dein\apnota{14}{M U Deinde.}, quod auctoritate concilii Constantiensis non vult se artari, viderit ipse, videant iudices, quam iuste fecerit. Nicenam synodum

non respuo in his\apnota{16a}{U iis.}, quae sunt fidei, at quae\apnota{16b}{U atque (statt at quae).} mores respiciunt ut ordinationes episcoporum, pro qualitate temporum et locorum mutari possunt\apnota{17}{M posse.}.

Postulat ex me, ut ius divinum ei afferam, quod etsi saepius fecerim, conclusio tamen mea illud non habet, sed dumtaxat, quia reverendus pater in resolutorio dixit, Rom. ecclesiam ante tempora Silvestri non fuisse aliis superiorem, quod in conclusione mea nego, tamen fateor, primatum summi\apnota{21}{U Romani.} pontificis esse de iure divino secundum allegata per me.

Dein\apnota{23}{U Deinde.} reprobaturus solutionem dicit, distinctionem meam humano arbitrio inventam, quasi non liceat in theologia distinguere, sicut Arius voluit Athanasium irridendo, nam nimis literae\apnota{25}{M litera.} tenax dixit filium esse creaturam, per illud: Pater maior me est [Joh. 14 v. 28]. Quod diluens\apnota{26}{W. A. verbessert diluit.} Athanasius distinguendo: Pater maior me est, secundum humanitatem, sed non secundum divinitatem; quia sacra scriptura\apnota{28}{U scripturas sacras.} non distinxit, Arius eam non recepit et ita in errore pertinaciter mansit. Sic in praesentia, cum qui elucidant sacram scripturam, vitam aeternam habebunt. diversos passus\apnota{30}{M U passus scripturae.} conferendo, ne in sacris literis contrarietas admittatur, distinctiones dandae sunt, ne literae nimis tenax occidatur, sed magis a spiritu vivificetur. Et quia Petrus aliis praefertur\apnota{32}{U praeferter, Uv praefertur.} saepius in sacris scripturis, saepe eis aequalis reperitur, ad contrarietatem vitandam necessarium est, distinctionem ponere sacrae scripturae conformem.

Sic et Hieronymus innuit et Cyprianus per me\apnota{35a}{M a me.} allegati\apnota{35b}{U allegati et epistola 9, W. A. allegati lib. 4 epistola 9.}, qui et aequalitatem apostolatus fatentur et tamen primatum Petro tribuunt, quod fieri non potest, nisi ad diversa referas. [316] Placet mihi apostolatum dicere officium, ita tamen quod etiam dicat honorem. Sunt enim ista nomina et honoris et oneris, ut dixi, quod apostolus ad Romanos indicat, cap. 10 [v. 15]: Quomodo evangelisabunt, nisi mittantur?

De administratione iterum assentior\apnota{4}{M assentio domino.} reverendo patri, illam scilicet in temporalibus esse multo inferiorem, quam illa\apnota{5}{W. A. illam.} in spiritualibus. Deum optimum rogemus, ut episcopi nostri illud agnoscant et spiritualia non per vicarios et temporalia per 'se ipsos' administrent\apnota{7}{M administrant.}. Denotavi ergo administrationem regiminis in spiritualibus potestative, nec obstat, quod reverendus pater dicit, sic se posse fingere Johannem cancellarium, Marcum sigilliferum etc., quoniam, ut Augustinus ait, liberis verbis utuntur philosophi, nobis vero\apnota{11}{M autem.} ad certam regulam loqui fas est, ideo talia oportet afferre, quae fundamentum recipiant\apnota{12}{M recipiunt.} ex sacra scriptura,

sicut primum membrum apostolatus in apostolorum aequalitate, alterum membrum regiminis in Petri primatu. Unde fateor robustissimam esse auctoritatem Pauli, sed quae mihi nec in una litera adversetur.

Contempsit reverendus pater simile per me adductum de schismatibus in provincia praeside carente, quoniam iam ita sit in regnis fidelium Hispaniae, Galliae, Angliae, Hungariae, quae sint\apnota{18a}{U sunt.} sine capite, et tamen stet pax et\apnota{18b}{U stat pax, quia.} concordia inter illa. Miror, quomodo reverendus pater tantarum caedium, bellorum, proeliorum iam fuerit oblitus, quae non potest ignorare. Dicere nolo\apnota{20}{M volo.} tam diuturnum dissidium et perpetuum odium inter Gallos et Anglos; odium etiam inveteratum inter Gallos et Hispanos, qui toties pro regno Neapolitano cum maxima Christiani\apnota{23}{M maximae Christianae.} sanguinis effusione bella gesserunt, et utinam tandem pax esset stabilis et permanens. Unam fidem, unum dominum Christum agnosco cum apostolo, sed Rom. pontificem ut Christi veneror vicarium. Stat ergo prior solutio invicta et valida meo saltem exili iudicio, de quo doctiores iudicent.

Praeterea observandus pater dixit\apnota{25}{M U dicit.}, verba Pauli per me non tacta: Deus personam hominis non accipit\apnota{29}{U a. R. quia omnes sumus caeci, dixit Eccius (soll heissen Martinus?); W. A. fügt diese Randnote schon 2 Zeilen früher nach indicent ein.}, et sic minus soluta. Forte obaudivit, quoniam tam ex Ambrosio quam Hieronymo adduxi eorundem verborum sententiam, Johannes enim et Petrus idiotae fuerant et analphabetici\apnota{31}{M alphabetici.} in lege; Paulus contra doctus fuerat\apnota{32a}{U fuit.}, proficiens supra multos coaetaneos suos, ut de se\apnota{32b}{M suos coetaneos, ut se (Mc de) se.} testatur [Gal. 1 v. 14]. Noluit ergo Paulus reputare, quales fuissent Petrus et Johannes, quia personam hominis deus non accipiat\apnota{34a}{U accipit.}, ita enim elegit piscatorem, sicut doctum\apnota{34b}{M doctam.} scribam. Sed per hoc primatus non tollitur, sed acceptio personarum\apnota{35}{M personatur.} deo non imputatur, etiamsi unum alteri praeelegerit, ut de Cornelio Petrus fatebatur [act. 10 v. 34]: In veritate comperi, quia non est acceptio personarum apud Deum.

Veniamus ad inducta per reverendum patrem. [317] Apostolum citat, 1. Corinth. 12 [v. 28], qui corpus mysticum describens referat apostolos, prophetas, doctores et novissime administrationem. Dico, egregie quidem apostolum descripsisse, sed primatum ibi nec negasse, nec astruxisse, invalidum autem\apnota{4}{U est ergo statt autem est.} est argumentum: Hic Paulus huiusmodi primatus non meminit, ergo primatus non est in ecclesia, tum quod pueri sciunt locum ab auctoritate negative non valere, tum quod etiam in sacra scriptura uno loco non expressum, iuxta Augustini doctrinam, lib. de doct. Christiana, ex alio loco accipi debet, sicut nos fecimus.

Pariformiter respondemus ad illud Apoc. 21. Fateor, duodecim sunt fundamenta, sint illa etiam figurata per duodecim boves, per duodecim leunculos\apnota{10}{P lenunculos.}, per duodecim lapides Jordanis, sed in iis\apnota{11}{M his.} numerus duodenarius apostolorum asseritur, primatus ob hoc non negatur. Ego autem eum locum ob hoc induxi, quod oporteat passus\apnota{12}{M oportet fassuss.} sacrae scripturae concordes ostendere, et cum apostolus unum fundamentum, Christum, posuisset et nemo aliud ponere posset, Petrum per hoc fundamentum ecclesiae negari\apnota{15}{fehlt in M.} non debere, quia Johannes et ipse sanctus duodecim posuerit fundamenta, quod nisi distinctionem acceperis de fundamento\apnota{16}{M susceperis de [fundamento fehlt], U funnamento.} fundamentorum, id est principali, et de fundamento substituto seu secundario\apnota{17}{M de [fehlt in P] fundamento subsecunda [Mc substituto] et secundario.}, Paulum et Johannem non concordabis.

Secundo\apnota{19a}{M U secundo loco.} induxit, Petrum non potuisse ordinare\apnota{19b}{M ordinari.} Matthiam, sic spiritum sanctum segregasse Paulum et Barnabam. Respondeo: Cum apostoli fuerint episcopi, non memini me reperisse in sacra scriptura ordinatos a Christo episcopos, cum in ultima coena primum eos ordinaverit sacerdotes. Quare et illud in Petrum referre possum, quoad episcopatus ordinationem; attamen, quia apostolatus officium a solo deo tributum admitto, non mirum, quod divinum postulaverunt praesidium, sortes enim miserunt non vulgari more\apnota{25}{M modo, gestrichen, also vom Disputator widerrufen.}, sed ut sacer Dionysius testatur, sors illa fuit spiritus sancti visibilis apparitio. Unde locus ille pro primatu Petri facit, quia primatum sibi concessum executus est: In diebus illis exurgens Petrus in medio fratrum\apnota{28}{M fratrum dicens.} etc. Sicut et de Saphira et Anania, sicut in excusatione apostolorum, quod non essent musto repleti, sicut in responsione et defensione coram concilio\apnota{30}{M U consilio.} Judaeorum [act. 1—5] et aliis in Actis apostolorum contentis, ubi regimen superioritatis exercuit. Unde Matthiam credo apostolum a domino factum et a s. Petro episcopum ordinatum. Sic et de Paulo dudum concessi\apnota{33}{M concessit.}, quod apostolatum a deo acceperit\apnota{34}{M acceperat.}, abundantius eum laborasse aliis fateor.

Sed ad tertium veniamus, quod videtur facere magis\apnota{35}{U magis facere.} ad propositum inter alia per reverendum patrem adducta: Qui cooperatus est Petro, ad Gal. 2. [v. 8] Et dico, Paulum non promulgasse ibi ius divinum, sed narrasse factum, quoniam et Petrus in Actis apostol.\apnota{38}{M dixit ad Judaeos: ecce convertimur ad gentes gestrichen, also vom Disputator widerrufen.} Cornelium gentilem iussus est in fide instruere et gentibus in Antiochia, ubi praesedit, et similiter Romae, ubi tandem martyrio a Nerone gentili coronatus est, praefuit et docuit\apnota{40}{U a. R. non est verum, dixit Martinus.} simul [318] cum Judaeis ad fidem

conversis et convertendis. Unde clare liquet, Paulum factum narrasse\apnota{2}{M U tunc narrasse.}. Praetereo, quod s. Paulus non meminit officii regiminis vel primatus praesidentiae, sed executionis apostolatus, secundum quem apostoli diversas sortiebantur provincias, in quo et Petrus aliis apostolis aequalis fuit secundum Anacletum can. In novo, secundum Hieronymum contra Jovinianum, et super epistola ad Galatas secundum Cyprianum, in epistola 9. ad Cornelium papam lib. 4. Qua de re sacra scriptura intelligenda est cum effectu, ut advertamus, quando de officio apostolatus loquatur, et quando de regimine potestatis seu praesidentiae.

Non refert, quod reverendus pater adhuc negat, me adduxisse ius divinum pro primatu et non dissolvisse auctoritates per eum inductas, quia idem dicerem de eo. Attamen nullus ex s. patribus ita intellexit auctoritates per reverendum patrem inductas ad negationem primatus Petri super alios, sicut ipse induxit. Contra vero plurifariam\apnota{14}{M plurivariam.} sancti patres intellexerunt auctoritates per me inductas loqui de primatu Petri. Unde circa illam conclusionem, si aliquando acerbius et durius quicquam locutus sum contra reverendum patrem, nolim\apnota{17a}{U nollem.} arbitremini in personam suam dictum\apnota{17b}{M dictam.}, sed causa ipsa, quae vertitur, tam praegnans et necessaria mihi videbatur, ut verba aliquando acriora essent utenda\apnota{19}{M uterquam.}. In reliquis conclusionibus Eccianam semper experiemini\apnota{20}{M reperiemini.} modestiam.

Unde et ego istam conclusionem volo finire in praesentia, quod reputem\apnota{21}{M repetens, U repetam.} s. Petrum primatum totius ecclesiae a Christo tenuisse ex promissione ei facta Matth. 16, sicut intelligit Hieronymus, Cyprianus, Augustinus\apnota{23}{U Hieronymus, Ambrosius, Cyprianus.} et alii; ex nominatione, qua\apnota{24}{M U quia.} primus nominatur Matth. 10. secundum glossam ibidem; ex solutione tributi, Matth. 17., ubi solus ipse fuit aequiparatus Christo, secundum Augustinum et Ambrosium; Ex eo, quod Christus rogavit pro fide eius indefectibili et iussit eum confirmare fratres Luc. 22., secundum Chrysostomum et s. Leonem; quod pastor ovium Christi\apnota{28}{M Christi, Mc a Christo.} constitutus est, Johan. ult., [c. 21] secundum Chrysostomum et Gregorium; quod Christus ait ad Petrum: Sequere me, non solum scilicet genere martyrii, sed et ordine magisterii, secundum Theophilum; quod solus iussus est Petrus venire super mare ad Christum [Mat. 14], ut sic mundus per mare designatus totus esset\apnota{32}{fehlt in U, in Uv nachgetragen.} ei subiectus secundum Bernhardum de consideratione ad Eugenium, et per alia superius adducta, reputo conclusionem illam a tot sanctis patribus et novissime per concilium Constantiense comprobatam\apnota{35}{U a concilio Constantiensi probatam.} esse veram. Nolo tamen inniti propriae prudentiae, paratus

captivare intellectum meum in his\apnota{36a}{U et in iis.} et quibuscunque aliis iudicio iudicum\apnota{36b}{fehlt in M.} ordinandorum sedis\apnota{37a}{U sedi.} apostolicae et aliorum prudentium et bonorum virorum\apnota{37b}{M vinorum.}.

Soli deo gloria. Eccius\apnota{38}{fehlt in U.}.

Hora secunda pomeridiana continuabitur disputatio.

[319] Hora secunda pomeridiana eiusdem diei\apnota{1}{Die Zeitbestimmung [2 Zeilen] fehlt in M ganz, U hora secunda.}.

D. D. Martinus.

Cogimur adiicere coronidem, quia egregius D. D. diluens obiecta mea simul plurima opposuit et rursus me respondere coegit. Et quanquam id voluissem praeterire, tamen, ut ordinandis iudicibus per omnia meam sententiam declarem, paucis respondeo.

Primum opponebat, quod Petrus praefabatur electioni Matthiae\apnota{7}{M Matthei.}, exhortans apostolos etc. Credo autem nec ipsi\apnota{8}{M ipso.} D. D. videri per hoc probari primatum Petri, si Petrus praedicat, cum in superioribus ipse mihi responderit, aliud esse apostolatum et aliud administrationem regiminis. Non recte ergo dixit, Petrum sic praedicando exercuisse administrationem regiminis\apnota{11}{non recte ergo dixit, Petrum sic praedicando exercuisse administrationem regiminis fehlt in P, offenbar durch ein Versehen (Homoioteleuton).}, alioquin et Stephanus erit primus et papa, qui\apnota{12a}{M U quia.} et ipse praedicavit et signa\apnota{12b}{fehlt in U, in Uv nachgetragen.} multa fecit in populo\apnota{13}{in populo fehlt in M.}.

Et quod credit, Matthiam apostolum esse episcopum ordinatum a Petro, sum contentus, ut credat, quidquid volet, ego non credo, nisi probet. Item quod Petrus exercuerit\apnota{15}{M excercuit, Mc exercuerit.} primatum in percussione Ananiae et Saphirae, potest dici, sed non probatur\apnota{16a}{M probari.}, nam suscitavit et mortuos\apnota{16b}{et fehlt in U; M et resuscitavit et mortuos.}, quod et alii apostoli fecerunt. Item quod excusavit apostolos, ne viderentur esse ebrii, nec hoc indicium est primatus, cum, ut dixi, egregius D. D. semper distinxerit primatum ab omnibus operibus et administrationibus communibus ceteris. Item quod Petrus ambulavit super mare et Bernhardus per\apnota{20a}{in W. A. weggelassen.} mare interpretetur\apnota{20b}{U interpretatur.} mundum, concedo, sed nihil ad primatum, cum et nos calcare debeamus mundum.

Illud libens\apnota{22a}{M liber.} transeo de fidei indefectibilitate rogante\apnota{22b}{M probante.} Christo promissa, concedo enim, fidem Petri nunquam cecidisse, etsi ipse ceciderit\apnota{23}{M ceciderat.}

a fide. Tunc enim latro credidit fide Petri, quando Petrus negavit, ut Augustinus ait. Longe aliud est fides, quam primatus\apnota{25}{M aliud fides et primatus.}.

Item et illud, quod Petrus dedit staterem pro Christo et se, in quo aequiparatus dicitur Christo, concedo, sed nihil ad primatum, imo contra primatum, nam mox sequitur, ut et Hieronymus interpretatur\apnota{28}{M etiam interpretatur.}, ceteros apostolos ideo offensos, quod Petrum suspicarentur fore maiorem, coeperunt disputare, quis esset maior. Ideo Christus advocato parvulo compescuit eorum ambitionem, et sic sequentia textus fortius pugnant contra primatum quam praecedentia pro primatu.

Et\apnota{32}{fehlt in U.} illud. Johan. ult. [c. 21], quod soli Petro dictum est: Sequere me, interprete Theophilo, sit intelligendum non modo de genere martyrii, sed et\apnota{33}{fehlt in P.} ordine magisterii, nego Theophilum, tum quod\apnota{34}{U cum quod, W. A. cum.} Augustinus melius videatur sentire dicens sic: Sequere me temporalia mala perferendo. Et si Augustinus non diceret\apnota{35}{M dixerit.}, ipse textus evangelii convinceret\apnota{36}{M convincit.}. Exiit enim sermo inter discipulos, quod discipulus ille non moreretur, cum audissent Christum dicentem: Si\apnota{37}{M sic modo.} volo illum manere, quid ad te? Tu me sequere. Ideo manifestissimum est, quod de genere martyrii et passionis imitatio ista intelligenda sit.

[320] Quare ego rogo, D. D. omittat nova adducereet praesertim ea, quae non faciunt adrem, cum et ego possem inducere illud Act. 8., [v. 14] ubi apostoli miserunt Petrum et Johannem tanquam inferiorem, et Act. 15. [v. 13 ff.] Jacobus sermonem Petri et confirmavit et mutavit. Sed nolui haec et similia inducere, cum iamdudum concesserim primatum honoris Petro solumque negaverim primatum potestatis super ceteros apostolos. Haec est enim\apnota{6}{M enim est.} illa persona, imo et primatus honoris, quam deus non respicit, ut superius induxi, non autem illa persona piscatoris vel\apnota{8}{U et.} paupertatis, ut Ambrosius referente doctore exposuit, Hieronymus enim melius de magnifica qualitate Paulum interpretatur\apnota{9}{M interpretavit.}, qua pseudoapostoli Galatas\apnota{10}{U quam pseudoapostoli ad Galatas.} tanquam ab auctoritate Petri subvertere voluerunt. Sum ergo\apnota{11}{U enim.} contentus optime, quod D. D. suam propositionem per auctoritates Hieronymi, Ambrosii, Bernhardi, Leonis confirmavit\apnota{12}{M confirmavit etc.}, praecipue super hanc auctoritatem: Tu es Petrus, modo mea sententia, ut in praecedentibus dictum est, eisdem et maioribus auctoritatibus et magis ad intentionem loquentibus, confirmatior habeatur, referens haec quoque\apnota{15}{U quaeque.} ad iudicium, ut supra.

D. D. Eccius.

Reverendus pater recludit hodie conclusa et ea, quae velut accessoria et emblemata\apnota{18a}{M emplamata.} in conclusionis adieci\apnota{18b}{M adiecit, U in conclusione adiecit.} coronidem, repercutit ut principalia et adeo metamorphosim curat, ut ignoret\apnota{19}{M curam ut ignorem, U curat ut ignorem.}, an in opponentem vel respondentem sit transformatus. Ad rem brevibus\apnota{21}{M brevius.}.

Nunquam somniavi illam consequentiam: Petrus praedicavit, ergo fuit primus. Deinde credulitatem meam de Matthia improbat et tamen rationem meam non diluit. Matthias et ceteri apostoli fuerunt episcopi et non fuerunt ordinati a Christo, nec se ipsos ordinarunt, quare non nisi a Petro fuerunt ordinati, quem Christus pastorem universalem constituit pascendo oves\apnota{27a}{universalem und pascendo oves fehlt in U, ersteres in Uv nachgetragen, M pasce oves.}. Quare mecum credat, vel\apnota{27b}{W. A. vel aliud.} ad illud respondeat.

De excusatione Petri imponit\apnota{29a}{M iterum imponit.} mihi nunquam cogitata\apnota{29b}{U cogitatum.}, quod primatum separaverim ab operibus administrationis aliorum apostolorum, quod\apnota{30}{fehlt in U.} nunquam fuit mentis meae, imo allegavi ei\apnota{31}{fehlt in U.} b. Gregorium, alios episcopos vocatos in partem sollicitudinis, non in plenitudinem potestatis.

De Bernhardo super transitu Petri super mare, dicit nihil facere ad primatum. De quo miror, quomodo dicere possit, si Bernhardum legat\apnota{34}{M U legit.}. Nam ex intentione Bernhardus voluit probare ex hoc primatum Eugenii\apnota{35}{M a. R. Petri.} super alios, et totum mundum ei debere esse\apnota{36}{M U debere ei esse (esse fehlt in M).} subiectum, non sic aliis episcopis et apostolis\apnota{37}{M apostolis, vel episcopis, U episcopis vel apostolis, Uv = P.}.

De statere ex litera sequenti vult elicere oppositum, quod tamen non fecit\apnota{39}{M feci.}, sed Christus murmur apostolorum compescuit. At ego non meo capite, [321] sed sanctorum patrum verba sacrae scripturae accipio, nam ex hoc b. Ambrosius voluit denotatum, Petrum reliquis superiorem futurum, quod et Augustinus in quaest. evangeliorum\apnota{3}{M a. R. August. q. 69.} testatur expressissime. Quare\apnota{4}{U quaestio mihi non occurrit quare.} nostro sensui non innitamur, sed s. patribus.

De fidei indefectibilitate transiit\apnota{5}{P M U defectibilitate transiit pater (pater fehlt in P).}, adducens tamen illud Augustini, quod fidem\apnota{6}{M U fides.}, quam Petrus negavit, latro in cruce servavit. Et ego de illo dicto transeo. Sciat tamen reverendus pater, quod s. Cyprianus verbis Christi innixus voluit nullas haereses Romae velut\apnota{8}{U vel.} in aliis ecclesiis exortas, sicut pridie dixi, quod in Romana sede fuerunt ducenti

et duodecim papae et tamen ante eorum amotionem nullus\apnota{10}{M nullus sensit.} definitive et sententiam ferendo in fide erraverit.

Ad Theophilum, quem suo nutu\apnota{11a}{M U pro suo nutu.} negat sine\apnota{11b}{U sive.} alicuius vel auctoritatis vel rationis inductu, miror, quod ecclesiasticos patres ita contemnit, qui sophistis est admodum infensus. Augustinum induxit, qui Theophilo non adversatur, Christum locutum de genere martyrii. Sed doceat, ubi Augustinus ordinem magisterii negaverit\apnota{15}{P negagaverit (so!)}, cum ibidem tract. 124. de Petro fateatur, eum fuisse abundantione gratia, unum eundemque primum apostolum. Et in epistola 162.: In Romana ecclesia, inquit, semper apostolicae cathedrae viguit principatus. Quare non me\apnota{18}{M me non.} roget, ut conclusionem nostram, quam reputo veram, non defendam.

Sua, quae\apnota{20}{M U Suaque.} adduxit de Petro misso in Samariam, et per hoc primatum ab eo ablatum, quasi nesciamus aliquem mittere se ipsum, ut filium dei a patre et a se ipso missum, non noviter ista ventilari\apnota{22}{M ventilare.} coeperunt, sed omnia dudum fuerunt soluta per Alpharum, Johannem de Turre cremata\apnota{23}{M per Johannem de Turre crematae.} et alios. Quod si Petrus Jacobo cessit Hierosolymis in concilio, et loci et senectutis rationem habens, primatum Petro non aufert. Sed ut videatur aliquid facere, reverendus pater\apnota{26a}{reverendus pater fehlt in U.} tribuit primatum honoris Petro, forte velut\apnota{26b}{M velut forte velut.} ambitioso, cum ex me audierit Cyprianum, Hieronymum et Anacletum, apostolos aequalis fuisse honoris. Ergo primatum honoris reverendus pater tribuit Petro, quem s. patres negant, contra primatum regiminis negat Petro, quem concors s. patrum sententia et concilii Constantiensis Petro tribuit.

Desinat ergo, oro, gloriari se contra tam s. patres, contra tam celebre concilium maioribus\apnota{32}{M U se maioribus.} auctoritatibus nixum, quod saepe ostenderim, auctoritates per eum\apnota{33}{M eam.} inductas primatum non negare alicuius s. patris sententia. Quare potius cum Graecis sero sapientibus, velut de Troianis dicitur, tandem sapiat, Rom. pontificem verum Christi vicarium primum in orbe tenere locum, sicut Graeci ipsi, imperator, patriarcha et primates hoc confessi sunt, et obedientiam fecerunt Rom. ecclesiae anno 1439. 22. Novemb. et sic fiat pax in diebus nostris. Quae omnia, ut supra, committo iudicio eorum, quorum interest et intererit. Soli deo gloria. Eccius\apnota{39}{fehlt in M U.}.

[322] D. D. Martinus.

Non displicent omnia ferme, quae locutus est egreg. D. D., praecipue illud Aug., quod apostolicae cathedrae principatus\apnota{3}{M catholicae (Mc apostolicae) principatu.} viguit semper

in ecclesia Rom., si addidisset unum verbum, scilicet principatus potestatis super\apnota{4}{U et unum verbum; M si addisset adhuc verbum unum principatus super.} omnes episcopos, et Bernhardum credo torquere textum de Petro ambulante\apnota{5}{M U ambulante Petro.} super mare, cum consequentia manifeste fidei tribuat\apnota{6}{M U sequentia .... tribuant.} ambulationem illam, non primatui\apnota{7}{M primatum.}, dicente Christo, cum submergeretur Petrus: Modicae fidei quare dubitasti [Mat. 14 v. 31]? Et quod cupit, ut\apnota{8}{M quod.} credam secum apostolos ceteros fuisse ordinatos episcopos a Petro apostolo, modo non placet obsequi, cum omnes fuerint\apnota{10a}{U fuerunt.} episcopi aequaliter Petro, ut textus per Petrum inductus\apnota{10b}{M induxtus.} probat de Juda\apnota{11a}{de Juda fehlt in M.}: Et episcopatum eius accipiat\apnota{11b}{M accipiet.} alter [act. 1 v. 20; Ps. 109 v. 8]. Haec de isto\apnota{11c}{M ista.}.

D. D. Eccius.

De episcopatu, quod etiam Judas fuerit episcopus iuxta illud Psal.: Et episcopatum eius accipiat\apnota{14a}{M accipiet.} alter, dico episcopatum ibi usurpari\apnota{14b}{M ibi usurpari ibi.} pro apostolatu et non pro ordine episcopali. Nam apud plerosque dubium est, an Judas etiam fuerit\apnota{16}{M etiam Judas fuerat.} sacerdos. Sed hoc nihil ad rem. Exiit tamen Judas buccella\apnota{17}{M porcella.} sumpta, quare dum in coena Christus ordinaverit discipulos sacerdotes, Judas non fuit\apnota{18a}{M fuerat.} episcopus consecratus. Committo iudicand\apnota{18b}{U haec iudicanda.}, aut alia.

Instaurata fuit disputatio mox et continue assumpta materia de purgatorio circa nonam positionem.

D. D. Eccius ut opponens orsus est\apnota{20}{Die Zwischenbemerkung lautet in M Eccius contra nonam conclusionem doctoris Martini sexta feria in die Kiliani. De animabus in purgatorio; in U De purgatorio. Contra nonam conclusionem argumentatur Eccius.}.

Non solum apud theologistas, sed etiam in sacra scriptura et\apnota{21}{M U in (statt et).} s. patribus habetur, quod vita haec praesens sit status et terminus merendi et demerendi. Quare existentes in purgatorio amplius mereri non possunt\apnota{23}{U posse, Uv possunt.} et sic etiam maiorem gratiam non accipere\apnota{24}{etiam fehlt in U; M U non maiorem gratiam accipere.}. Haec\apnota{25}{M U hoc.} ita esse probo per Jer. 25 [v. 14]: Reddam eis secundum opera eorum et secundum facta manuum eorum, ita etiam Paulus ait: Omnes nos manifestari oportet ante tribunal Christi, ut referat unusquisque propria corporis, prout gessit, sive bonum, sive malum [2. Cor. 5 v. 10]. Scientes ergo timorem Domini hominibus suademus.

Addo Hieronymum, ut gradatim descendam, eo loco apostolum intelligentem: Hoc dicit, inquam, quamdiu in mundo positi peregrinamur

a Domino, conversatione bona id agamus, ut ei in futuro placere possimus, non, ut quidam putant, quod posteaquam excesserimus a corpore\apnota{33}{U e corpore.} ibi aliquid operantes promeremur. Idem fatetur glossa ordinaria ibi\apnota{34}{fehlt in U.}: Frustra ergo homo post corpus sibi hoc promittit, quod hic comparare neglexit. Ita etiam intellexit Ambrosius, solum facta corporis alibi praemiari.

Ex quibus liquet, cum gratia augeatur\apnota{37}{M agatur.} ex meritorum respondentia et animae in purgatorio non poterunt\apnota{38}{M poterint.} aliquid operari meritorie ex allegatis, non [323] poterit etiam in eis augeri gratia. Subscribit apertissime b. Augustinus in ench. cap. 109. et 110. et repetit eadem verba principaliter quaest. 2. de octo Dulcitii\apnota{3}{M stulcicii.} quaest. post multa. Quocirca, ait, hic omne meritum comparatur, quo possit post hanc vitam relevari\apnota{4}{M revelari.} quispiam vel gravari. Nemo autem se speret, quod hic neglexit, cum obierit, apud deum promereri. Cessat\apnota{5}{M cesset.} ergo post mortem meritorii operis auctio, cesset\apnota{6}{U cessat.} et gratiae auctio meritis respondens.

D. D. Martinus.

Opposuit egregius D. D. duas divinae scripturae auctoritates et quatuor s. patrum. Prima Jeremiae, ubi deus dicit: Reddam eis secundum opera eorum. Miror satis egregium D. D., quod haec et similia arbitretur contra me pugnare, cum ego in resolutorio meo multo plures adduxerim in hanc formam et exposuerim, quomodo nihil contra me pugnent\apnota{14}{M pugnet.} aut aeque contra purgatorium pugnent\apnota{15a}{M pugnet, U pugnant.}. Cum enim scriptura divina\apnota{15b}{U divina scriptura.} tota prorsus nihil habeat de purgatorio, sed omnia loquatur vel de inferno vel caelo, volo hac\apnota{16}{M U hoc hac.} una responsione ad omnes auctoritates respondisse, quod non faciunt\apnota{17a}{U faciant.} ad propositum, sequenter\apnota{17b}{Uv sequentur.} et omnes\apnota{18}{fehlt in M U.} expositiones omnium patrum, nisi expresse meminerint purgatorii. Nam credibile est, quod, quando tractant scripturas sanctas, sensum etiam scripturae sanctae\apnota{20}{fehlt in M.} sequantur. Ideoque aliis rationibus et auctoritatibus doceri oportet, animas esse certas et non augeri in eis charitatem\apnota{21}{U eas in charitate.}.

D. D. Eccius.

Cum reverendus pater plures se dicat allegasse\apnota{23}{U allegasse dicat.} auctoritates in resolutorio et eas\apnota{24}{fehlt in U.} exposuisse, vidimus. Vidimus quoque glossam suam ad textum nihil facere, quod iam ostendemus. Hoc quoque non

accipimus, in sacra scriptura nihil haberi de purgatorio, quod dictum quidem Graecis et Picardis esset favorabile, verum Christianae fidei adversum, ut recte meminit in resolutorio, sed ponamus solutionem suam unicam ad aciem\apnota{28}{Löscher „in usto legitur arcem" beruht auf Irrtum.}, qua dicit, adducta nihil facere ad propositum, quam concinne respondit\apnota{29}{M respondeat.}! Apostolus dicit: Unumquemque recepturum secundum quod in corpore gessit, D. pater dicit, hominem recipere etiam secundum quod in purgatorio gessit. Augustinus dicit, omne meritum hic comparari, et D. pater contra dicit, etiam in purgatorio meritum comparari. Hieronymus dicit, postquam e corpore excesserimus\apnota{34}{M excesseribus, Mc excesserimus.}, nihil nos promereri\apnota{35}{M promeremus (nos fehlt).} apud deum, D. pater contra dicit, postquam e corpore excesserimus, adhuc in purgatorio nos promereri, ita de Ambrosio. Judicent iam illustrissimae dominationes, excellentiae et praestantiae vestrae\apnota{36}{fehlt in P.}, an per me adducta\apnota{37a}{M aducta, U inducta.} nihil faciant\apnota{37b}{M faciunt.} ad propositum, quae directo dictis patris\apnota{37c}{M directe dictis patrum, U directe nihil dictis patrum.} contrariantur. Peto ergo, quod vere solvat vel ostendat, quomodo ad pro-[324]positum non faciant, nec miretur, quoniam\apnota{1a}{U nihil faciat nec meretur (Walch XV Sp. 1198 moretur) quomodo; Uv verbessert faciant und miretur.} et ego plures sum ei\apnota{1b}{fehlt in M.} adducturus auctoritates, ut\apnota{2}{M quod.} manifestarium sit, 'hoc' non esse theologistarum, sed verorum theologorum.

D. D. Martinus.

Ego, qui credo fortiter, imo ausim\apnota{5}{M ausum.} dicere, scio purgatorium esse, facile persuadeor in scripturis de eo fieri mentionem, quemadmodum illud Matth. [12 v. 32] inducit\apnota{6}{M illud inducit illud Matthaei.} Gregorius in dialogo: Non remittetur\apnota{7}{U remittittur.}, neque in hoc saeculo neque in futuro, volens\apnota{8}{M nolens.} peccata quaedam remitti in purgatorio.

Admitto et illud Macha. 2. [12 v. 45 f.]: Sancta et salubris est cogitatio pro defunctis exorare\apnota{9}{M U sancta est et salubris cogitatio pro defunctis orare (U exorare etc.)}. Sed hoc volo, quod in universa scriptura non habeatur memoria purgatorii, quae posset stare in contentione et convincere, nam et liber Machabaeorum, cum non sit in canone, pro fidelibus potens est, contra pertinaces nihil valet\apnota{12}{M U facit.}. Et dictum Gregorii facillime eluditur\apnota{13}{U eliditur.}, quod neque hic neque in futuro peccatum remittatur\apnota{14}{M remittatur peccatum.}, id est, nunquam. Ideo nolo mihi suspicionem fieri, quod faverem Bohemis et Graecis. Hoc pro primo\apnota{15}{Hoc pro primo fehlt in M.}.

Ad replicam domini D. dico\apnota{16}{U dicam.} nihil aliud, quam quod iam prius dixi. Idem enim repetit, ideo et idem respondeo. Dixi enim, in

sacris literis nihil haberi de purgatorio, ideo non posse eas neque expositiones earum ad negotium purgatorii aptari. Redditus enim\apnota{19}{U autem.} aut bonum bonis aut malum malis post mortem. Medii sunt, qui in purgatorio sunt, de quibus dictae auctoritates non loquuntur, quibus neque redditum est bonum sive\apnota{21a}{U neque.} malum, et\apnota{21b}{fehlt in M.} sic purgatorium per medium auctoritatum\apnota{22a}{Mc 'illarum' auctoritatum.} semper transit\apnota{22b}{U transiit, Uv transit.}.

D. D. Eccius.

Quod se excuset\apnota{24}{M U excusat}, credere se et scire purgatorium esse, bene accipio et prius novi. Sed quod negat in sacra scriptura purgatorium probari, hoc tamquam falsum et Graecanicum respuo, in proposito manens, quamquam nec illa\apnota{27}{M U illud.} nec alia multo maiora contra protervos et pertinaces convinci possint, ut de Ario et aliis patuit\apnota{28}{U patribus, Uv patuit.}. At quod propterea illud dicatur non reperiri in scriptura, quia contra pertinacem convinci non possit, est omnino captiosum et religioni nostrae adversum.

De libris Machabaeorum, quos dicit facere pro fidelibus, sed non esse in canone, et hoc, inquam, falsum est, quamvis enim apud Hebraeos in canone non fuerint, tamen ecclesia recepit eos in canonem, ut pater domini patris Augustinus libr. de civitate dei testatur\apnota{34}{U testatur lib. 8. [W. A. lib. 18].}, et s. Ipho in suis decretis constitutionem inserit, qua ecclesia libros illos in canonem recepit.

Quod in suo proposito persistere vult reverendus pater, allegata per me non facere ad propositum, quae tamen de directo ei contrariantur. Primo\apnota{38}{M U primo enim.} quod dicit auctoritates illas respicere caelum vel infernum, non iuvat, quia de directo faciunt mentionem de statu merendi et demerendi et hunc [325] statum definiunt\apnota{1}{U diffiniunt.} morte terminari. Contra reverendus pater nullo s. doctore\apnota{2a}{M doctori.} innixus audet animabus a corpore exeuntibus\apnota{2b}{U existentibus.} terminum merendi prorogare et merita purgatorii in caelo remunerari, et apostolica clamat sententia, solum remunerari opera, quae homo in corpore gessit. Quare dicta in resolutione sua pervidi\apnota{5}{U perpendi.}, sed tamquam insufficientia acceptare non potui, ut illas contrarietates prius aperui\apnota{6}{M aperuit, U aperiat.}.

D. D. Martinus.

Primo\apnota{8}{U primum.}, quod egregius D. D. dicit, non ideo aliquid negandum esse de scriptura, quia pertinaces convinci non possint\apnota{9}{M possunt.}, optime et verissime dicit, sed loquor\apnota{10a}{M loquar.} ego de iis pertinacibus, qui nos propria\apnota{10b}{M U de his ..... nos nostra.}

auctoritate et proprio iaculo 'confodere' possunt. Evidens enim\apnota{11}{M convincere (Mc a. R. confodere); enim fehlt.} est librum Machabaeorum 'pertinere' ad vetus 'testamentum'\apnota{12}{M ad vetus 'testamentum pertinere'.}. Quando ergo 's. Hieronymus' canonem Hebraeum descripserit\apnota{13a}{M scripserit, Mc descripserit, U conscripserit.} et eos solos libros valere in\apnota{13b}{M de.} contentione, qui de canone sunt, definiat sitque in hac sententia sua\apnota{14}{U sua sententia.} receptus, facile nostro telo perforabimur, nisi fidelibus\apnota{15}{M U telo verberabimur (Uv perforabuntur) nisi fidelibus (M videbimus).} persuadeamus.

Secundo probat\apnota{16}{M probat.}, librum Machabaeorum esse receptum in canonem, contendit ad aequivocationem, et facile concordabimur Scio, quod ecclesia recipit hunc librum, et hoc\apnota{18}{fehlt in M.} dixi, sed non potest ecclesia plus tribuere auctoritatis aut\apnota{19}{M vel.} firmitatis libro, quam per se ipsum habeat, sicut et ceterorum patrum opuscula\apnota{20a}{fehlt in M.} approbat et reiicit\apnota{20b}{M U recipit.}, sed non ideo omnia confirmat aut meliora reddit, quae in his dicuntur. Transeat ergo ista aequivocatio canon et canon\apnota{21}{M approbat statt omnia confirmat; igitur statt ergo; et canonem 'et canon'; U non ideo confirmat aut meliora reddit. Transeo ergo ista, quae in multis dicuntur canon et canon; Uv = P, nur omnia fehlt.}.

Tertio dicit, auctoritates inductas directe facere mentionem\apnota{22}{M intentionem.} de statu merendi et eundem\apnota{23}{fehlt in M.} morte terminari. Respondeo sicut prius: Quia nihil de purgatorio meminerunt, sed solum de caelo aut\apnota{24}{M U tantum de caelo et.} inferno, ideo moriens ad infernum non meretur et moriens ad caelum non meretur\apnota{25}{M U haben die beiden mit moriens ad anfangenden Sätze umgestellt, M beidemal nihil statt non, et fehlt in M U.}, et sic terminatur utriusque status merendi, alioquin possit eisdem auctoritatibus evidentissime probari purgatorium.

Quarto dicit, me nullo doctore\apnota{27}{U auctore.} inniti et prorogare terminum merendi, item merita purgatorii remunerata asserere. Fateor, hoc enim feci, ut disputarem et audirem meliora, quam ego noverim. Ego enim nihil de purgatorio novi, nisi animas ibi pati iuvandas operibus et orationibus nostris, paratus humiliter doceri, si quae plura de ipso\apnota{31}{M ipsis.} possunt tradi. Quare apostolus contra me adductus\apnota{32a}{U inductus.}, quod solum opera in corpore gesta remunerentur\apnota{32b}{M in corpore gesta opera remunerabuntur.}, relative ad purgatorium acceptus est a domino\apnota{33}{fehlt in M.} D., quomodo iam dixi non posse accipi, sed solum ad caelum vel\apnota{34}{M et.} infernum. Ista relatio si observetur, patet, quod non pugnent contra me, aut demonstrabunt\apnota{35}{M U pugnet contra me aut demonstrabit (M demonstabit) purgatorium.} purgatorium.

Crastina hora septima continuabitur disputatio.

Die Saturni nona mensis Julii hora septima antemeridiana\apnota{36}{Die Zeitbestimmung lautet in M Eccius sabbato die mane, Löscher liest falsch D. Marci, W. A. vermutet 8. Marie; U Nona Julii.}.

D. D. Eccius.

Primo, cum D. pater dicit, auctoritates adductas\apnota{38}{U inductas.} non intelligi de purgatorio, placet, quoad hoc, quia meritum ponunt\apnota{39}{fehlt in U, in Uv nachgetragen.} solum in hac vita, ergo meritum [326] excludunt in purgatorio\apnota{1}{U ergo purgatorium excludunt in hac vita, Uv = P.}. Deinde certum est, Augustinum in Ench. loqui etiam de purgatorio, quia ibidem probat animas suffragiis vivorum relevari. Et Hieronymus increpat eos, qui putant animas exutas mereri, qui purgatorium asserit, Esa. ult. [66 v. 24].

Quod vero dicit, conclusionem suam esse disputationem, quasi non habens de ea magnam fiduciam, miror plurimum, cum iam post annum illam ventilaverit et magnifico titulo disputationem inscripserit contra novos et veteres errores, sententiam quoque recentiorum\apnota{8a}{M recensiorum.} theologorum nec stulto\apnota{8b}{M stultam.} verisimilem affirmat ac contemptim eos theologistas nuncupat\apnota{9}{M nuncupans.}. Quod vero divinatur canonis aequivocationem, non patior\apnota{10}{M patio.}, quoniam Augustinus in illo lib. 18. de civit. dei eundem terminum, maxime in puncto adversativo, non potuit aequivocare, dicendo, quod non fuerit in canone apud Hebraeos, sed apud ecclesiam. Deinde exploratum est, cum plura essent evangelia conscripta\apnota{14}{U scripta.}, auctoritate ecclesiae quatuor in canonem recepta, et sic libros Machabaeorum receptos testatur prologus, tamen ab ecclesia inter divinorum voluminum annotantur historias. — Sed melius ad rem.

Quia se fundat in hoc, quod purgatorium non sit in sacris literis expressum, contra quod est concilii\apnota{18}{U consilii.} Florentini decretum, quod et Graeci abnegato errore assumpsperunt, dein\apnota{19}{U deinde.} plures loci sacrae scripturae, ut illud: Transivimus per aquam et ignem, et eduxisti\apnota{20}{M reduxisti, U perduxisit.} nos in refrigerium [Ps. 66 v. 12]; illud Ecclesiastis\apnota{21}{U Ecclesiastes.} 4 [v. 14]: De carcere et catenis egreditur ad regnum; illud Matth. 5 [v. 25]: Dixit dominus: Esto consentiens etc., ne in carcerem mittaris. Amen dico tibi, non exies inde, donec reddas novissimum quadrantem. Ubi per carcerem intelligit locum purgatorii, ex quo homo non liberatur, nisi plene satisfecerit, sicut interpretatur b. Ambrosius\apnota{25}{U Gregorius, Uv = P.} in loco respondente super Lucam. Ratio tamen Gregorii, quam reassumit b. Bernhardus serm. 66. super Canticis, est valida, quod Christus innuat aliqua peccata remitti in futuro saeculo.

Sed apertissimus textus est 1. Cor. 3 [v. 12—15]: Si cuius opus arserit, detrimentum patietur, ipse autem salvus erit, sic tamen quasi per ignem. Et dies Domini declarabit, quia in igne revelabitur, et uniuscuiusque opus quale sit, ignis probabit, et ibi de stipula, foeno\apnota{31}{M veno.} etc., ubi b. Ambrosius dicit, apostolum esse\apnota{32a}{M U expresse.} locutum de purgatorio. Item\apnota{32b}{M U idem.} testatur glossa interlinearis, glossa quoque ordinaria. Subscribit b. Hieronymus lib. 2 contra Jovinianum, b. quoque Gregorius de purgatorio accipit, quod et b. Bernhardus super Canticis facit. Et ne quis obiiciat mihi Augustinum in Enchir. cap. forte\apnota{36}{fehlt in M U.} 67., qui et de igne praesentis saeculi verba illa posse intelligi ait, legat eundem completius ibidem et quaest. 2 de octo Dulcitii quaestio., et reperiet Augustinum, illa verba apostoli, sicut sunt manifestissima, accipere\apnota{38}{M accipi.} quoque de purgatorio.

[327] Quare nedum in libris Machabaeorum, quod utique ecclesiae sufficeret, sed et in aliis scripturis purgatorium asseritur, nisi Gregorius, Ambrosius, Augustinus, Hieronymus, Bernhardus sint de numero theologistarum; nescio alioquin, quomodo reverendus pater possit fateri, se scire, esse purgatorium, cum literae sacrae scripturae ita valide velit inniti, et cum clarissima sint scripturae testimonia pro neotericis theologis. Ipse nec unum locum in scriptura\apnota{7a}{M scripturae statt in scriptura.} assignare poterit, ubi merita vel gratiam in purgandis\apnota{7b}{M purgantis.} augeri asseratur, cum contra Sapiens dicat, Eccles.\apnota{8}{U Ecclesiastes.} 11 [v. 3]: Et si ceciderit lignum ad austrum aut aquilonem, ad locum ubi ceciderit, ibi erit. Glossa ordinaria: Id est, locum quem hic tibi praeparaveris, tunc habebis. Et\apnota{10}{fehlt in U.} quia in domo patris mei, ait Christus, mansiones multae sunt, certa mansio morienti deputatur, ultra quam suis meritis in purgatorio ascendere nequit. Sic Damasceen. lib. 2. cap. 4 testatur: Scire autem oportet, quod hoc est in hominibus mors, quod in angelis casus, et declarat hoc quoad terminum merendi. Sic Sapiens Ecclesiastis\apnota{15}{U Ecclesiastes.} 9 [v. 10]: Quidquid potest manus tua instanter operare etc.\apnota{16}{fehlt in M U.}, sicut ad illum sensum allegavit devotus ille et dignus ecclesiae minister, Johannes Capistranus\apnota{17}{M cupristanus.}. Ita hoc apostolus docuit ad Galat. 6 [v. 7 ff.]: Nolite errare, deus non irridetur. Quae enim homo seminaverit, haec et metet. Quoniam qui seminant in carne sua, de carne metent corruptionem etc.\apnota{19}{M metet; etc. fehlt in M U.}, bonum autem facientes non deficiamus, ergo dum tempus habemus, operemur bonum\apnota{20}{M oparemur; bonum fehlt in M, U tempus habentes operemur bonum.} ad omnes. Glossa: quod homo seminaverit\apnota{21}{M seminaveris.}, scilicet praesentis vitae labore. Non ergo in purgatorio

seminantur merita, sed in praesenti vita. Quod et Chrysostomus confirmat per illud Joh. 9 [v. 4]: Me operari oportet opera eius, qui misit me, donec est dies, venit enim nox, quando nemo operari potest. Dies, ait Chrysostomus, est praesens vita. Ideo Augustinus adhortatur, nos debere\apnota{26}{M U hortatur nos debere (M deberi).} operari, dum vivamus, ne per noctem praeveniamur. Accedat testis etiam b. Ambrosius, volens apostolum loqui de tempore praesentis vitae nobis concesso\apnota{28a}{M U concessae.}, ut iuste conversemur. Quare egregius psaltes ait, Psal. 104 [v. 23]\apnota{28b}{M 33 statt 103.}: Exibit\apnota{29}{P M U exivit.} homo ad opus suum et ad operationem suam usque ad vesperam.

Postremo, expressa est illa Augustini sententia praeter superius allegata\apnota{30}{M allegatam.}, ubi tamen purgatorii meminit, sententia est eius cap. 3. de fide ad Petrum: Tempus acquirendi vitam aeternam in ista tantum vita deus hominibus dedit. Hieronymus quoque lib. 3. super illud ad Galat. 6 [v. 5]: Unusquisque onus suum portabit, ostendit, quomodo homines meritis suis iuventur in hac vita, sed vocati ante tribunal, quod fit in morte, tunc nec Job nec Noe nec Daniel possint stare pro quoquam, sed unusquisque portabit onus suum. Nam et\apnota{36}{M U ut.} pulchre b. Bernhardus docet, quod\apnota{37a}{fehlt in M U.} in cuiusque hominis morte particulare fiat\apnota{37b}{M U fiat (U fiet) de eo.} iudicium, quod fieri non posset, si adhuc terminum haberet merita sua in purgatorio accumulandi\apnota{39}{U cumulandi.} et gratiam augendi, quia, ut per prophetam Zachariam [4. v. 7] inquit, adaequabit gratiam gratiae, gratiam gloriae reddet in proportione ad gratiam fidei. [328] Quodsi etiam testimonia illa non essent ita clara, quae tamen sunt apertissima, hoc unum deberet reverendum patrem in proposito absterrere, quod theologus absque theologico fundamento pro arbitrio suo\apnota{3}{M U suo arbitrio.} non existens iudex nec iudicis assessor, animabus exeuntibus terminum merendi vult prolongare\apnota{5}{M promerendi vult prorogare.}. Quare quoad istam particulam concludendum videtur, animas in purgatorio satis pati et purgari a delictis, sed maiorem gratiam non mereri.

De secundo principali, quia\apnota{7}{M qua, U quo.} negat, animas in purgatorio esse certas de salute, hoc iterum arbitror nulla sacrae scripturae auctoritate fundatum. In praesentia tamen, quod sint certae de salute, adduco primo\apnota{9}{fehlt in U.} illud Apoc. 5 [v. 1 ff.]: Vidi in dextra sedentis super\apnota{10}{U supra.} thronum Et nemo inventus est dignus aperire librum, nec in caelo nec in terra nec subtus terram. In inferno indubitate nulla est dignitas librum aperiendi, quare de purgatorio loquitur: In quo etiam sunt aliquando s. viri miraculis coruscantes, sicut de s. Paschasio testatur b. Gregorius 4. dialogo, et historia de s. Severino Coloniensi episcopo, natione tamen

Suevo, illud idem asserit\apnota{15}{P M U asserit, Uv accidit.}. Et infra: Et omnes cantabant, in caelo, in terra et subtus terram\apnota{16}{M terra.}. Sed de salute quasi desperantes, non cantant, quare eo ipso, quod sacra anima cantum\apnota{17}{U tantum, Walch XV. 1206 und nach ihm W. A. richtig cantum.} ibi esse asserit, animas certas\apnota{18}{M certas esse.} de salute affirmat.

Praeterea in canone missae, ubi pro defunctis in purgatorio oratur, ita dicimus: Memento, Domine, famulorum tuorum, qui nos praecesserunt cum signo fidei et dormiunt in somno pacis etc., ipsis et omnibus in Christo quiescentibus etc.\apnota{21}{fehlt in M U.} Si ergo iuxta sacratissimum missae canonem animae quiescunt in Christo, quomodo possunt esse in tali turbatione et horrore quasi desperationis, sicut declarat reverendus pater in suo resolutorio concl. 14. usque ad 20.? Non enim video, quomodo in tali horrore, tremore, turbatione et quasi desperatione, quae omnia maximam dicunt inquietudinem, animae purgandae possint dici dormire in somno pacis, quod praedicta omnia paci adversentur. Sed tunc a veris Christianis recte censebuntur in pace quiescere, dum securae exspectant terminum suae purgationis. Haec sunt partim, quae me moverunt et hodie movent ad dissentiendum reverendo patri, paratus tamen stare iudicio et informationi\apnota{30}{M informatione.} aliorum.

D. D. Martinus.

Tria per ordinem egregius D. D. copiosissime contra me prosecutus est, primum canonem scripturae\apnota{33}{Mc 'sacrae' scripturae.} comprehendere libros Machabaeorum, deinde purgatorium probari etiam aliis locis scripturae, tertio conatus ostendere, animas esse certas de salute sua.

Respondeo: De primis duobus nulla\apnota{36}{M U nulla est.} inter me et dominationem suam controversia. Ideo non erat necesse tot auctoritates colligere ad eam [329] rem, quam forte constantius affirmo quam ipse, ut qui me\apnota{1}{fehlt in M.} scire professus sim, esse purgatorium. Sed hic vertitur scopus, an possit\apnota{2}{U posset.} probari, animas in purgatorio mereri et gratiam in eis\apnota{3}{M eius statt in eis.} augeri, tamen transeamus per ordinem.

Primo dixit, auctoritatibus heri adductis statui meritum solum\apnota{4}{M solum meritum.} in hac vita, ergo excludi in futura vita, quod ipse intelligit purgatorium, et sic aptat auctoritaes ad purgatorium. Ego autem concedo de futura vita vel inferni vel caeli. Adduxit etiam Augustin. in Ench. loquentem de purgat. et Hieron. increpantem eos, qui exutas animas mereri asserant. Utrumque concedo. Me etiam reprehendit, quod dixerim,

me disputare\apnota{10}{M quod disputem statt quod dixerim me disputare.} hanc propositionem, quasi non habens fiduciam de ea, qui magnifico titulo inscripserim\apnota{11}{M U scripserim.}, contra novos et veteres errores me disputaturum, sententiam recentiorum\apnota{12}{veriorem.} nec stulto verisimilem affirmarim ac contemptim 'eos' theologistas nuncuparim\apnota{13}{U nominarim; M similem (Mc a. R. veri) affirmari ad contemptim 'eos' theologistas nominari (Mc a. R. appellari).}.

Dico, sicut\apnota{14a}{M quod et statt sicut.} heri dixi, me adhuc nihil scire 'de negotio animarum et purgatorii'\apnota{14b}{Mc de negotiis animarum in purgatorio.} et hoc appellem\apnota{15a}{M appello.} errorem, quod quidam audent asserere\apnota{15b}{M assere.}, quod ignorant, et opiniones patrum ut b. Thomae, Bonaventurae et similium\apnota{16}{et similium fehlt in M.}, quas ipsi non asseruerunt, ipsi tamquam certos articulos fidei pronunciant. Hos ego appello theologistas et non theologos. Opiniones tractari debent in scholis, in populum praedicari debent\apnota{19a}{fehlt in U.} verba et opera dei. Psal. 19 [v. 2]: Caeli enarrant gloriam dei et opera manuum eius\apnota{19b}{M psal. 108, U psal. 8; M U etc. statt et opera manuum eius.}. Non ergo ego\apnota{20}{fehlt in M.} damno opiniones optimorum patrum, sed resisto fabris illis, qui ex opinionibus hominum nobis conflant articulos fidei, quod non est boni theologi officium.

Quod canonem ego aequivocaverim contra Augustinum lib. 18. cap. 36.\apnota{23}{M de trinit. c. 26, U c. 26.} coegit me divus Hieronymus et\apnota{24}{M U item.} Eusebius in historia ecclesiastica recensens antiquorum\apnota{25a}{M U et antiquorum.} auctoritates. Ideo stat aequivocatio, dum\apnota{25b}{U cum.} aliter Augustinus, aliter Hieronymus cum\apnota{26}{U de.} antiquioribus de canone sentiunt, et per consequens nullum robur argumenti in contentione reliquum est\apnota{27}{U relictum est.}. An prologus Hieronymi inter divina volumina libros Machabeorum numeret\apnota{28}{U enumeret, M enumerat.}, non memini, transeo illud, quod quatuor evangelia auctoritate ecclesiae recepta\apnota{29}{M accepta, gestrichen, recepta.} sunt. Haec enim erit alia materia.

Post hoc dicit, contra concilium Florentinum esse, quod purgatorium non\apnota{32a}{Mc quod non, letzteres Wort [rot] durchstrichen.} sit in scriptura expressum. Respondeo, quod\apnota{32b}{fehlt in U.} concilium non potest facere esse de scriptura\apnota{33}{M scripturae facere, U facere de scriptura esse.}, quod non est de scriptura natura sua, sicut nec ecclesia potuit facere evangelia, etiamsi approbavit\apnota{34}{M approbaverit.} evangelia.

Quare videamus auctoritates. Primo illam Psal. 66 [v. 12]\apnota{36}{M prima illa psal. 165.}: Transivimus per ignem et aquam. Respondeo: Non valet ad purgatorium, loquitur de persecutionibus sanctorum, sicut in multis aliis locis: Psal. 17 [v. 3]: Igne me examinasti, item 26. [v. 2]: Ure renes meos

et cor meum, et Petri primo [1. v. 6 f.]\apnota{39}{U 1. Petr. 1.}: Modicum, si nunc oportet contristari in variis [330] tentationibus, ut probatio fidei vestrae multo sit pretiosior auro, quod per ignem probatur\apnota{2a}{M probetur.}. Et breviter, tropus est ille vulgatissimus\apnota{2b}{M 'iste' est vulgatissime, U est iste vulgatissimus.} scripturae, per ignem et aquam\apnota{3}{M per aquam et ignem.} intelligi tribulationes. Ideo nimis haerent in literis et syllabis, quod mihi imponit egregius D. D., qui ignem pro purgatorio accipiunt\apnota{4}{U ideo nimis haeret . . . . accipit.}.

Item et illud Ecclesiastis: quod aliquis de carcere et catenis\apnota{5}{M catena.} egreditur ad regnum\apnota{6a}{M etc. statt egreditur ad regnum.}. Satis clarus est\apnota{6b}{U est clarus.} textus de vanitate huius mundi disputantis, quod casu vertente is, qui nunc servus est, fit rex, et qui rex\apnota{7}{M his statt is, qui rex statt et qui rex, U et qui rex est.} fit servus. His et similibus auctoritatibus, si contra negantes pugnaremus, nihil nisi ludibrium de nobis et ecclesia adversariis faceremus, quamquam ego libentissime haec omnia admittam\apnota{10}{M haec omnia libentissime admitto, U libentissime omnia haec admitto.}, et si qua sunt similia.

Tale et illud est Matth. 5 [v. 25]: Esto consentiens adversario tuo, ne tradat te tortori\apnota{12}{M U tortori etc.}, ubi per carcerem dicit intelligi per Ambrosium locum purgatorii, libenter admitto. Sed quia alii patres in diversum exposuerunt, praesertim Augustinus de inferno et nunquam egressuro exponit, fit auctoritas dubia, fidelibus suadens, resistentes\apnota{15}{U resistens.} non convincens, ut id taceam, quod textus consequentia ne patiatur quidem de purgatorio intelligi, dicit enim de consentiente adversario et dissentiente, volens dissentientem damnabiliter et contra Christi praeceptum peccare, quae\apnota{18}{P qui.} pertinent ad infernum, non ad purgatorium.

Rationem Gregorii consentiente Bernhardo dicit esse validam, quod deus remittat\apnota{21}{U remittit.} aliqua peccata in futuro. Respondeo: Verum est apud fideles, sed invalida apud resistentes.

Post haec dicit, esse apertissimum textum prima ad Cor. 3 [v. 15]: ipse\apnota{23}{U si cuius, ipse.} salvabitur, sic tamen, quasi per ignem, et Ambrosium et Hieronymum, glossam\apnota{25}{U glossa.} ordinariam et interlinearem, Gregorium, Bernhardum intellexisse apostolum de purgatorio, licet Augustinus alicubi\apnota{26}{M alicui.} de praesentis saeculi igne eadem verba posse intelligi dicat. Respondeo: Adeo non est apertissimus hic textus, ut ego hodie fatear, cum multa excusserim, me adhuc nescire\apnota{28}{M non scire.} germanum sensum Pauli, ita variantibus interpretibus, quamquam pro me libens admittam de purgatorio. Verum cum apostolus apertissime dicat, opus uniuscuiusque per ignem probari, in quo igne dicit revelandam diem Domini, quae declaret uniuscuiusque

opus, ut sunt aperta verba, meo tenui iudicio videtur loqui de igne conflagrationis et extremi iudicii, aut ut Augustinus tropologizat\apnota{34}{M tropologisat, U tropologice.} de igne persecutionis temporalis, qua maxime probatur fides et fidei doctrina, et quaecunque super haec aedificantur.

Quare adhuc nihil habetur ex sacris literis\apnota{36}{M ex sacra scriptura.} manifestum de purgatorio, quod in contentione valeat. Non ergo beatissimos patres de numero theologistarum habeo, quia purgatorium cum eis confiteor, nec ipsi opiniones et ignorantias suas de statu animarum pro articulis sanxerunt, sicut faciunt theologistae.

[331] Miratur D. D.\apnota{1}{fehlt in P.}, quomodo possim scire, esse purgatorium, et ego nec unum locum habeam\apnota{2}{M habeo.} pro me in scripturis, cum haec, ut dicit, clarissima scripturae 'testimonia' pro neotericis theologis\apnota{3}{M pro theonericis scripturis (so!).} facere putet. Respondeo: Non est necesse confiteri, qua via noverim purgatorium, vel quaecunque. Deinde haec\apnota{4}{M dehinc statt deinde haec.} testimonia clarissima pro purgatorio sunt adducta, non pro statu animarum in purgatorio.

Jam illa auctoritas Eccles. 11 [v. 3]\apnota{7a}{P M decimo.}: In quocunque loco\apnota{7b}{fehlt in P M.} ceciderit lignum, sive ad austrum sive ad aquilonem\apnota{8}{M sive austrum sive aquilonem.}, ibi erit; ubi glossa ordinaria intelligit: locum, quem hic tibi praeparaveris, tunc\apnota{9}{M tibi hic praeparaveris tunc; tunc fehlt in U, in Uv nachgetragen.} habebis, nescio, quo ingenio ad propositum ducatur; si enim per locum praeparatum et habendum\apnota{10}{U et humidum statt et habendum (W. A. et habitum).} intelligit purgatorium, manendum erit in aeternum in purgatorio, si autem intelligit, quod hic\apnota{12}{fehlt in M, U quod si hic.} praeparaveris, id est merueris, tunc non facit ad auctoritatem Eccles. 11, qui nihil de merito loquitur, sed de morte hominis.

Transeo illud, quod Christus dicit, in domo patris sui esse multas mansiones, et\apnota{15}{M in domo patris mei sunt mansiones multae; 'et' et (so!)} quod certa mansio morienti deputatur, ultra quam ascendere nequit. Totum pro me est. Scio, quod unicuique mansio deputatur post mortem, sed 'non' statim pervenit ad mansionem post mortem, nisi mansionem iterum pro purgatorio accipiat et sic morientes ab aeterno regno excludat\apnota{18}{M accipit et sic morientes a regno coelorum (Mc 'ab regno aeterno') excludat; U excludit, Uv excludat.}.

Item et illud Damasceni: Hoc hominibus mors, quod angelis casus. Respondeo: Si ergo post mortem ceciderint\apnota{21a}{M U ceciderit.} in purgatorium, sequitur, quod in aeternum purgabuntur\apnota{21b}{U purgabitur, Uv purgabuntur.}, aut oportet Damascenum, sicut et omnia praecedentia, intelligi de duabus ultimis mansionibus, damnationis et beatitudinis, et non de purgatorio.

Ita et illa Eccles. 9 [v. 10]: Quodcunque poterit manus tua instanter operare, quia nec ratio nec sensus apud inferos est, quo tu\apnota{25}{M quod tu, Mc quo tu.} properas. Si de purgatorio intelligit, iterum purgatorium erit infernus, ergo solum de inferno, citra memoriam purgatorii, loquitur.

Johannem Capistranum libenter admitto, sed extra contentionem\apnota{28}{M conventionem.}.

Jam quod apostolum Gal. 6. [v. 7] dicit contra nostram sententiam loqui: Quae homo semina verit in\apnota{30a}{fehlt in P.} carne, haec et metet, et illud Joh. 9 [v. 4]\apnota{30b}{U haec metet; P M U Johan. 8.}: Venit nox, quando nemo poterit operari\apnota{31}{M operare.}, quod Ambrosius de praesentis tempore vitae exponit, quo\apnota{32a}{M de praesentis temporis vitae exponit in quo, U de praesenti temporis vita [Uv tempore vitae] exponit, quomodo.} meremur, clarum est, quod nihil faciat\apnota{32b}{M U facit.} ad rem. Apostolus enim non de purgatorio, sed de extremo iudicio loquitur. Et Christus per noctem salva\apnota{34}{M U salva tamen.} reverentia Chrysostomi intelligit infidelitatem, ut clarum est ex proximo textu: quamdiu ego in mundo sum\apnota{35}{M sunt, Mc a. R. 'sum'.}, tam diu lux sum mundi volens, quod extra fidem Christi nemo possit\apnota{36}{M potest, Mc 'possit'.} bene operari, tamen Chrysostomum admitto, quia respicit ad extremum iudicium\apnota{37}{M indicii.}.

Auctoritas Psal. 104 [v. 23]: Exibit homo ad opus suum usque ad vesperam, permittitur trahi figurativo sensu ad vitam et mortem hominis, sed genuino sensu, et qui pugnet in contentione, loquitur de admirabili dispositione dei, [332] quod hominem sic disposuerit, quod mane exeat usque ad\apnota{1}{U ita disposuerit quod homo exeat ad; usque fehlt in M.} vesperam diei naturalis ad opus suum\apnota{2}{M U operationes suas.}. Et Augustinus\apnota{3}{P Augustinus.} recte sentit 'capite 3. de fide' ad Petrum, tempus acquirendi vitam\apnota{4a}{U viiam, Uv vitam.} aeternam esse datum hominibus tantum\apnota{4b}{U tamen.} in ista vita, quia ut dixi, semper ad futuram vitam respiciunt\apnota{5}{M Lücke hinter ad, Mc 'praesentem (gestrichen) vitam futuram respiciunt'.}, non ad purgatorium.

Et Hieronymus ad Galat. 6 [v. 5]: Unusquisque onus suum portabit, bene\apnota{6}{M Et quod Hieronymus, sonst = P; U Et Hieronymus super Galat. 6. Unusquisque portabit onus suum etc. et bene.} ostendit, quod homines meritis suis iuventur in hac vita, sed vocati ante tribunal, quod fit in morte, tunc nec Job nec Noe\apnota{8}{nec Noe fehlt in U.} nec Daniel possunt stare pro quoquam. Respondeo, haec auctoritas robustissima est\apnota{9}{M quod haec auctoritas est robustissima est (so!), U quod auctoritas est robustissima.} contra egregium D. D., quia si verum est, hominem in morte vocari ante tribunal, et tunc nec Job nec Noe\apnota{11}{nec Noe fehlt in M U.} nec Daniel possunt stare, frustra orat ecclesia pro mortuis. Et sic negatur purgatorium. Quare de tribunali novissimo loquitur Hieronymus, post finitum purgatorium.

Bernhardum etiam induxit, quod in hominis morte agatur particulare iudicium, quod non fieret, si haberet terminum adhuc merendi. Respondeo: quidquid sit de particulari iudicio, consequentia non valet.

Transeo auctoritatem Zachariae [4 v. 7]: Adaequabit\apnota{17}{P adaequavit.} gratiam gratiae, quia tropologico et bono sensu exponitur, sed non proprio.

In fine huius articuli dicit\apnota{19}{Mc fügt bei 'Eccius'.}: Etiamsi non essent tam clara haec testimonia, hoc tamen deberet me absterrere theologum, quod sic\apnota{20}{M testimonia haec, hoc tamen me deberet abstinere (Mc 'absterrere') theologum, quod (sic fehlt, Mc 'si').} absque fundamento pro meo arbitrio, cum nec iudex nec assessor fuerim\apnota{21}{M fuerit, verbessert fuerim; U fuerim et tamen.}, animabus 'exeuntibus' terminum 'prorogarim\apnota{22}{Mc 'prorogarem'.} merendi'. Idem retorqueo in D. D.\apnota{23a}{U ipsum D. D.}, cum nec ipse iudex nec assessor\apnota{23b}{M assessor nec iudex.} fuerit, et tamen animabus certum statum iuxta opiniones suas definiat\apnota{24}{M diffiniat.} absque fundamento, praesertim cum id nolit esse opinionem\apnota{25}{M opinionem esse.} dumtaxat (quod libentissime tolerarem), sed certam scientiam\apnota{26}{M Lücke, Mc 'tolleraret sed certam sententiam'.}.

Hora secunda pomeridiana continuabitur disputatio.

Hora secunda pomeridiana eiusdem diei.

D. D. Martinus continuavit coepta sua, ut sequitur\apnota{29}{M die Zeitbestimmung fehlt, U Hora secunda continua: disputatio per eundem d. Martinum, qui dixit se declaraturum recte et candide de scripturis sentire.}.

Et quia mea sententia non satis intellecta videtur egregio D. D., quantum possum apertius 'me' exponam\apnota{31}{M U expono.}, quod scripturae et s. patres antiqui in suis sententiis respiciunt futuram vitam, qua vel salvantur\apnota{32}{M samnantur.} vel damnantur animae, nihil de purgatorio interim cogitantes. Ideo multae illae auctoritaes, quae abnegant statum merendi post hanc vitam, non pertinent ad purgatorium. Quod si adhuc non intelligitur, fingatur purgatorium nullum esse, sicut non fuit eius cogitatio in mentibus eorum. Tunc adhuc stabunt et verae erunt omnes auctoritaes inductae, quod ante mortem est vita merendi. Deinde non sic de meritis loquor animarum in purgatorio, quod aliquid operentur, sed quod recipiant gratiam ampliorem. Cum apud omnes constet, [333] culpam non remitti, etiam venialem sine augmento gratiae et in purgatorio culpas remitti scribat divus Gregorius in dialogorum 4. et recitatur dist. 25.\apnota{2}{M dist. 24.} cap. Qualis, sic tamen haec\apnota{3a}{U hoc.} assero, ut pro opinione habeam,

imo ut fatear\apnota{3b}{M habeant imo ut fateor.}, pro ignorantia. Soli deo credo cognitum esse statum animarum in purgatorio.

Ad tertium principale de scientia salutis earum, ubi induxit\apnota{5}{M adduxit.} primo Apocalypsim 5. capitulo [v. 3]\apnota{6}{M quintum capitulum, U quintum caput.}, quod nemo inventus est nec in caelo nec in terra nec subtus terram, qui dignus esset aperire librum, volens per subtus terram intelligere purgatorium, ut in quo 'aliquando' et\apnota{8}{U et aliquando.} sancti viri fuerunt, sicut Paschasius et Severinus:

Dico quod haec glossa auctoritatem non habet, ideo eadem facilitate contemnitur, qua probatur. Fateor quidem, sanctos in purgatorio fuisse, imo nullum in purgatorio esse\apnota{11}{M esse in purgatorio.}, nisi sanctum. Posset ergo contentiosus dicere, subtus terram esse idem, quod infernus vel quodlibet aliud, ut quod nec daemones nec homines nec angeli possint\apnota{13}{M quod quod nec daemones nec angeli nec homines possit (so!)} aperire librum, sicut dicitur, trina rerum machina colere deum: et apud apostolum tam caelestia, quam\apnota{15a}{fehlt in P, in Uv gestrichen.} terrestria, quam inferna genu flectere\apnota{15b}{M flectunt.}; siquidem et daemones contremiscunt, imo, ut in eodem libro modus est, subtus terram significat etiam mortuos\apnota{17}{U mortuus.}, ut dicit in alio cap. [20 v. 13]: Terra dedit mortuos suos, et infernus dedit mortuos suos\apnota{18}{M infernus dedit . . . . et terra dedit.}.

Quod autem sequenter inducit: omnes cantabant\apnota{19}{M contabant, Mc cantabant.} in caelo et in terra et subtus terram, desperantes autem non cantare posse, respondeo: Non dixi animas desperare unquam, sed quando una ecclesiae sententia eadem est poena inferorum et purgatorii, dixi eas esse similes\apnota{22}{U similes esse.} desperatis, quomodo leguntur et in hac vita quidam tentati fuisse desperatione, ut in Psalterio\apnota{23}{M temptatos . . . ut psalterium.} in multis locis: Ne avertas faciem tuam a me, et assimilabor\apnota{24}{M U assimilabor.} descendentibus in lacum [27 v. 9; 28 v. 1]. Deinde cantare\apnota{25}{M cantare, Mc cantari.} Domino non semper est laetari et gaudere, imo canticum novum\apnota{26}{U nonum.} est canticum crucis, hoc est, laudare et portare deum in mediis tribulationibus atque adeo in morte.

Quod tertio adduxit canonem missae, ubi pro defunctis orantes dicimus: Dormiunt in somno pacis et quiescunt in Christo, non videt, ut dicit, quomodo possint\apnota{30a}{M possunt.} 'dici' dormire in somno pacis, si sint in tali turbatione\apnota{30b}{fehlt in M.}, tremore et inquietudine\apnota{31}{U inquitudine, Uv inquietudine} maxima. Ideo interpretans hanc quietem pacis dicit, quod exspectant secure terminum suae purgationis. Respondeo: Non satisfacit mihi haec glossa, et idem per idem probatur. Nam esse eos inquietos probat id, quod sequitur in canone:

Ipsis, Domine, et omnibus in Christo quiescentibus locum refrigerii, lucis et pacis indulge, et illud, quod omnes oramus: Requiem aeternam dona eis, Domine. Item: Dona eis pacem, quod non potest intelligi, dona eis, exspectare secure terminum suae purgationis, quod D. D. interpretatus est esse pacem. Rectius ergo meo iudicio requiescunt in pace, quoad corpus. Dormire enim in sacris literis in pace significat quiescere in sepulcro. [334] Et sic patet, quod egregius D. D. non probavit scientiam salutis eorum, imo quodammodo tollit poenas purgatorii, dum\apnota{2}{M cum.} eis tribuit pacem.

D. D. Eccius.

Quia nobis constitutum est artum tempus hanc materiam hodie finiendi, ut pro qualitate materiae solutiones nihili refutare\apnota{5}{M nihil non reputari.} non possim et alia puncta purgatorii attingere, paucula quaedam afferam, ut facile agnoscatur vitasse reverendum patrem materiae nostrae medullam et diverticula quaesivisse. Unde primo cum Augustinum induxissem, omne meritum hic comparatur\apnota{9}{M comparatum.}, et quod nemo speret, quod hic neglexit, cum obierit, apud deum promereri, similiter et Hieronymum, dicit se utrumque concedere. Sed quae tunc est ista obstinatio\apnota{11}{M est tunc; U obstinatia.}, si Augustinum concedit, omne meritum hic comparari, ut dicat gratiam in animabus, etiam in purgatorio augeri?

Neque valet in praesentia collyrium suum ad omnes auctoritates, quas dicit non loqui de purgatorio, nam b. Augustinus illud idem dixit, cum de animabus in purgatorio tractaret, et in Enchiridio et in libro de cura pro mortuis agenda c. 1. Nam post mortem non possumus agere, quod prosit\apnota{16}{U prodest.}, sed recipere, quod egimus. In hanc enim inciderat\apnota{17}{M inciderit.} Augustinus difficultatem, quomodo suffragia mortuis prodesse possent, cum nihil mererentur?\apnota{18}{M possint quia nihil mererent, U possunt, quia nihil mererentur.} Recurrit ad hoc, ut in vita meruerint, quod eis post mortem\apnota{19}{M pro mortem.} prodesset, quare frustra nitebatur me instruere de sensu\apnota{20}{M U de sensu suo.}, quem satis bene accepi. Nec persuasiuncula sua quicquam momenti habet apud Christianum, quia venialia non remittantur sine augmento gratiae. Probet hoc reverendus pater. Nulla auctoritate hoc est fulcitum, imo est falsum: cum peccatum veniale dei offensam non faciat ex dei misericordia, non opus est nova gratia in eius deletione, sed sufficit, quod aliquis pro eo satis patiatur.

Quod vero secundo loco excusat se de errore, quia non video unam literam adductam per D. patrem pro sua nova doctrina: ideo merito

debet reputari\apnota{28a}{M suputari, Mc reputari.} suspecta, cum recentiorum theologorum\apnota{28b}{M neotericorum doctorum.} sententia tot sacrae scripturae et sanctorum patrum nitatur auctoritatibus. Nec hoc excusat, quod aliquibus imponit facere\apnota{30}{P M facere se.} opiniones Thomae vel Scoti sicut articulos fidei, de quo mihi non constat. Nam\apnota{31}{fehlt in U, in Uv nachgetragen.} ego non Thomam vel Scotum, sed clarissima ecclesiae lumina cum sacra scriptura adduxi, viderit ipse, quos theologistas incuset\apnota{32}{M incusat.}.

Tertio ad Augustinum dicit de libris Machabeorum, fortiorem opponendo b. Hieronymum, at Hieronymus nullibi negat libros Machabaeorum apud ecclesiam esse de canone bibliae, quin in prologo hoc constanter asserit fuisse annotatum inter divinorum voluminum historias. Quare opusculis s. patrum in can. Sancta\apnota{37}{fehlt in P M U.} Rom. 15. dist. non debuit aequiparari.

Cum vero ei opposuissem de concilio Florentino, respondit\apnota{38}{U respondet.}, concilium non posse [335] facere aliquid esse de scriptura, quod non sit. Hoc quidem verum, sed quid hoc est, concilium\apnota{2}{U consilium.} tam laudabile tanta temeritate contaminare, ut hoc absurdum decernat? Cum vero doctissimi fuerint in eo concilio viri, malo credere concilio, quod a spiritu sancto regitur, quam domino Luthero, non quod concilium faciat aliquid esse\apnota{5}{M concilium faciat esse aliquid, U consilium faciat aliquid (esse fehlt).} de scriptura, quod non sit, sed quod credam concilium\apnota{6}{M concilium melius, U consilium melius.} habere sensum et intelligentiam scripturarum decernendo hoc esse de scriptura, quod in scriptura reperitur.

Evanida\apnota{8}{M evanida, Mc evalida.} est sua excusatio, posse reperiri semper expositionem, ita, quod textus allegatus non valeat in contentione contra pertinaces, nam illud esset omnium haereticorum labitulum, qui semper aliquem expositionis fucum possent afferre, quo\apnota{11}{M U quibus.} contenderent catholicas veritates in scriptura sacra non esse expressas. Ita hodie adhuc perfidi Arii duraret haeresis, quod homousia\apnota{12}{M omusia.} ex sacris literis tam expresse probari non posset, quin in contentione pertinax qualicunque fuco evaderet\apnota{14}{M U non evaderet.}.

Similiter bene novimus, ecclesiam non posse facere evangelia, tamen ecclesia facit, ut relictis Nicodemi, Bartholomaei, Thomae et aliorum evangeliis, quatuor dumtaxat indubitatam fidem adhibeamus, in quo ecclesiae iudicio standum in evangeliorum electione\apnota{18}{M U standum est in evangeliorum acceptatione.}. Ita et in sacrarum scripturarum intelligentia et expositione.

Porro assumens auctoritates per me inductas ab initio, voluit praeter necessitatem hoc factum, quia et ipse crederet purgatorium esse,

non meminit, quod ex tot\apnota{22}{M ex hoc.} auctoritatibus terminum merendi in praesenti vita statuentibus voluit se evolvere, quia purgatorium ex sacris literis non probaretur, quod dictum, ne Picardi et alii schismatici pro se arriperent, qui non sciunt purgatorium esse, sicut d. pater, merito fuerat a me confutandum. Dimissis aliquibus duos saltem conspiciamus locos.

Matth. 5. [v. 25] noluit per carcerem intelligi purgatorium, quia et Augustinus, quod minime me latebat, per carcerem infernum intelligat\apnota{28}{M U intelligit.}. Addo ultra D. patrem, Chrysostomum intelligere praesentem vitam, sed rectiorem esse Ambrosii sententiam vel saltem non contemnendam, qui per carcerem purgatorium intelligit. Verba Christi hoc indicant, debere illum solvere usque ad novissimum quadrantem\apnota{32}{P U dodrantem.}, sed in inferno nulla est solutio, sicut nulla est redemptio. Neque fucus additus Ambrosii expositionem repellere\apnota{33}{U ad Ambrosii expositionem refellere.} potest, quod in carcerem trudendus peccaverit mortaliter, quoniam et propter peccata venialia et propter peccata mortalia, tamen contrita, in purgatorio puniuntur, quod ex Machabaeis accipimus, dum inquit: Sancta et salubris ergo est cogitatio pro mortuis exorare\apnota{37}{M U sancta (M sanctam) et salubris est ergo (ergo fehlt in M) cogitatio pro defunctis exorare (Mc 'orare pro defunctis', M Lücke).}, ut a peccatis solvantur. Hi enim, qui occisi fuerant, et pro quibus Judas Machabaeus oblationes fecit\apnota{38}{M fecerat.}, peccaverant mortaliter propter spolia idolorum\apnota{39}{M spolia ablata.}, quamvis credantur poenituisse in [336] ipsa caede, 'iuxta' glossam ordin. ibidem, et illud Psal. [78. v. 34]: Cum occideret eos, quaerebant eum.

Quarto, dum apertissimum Pauli apostoli locum 1. Cor. 3. [v. 12ff.] adduxissem, quem Augustinus, Ambrosius, Hieronymus, Gregorius, Bernhardus, Isidorus cum glossa ordin. intelligunt de igne purgatorii, fatetur reverend. pater, se germanum sensum huius loci non habere, ideo contemptis tot patribus novam attulit sententiam, apostolum loqui de igne conflagrationis, quia apostolus meminerit diei Domini, et quod uniuscuiusque opus ignis probabit; ego, qui semper audivi, antiquioribus maxime sanctis credi debere, s. patrum exosculor sententiam et novam illam glossam et verba D. Lutheri non accipio, nisi probet sacrae scripturae auctoritate. Nec iuvant ea, quae ponderat\apnota{12}{M ponde dat, Mc ponderat.} in praesentia, quae et s. patres bene legerunt. Quamvis enim peculiariter dies iudicii extremi dies Domini dicatur, ut Bernhardus quoque meminit, tamen in cuiusque hominis morte, cum iudicium fert\apnota{14}{U fit.}, dies Domini dici potest, neque tam exactam et captiosam accipiat distributionem, ut sophistae facere solent\apnota{16}{M solem, Mc solent.}, in vocula uniuscuiusque. Sed de his accipiat, qui stipulam, lignum et foenum\apnota{17}{M venum.} superaedificarunt, sicut et August. complicat distributionem Johan. 1 [v. 9]: Illuminat\apnota{18}{P illuminet.} omnem hominem venientem in hunc mundum. Unde invicta est apostolica sententia, eos, qui super fundamentum aedificant lignum, foenum\apnota{20}{M venum.}, stipulam, salvos fieri, sed per ignem purgationis, ut sic nullum malum maneat impunitum.

Ad illud Ecclesiastis\apnota{22}{P U Ecclesiastes.} de casu ligni respondet: Si loquatur de purgatorio, sic probari perpetuo animam manere in purgatorio. Dico optime sensisse glossam, per lignum intelligi hominem, sicut\apnota{24}{U hinc Uv sicut.} et Ezech. 31. [v. 8] rationalem creaturam significat: Omne lignum paradisi etc.\apnota{25}{fehlt in M.} et per casum mortem, nec tamen sequi ideo perpetuo remanere in purgatorio, sed, ut egregie exponit b. Hieron., per austrum bonum, per aquilonem malum et damnatio significatur\apnota{28}{M damnationem significari, U damnationem significare.}. Quare sapiens non loca illa intellexit, sed statum, quod in casu vel est bonus et ita perseverabit sine augmento gratiae ob peccatorum deletionem, si est malus, perdurabit in malitia\apnota{30}{U malignus statt in malitia.}.

De mansione apud Joh. [14 v. 2] dixit, statim post mortem certam animae deputatam mansionem. Quomodo ergo erit certa, si fiat accessio maioris gratiae? Nam tunc necessario accedet\apnota{33}{M non tunc necessario accedat.} altior mansio, consectarium quoque esset, eum, qui ad purgatorium descenderet, cum pluribus venialibus melius habere quam descendentem cum paucis vel morientem cum nullis, quia maior fieret ei gratiae accessio ad plurium venialium deletionem. Addo, quod noxium et damnabile esset orare pro mortuis et plus prodesset eis, si in poenis purgatorii perdurarent propter maioris gratiae accessionem, sicuti alicubi scripsit reverendus pater, animam alienis suffragiis\apnota{39}{M suffragii.} liberatam minus beari, quam si per se in purgatorio satis pateretur. Quod est non solum contra sacram [337] scripturam: Sancta ergo et salubris est\apnota{1}{fehlt in U.} cogitatio etc., sed et contra observantiam et pietatem totius ecclesiae, quam mortuis impartiri debemus, ut b. Augustinus pulcherrime libro de cura pro mortuis agenda\apnota{3}{M habendam.} et lib: Enchiridio explanat.

Ad Damascenum porro respondit, eum non loqui de purgatorio, quoniam sic perpetuo durarent in purgatorio. Et ego sentio, nec Damascenum nec alias auctoritates sonare de purgatorio, quoniam alioquin mihi obessent et D. patri prodessent. At cum de praesenti loquantur vita, in qua terminum ponant\apnota{9a}{M ponant.} merendi et demerendi, non potest porrigere\apnota{9b}{U porrigi.} tempus merendi ultra terminum constitutum, scilicet mortem. Quare non aliam accipio expositionem, quam auctor ipse velit, qui de termino merendi in loco allegato loquitur.

Ita et de aliis dicere possum neque simili obiurgatione increpari possum sicut reverendus pater, quod\apnota{14}{M quo.} auctoritate s. patrum et sacrae scripturae cursum et merendi terminum cum morte finiam. Ipse autem nulla\apnota{15}{M finiant .... nullo.} fulcitus sacrae scripturae auctoritate eis terminum porrigit et prorogat, contra Augustinum et contra Hieronymum, quamvis suae sit modestiae, ut hanc dicat sententiam suam esse opinionem, sed tunc non debuit oppositum vocare errorem.

De Hieronymi auctoritate\apnota{19}{M U De auctoritate Hieron.} conatus est eam in me retorquere, nec Noe nec Job nec Daniel stare pro quoquam post mortem, quia ibi de extremo loquatur iudicio, quod\apnota{21}{M U quod tamen.} ex litera non convincit. Non solum enim in extremo iudicio, sed statim etiam\apnota{22}{M etiam statim.} post mortem unusquisque onus suum portabit.

Postremo ad aliam accedens particulam, respondit ad illud Apoc. [5 v. 3ff.], quod per sub terra intelligitur\apnota{25}{M U sub terram intelligatur.} infernus. At tunc sequens de cantico stare nequit, quia daemones et damnati non cantant, sed blasphemant et eiulant\apnota{26}{M plasfemant et egulant.}. Quare per sub terra\apnota{27}{M U sub terram.} necessario purgatorium et non inferos intelligere potest. Quamvis non negem in\apnota{28}{M quin.} sacra scriptura etiam daemonum fieri mentionem in inferis, ut in apostolo caelestium terrestrium et infernorum\apnota{29a}{M in inferis ut coelestia etc., U in inferis in apostolo coelestia terrestria, Uv = P.}. Sed hoc nihil ad rhombum\apnota{29b}{M romphum.}.

Apportavit autem glossam super illo verbo »cantat«, quod sit crucem portare et in tribulationibus deum laudare, quam glossam, ut a D. patre hoc loco fictam, eadem\apnota{32}{P M U ea.} facilitate contemno. Imo Johannes non patitur sic se exponi, qui\apnota{33a}{P M U quia.} inquit: Cantabant in caelo, non in terra\apnota{33b}{non fehlt in P, M cantabo in coelo et in terra; U cantabunt (Uv cantabat) in coelo in terra, wie oben zuerst bei Walch, 1223.}, modo in caelo non portant crucem nec in tribulationibus laudant, quoniam abstersit deus omnem\apnota{34}{M U laudant deum quoniam abstersit (U abstergit) deus omnem (U sic omnem).} lacrimam ab oculis sanctorum et iam non est labor etc\apnota{35}{M ab oculis eorum etc. (Die Worte von et iam an fehlen), U ab oculis eorum iam non est nullus labor et dolor etc.}.

Quodsi etiam daretur expositio sua, ut\apnota{36}{M U et.} esset vera, sicut mtaen literae repugnat, adhuc hoc propositum non rumperet, nam laudantes deum in tribulationibus post mortem certissimum haberent futurae beatitudinis signum, cum damnati contra consummato odio dei in perpetuam prorumpant\apnota{39}{U prorumpunt.} creatoris blasphemiam. [338] Quare animae purgandae\apnota{1}{fehlt in P.} de salute securae cantant et deum laudant, terminum purgationis expectantes.

Cum vero induxissem canonem missae, quo eos dormire in somno pacis asseritur, non patitur illationem meam: si in somno pacis, ergo sunt securi\apnota{4}{U securae.} de salute et non in horrore, 'tremore', pavore et quasi desperatione. Rationem duplicem adduxit\apnota{6a}{U adduxi.}, primo, quia sequatur, ut eis locum refrigerii, lucis\apnota{6b}{P et lucis.} et pacis indulgeas\apnota{7}{M indulgeas etc.}, secundo, quia frequenter oramus: Requiem aeternam dona eis, Domine, et: Dona eis pacem. Ideo aliam glossam dedit, quod dormiant somno\apnota{8}{U somnum.} pacis, scilicet quoad corpus. Non contentor\apnota{9}{M contento.} nec obiectione nec nova sua glossa\apnota{10}{M glossam, Mc glossa.}, quae nullam habet apparentiam, nullum quoque fundamentum. Nam dicit canon: Memento eorum etc., qui dormiunt in somno pacis, et infra\apnota{11}{U eorum, qui dormiunt (in somno pacis fehlt); et infra fehlt in M.}: Et omnibus in Christo quiescentibus. Non dicit: In sepulcro quiescentibus. Porro quid refert ad animam, corpus habere pacem vel minus, sive in mare iactetur vel\apnota{13}{M sive mari iactetur ac, U sive a mari iactetur sive.} a feris laceretur, quemadmodum beatus\apnota{14}{fehlt in M, U d. statt b.} Augustinus testatur? Ecquis hominum patitur corpus exanime et exsangue dormire in pace? Unde distinctiones in neotericis theologis non libenter admittit, et tamen egregie hic de pace distinguit corporis et animae.

Quare 'verior amplectenda est\apnota{17}{M esse.} sententia, animas purgandas et dormire in pace et quiescere in Christo. Unde quies illa et pax turbationem animae, horrorem et quasi desperationem excludit, cum inquietissimae sint\apnota{19}{U sunt.} animae, quae hac turbatione mentis agitantur, ut quasi desperent.

Ad inducta vero non impendio respondemus: Orare nos et obsecrare deum, ut det eis pacem, non pacem, quam\apnota{22}{M quem.} modo habent, securitatis scilicet de salute, sed pacem aeternam, quoniam pax est, ut Bernhardus testatur, et a culpa et a miseria. Det ergo deus animabus in purgatorio pacem a poenis, sicut contulit eis gratiam contra peccatum\apnota{25}{U peccata.}.

Quare adhuc vera et inconcussa est veritas, animas in purgatorio non quasi desperare de salute.

D. D. Martinus.

Contra solutiones replicat egregius D. D. ex Augustino dicente\apnota{29}{M dom. egreg. doctor ex Aug. dicentem.}: Omne meritum hic comparatur\apnota{30}{M comparatum.}, atque agente de purgatorio, quod post mortem solum recipiamus, quod egimus. Respondeo brevissime: Ergo non est purgatorium, aut purgatorium est\apnota{32}{M U erit.}, quod recipitur pro vita praeterita. Quare patet, Augustinum non posse de purgatorio intelligi. Non enim hoc egerunt in vita praeterita, ut purgatorium reciperent. Quamquam et haec auctoritas mihi patrocinetur, quod meritum, quo merentur iuvari\apnota{35}{M iuventur iuvari.}, hic meruerunt. Ego autem fateor intelligere me non posse, quomodo iuventur animae in purgatorio et liberentur absque ullo munere gratiae, per solam ablationem poenarum: intelligat, qui potest.

Deinde dicit, hoc falsum\apnota{38a}{M falum, Mc falsum.} esse, venialia 'non' remitti sine\apnota{38b}{M sine, Mc absque.} augmento gratiae.

Quando ergo vult D. D., quod venialia possint\apnota{39}{M U possunt; quod venialia fehlt in M.} remitti absque 'augmento gratiae', non concedo, donec probet. Imo est contra expressum\apnota{40}{fehlt in M.} textum, ad Roma. 7 [v. 14ff.]\apnota{41}{P Rom. V.}, Gal. 5 [v. 17], ubi apostolus dicit: Venumdatus sum sub peccato, quis liberabit me de corpore mortis\apnota{42}{M de corpore mundi, U de morte corporis.} huius? Respondet: Gratia per Jesum [339] Christum\apnota{1a}{M respondit: gratia dei Jesu Christi, U = P, nur gratia dei.}. Certum est autem, apostolum non fuisse in peccato mortali\apnota{1b}{M esse, verbessert fuisse, M U in peccatis mortalibus.}. Cetera relinquo disputationi futurae\apnota{2}{M ad disputationem futuram, U futurae disputationi.} inter egregios D. doctores Andream et Iohannem Eccium\apnota{3}{M Johannem et Andream.}. Et hoc erroneum est, quod veniale peccatum deum non offendat, cum displiceat deo omne immundum et reprobet eum, qui minima solverit, Matth. 5 [v. 18f.]\apnota{5}{Matth. 5. fehlt in M.}. Nec probavit D. D. sufficere, quod satis patiantur.

Secundo criminatur, quod nec unam literam pro mea sententia induxerim, cum ipse pro recentiorum theologorum\apnota{7}{M recensiore; theologorum fehlt in M U.} sententia tot scripturae et patrum induxerit auctoritates. Respondeo: Eo magis mihi suspecta est\apnota{8}{fehlt in U.} eius sententia, quo plures induxit, quia per capillos et obtorto\apnota{9}{M optorto.} collo adduxit, sicut in praecedentibus satis dixi. Minus peccat, qui dubitat in cogitationibus suis, quam qui dubia sua verbis divinis nititur statuere\apnota{11}{M statuere, verbessert constituere (construere?)}.

Dicit etiam, non constare sibi, qui pro articulis fidei opiniones Thomae vel aliorum statuerint, quia ecclesiae lumina cum scriptura adduxerit. Monet ergo, quod videam, quos appellem\apnota{14}{M U ut statt quod; M appellam.} theologistas. Dixi hodie et iterum dico: Permitto opiniones esse opiniones neminique hoc vitio verto, sed trahere repugnantem scripturam in aliorum sensum pro firmandis\apnota{16}{M U in alienum sensum pro formandis (U confirmandis).} opinionibus et in hoc pertinaciter stare, hoc dico theologistarum esse officium.

\noindent\textit{[nota editorial sin llamada en el original, al inicio del aparato: coepit d. Martinus latine prosequi. W. A. 299, 3–11 fehlen in M, a. R. [rot] fecit sermonem vulgarem.]}

\noindent\textit{Notas del aparato sin llamada en el cuerpo:} [15b] M cephale, U Kephale, a. R. κεφαλη.; [15] M dici.

Ad aliud, de canone librorum, ubi nixus Hieronymi et concilii Florentini\apnota{19}{fehlt in M.} auctoritate mavult credere concilio, quod a spiritu sancto regitur, quam mihi, et\apnota{20}{W. A setzt dafür ei.} gratias ago. Pie enim sapit, nunquam volui mihi credi. Sed respondeo breviter: Conciliet ipse primo\apnota{21a}{U primum.} Hieronymum sibi, qui\apnota{21b}{M quia.} in prologo galeato\apnota{22a}{P Uv allegato.} Machabaeorum libros et nonnullos alios manifeste inter apocrypha\apnota{22b}{M apographa.} recenset, qua auctoritate fit, ut mihi liber Machabaeorum sit gratus et probatus, sed contentiosis pateat ad\apnota{24}{M pati ad. Mc schreibt über 'fateor'.} repulsam. Satis hodie de concilio dixi, imo, ut ipsemet D. D. dicit: Concilium non errat, si autem errat, non est concilium, et ut meo sensu loquar, credo concilium et ecclesiam nunquam errare in his, quae sunt fidei, in ceteris non est necesse non errare.

Quod vero exclamat, latibulum hoc esse haereticorum, qui hac fiducia quaslibet expositiones auctoritatum refutarent vel acciperent\apnota{29}{M quemlibet librum negarent, gestrichen vor quaslibet (vom Disputator widerrufen); M U vel reciperent, Mc a. R. vel acceptarent.}, respondeo: Quare hoc non sunt conquesti\apnota{30a}{M non sunt conquesti hoc.} s. Augustinus, Hieronymus aliique victoriosissimi\apnota{30b}{M victoriosissimi.} haereticorum triumphatores, sed donatis ambiguis locis studuerunt certis et apertis locis pugnare? Hoc faciamus et de animabus in purgatorio.

Ad rem et auctoritates\apnota{33}{et auctoritates fehlt in M U.}.

Ad solutiones auctoritatum\apnota{34}{U ad auctoritatum solutiones.} dixit, a me contemptos esse Augustinum, Hieronymum, Ambrosium, Bernhardum, Gregorium, Isidorum\apnota{35}{M Ambrosium Hieronymum; Isidorum fehlt; U Aug. Ambros. Bernh. Hieron. Gregor. Isid.} cum glossa ordinaria\apnota{36}{M ordinariis.}. Hoc pro modestia, imo molestia Ecciana dixit, nimis cupidus movendae invidiae\apnota{37}{M invidii, Mc invidiae.}. Dixi hodie, non esse contemptos a me, et iterum dico.

Quod autem meam glossam non accipit, non curo. Probet autem ipse et suam: cum textus sit clarus de die Domini et igne, in\apnota{39}{fehlt in M.} quo revelabitur dies Domini. Quod etsi potest trahi ad purgatorium, ut dixi et sic sententiam D. doctoris [340] non reprobavi, purgatorium scilicet esse, non tamen potest hac intelligentia oppilari os contentiosum.

Item et illud Matth. 5. [v. 25] de carcere non contempsi pro Ambrosii sententia. Sed\apnota{4}{fehlt in U.} quod addit D. D. verba Christi indicare, debere solvi usque ad novissimum quadrantem, ideo de inferno intelligi non debere: hoc\apnota{5}{M posse, übergeschrieben debere, beides gestrichen, daneben deberi; M U respondeo hoc.} satis eluitur per divum Hieronymum, qui adversus

Helvidium\apnota{6}{P Uv Helpidium.} hanc dictionem »donec« recte exponit, quod\apnota{7}{M U ut.} non cogatur significare, quod Ambrosius cum D. D. hic voluit. Nam sic et Matth. 1. [v. 25] scribitur\apnota{8}{U hic ponit nam et sic . . ., M U scribitur de Joseph.}: Nec cognovit eam, donec peperit filium suum, et tamen non sequitur, quod post partum cognoverit eam, ita hic non sequitur, ut\apnota{10}{M U quod.} sit exiturus post solutionem, sed solvet et non exibit.

Quod autem confutavit, quod non peccaverit mortaliter dissentiens adversario\apnota{12}{M dissentio ab adversario.}, atque quod et mortalia, modo\apnota{13a}{U non.} sint contrita, et venialia purgentur\apnota{13b}{M U purgantur.} in purgatorio: dico, quod hic textus non loquitur de contritis mortalibus, nam contritum mortale iam\apnota{15}{M nam, Mc iam.} non est mortale, et poenitens iam non est dissentiens adversario\apnota{16a}{M ab adversario.} suo, ergo de dissentiente et sic permanente loquitur Christus\apnota{16b}{M deus.}.

Transeo illud de casu ligni apud Ecclesiasten, ubi Hieronymum inducit\apnota{17}{M U dicit.} intelligere per austrum statum bonum et per aquilonem statum malum, concedo utrumque, et nihil ad purgatorium.

De mansione apud Iohannem certa, ubi argutatur in hunc modum: Quomodo fiet certa, si fiat accessio gratiae, cui debeatur altior mansio? potest idem argui de quolibet fideli post baptismum\apnota{22}{M baptisma.}, cui est certa mansio ab aeterno praedestinata.

Deinde in purgatorio existentes sic ordinati sunt, ut hoc modo ad certam mansionem veniant. Illud etiam humanum argumentum nihil concludit, quod sequeretur descendentem cum pluribus venialibus\apnota{25}{fehlt in M U.} peccatis melius habiturum esse quam eum, qui cum paucis. Si haec ratio bona est, timendum est, ne meretrix melior fiat\apnota{27}{P fieret.} quam b. virgo, quasi non sint differentes gradus animarum in purgatorio, cum ipse hodie sanctos velut prae ceteris praestantiores in purgatorio asseruerit.

Deinde Silvestrinum argumentum inducit, quod prodesset eis, si perdurarent in poenis, nam sic et martyribus expediret usque in diem iudicii mori et pati, quasi ignoret D. D. ad certam mensuram poenas esse animabus constitutas. Non enim\apnota{33}{M U ergo.} damnabile est, orare pro mortuis, sicut nec damnabile fuit, quod apostolus pro se oravit et orari petiit, cum tamen cresceret virtus eius in infirmitate. Ita quilibet fidelis pro qualibet necessitate cuiuslibet fidelis debet orare, debet adiuvare, non obstante, quod ille per necessitatem istam magis ac magis mereatur.

Per idem volo solutum, quod contra observantiam ecclesiae et\apnota{38}{M et ecclesiae} pietatem sit, quod dixi animam alienis suffragiis liberatam minus beari, quam si per [341] se satis\apnota{1}{fehlt in M.} patiatur in purgatorio. Hoc velut corollarium sumpsit ex dictis meis, ubi ego sensi animabus esse succurrendum, quantumlibet perfectius genus meriti esset iustitiae divinae omnibus modis satisfacere.

Damascenum transeo cum termino meriti et remitto ad praecedentia.

Item, iterum me carpit, quod nullo fulcitus testimonio prorogo terminum meriti, ipse autem nitatur multis et neget. Respondeo: Nullis nititur et ipse nisi violenter intellectis\apnota{7}{M intellexisset.}, ut satis dictum est.

Capere etiam me in verbis meis volens dicit, me appellasse sententiam meam opinionem, ergo male oppositum errorem appellaverim\apnota{9}{M U appellaverim errorem.}. Dico sicut prius. Non solum opinionem, sed et ignorantiam appello meam sententiam, errorem appellavi non contrariam opinionem, sed quod opinionem pro veritate statuunt\apnota{12}{M statuunt proferunt.}.

Deinde auctoritatem Hieronymi de Noe, Job et Daniele\apnota{13a}{Hieronymi fehlt in M U, de fehlt in M; M U Daniel.} dicit valere non\apnota{13b}{U non valere.} tantum pro extremo iudicio, nec possim hoc convincere ex litera. Relinquo hoc iudicio melius sentientis, cum manifeste textus\apnota{15a}{U textus manifeste.} habeat, quod pro homine ad tribunal\apnota{15b}{M tribunal.}, quod in morte fit, rapto nec Noe nec Daniel nec Job\apnota{16}{M raptum statt rapto, sonst = P; U rapto nec Job nec Daniel.} orent, quod omnino de iudicio extremo intelligi oportet, aut saltem non de purgatorio, quia pro mortuis in purgatorio orant Noe, Job, Daniel et omnis ecclesia\apnota{18}{M U Daniel et (M etc. et) omnes ecclesiae.}.

Ad auctoritatem Apocal. recte dicit, quod daemones sub terra\apnota{19}{M ad tertium (?), Mc sub terra, U in inferno, Uv = P.} 'non cantant, sed blasphemant etc.'\apnota{20}{mit sub terra non cantant beginnt in M der Passus, der aus dem Notariatsprotokoll abgeschrieben ist, am Rande Mc [rot] verba notarium iuratorum. M liess eine Seite dafür frei; da sie nicht reichte, sind 4 Blätter (Signatur post o ante p) später eingeheftet. U non laudent (Uv non cantant), sed blasphemant deum.}.

D. D. Martinus exclusus tempore, cum iam duae horae transissent, obtulit se ad reliqua in scriptis tradenda notariis et primum D. doctori Eccio exhibenda. Obtulit autem\apnota{22}{M autem sequenti luce.} infrascripta per notarium suum\apnota{23}{in U lautet die Zwischenbemerkung: Reliqua exclusus tempore doctor Martinus vult signare in scheda et ostendere domino doctori et dominis notariis. Obtulit autem sequenti loco (Lö sequenti luce) infra scripta.}.

Quod daemones in inferno non laudant, sed blasphemant deum, scio, sed quod ideo subtus terram significet purgatorium, in quo cantent animae, non sequitur. Primum, quod sunt in tribulationibus et poenis necdum abstersit deus omnem lacrimam ab oculis earum\apnota{27}{M U eorum.}. D. D. autem negat, cantare esse in tribulationibus laudare deum. Ideo

contra se ipsum loquitur, dum subtus terram in purgatorio\apnota{30}{M subtus terram (U motus terra, Uv = P) in purgatorio subtus terram.} laudem ponit simul et poenas, in quibus non laudent. Nego autem et ego, quod animae ideo cantent, quia habent certissimum signum beatitudinis. Hoc enim signum certissimum fuit probandum. Deinde ego sub terram non dixi significare solum infernum, sed et alia, quaecunque sub terra sunt, in quibus omnibus laus dei abundat, omnia enim sunt plena gloria et laude dei, qui laudabilis est in omnibus operibus suis. Alioquin oportet D. doctorem non solum purgatorium per subtus terram intelligere, sed etiam aliquod quartum per subtus mare, cum eodem Apocal. 5. etiam subtus mare recitetur et omnia, quae in eo sunt. Rectius ergo creaturas omnes hoc quaternario intelligimus significatas, quam ea, quae D. doctor\apnota{38}{fehlt in M U.} Eccius imaginatur.

[342] Quod\apnota{1}{hier beginnt die eingeheftete Lage von 4 Blättern, s. o.; Me (rot) quia per schedam sunt haec notariis praesentata.} autem in tribulationibus laudare deum sit cantare, probo non esse a me fictum, ut D. D. dicit, quia Rom. 5. [v. 3.] dicit: Gloriamur\apnota{2}{M gloriamini.} in tribulationibus, et Psal. 42 [v. 9]: In die mandavit Dominus misericordiam suam et nocte (id est, tribulatione), canticum eius, et Jac. 5 [v. 13]: Tristatur aliquis, oret aequo animo et psallat; item: Benedicam Dominum\apnota{5}{P domino.} in omni tempore [Ps. 34 v. 2]. Omni, inquit tempore, etiam adverso. Contra vituperatur ille, de quo Psal. [49 v. 19]: Confitebitur tibi, cum benefeceris ei. Fateor autem hoc me non dixisse de damnatis subtus terram, sed ut excluderem sententiam D. doctoris volentis, quod cantare sit tantum de beatitudinis securitate in purgatorio.

De canone missae non est contentus, quod dormire in pace ad corpus retuli, quia dicit canon: In Christo quiescentibus, non in sepulcro quiescentibus. Porro quid referat ad animam, corpus in mari aut aëre versari ac non quiescere? Item, quis (inquit) corpus exanime dicat dormire in pace? Respondeo: Nihil ista argutia valet, nam quiescentes in sepulcro aut ubilibet secundum corpus iactati vere quiescunt in Christo, hoc est ut Apocal. 14 [v. 13]: requiescunt\apnota{16}{U requiescant.} a laboribus suis, videlicet quod separata anima a corpore iam non laborat in corpore in variis huius mundi molestiis. Alioquin quomodo Christi corpus in sepulcro negabit dormire in pace et requiescere\apnota{18}{M U quiescere.} in deo, cum in tot locis de eo sic loquatur\apnota{19}{M loquitur.} scriptura et cantet ecclesia: In pace in id ipsum dormiam et requiescam [Ps. 4 v. 9]; et Psal. 16 [v. 9]: Et caro mea requiescet in spe. Item Apocal. 14 [v. 13]: Beati qui moriuntur in Domino. Item Stephanus Act. 7 [v. 59]\apnota{21}{P M U act 8.}: Obdormivit in Domino, quod sine dubio de corporali morte tropo scripturae dicitur. Si Eccio

non dormit corpus exanime, sinat, quaeso, illud dormire suscitatori deo et scripturae auctoritati\apnota{24}{M U suscitaturo deo: M auctoritatem.}.

In fine dicit, orari\apnota{25}{M orare.} a nobis pro pace, non quam habent, sed quam habere a miseria poenarum contingat purgatis. Respondeo: Haec\apnota{26}{M U hoc.} est petitio principii, quia pax illa securitatis, quam habeant, nondum est probata. Relinquitur ergo, quod in poena sunt, et ut ab hac solvantur et pacem habeant, oremus.

Haec omnia sic volo dixisse, ut ostendam me ignarum esse eorum, quae in purgatorio agat deus, et nimio\apnota{30}{U nimis, Uv et nimis.} eos audere, qui huius ignorantiae impatientes potius fingere volunt sua, quam confiteri sese\apnota{31}{U se.} ignorare. Mea probare non possum, sed nec illi sua. Quare recte mea conclusio id tantum voluit, esse nec scripturis nec rationibus probatum, animas certas esse saltem omnes de salute sua et gratiam in eis non augeri, sufficiat\apnota{34}{M sufficimus, U sufficit.}, quod scimus, eas pati et nos debere eis succurrere, cetera deo relinquenda soli.

Acta sunt haec in praesentia honorabilium virorum Johannis Teuber Weysmann et Bartholomaei Schaller Ernfriedsdorfensis, Bambergensis et Misnensis dioecesis, laicorum testium ad praemissa requisitorum specialiter et vocatorum. Die solis decima mensis Julii in solito disputationis loco, hora meridiana vel quasi\apnota{36}{M Wassmann statt Weissmann, ad praemissis statt ad praemissa, am Ende fehlt vel quasi: U hat statt der ganzen Angabe nur Acta sunt haec 10. Julii praesentibus etc.}.

D. D. Eccius sequentia obtulit in scheda quadam Die lunae undecima mensis Julii mane hora septima\apnota{37}{U Eccius per schedam respondit 11. Julii. Löschers Anm. p. 436 „in M Sto additur hora prima" beruht auf Irrtum.}.

Solutiones reverendi patris apparent mihi insufficientes, potissimum quod Augustinum dicat non loqui de purgatorio, cum maxime hoc\apnota{39}{U id.} agat, scilicet [343] de purgatorio. Neque negandum est, animas defunctorum pietate suorum viventium relevari, cum pro illis sacrificium mediatoris offertur vel eleemosynae in ecclesia fiunt, sed haec eis prosunt, qui, cum viverent, ut haec sibi postea prodesse possent\apnota{4}{U possint.}, meruerunt. Est enim vivendi modus nec tam bonus, ut non requirat illa\apnota{5}{U ista.} post mortem, nec tam malus, ut ei non prosint post mortem. Quocirca hic omne meritum comparatur, quo possit post hanc vitam relevari quispiam vel gravari. I nunc, Luthere, et dic Augustinum hic non de purgatorio locutum.

Praeterea non probat, ad deletionem peccati venialis requiri augmentum gratiae. Fateor, augmentum gratiae tollere aliquando veniale\apnota{10}{U venialia.}, sed etiam veniale tollitur satispassione. Unde apostolus allegatus Rom. 7. [v. 14 ff.] non loquitur de veniali peccato, porro si loqueretur, satispassionem non excluderet. At de eius consensu, velut impertinente, nihil dico in praesentia. Veniale, fateor, offendit, sed venialiter, quia hominem in odio dei non constituit nec inimicum dei facit, quod alioquin cum gratia non\apnota{15a}{fehlt in U, in Uv nachgetragen.} staret, ideo in purgatorio quidem habent\apnota{15b}{M U habent quidem.} peccata venialia citra inimicitiam dei.

At quod causatur, me torquere auctoritates per capillos ad propositum, dispeream, si tota disputatione vel unam adduxerit auctoritatem tam pertinentem, sicut ego hic adduxi. Sed hoc iudicent iudices, non Lutherus.

Indignatur, cur non et Augustinus et Hieronymus causati sint hoc latibulum haereticorum, quod dixi; bone deus, si\apnota{21}{W. A. verbessert quasi.} non ubique clamet Augustinus\apnota{22}{M U Augustinus et Hieronymus.} contra Pelagianos, Cyprianus et Ambrosius contra Arianos et Novatianos, eos pro arbitrio interpretari sacram scripturam!

Petit, quod probem meam sententiam, cum suam probare non possit, meam vero probatam iudicum relinquo iudicio.

De »donec« alia est ratio Hieronymi ad Helvidium, quia aliunde probatur ex scriptura, Mariam semper mansisse virginem, igitur donec ibidem non dicit consummationem, quod hic non potest facere.

Porro de rationibus meis gratia brevitatis non moror. Sed quod dixit\apnota{29}{U dicit.}, me certum fecisse ex dictis suis, quod animae suffragiis viventium liberatae minus beentur, male mihi imponit, nec enim mihi\apnota{31a}{M nisi.} somnia fingo. Refert\apnota{31b}{U reperietar, Uv refert.} enim in dilutione adnotationum mearum, noluisse Paschalem sua sibi minui praemia, ideo maluisse ardere, sed missa haec facio et iudicibus committo.

Porro, non nego, quin in tribulationibus quis laudet deum\apnota{34}{M eum}, sed hoc loco accipi non potest, quod et in caelo cantare eos dicit, cum ab illis absterserit\apnota{36a}{M et coelo cantare eos dicit, ab illis abstergcrit (so!), U = P, nur et statt cum.} deus omnem lacrimam\apnota{36b}{U lacrimam etc.}. Unde liquet lectori, quomodo sententiam meam\apnota{37}{M suam, verbessert meam.} invertat. Unde animae in purgatorio cantant, secundum Johannem, deo, et quaero, causam explicet incertitudinis\apnota{38}{M U nisi certitudinis.} gloriae; quare falsissime comminiscitur\apnota{39}{M comminiscatur.}, quod contra me ipsum loquar. Unde non nego, bonos gloriari in tribulationibus et cantare, sed negavi illam glossam, quam ipse finxit, quod [344] cantare sit in tribulationibus deum laudare, alioquin angeli non canerent deo et beati.

Quod errorem patentissimum defendere conatur, scilicet animas dormire somno pacis quoad\apnota{4}{U scilicet quoad.} corpus etc., quoniam anima separata a corpore

non laboret\apnota{5a}{M U laborat.} in corpore variis molestiis: vide vafritiem, quia\apnota{5b}{M vabriciem; M U qui.} cum solvendo ad corpus retulisset in sepulcro quiescens\apnota{6}{M quiescente, U quiescentes.}, iam se transfert ad quietem animae a laboribus corporis. Sed heus bone vir, est hoc quiescere a laboribus, quia videlicet non in corpore molestias patiatur, tamen incomparabiliter maiores molestias sentiat extra corpus in anima? quid iuvat, si quiesco a febri et infestor calculo?

De Christi quiete\apnota{11}{M quete.} aliud est, resurrectionem absque corruptione expectante. Beati ergo, qui moriuntur in Domino. Sed quis dicturus est corpus beatum nisi per redundantiam animae in corpus? Sic et\apnota{13}{fehlt in U.} de b. Stephano, qui in Domino obdormivit.

Deinde\apnota{15}{M U demum.} inquit me petere principium: Sunt in pace, et tamen eis petimus pacem. Et\apnota{16}{W. A. verbessert sed.} quia ridiculum est, referre pacem habitam ad corpus, sequitur, illam esse in anima scientiam\apnota{17}{M iam statt illam; P U scientia.} securitatis.

Quod ait, se sua probare non posse, reputo verissimum, contra alii sua probant, quamvis pertinax et cervicosus etiam demonstrationibus non acquiescat, quoniam semper contenderet\apnota{20}{M U quin statt quoniam; M contenderit.}, non esse demonstrationem, velut Gregorius Arimin. et Petrus Aliacensis Aristoteli faciunt in rationibus de primo motore in 7. et 8. Physicorum. Quare cum communi concludo sententia, et in purgatorio gratiam non augeri, et eas\apnota{24}{W. A. verbessert animas.} certas esse de salute.

Die lunae undecima mensis Julii hora septima coepit disputari materia de indulgentiis\apnota{26}{U De indulgentiis. Die undecima Julii, quae fuit dies lunae. Lö. p. 438 Anm. „In M Sto additur hora septima".} D. D. Eccius\apnota{27}{M Egreg. D. D Eccius Damit endet der aus dem Notariatsprotokoll abgeschriebene Passus in M. Es folgt eine leere Seite mit der Notiz vacat. Die nun folgende Lage von 8 Blättern müsste die Signatur „P" tragen; da aber das 1. (leergelassene) Blatt fehlt, so ist auch die Signatur weggefallen. Das 2. Blatt beginnt: De indulgentiis Eccius.}.

Imprimis, antequam descendam in hanc disputationem indulgentiarum, protestor, non esse mentis meae aut propositi\apnota{29}{U propositum.} velle hac mea disputatione contravenire mandatis summi pontificis, quibus mandat in decreto incipiente: Cum postquam, non debere defendi vel praedicari contrarium certis punctis de indulgentiis sub poena excommunicationis, sed magis esse intentionis meae, defendere veritates in eodem decreto approbatas\apnota{33}{M approbare, U veritatem . . . . approbatas (so!).}. Quo sic praefato, accedo ad impugnandum\apnota{34}{M U impugnanda m.} conclusionem undecimam.

Dicere indulgentias esse utiles Christianis, est verum et pium, nec indulgentiae sunt vitium boni operis, quare non videtur bene dictum, quod [345] hoc affirmantes insaniant. Hoc probatur, quia ecclesia in iis\apnota{1}{M his.}, quae sunt fidei, non errat, et quae respiciunt animarum salutem. Imo, ut s. Cyprianus ait, deus non permittit maiorem partem cleri errare, sed in conciliis generalibus et per universitatem cleri iam a trecentis annis indulgentiae reputatae sunt utiles et piae Christianis. Quod patuit in concilio Viennensi\apnota{5}{M binensi, U Wienensi.}, ubi approbatae sunt indulgentiae\apnota{6a}{fehlt in U, in Uv nachgetragen.} datae per Urbanum IV. pro venerabilis eucharistiae veneratione, nam concilium\apnota{6b}{U consilium.} hanc addidit rationem motivam, ut Christi fideles essent magis parati ad obsequia venerationis et honoris sacramento impendenda.

Sic Innocentius III. in concilio generali, ubi utilissimam constitutionem: Omnis utriusque a tota ecclesia receptam\apnota{10}{M recepta.} edidit, de indulgentiis quoque disposuit hospitaliorum, dicens: Iniungimus vobis in remissionem peccatorum, quatenus de bonis vobis collatis grata eis charitatis subsidia\apnota{12}{M de bonis collatis grata charitate subsidiis, U = P, nur subsidia charitatis.} erogetis, ut per subventionem\apnota{13a}{M subfensionem.} vestram ipsorum inopiae consulatur et vos per haec bona\apnota{13b}{U et vos praeponat.} et alia, quae Domino inspirante feceritis, ad aeterna possitis gaudia pervenire. In eodem quoque concilio limitata fuit inferiorum praelatorum indulgentias dandi\apnota{16a}{M danda, verbessert dandi.} potestas, quae in Lugdunensi concilio\apnota{16b}{U in Lugduniensi consilio.} deinde fuit approbata. Sed si indulgentiae essent vitium boni operis et Christianis inutiles, quid opus esset collationem vitii operis minuere in inferiori praelato et conferre in plenitudinem potestatis?

Accedit, quod laudabile Constantiense concilium\apnota{20}{U consilium.}, quod et inter errores damnavit indulgentiarum contemptum\apnota{21}{M contemptor, Mc contemtum.}, contulit indulgentias omnibus in concilio existentibus et feria sexta ieiunantibus, quod fuisset bona opera tantorum patrum inquinare et inficere. Et Gerson, quem reverendus pater illustrem theologum appellat, sicut fuit veri et honesti studiosissimus\apnota{25}{M honesti studii honestissimus.}, is decidit, indulgentiarum concessionem non esse parvi pendendam vel\apnota{26}{M seu, U pervidendam seu.} contemnendam, sed devote amplectendam in fide, spe et charitate Domini nostri Jesu Christi, qui talem potestatem clavium ecclesiasticarum dedit hominibus. Constat enim, quod fructuosior est et deo acceptabilior operatio talibus indulgentiis innitens, quam altera et hoc\apnota{29}{M et ob, U et.} ceteris paribus non innitens. Haec Gerson.

Sic et alii sancti et optimae\apnota{31}{M optimi.} existimationis viri tempore datarum indulgentiarum voluerunt, quae a b. Gregorio quoque datae sunt ante

nongentos annos et Paschasio\apnota{33}{M Paschasium.} ante sexcentos annos. Unde sola tanti patris auctoritas debet quemlibet Christianum movere, ut credat indulgentias esse utiles\apnota{34}{M utiles esse.} Christianis, quamvis non ignorem\apnota{35a}{M ignoro.}, quod Vuilhelmus Altisiodorensis\apnota{35b}{M Altioridensis, U Altisidorensis.} referat non defuisse, qui etiam vivente b. Gregorio sancto patri in indulgentiis contradixerint.

Porro hoc addendum\apnota{38}{M addendo, U addam.}: Si indulgentiae sunt vitium boni operis, hoc potissimum meo arbitratu ex eo suspicari posset\apnota{39}{U possem.}, quod sunt satisfactoriae, sed [346] per opus satisfactorium non minus meremur, quam si non esset satisfactorium, alioquin praestaret nos nulla facere opera satisfactoria, ne merita nostra pro isto statu minueremus. Sed in benedictionibus seminantes de benedictionibus quoque meteremus\apnota{4a}{U metemus.}. Sed hic audiam\apnota{4b}{M audeam.} reverendum patrem. Accedat tantus totius Christianitatis consensus in iubilaeis per summos pontifices celebratis, per Bonifacium\apnota{6}{M celebratos statt celebratis; W. A. per Bonifacium 8.}, Clementem VI., Urbanum VI., Nicolaum V., integerrimum et doctissimum pontificem, Sixtum IV.\apnota{7}{M Sixtum IV etc.}, et quod communi consensu Christianissimorum\apnota{8}{M christianorum.} regum et principum pro passagiis et cruciatis summus pontifex plenissimas\apnota{9}{fehlt in M.} saepe dedit indulgentias.

Cum ergo fides ecclesiae sit indefectibilis, pro qua Christus rogavit, ut non deficeret, Lucae 22, et pollicitus est, se nobiscum esse omnibus diebus usque ad consummationem saeculi Matth. 28. [v. 20], non est admittendum, ecclesiam cum animarum pernicie tanto tempore errasse, et quia decreta summorum pontificum ab omnibus communionem ecclesiae sanctae habentibus sunt acceptanda, ut ait Gregorius can. Praeceptis, 12. distin., item can. Omnia decretalia, 25.\apnota{16}{P M U 24.} quaest. 1. cum multis concordantibus, et pontifex modernus sub poena excommunicationis latae sententiae definivit, indulgentias esse utiles ad poenam pro peccatis debitam his\apnota{18}{M U poenam peccatis debitam his (U iis).}, qui charitate iungente membra sunt Christi, delendam, et quod huiusmodi remissio ex merito Christi et sanctorum, auctoritate summi pontificis, quasi compensetur: his omnibus persuasum volo Christiano fideli, indulgentias non esse vitium boni operis, nec dicere indulgentias Christiano utiles esse\apnota{22}{M U esse utiles esse.} insanire, quod nihil aliud est, quam dicere, totum clerum in ecclesia tanto tempore insanivisse.

D. D. Martinus\apnota{24}{P D. D. Eccius.}.

Duodecim argumentis impugnat egregius D. D. conclusionem meam undecimam. Ad quae priusquam respondeo\apnota{26}{M U respondeam.}, primo diluam illud,

quod insaniam esse dixi, habere\apnota{27}{M U haberi.} indulgentias bonas Christiano esseque eas vitium boni operis. Dico ergo, quando propheta Psal. 40. [v. 5] audet insanias falsas appellare, et Psal. 119. [v. 85] iniquorum fabulationes ipsas traditiones hominum, quae de se non erant malae, nisi quod malam fiduciam adversus iustitiam dei praestabant insipientibus: quanto magis insania est indulgentias, quae neque praeceptae neque consultae\apnota{32}{M U consultae sunt.} neque necessariae ad salutem, imo remissiones sunt bonorum multorum operum, haberi pro bono Christianorum? cum omne bonum nostrum et necessarium ad salutem comprehensum sit in praeceptis et consiliis. Imo amplius dico: Quando saluberrima dei lex, teste Ezechiele et apostolo Paulo Rom. 7. [v. 11 ff.], non est bona homini, imo occasio peccati, sola autem gratia bonum est Christiano, quanto minus\apnota{37}{P M U magis, in W. A. verbessert.} indulgentiae, quae nedum cum gratiae, sed nec cum legis ullius bonitate ullo modo comparari possunt\apnota{38}{U possint.}, [347] cum sint nihil et privatio quaedam bonorum operum, haberi debent pro bono? Non quod noxias aut perniciosas arbitrer\apnota{2}{M arbitremur.}, ut dicemus.

Quando ergo egregius D. D. dicit primo, ecclesiam 'in his'\apnota{3}{U iis.}, quae sunt fidei, non errare, nec in his\apnota{4}{U iis.}, quae respiciunt animarum salutem, recte dicit, sed tales non sunt indulgentiae, ut dictum est.

Quod autem Cyprianus dicit, deum non permittere, ut\apnota{6}{M U quod.} maior pars cleri erret, viderit ipse, quid dixerit. Certum est, quod tempore Arianae\apnota{7}{M Arriani.} perfidiae maior pars et eloquentiorum episcoporum et doctiorum virorum erravit in tanto articulo fidei, adeo ut nullus catholicus episcopus ferme resideret in cathedra sua.

Quod dicit, in conciliis generalibus et a trecentis annis habitas\apnota{11}{U iam statt et a; habitis.} indulgentias utiles et pias Christianis et in concilio Viennensi\apnota{12}{M binensi, U Wienensi.} approbatas, quas Urbanus IV. pro reverentia venerabilis eucharistiae in fidelibus\apnota{13}{M infidelibus statt in fid.} excitanda largitus est, respondeo: Nunquam negavi esse utiles indulgentias, sed non Christianis, hoc est, fervide agentibus, Christum secundum nomen suum amantibus et quaerentibus, his\apnota{16}{U iis.} enim ingrata est remissio operum, grata vero impositio operum. Secundo dico, quod in his\apnota{17}{U iis.} rebus non est respicienda persona ullius\apnota{18}{M unius.} sive concilii sive papae, sed quid dicatur, praesertim quando hic error de indulgentiis, ut in re non necessaria, non est periculosus, ceteris tamen paribus, hoc est, si sciant esse\apnota{20a}{U eas esse.} non necessarias nec in eas confidant\apnota{20b}{M perfidant.}.

Ad secundum, de concilio generali et Innocentio III.\apnota{21}{M de concilio generali Innocentii.}, qui iniunxit pro dispositione hospitaliorum, ut in remissionem peccatorum conferrent

grata subsidia etc., dico me adhuc non intelligere hoc verbum papae, cum remissio peccatorum mihi videatur esse aliud quam indulgentiae\apnota{24}{U indulgentias.}, atque si esset idem, dico sicut prius.

Ad tertium, de Lugdunensi\apnota{26}{M luculenso, gestrichen, a. R. luetnensi.} concilio, in quo approbata dicitur et limitatio potestatis conferendi indulgentias inferiorum praelatorum, admitto totum. Quod autem sequitur: Si essent inutiles, non fuisse necessarium collationem vitii operis minuere in inferioribus praelatis\apnota{29}{M in inferiori praelato.}; dico indulgentias esse vocatas a me vitium operis\apnota{30}{Mc 'boni' operis.}, secundum quod impugnatus sum a D. D. conclusione mea 42. in Resolutorio, ubi sic dixi: Non quod veniae sint malae et noxiae, sed quod abusus perversus nocet, dum tale opus non facerent, nisi veniae essent. Sic enim finis operis huiusmodi fit ipsa venia, et clarissimum est, quod raro homines contribuerent, nisi indulgentiae promitterentur. Ideo semper est ibi\apnota{35}{M U est ibi.} periculum, saltem vitiosi operis.

Quarto induxit laudabile Constantiense concilium, quod inter errores ceteros damnavit contemptum indulgentiarum. Respondeo: Nunquam ego contempsi aut contemnendas esse docui, nisi contemptus intelligatur hoc modo, quod incomparabiliter meliora possimus\apnota{39}{M U possumus.} eisdem expensis facere, quibus indul-[348]gentiae redimuntur, vel quam sunt indulgentiae ipsae. Quare praerogativa melioris non est contemptus deterioris, sicut aurum non est contemptus ligni aut foeni\apnota{2}{M veni, verbessert feni.}, etsi praestantius.

Per idem ad quintum, ubi Gersonem inducit, indulgentiarum concessionem non esse parvi pendendam, addo, in suo genere, seu sicut ipse dixit\apnota{5}{M dixerit.}, ceteris paribus. Nec credo hanc sententiam D. doctori satis placere, quod indulgentiae sint suscipiendae in fide, spe et charitate, cum tenuerit satisfactionem impleri et remitti posse etiam his\apnota{8}{U iis.}, qui sunt extra gratiam. Admitterem, quod fructuosior sit operatio innitens indulgentiis quam altera non innitens ceteris paribus, modo cetera illa paria recte comparentur, sicut non nego, meliorem esse posse unam orationem dominicam laici, quam omnes horas canonicas unius sacerdotis.

Sexto\apnota{13}{U Sexo, Uv Sexto.} inducit b. Gregorium ante nongentos annos, Paschasium ante sexcentos\apnota{14}{M sexcentum.} annos dedisse indulgentias, quamvis, recitante Altisiodorensi\apnota{15}{M 'referente' Wilhelmo Altiodorisiensi, U recitante Wilhelmo Altisidiorensi.}, aliqui contradixerint s. Gregorio. Respondeo: Nondum habeo fide\apnota{16}{M vide, verbessert fide.} dignam historiam de indulgentiis Gregorianis. Tamen quidquid sit, non ideo\apnota{17}{fehlt in M.} sequitur, indulgentias esse aliud quam indulgentias, hoc est, remissiones bonorum operum, quas nemo bonas poterit appellare Christianis saltem\apnota{18}{fehlt in M U.}, quantumlibet sancti, concilia vel quicunque\apnota{19}{U quantumlibet sancta concilia vel quamlibet sancti.} eas dederint vel approbarint. Ideo D. D. non contra me agat solum\apnota{20}{M solis, U non solum contra me agat solis.} nominibus auctoritatum, sed et ipsius rei veritate.

Septimo suspicatur, ideo esse eas\apnota{22}{M ideo esse, U eas esse ideo.} boni operis vitium, quia sunt satisfactoriae. Dico, quod non sunt satisfactoriae, sed remissiones satisfactionis nec habendae\apnota{24}{M satisfactorium nec habendi.} ut opus satisfactorium. Ideo male hoc argumento comparantur operibus satisfactoriis, quasi ideo merita nostra minuantur, si indulgentiae relinquantur, imo contrarium verius est\apnota{26}{fehlt in M.}: merita augentur, dum indulgentiae relinquuntur.

Octavo, dicit consensu fidelium receptum annum iubilaeum Bonifacii VIII., Clementis VI., Urbani VI., Nicolai V., Sixti IV. etc.\apnota{29}{Urbani VI., Nicolai V. fehlt in M U, Sixti IV. fehlt in M; U Sixti statt Sixti IV.}. Respondeo: Quis scit, si consensus fidelium eum receperit? Nec est Bonifacius tantae opinionis, ut, quidquid gesserit, statim pro re bene gesta habendum sit, praesertim in re ista\apnota{31}{fehlt in U.} non bona nec necessaria ad salutem, qui\apnota{32}{M quia.} et alia monstra perpetravit, quae in rebus necessariis malum exemplum praebuerunt.

Nono, quod dederunt Romani pontifices consensu Christianissimorum\apnota{34}{M roma, gestrichen, christianorum, verbessert christianissimorum.} regum et principum indulgentias pro passagio etc.\apnota{35}{M passaio, verbessert passagio; etc. fehlt in U.}, dico sicut prius, non est miraculum, quod deus haec contemnat et sinat agi et agere, quae impertinentia sunt ad salutem, qui in hoc super nos vigilat, ut doceat utilia, ut apostolus ad\apnota{38}{fehlt in U.} Titum 3. [v. 8] dicit, hoc est, praecepta sua et consilia.

Decimo arguit, fidem ecclesiae esse indefectibilem et Christum nobiscum, ideo ecclesiam non potuisse errare tanto tempore cum animarum pernicie. [349] Dico: non tota ecclesia erravit. Deinde si etiam erraret in his\apnota{1}{U iis.} rebus nihili (semper de Christianis loquor), non est periculum salva fide. Variari opiniones et errare in temporalibus, non tollit ecclesiam Christi.

Undecimo decreta Romani pontificis\apnota{4}{W. A. Romanorum pontificum (U Ro: pon:).} acceptanda esse contendit. Respondeo: Sine dubio sunt acceptanda, sed tamen cum iudicio, ut habetur dist. 19. cap. Anastasius\apnota{6}{P dis. XXII. vel XXI. ca. Anastasius, a. R. Eccius invenit esse XIX. dis.; M U dist. 22. ca. Anastasio, Mc verbessert im Text 22 in 19, a. R. 19 vel ut alii 22.}; quia Romani pontifices homines fuerunt et homines circum se habuerunt, peccare\apnota{7}{M U errare.} potuerunt. Deinde nullum est decretum, quod praecipiat indulgentias redimere, nec est usque hodie, quod indulgentiarum valorem aperte declaraverit.

Ultimo, de moderni pontificis definitione, esse videlicet utiles indulgentias et de merito Christi et sanctorum acceptas, respondeo: Nec ipsa\apnota{11}{U ipse.} satis exprimit nec probat vel una\apnota{12a}{M pro una (statt vel una), Mc per unam.} syllaba, quae dicit: de qua re copiosius in actis\apnota{12b}{Mc actibus, verbessert actis.} meis Augustensibus 'dixi'. Non sequitur ergo, totum clerum in ecclesia tanto tempore\apnota{14}{M tanto tempore in ecclesia.} insanivisse, cum interim semper fuerint, quibus et displicuerint, qui et contradixerint saltem privatim indulgentiis, ut testatur vulgatissimum\apnota{16}{Mc 'pro' vulgatissimum.} proverbium: Indulgentiae sunt piae deceptiones.

D. D. Eccius.

Accepi reverendi patris sententiam in multis mitiorem et aequabiliorem, quam verba conclusionum praetenderent aut alia sua scripta mihi videbantur significare\apnota{21}{M significabantur, verbessert significare videbantur.}. Neque mei est instituti, indiscretas indulgentiarum proclamationes aut abusus defendere, sed veritatem et quantum indulgentiae rationabiliter datae possint prodesse, tutari\apnota{23}{U possunt statt possint, M putari statt tutari.}. Et ut sim brevis, non discutio de excusatione suae conclusionis. At cum assumit indulgentias non praeceptas\apnota{24}{M praeceptas fore.}, non necessarias esse, fatemur, Sed in hoc toto meo exili iudicio apparet esse difficultas nostrae dissensionis principaliter, quia reverendus pater dicit, indulgentias non esse, nisi remissiones bonorum operum, quam interpretationem apud nullum indulgentias vel dantem vel admittentem memini me reperisse, quoniam sic indulgentiae essent vere noxiae, periculosae et damnabiles, cum sic concilia et summi pontifices remitterent hic\apnota{30}{U haec.} bona opera, ut vel sic homo non mereretur tantam beatitudinem, quod est contra animarum salutem, vel homines decepti remissione bonorum operum ad purgatorium mitterentur\apnota{32}{M remitterentur.}, ut poenas hic non dilutas ibidem sustinerent, quod revera damnabile est\apnota{33}{M dandabile esset.}. Unde omnium vel dantium vel admittentium indulgentias unus consensus est, indulgentias non esse bonorum operum remissiones, sed remissionem esse poenae temporalis pro peccato contrito adhuc debitae et non solutae. Sicuti Leo modernus papa expresse in praefato decreto definivit, quod Sixtus papa\apnota{37}{M diffinit, U definit; papa fehlt in M U.} IV. ante eum expressissime fecit tempore iubilaei sui.

Unde concordare non possum, quod reverendus pater in suis resolutionibus dicit, indulgentias solum datas pigris, quod [350] etiam in sermone vulgari affirmat, non volentibus implere canones poenitentiales. Et tamen in eodem resolutorio dicit canones poenitentiales pigris impositos, quod viri alioquin strenui totam vitam\apnota{3}{P totam vitam vel omnem, M totam vitam, U omnem vitam.} poenitentiam agerent. Unde infelicissimae essent indulgentiae, si essent bonorum operum remissiones.

Soluturus ad primum respondit, indulgentias non attinere ad animarum salutem, cum profecto maxime intersit salutis animarum, an salutem promoveant vel impediant. Interest etiam fidei nostrae, an papa, ut Christi vicarius, poenas peccato debitas\apnota{8}{M debitae.} facta recompensatione per merita Christi et sanctorum possit\apnota{9}{M posset.} remittere et delere.

De Cypriano obiecit mihi, quod occurrebat Cyprianum legenti, quamvis non plene mihi constet, an sicut in Graecia, ita et in aliis regionibus ecclesiae maior pars cleri\apnota{12}{P ecclesiae.} fuerit infecta.

Deinde ad concilia\apnota{13}{M dein; U consilia.} respondit, indulgentias non esse bonas Christiano, id est, qui fide et charitate Christum sequatur. Hoc certe dicta conciliorum destruit et indulgentiarum largitionem, quia non nisi Christi fidelibus et contritis et confessis sub communi forma conceduntur. Quare communi sententia indulgentiae non sunt bonae malis, quia eas non consequuntur, sed bonis et vere\apnota{18}{M veris.} Christianis.

Non accepto, quod reverend. pater\apnota{19}{U quod 2. p[ater].} dictum Senecae hic usurpat: Non attendendum, quis, sed quid dicatur, quoniam maxime in his\apnota{20}{M quoniam quam maxime in his; U in iis.}, quae sunt fidei, attendere oportet etiam, quis dicat, et concilium\apnota{21}{M U consilium.} legitime congregatum est ab omni Christiano audiendum.

Dicit errorem non periculosum, si\apnota{23}{M sic.} hic de indulgentiis erretur. At mihi contrarium videtur, cum periculosi errores fuerint\apnota{24}{U fuerunt.} iudicati, qui sacris ordinibus vestris fratrum mendicantium ab initio detraxerunt.

Ad Innocentii\apnota{26}{U Innocentii, Uv Innocentii.} in generali concilio constitutionem dicit, se non intelligere, quomodo valeant, vel indulgentiae sint remissio peccatorum. Dico, ut semper, non dabo, tam solenne et legitimum concilium potuisse errasse\apnota{28}{M U concilium (U consilium) et legitimum potuisse errare.}, cum a spiritu sancto non derelinquatur et remissio peccatorum intelligenda sit quoad poenam, quia nomen peccati non semper pro culpa, sed etiam pro poena culpae debita usurpatur, sicut apud Machabaeos, ut a peccatis solvantur, quod de culpa nequit intelligi mortali, sed de poena culpae debita. Et si indulgentiae dicantur dari, ut absolvantur a poena et culpa eas consequentes, nullus existimare debet, indulgentias remittere culpam, sed quia papa concedit, ut a culpa, mediante sacramento poenitentiae, a deputatis commissariis absolvantur et dein\apnota{36}{U deinde.} indulgentias, id est, poenarum remissiones consequantur.

[351] Quod\apnota{1}{Vor quod hat W. A. die in U a. R. stehende Zeitangabe Hic continuavit Eccius hora secunda die 11. Julii 1519. Uv deutet die richtige Stellung dieser Angabe hinter facerent an; vgl. unten Z. 5.} vero se excusat, quod dixit\apnota{2}{M dicit.} indulgentias esse vitium boni operis, quia et ipse voluerit ceteris paribus opus bonum cum indulgentiis esse melius quam sine indulgentiis, sed sicut fieri soleat, quia ita homines indulgentias redimant, ut bona opera illa sine indulgentiis non facerent.

Exactis iam duabus horis reliqua in horam secundam pomeridianam dilata sunt; qua veniente prosecutus est, ut sequitur\apnota{5}{Die Zeitangabe fehlt in M, doch ist der neue Anfang an der Schrift deutlich zu erkennen; über U und Uv vgl. oben Z. 1.}.

Quamvis pleraque super illo dicto afferri\apnota{6}{M illo verbo 'inferri'.} possent, et in simili, quando vicini propter amicitiam vicinorum conveniunt ad oblationes et offertoria plus quam intuitu mortuorum, attamen mihi apparet, meliori iudicio semper salvo, hoc non esse vitium boni operis. Primum\apnota{9a}{M primo.} ex eo, quia sacrum concilium\apnota{9b}{U consilium.} per indulgentias vult excitare homines ad bona opera. Secundo, quia esto, quod tale quid faciant vel exhibeant propter indulgentias consequendas, tamen ista opera ex illo fine non inficiuntur, quoniam finis iste indulgentiarum non discrepat a fine ultimo, ita quod poenarum remissio congruenter in deum ordinatur, alioquin praescindendo subordinationem finium paucissima vel nulla essent bona opera.

De contemptu indulgentiarum per Constantiense concilium improbato\apnota{15}{M concilium improbatum, U consilium improbatum.} dicit reverendus pater, se illas non contempsisse, fateor ingenue, si eo semper fuisset usus moderamine, sicut se hodie exposuit, et forte eius interpretationi\apnota{18a}{M interpretatione.} est standum et in conscientia sine 'foro', tamen scriptus\apnota{18b}{M scripta.} eius sermo vulgaris aliam exhibet intelligentiam laicis. Sed hoc nihil ad me.

Ad Gersonem respondit admittendo, tamen arbitratur Gersonem mihi adversari, qui indulgentias\apnota{21}{M gratiam, gestrichen, indulgentias.} tunc utiles reputat, quando fide, spe et charitate percipiuntur. Quod ego alias tenuerim posse fieri satisfactionem in peccato mortali, dico et ipsum Christianissimum cancellarium hoc tenere, aliquem posse satisfacere de iniuncta poenitentia in peccato mortali, alioquin imprudentissime agerent confessionum auditores, durantem imponendo poenitentiam. Sed an quispiam\apnota{26}{M quaespiam (Verwechslung der Abbreviatur).} indulgentias in peccato mortali consequatur, nihil ad rem, ratio diversitatis est obvia, quoniam indulgentiae non solent concedi, nisi contritis et confessis,

ideo non mirum, si peccatores\apnota{28}{M peccatores existentes.} harum non sunt capaces, quia indulgentiae tantum valent, quantum sonant\apnota{29}{M sunt, Mc sonant.}.

Subiunxit reverendus pater, orationem dominicam\apnota{30}{M dominici.} laici tantum valere, sicut horas canonicas clerici, si cetera sunt paria, illud non admitto.

De b. Gregorio dicit, se nondum receptam vidisse historiam. At tam constantissima fama, sicut b. Augustinus lib. de trin. de Alexandria inquit, hoc negari\apnota{34}{U negare.} non permittit. Negat autem indulgentias esse satisfactoria, sed solum sint remissiones bonorum operum. In hoc credo principali nos dissentire et magnam partem controversiae hic fundari. At prius memini, indulgentias esse non remissiones bonorum operum, ut credit reverendus pater, [352] sed remissiones poenarum peccatis debitarum\apnota{1}{M debita.}. Unde homo vere Christianus accipiens indulgentias non minora, sed plura facit bona opera, quam si non consequeretur indulgentias. Sic autem intelligo indulgentias\apnota{3}{U indulgentis.} esse satisfactoria, sicut papa Sixtus IV., Clemens VI. et papa modernus declarant\apnota{4}{IV. und VI. fehlen in P, M U Sixtus Clemens (U Clem. 8) papa modernus declararunt.}, non quod consecutus indulgentias satisfaciat, sed quia nullum malum manet impunitum et nullum peccatum manet inultum (can. Sicut primi de poenit. dist. 1.). Ita vicarius Christi dispensator eius thesauri, ubi homo pro poena peccati non reddit aequivalens, ipse ex thesauro sibi commisso merito Christi sufficientissimo illud recompensat et pro eo satisfacit.

Quare concessi in prioribus, indulgentias de se non delere culpam, sed poenam, et quod sit talis remissio poenarum in\apnota{11}{fehlt in M.} indulgentiis et non remissio bonorum operum, patet per illud Esaiae: Spiritus Domini super me, eo quod unxerit me, ad annunciandum mansuetis misit me, ut mederer contritis ac praedicarem captivis indulgentiam et clausis apertionem etc.\apnota{14}{fehlt in U.}. Christus eadem verba de se assumens pro verbo indulgentiam\apnota{15}{U indulgentias.} remissionis nomen accepit. Sic Theophilus ait, ut b. Thomas eum enarrat\apnota{16}{M eum renarrat, U enim renarrat, Uv eum renarrat.}: nam prius ab inferis animae omnium tenebantur, donec venit, qui praedicaret captivis remissionem.

Deinde contendit reverendus pater, merita augeri relictis indulgentiis. Hoc ego non video\apnota{19a}{M vidi.}, quoniam, ut dixi, relinquentes\apnota{19b}{M non, gestrichen, relinquentes, U non redimentes vel reliquentes (so!)} indulgentias aliquando istam pecuniam in bonis conviviis volunt consumere\apnota{20}{P absumere}, non merita augere. Fateor, quod aliquando posset\apnota{21}{M non merita (Mc 'indul.') redimere fateor quod aliquando possit.} quis facere maius

opus meritorium cum illa pecunia\apnota{22}{U meritoritorium, in Uv verbessert; M peccata, Mc pecunia.}, quam si redimeret indulgentias, casus non oportet specificari\apnota{23}{M specificare, U specificari.}.

De iubilaeis celebratis non vult admittere, accessisse consensum fidelium, et suspecta Bonifacii VIII. vita posset suspectum\apnota{25a}{M Bonif. octava vita posset suspecta, U = P, nur potest statt posset.} facere eius iubilaeum. Dico\apnota{25b}{M dicam (?).}, satis magnum esse populi consensum, quod ex toto Christiano\apnota{26}{M Christiani.} orbe in iubilaeo Christi fideles confluunt; quomodo alias consensum fidelium probare\apnota{27}{M probari.} possit, non video.

Quod consensu principum dixi pro passagio et cruciatis a summo pontifice indulgentias datas, respondet parum conficere, si hic erraverint in re nihili, quasi nihil esset, Christi fideles decipere et unanimem Christianorum principum et regum consensum in tam fallacem descendere fidelium deceptionem.

Praeterea, nullum, ait, esse decretum expressum et clarum super indulgentiis. Oppono ei declarationem Sixtinam et iam novissime declarationem moderni pontificis, ubi inquit, Romanam ecclesiam, quam reliquae tamquam matrem sequi tenentur, tradidisse Romanum pontificem, Petri clavigeri\apnota{37}{M clafigeri.} successorem et Jesu Christi in terris vicarium, potestate clavium, quarum est aperire tollendo illius in Christi fidelibus impedimenta, culpam scilicet et [353] poenam pro peccatis actualibus debitam. Culpam quidem\apnota{1}{U deletam statt debitam, M scilicet statt quidem.} mediante sacramento poenitentiae, poenam vero temporalem pro peccatis actualibus secundum divinam iustitiam debitam, mediante ecclesiastica indulgentia: posse pro rationabilibus causis concedi eisdem\apnota{4}{M U concedi fehlt; M eiusdem.} Christi fidelibus, qui charitate iungente sunt membra Christi, sive in hac vita, sive in purgatorio\apnota{5}{M purgatoria, Mc purgatorio.}, indulgentias ex superabundantia meritorum Christi et sanctorum ac tam pro vivis quam defunctis\apnota{6}{M = P, nur et pro statt quam, U et tam pro vivis quam pro defunctis.} apostolica auctoritate indulgentiam concedendo, thesaurum meritorum Christi et sanctorum dispensare\apnota{8}{M U dispensare etc.}. Ubi clare determinat cum sequentibus, indulgentias non esse remissiones bonorum operum, sed remissiones poenarum, quae fiant\apnota{9}{U fiat.} compensatione facta meritorum Christi.

Et quod ratiocinabar\apnota{11}{P rationabar; M quo statt quod.}, ecclesiam tot annis totam non errasse, respondit, non totam ecclesiam errasse, quia semper interea\apnota{12}{M interim.} fuerint, quibus hoc negotium fuerit suspectum. Hinc vulgatum illud: Indulgentiae sunt piae deceptiones fidelium. Ad quod dico: Quando alicuius rei habenda est ratio, non quod quisque dicat, sed quod boni, docti, honesti et optimi statuunt, audiendum est. Tot\apnota{16a}{U quot.} enim fuerunt

probatissimi et honestissimi viri per totum Christianum\apnota{16b}{M probatissimi viri per totam Christianam.} orbem in trecentis annis et horum aliqui in catalogum sanctorum relati, non tamen indulgentias hoc pacto improbaverunt\apnota{18}{U probaverunt.}, etiamsi abusus eis sicut in omni re displicuerit. Thomas sanctus est, Bonaventura, Albertus Carmelita, Bernhardinus\apnota{20a}{M Bernhardinus Antonius de Badna (so!)} et alii magni viri, Ales\apnota{20b}{fehlt in M.}, Gerson, Capistranus etc., qui tamen indulgentiarum concessionem non reprobarunt. Quare nihil refert, quod Johanni Wikleff et Ulrico Kalteisen in Anglia, domino Johanni de Wesalia\apnota{23a}{M Johanni bicleff (Mc 'bicler') et Ulrico Calte in Anglia, d. Joh. de Besalia.}, qui tamen hoc retractavit Maguntiae, et Bohemis\apnota{23b}{et fehlt in U; P Boiemis.} indulgentiarum concessiones non placuerunt.

Quare concludamus, indulgentias rationabiliter datas ad dei honorem et fidei exaltationem non esse inutiles Christi fidelibus devotis, non praesumentibus, sed eas in humilitate et gratiarum actione\apnota{27}{M U cum humilitate et (M et cum) gratiarum actione.} accipientibus, abusibus resectis et repulsis, ut sic potius de tanto munere deo gratias agamus, quam ut id\apnota{29}{M U illud.} superbe contemnamus, semper in animo habentes, optimam indulgentiam\apnota{30}{M contritionem gestrichen, indulgentiam.} esse veram contritionem. Tamen omnino, si hodie exposita sententia fuit mens reverendi patris, maluissem, ut illud clare dixisset et simplicibus non dedisset occasionem aliud cogitandi, quia nemo est omnium auditorum, qui sensum hodie datum intellexerit in conclusione undecima: Dicere, indulgentias 'esse bonum Christiano, est insanire, sunt enim verissime operis boni vitium. At sum bene contentus de sua expositione, quantum ad me attinet.

D. D. Martinus.

Contra solutiones meas opponit egregius D. D. primo, indulgentias non esse remissiones bonorum operum, nec se vidisse, qui hoc assereret. Sequi [354] etiam ad hoc, indulgentias esse vere periculosas, noxias\apnota{1}{M obnoxias.} et damnabiles, esse potius eas uno consensu dantium et 'admittentium\apnota{2}{M Lücke, Mc 'admittentium esse'.} remissiones poenarum pro peccatis debitarum\apnota{3}{pro peccatis debitarum fehlt in M.}. Respondeo: Satis admiror, quid haec verba D. doctoris velint, cum sine dubio per poenas pro peccatis debitas intelligant omnes labores satisfactionis, quae tertia pars est poenitentiae. Inter quos certe sunt bona opera, eleemosynae, orationes, ieiunia, vigiliae\apnota{6}{M et vigilhae.} et similia. Ad quod est expressus textus lib. 5.\apnota{7}{Mc 'lib. 6.'} cap.: Cum ex eo, ubi papa dicit, per indiscretas indulgentias enervari poenitentialem satisfactionem, et omnes bullae indulgentiarum ferme hanc clausulam habent: De iniunctis poenitentiis relaxamus.

Sed age, dato, quod essent poenae et non opera, dico peius esse, remitti poenas quam opera, cum vita passiva et poenosa plus perficiat ac prosit quam mere\apnota{11}{U aut prosit quam mere; M mera.} activa, ut apostolus 2. Cor. 12. [v. 10]: Dum\apnota{12a}{P 2. Corinth. XI. forte dum, M 2. ad Corinth. 11 dum (Mc 'cum'), U 2 Corinth. 11 cum.} infirmor, tunc fortior sum. Purius\apnota{12b}{Mc a. R. durius.} est enim opus, quod solo deo operante, nobis patientibus, perficitur, quam quod nobis cooperantibus efficitur, iuxta illud Deut. 32. [v. 39]\apnota{14}{P M U Deut. 22.}: Ego percutiam et sanabo, occidam et vivificabo\apnota{15}{M occidam etc.}. Quare ad sensum egregii D. D. plura absurda sequuntur ex indulgentiis quam ad meum sensum, et sane nescio, an verissime dicam\apnota{16}{M dicam, Mc dicat.}.

Secundo, non potest concordare, quod dixi, indulgentias solum pigris esse datas et canones similiter tantum\apnota{18}{M tantum similiter.} pigris impositos. Dico: Non memini me dixisse, canones esse pigris impositos, si autem dixi, sensus meus est\apnota{19}{M esset.}, quod strenui et ferventes peccatores\apnota{20}{U ferventes sive poenitentes, Uv = P.} plus faciunt, quam canones exigunt.

Tertio\apnota{21}{M U Tertio contendit eas esse necessarias ad salutem, quia expedit nosse an promoveant, necne ad salutem, Mc streicht die Worte von expedit an, Uv streicht das Ganze; in P fehlt alles ausser tertio; vgl. Einleitung p. 7.} etsi fateatur, non esse necessarias indulgentias ad salutem, expedit tamen scire, an promoveant salutem necne\apnota{23}{U Etsi non fateatur necessarias esse ad salutem expedit tamen scire, promoveant salutem necne.}. Respondeo: nec hoc est necessarium scire, sed si nocerent saluti; non enim est necessarium scire non necessaria ad salutem\apnota{25}{M necessaria saluti.}.

Item, quid interest, an\apnota{26}{M U quod interest an (M quod, Mc 'an').} papa per merita Christi possit remittere? Nescio an intersit\apnota{27}{M an id 'intersit', U an inersit (Uv intersit).}. De meritis Christi postea\apnota{28}{M = P, Mc 'quae merita Christi postulant' (a. R. wiederholt).}.

Post haec dicit: Meam sententiam destruere dicta conciliorum et indulgentiarum largitionem, quae communi forma concedunt contritis et confessis. Respondeo: Non ideo destruit mea sententia dicta conciliorum, quia sunt contriti et confessi, qui pigri sunt ad solvendas satisfactiones, quos scriptura\apnota{32}{M in scriptura.} etiam infirmos vocat, et propter hos damnandae non sunt indulgentiae.

Reprehendit etiam me\apnota{34}{M U in me.}, quod usurpavi dictum Senecae: Non quis sed quid dicatur, attende\apnota{35}{M attendi, fehlt in U.}, quod in his, quae sunt fidei, attendere oporteat etiam, quis dicat. Respondeo: Hoc non est ad propositum, quia materia indulgentiarum non est materia fidei. Deinde, in materia fidei multo maxime opus est attendere, non quis, sed quid dicatur,

quia non est respectus personarum habendus, sicut manifeste praemonuit Christus\apnota{39}{fehlt in M U, von Mc Uv nachgetragen.} Matth. 7. [v. 15] et 24. [v. 5]: Attendite a falsis prophetis, qui veniunt ad vos etc. Item: multi venient\apnota{40}{U veniunt.} in nomine meo dicentes: Ego sum [355] Christus; et Joh. 10. [v. 27]: Oves meae vocem meam audiunt. Proinde\apnota{1}{M audient, P Deinde.} etiam summo pontifici non temere credendum est, ut stet regula Johannis apostoli [I. 4 v. 1]: Carissimi, probate spiritus, utrum ex deo sint.

Dicit deinde, periculum esse in errore indulgentiarum ex eo, quod periculosus error habitus sit eorum, qui ordinibus mendicantium ab initio detraxerunt. Quidquid sit de detractoribus\apnota{6}{U detractatoribus.}, per hoc non probatur error periculosus in indulgentiis, nec mox error est, quod ordines mendicantium vel sacerdotes ecclesiarum errorem iudicaverint. 'Per idem', ad id, quod dicit, concilium non errare, quia regitur a spiritu sancto, addo ego, scilicet in his, in quibus regitur a spiritu sancto, id est, in rebus fidei. Et ut aliquando dicam, non dixi errasse concilium in dandis indulgentiis, sed posse errare\apnota{12}{U errasse in concilio in dandis indulgentiis, sed postea (Uv. a. R. sed posse errare).}.

Transeo illud, quod remissionem peccatorum intellexit\apnota{13}{MU intelligit.} remissionem poenarum, cum saepe peccatum pro poena accipiatur, 2. Maccab. 12 [v. 46]\apnota{14}{P M Macc. 2, U Macc. 20.}. Ego non saepe invenio peccatum pro poena accipi, nisi forte, ubi de Christo scribitur, quod peccata nostra ipse portavit, quod tamen et ipsum non ausim\apnota{16}{M ausum.} dicere simpliciter pro poena accipi. Sed haec extranea.

Confutaturus illud, quod venias dixi\apnota{18}{U dixi esse.} vitium operis, dicit esse subordinationem finium in operibus\apnota{19}{M in bonis operibus.}, ideo si veniae referantur in deum, non vitiant opus. Respondeo: Da exemplum\apnota{20}{M da exemplum 'et da', die letzten beiden Worte gestrichen, U da exemplum et.} demonstrandi. Dixi enim hodie raros inveniri, qui tantum darent gratis et propter deum, quantum dant pro indulgentiis; ex fructibus ergo cognoscetis\apnota{22}{U cognosceris.} eos, imo si audiant melius posse se\apnota{23a}{M si audeam melius posse se, U si audiant melius se posse.} facere quam redimere indulgentias, probabile est, nullos\apnota{23b}{U nullas.} esse redempturos indulgentias.

De contemptu indulgentiarum dicit: Si fuissem usus hoc moderamine\apnota{25}{M moderamini.}, standum esset meae\apnota{26}{M esset eius, U est eius.} declarationi, verum vulgaris sermo aliam indicat laicis intelligentiam. Respondeo: Videat diligentius eundem sermonem egregius D. D. et inveniet, me expressis verbis dixisse, non esse contemnendas indulgentias, sed relinquendas liberas, neminem enim\apnota{29}{M U neminem esse.} absterrendum, sicut neminem urgendum, esse videlicet pigrorum

solatium\apnota{30}{M esset videlicet 'pigrorum' (ursprünglich virorum?) solatia, U = P, nur solatia.}. Quare sine mea culpa hanc declarationem et satis moderatam sententiam non intelligunt.

Transeo illud de Gersone, ubi dicit: Aliud esse satisfacere in peccatis et redimere indulgentias in peccatis: quia in his\apnota{33}{U iis.} obstat forma concessionis, quae solum confessis et contritis permittit indulgentias, illic vero et ipse Johannes Gerson secum sentiat. Dico: Relinquo\apnota{35}{U relinquendo, Uv relinquo.} hoc tamquam extra capacitatem intellectus mei, vere enim non capio, quomodo peccator possit satisfacere in peccatis et non\apnota{37a}{W. A. lässt dieses non fälschlich aus.} etiam possit recipere licentiam\apnota{37b}{M possit etiam recipere (dicentium, gestrichen, Mc a. R.) 'licentiam'.} non satisfaciendi, cum plus videatur requiri ad operum perfectionem quam ad\apnota{38}{fehlt in M.} omissionem.

[356] De Gregorio, etsi historia non habeatur recepta\apnota{1}{M etiam recepta.}, tamen famam dicit non permittere negari. Dico: Famae temere nemo credat, et transeo.

Tandem ad id redeo, quod indulgentiae sint satisfactoriae et non remissiones bonorum operum, sed poenarum, fortassis contendimus in aequivoco, quod satisfactorium D. D. appellat, quia permittitur virtute indulgentiarum, ut non satisfaciat, et haec permissio\apnota{6}{U promissio.} non satisfactionis pro satisfactorio habeatur. Ego satisfactorium appello, quod implet satisfactionem, non omittit\apnota{7}{in M durchstrichen, daher unleserlich, Mc. a. R. omissionem.}.

Adduxit etiam auctoritatem Esai. 61 [v. 1]\apnota{8}{P U Esai. 66, M Esai. 60.}, ubi Christus Lucae 4 [v. 18] legens ex Esaia\apnota{9}{M Esaie.}, dicit: Spiritus Domini super me, eo quod unxerit me, praedicare captivis indulgentiam, ubi Christus remissionem dixit\apnota{10}{fehlt in P M.}. Addidit Theophilum, qui exponit, quod animae tenebantur etc., donec venit\apnota{11}{etc. fehlt in M; M veniat.}, qui praedicaret captivis redemptionem. Quod illa auctoritas\apnota{12}{M dico quod illas auctoritas.} nihil faciat ad rem, ipsa verba et consequentia textus\apnota{13}{M verba textus et sequentia.} declarant, quia Christus ibidem dicit: Hodie impleta est scriptura haec\apnota{14}{M haec scriptura.} in auribus vestris; deinde, quia loquitur de vero anno iubilaei, hoc est, de plenitudine temporis, de corona anni benignitatis, et non de remissionibus poenarum, sed de remissione peccatorum, de qua tota scriptura loquitur, cum iubilaeus Bonifacianus tanto post tempore coeperit. Item quod neget augeri merita relictis veniis, transeo, quia non per hoc confutatur, quod aliquando pecuniam eandem volunt consumere in conviviis.

De iubilaeis celebratis dicit, satis magnum esse populi consensum, quod ad indulgentias tot\apnota{22}{fehlt in M U.} fideles in toto orbe confluunt, excepta

tamen Italia et Roma, nec videt, quomodo alias consensum possim probare\apnota{23}{M Nec videat quomodo alias possit probari consensus; U possim probare consensum, sonst = P.} fidelium Respondeo: Eum consensum ego probarem esse fidelium\apnota{24}{M ego fehlt; esse fidelium probarem.}, qui re patefacta et veritate indulgentiarum exposita fideles ad confluendum moveret. Nunc vero confluunt ignorantes, quid et quare, cum maior pars credat se\apnota{26}{fehlt in M.} rem necessariam et meritoriam facere.

Post hoc de passagiis et cruciatis arbitratur egregius D.\apnota{28}{fehlt in M.} D., non esse rem parvi momenti, tot fideles decipi, praesertim magnates. Respondeo: Nec est res magni momenti, quandoquidem et divus Hieronymus Matth. 23. [v. 35] de sanguine Zachariae prophetae dicit: Non damnamus errorem, qui de pietate fidei descendit\apnota{32}{M descendunt, verbessert descendit.}. Ita et hic error est sine damno animae, citra etiam lucrum pecuniae, quod deferunt pro honore s. ecclesiae rectoribus ecclesiarum.

Quando ego\apnota{34}{fehlt in M.} dixi, nullum esse decretum super indulgentiis, opposuit Sixtinam et Leoninam novissimam, ubi declarantur virtutes indulgentiarum. Dico: haec alia quaestio est. Ego dixi, nullum esse decretum, quod praeciperet indulgentias redimere, ideo non esse necessarias.

In fine, ad proverbium illud: Indulgentiae sunt piae fidelium deceptiones, per quod volui, non totam ecclesiam errasse, dicit non esse audiendum, quod quisque dicat, sed quod boni et honesti viri statuant. Nunc vero aliqui etiam\apnota{40}{M statuunt nunc etiam aliqui (sancti, gestrichen) etiam. U = P, nur statuunt.} [:357] in sanctorum catalogum relati non improbant indulgentias. Respondeo: Nec ego sane. Hoc tamen addo, quod ecclesia accipiens opiniones non facit ideo ex opinionibus veritates. Ideo ad illam Leonis X.\apnota{3}{M U istam Leonis X. (X. fehlt in M).} declarationem dico, sicut in actis dixi, non probari nudis verbis praesertim hominum, indulgentias esse de thesauro Christi sumptas et solvere poenas iustitiae divinae\apnota{5}{U iustitia divina.} requisitas, licet opinionem non damnem. Nam, ut dicam, quod sentio, cum Christi merita, sive accipiantur\apnota{7}{M accipientur.} ut suffragia, sive quocunque modo et nomine, nihilominus sunt merita Christi, nec propter usum vel applicationem metamorphosim subeunt. Si autem sunt merita Christi, gratia et veritas sunt, iuxta illud Psal. 25 [v. 10]: Universae\apnota{10}{M universi.} viae Domini misericordia et veritas, et Joh. 1 [v. 14ff.]: Gratia et veritas per Jesum Christum facta est; ibidem: Vidimus plenum gratiae\apnota{12}{P M gratia.} et veritatis. Quare etiamsi angelus de caelo aliud persuaserit, non credam esse in manu ullius\apnota{13}{M unius.} hominis gratiam et veritatem dare, hoc est, merita Christi dispensare, non tamen resisto pontifici. Sed ad declarationem in actis me refero.

D. D. Eccius.

Quia, ut dixi, magna pars dissensionis iacet in eo, an indulgentiae sint remissiones bonorum operum vel poenarum, ubi in meam partem dixi descendere ecclesiae doctores et Romanos\apnota{19}{W. A. Romani.} pontifices, qui hactenus de indulgentiis scripserunt. Sed hoc repellit\apnota{20}{M rescripserunt statt scripserunt; U refellit.} acute et docte reverendus pater: Quia remittere poenam 'esset remittere labores satisfactionis'\apnota{21}{Mc 'satisfactionum'.}, inter quos sunt bona opera, contra decretum Innocentii, can.: Cum ex eo, de poenitentiis et remis.\apnota{22}{et remis. fehlt in M U.}, 'ne poenitentialis satisfactio enervetur'. Dico bene remitti labores satisfactionis, sed non propterea bona opera, quamvis enim oratio, ieiunium etc. sint satisfactiones, tamen sunt bona opera, etiamsi his\apnota{25}{U iis.} non satisfaciamus. Accedit, quod satisfactio non solum fit\apnota{26a}{M fit, Mc 'satisfit'.} labore bonorum operum, sed etiam satispassionis\apnota{26b}{P M Uv satispassio.}. Quare plerique volunt, non debere cogi filium confessionis ad satisfactionem in se recipiendam; iuxta illud commune: Melius\apnota{28}{U melior.} est, ut confessor filium cum parva poenitentia mittat ad purgatorium, quam cum magna ad infernum. Nec mirum, si pontifex plenarias resecuit\apnota{30}{M resequit, Mc 'resecuit'.} indulgentias ad satisfactiones poenitentiales eludendas, si passim a praelatis fierent\apnota{31}{P Uv fieret; M fiat, U fiant; M U eludens statt eludendas, Uv = P.}. Sic particula bullae de iniunctis poenitentiis, id est, quae secundum iustitiam divinam iniungi deberent\apnota{32}{M deberet.}, ut doctores ecclesiastici\apnota{33}{M U ecclesiae.} de indulgentiis scribentes declarant.

Porro, quod argumentatur, peius esse remittere poenas quam opera, non accipio, quia pro tanto poenae prosunt, ut homo instruatur bene operari, et sic virtus in infirmitate perficitur, sed hoc ad iudices\apnota{36}{W. A. indices.}. Et quia errare in indulgentiis esset animabus 'periculosum', ut hodie induxi, ideo ecclesia non errat, si solum pigros contritos vult\apnota{38}{M U ideo ecclesiae non errant, si . . . . volunt (M vult).} recipere indulgentias, certe ad Carthusianos venire non deberent\apnota{39}{M deberet, U debent.}.

[358] In materia fidei dicit\apnota{1}{M dixit.} maxime attendendum, quid dicatur. Ego autem auctoritatem reputo in theologia potentissimam, et quod audire praecipimur sedentes super cathedram Moysi etc.\apnota{3}{M sedentem statt sedentes, U supra statt super; etc. fehlt in M U.}. Ideo quis dicat, attendere debemus, ut captivemus intellectum, iuxta praeceptum apostoli, in obsequium fidei, propterea\apnota{5}{M proptera.} concilia habita sunt, ut eorum auctoritate errores exterminarentur.

De mendicantibus transeo.

Quod vero ait, concilium in his, in quibus regatur\apnota{7}{M concilium in his, quae regantur, U consilium in iis in quibus regitur.} a spiritu

sancto, non errare, dico ego, antiquorum patrum\apnota{8}{fehlt in M.} secutus sententiam, semper debere praesumi concilium\apnota{9}{U consilium.} regi a spiritu sancto, quamdiu non constat de opposito, scilicet de illegitimitate congregationis.

Sed adhuc fortiter stat reverendus pater in vitio boni operis, quod dans pecuniam\apnota{12a}{M pecunias.} propter indulgentias alias non daturus\apnota{12b}{M daturas.}, quasi hoc sit vitium boni operis. Quaererem a 'reverendo patre', si ipse 'cum suis fratribus' habeat anniversarium alicuius 'principis' vel alicuius nobilis cum triginta missis, quia dantur bonae praesentiae, alioquin non habiturus anniversarium\apnota{15a}{Mc 'cum suis fratribus habet anniversarium' (a. R., die beiden letzten Worte auch im Text) 'cuiusdam principis' (Lücke ausgefüllt) eum 30 vel 20 missis non 'habiturus etc.' (statt anniversarium); U = P, nur habet statt habeat, nobilis statt alicuius nobilis.}: an hoc sit vitium boni operis? [Martinus dubitat. Eccius: Et sic non erit salva omnis caro tam in cuculla quam extra\apnota{15b}{Die Zwischenbemerkung in M U a. R.; M dubitet statt dubitat, Eccius und tam in cuculla quam extra fehlen.}.]

De indulgentiis ab eo non contemptis iudicent, qui sermonem vulgarem legerint\apnota{17}{M vulgarem suum sermonem legent.}.

Porro reverendus pater factus interpres meus dicit, indulgentias forte sic esse\apnota{19a}{U factas esse.} satisfactorias, quia propter ipsas\apnota{19b}{U istas.} aliquis 'satisfacere non teneatur'. Non iste est sensus meus, sed quia homo deberet pro peccatis etiam\apnota{20}{M etiam pro peccatis.} contritis satisfacere de suo et per indulgentias satisfacit de alieno, quia summus pontifex dat ei de thesauro ecclesiae, unde solvat, ut sic, iuxta communem patrum sententiam, etiam in indulgentiis peccatum non maneat\apnota{23}{P U manet.} impunitum, salvo eo, quod etiam de proprio satisfacere non posset, non accedentibus Christi meritis.

De consensu fidelium credo neminem esse tam stultum, qui credat, rem esse necessariam indulgentias in iubilaeo, et si existimet iter Romanum sibi esse meritorium, nec hic errat\apnota{28}{M esse sibi meritorium, nec hic erret, U esse meritorium nec hic errat.}, stat enim opus, per quod quis consequitur indulgentias, esse meritorium et tamen ipsas indulgentias non esse meritorias.

Porro, de verbis Sixtinis\apnota{30}{M Sixtini, U Sixtinae.} et moderni pontificis non facio vim, quia cathedram Petri sequor et eius sententiam, qui in ea sedet, quamdiu in haeresim, quod absit, prolapsus non fuerit. Scio eum non probare\apnota{32}{M Scio enim eum probare.}, sed definire, quem ut Christi vicarium definiendo\apnota{33}{M diffinire und diffiniendo.} arbitror esse indefectibilis fidei. Sed haec omnia eorum committo iudicio, quorum interest, paratus errata expungere, si quae ostendentur\apnota{35}{M demonstrentur.}.

Addidit D. D. Martinus: Et ego\apnota{36}{M Addidit D. D. fehlt, U Jussit Martinus addere: Et ego.}.

[359] Crastina hora septima ante meridiem inter eosdem E. D. doctores continuabitur disputatio tractabiturque materia de poenitentia.

Die Martis duodecima mensis Julii, hora septima antemeridiana.

D. D. Eccius\apnota{5}{M Die Zeitbestimmung von crastina an fehlt, Überschrift De poenitentia; U duodecima Julii argumentatus est Eccius respondente Martino de poenitentia. In nomine tuo, dulcis Jesu.}.

Contra conclusionem reverendi patris et partem sermonis sui de poenitentia pro defensione optimorum patrum praedicantium intendo probare, poenitentiam veram incipere a timore etiam poenae, et quod bene quis se praeparet\apnota{9}{M se quis praeparet.} discutiendo, recogitando etc.

Primo, quia talem modum praedicandi observat\apnota{10}{M U observavit; U a. R. Eccius.} Dominus Jesus et eius praecursor s. Johannes. Luc. enim 15. filius prodigus gerit typum poenitentis secundum Augustinum de quaest. evang., Ambrosium libro 2. de poenitent. cap. 3., Chrysostomum, Hieronymum et alios, sed eum\apnota{13}{M eam.} ita proponit nobis Christus: in se conversus dixit: Quanti\apnota{14}{U quod in se conversus dixit: quantum.} mercenarii in domo patris mei abundant panibus, ego autem hic fame pereo, surgam et ibo ad patrem meum et dicam illi: Pater, peccavi in caelum et coram te etc. Hic Dominus Christus modum poenitentis describens, exponit primo motum poenitentem praemiorum magnitudine, scilicet »abundant panibus«, et timore poenae, scilicet »hic fame pereo«, quibus gradibus evectus poenitentiam veram coepit meditari, scilicet »et dicam: Pater, peccavi« etc. Et ut\apnota{20}{U peccavi et hoc.} Basilius explanat: Tres sunt ibi poenitentiae gradus, mercedis spes, suppliciorum timor et bonitatis paternae sincera dilectio. Et sic, antequam fiat amicus dei, qui extitit inimicus, efficitur prius servus.

Sic Johannes incepit a timore: Progenies vel genimina viperarum, quis demonstravit\apnota{24}{M demonstrabit, U monstrabit.} vobis fugere etc.?

Id\apnota{25a}{U sic.} etiam in toto actum est veteri testamento\apnota{25b}{U instrumento.}, ut timore poenarum ad mandatorum obedientiam homines traherentur, ut in legalibus, historiis et\apnota{26}{fehlt in U.} prophetis hoc ubique reperitur. Sufficiat adduci\apnota{27}{M repetitur statt reperitur; M U adducere, Uv = P.} unum locum, Psal. 89. [v. 31ff.]: Si autem dereliquerint filii eius legem meam et in iudiciis meis non ambulaverint\apnota{29}{U ambulaverit, Uv ambulaverint.}, si iustitias meas profanaverint et mandata mea non custodierint, visitabo in virga iniquitates eorum et in verberibus peccata eorum; cum concordantiis\apnota{31}{M concordanti, U concordantibus, Uv = P.}. Quare praedicatores sacram imitati scripturam hactenus bonum habuerunt modum de poenitentia paranda, recogitando gravitatem peccatorum, aeterni supplicii etc.

Accedit ratio Dionysiaca: Quia sic deus res administrat, ut ima ad superiora per media ducat. Quare cum peccator sit in imo, quia non est dignus pane, quo vescitur, volens ascendere ad superiora gratiae hoc faciet\apnota{37}{U faciat.} per medium timoris.

Bene esset perfectionis, quod aliquis immersus luto [360] peccati posset se mero obtutu dilectionis dei et amore iustitiae erigere\apnota{1}{M erire, Mc erigere.} ad gratiam capessendam; sed quis est hic et laudabimus eum?

Imo b. Augustinus approbat modum praedicandi nostrae tempestatis, docens expresse, nullum pervenire ad charitatem et ad veram gratiam, nisi praecedente timore, non filiali, sed etiam servili. Aug. testatur tract. 9 super Joh.\apnota{5}{P M U Aug. est tract. 9; M Test., verbessert Augustinus; super Joh. fehlt in U, in Uv nachgetragen.}: Ergo incipiat timor, quia initium sapientiae timor Domini. Cum autem coeperit charitas habitare, pellitur timor, qui ei locum praeparavit. Quantum enim illa crescit, ille decrescit, et quantum illa fit interior, timor pellitur foras; maior charitas, minor timor; minor charitas, maior timor. Si autem nullus timor, non est, qua intret\apnota{9}{U intrat.} charitas (dat comparationem), sicut videmus\apnota{10}{M vitemus, Mc videmus.} per setam introduci linum, quo suitur, seta prius intrat, sed nisi exeat, non succedit linum: sic timor prius occupat mentem, quia ideo intravit, ut introduceret charitatem. Allegat illud psalmistae [30 v. 12]\apnota{13a}{M U psalmi.}: Convertisti luctum meum in gaudium\apnota{13b}{M gaudium etc.}. Et infra: sententia dicta est de scripturis, nam qui sine timore est, non poterit iustificari. Opus ergo est, ut intret primo timor, per quem veniat\apnota{15}{M veniet.} charitas. Timor medicamentum, charitas sanitas. Ex istis apertissime liquet, quia poenitentia est medicinalis et timor est medicamentum secundum Augustinum et charitas non intrat nisi per timorem, ideo bene praedicatur, poenitentiam a timore incipere.

Et non accipio, quod reverendus pater ab initio resolutorii refert reverendi patris Staupicii\apnota{20}{M stabicii.} vocem quasi caelitus demissam, poenitentiam incipere ab amore et dilectione. Nam huic expresse reluctatur b. Ambrosius, bonus animarum medicus, in epistola ad Studium: Ubi poena praescribitur, debet esse poenitentia peccatorum; ubi remissio donatur, gratia est; praecedit poenitentia, sequitur gratia. Neque ergo poenitentia sine gratia, neque gratia sine poenitentia. Debet enim poenitentia prius damnare peccatum, ut gratia possit id\apnota{26}{M U illud.} abolere. Illud

idem testatur Chrysostomus toto libro de cordis compunct. et Homilia 80. de poenit. et Sermone 29.

Addo etiam Isidorum libro 2. de summo bono. ca. 12.: Compunctio cordis est humilitas mentis cum lacrimis, exoriens de recordatione peccati et timore iudicii, et toto libro de contritione cordis et speculo peccatorum Augustinus nihil aliud agit, imo hortatur peccatorem volentem poenitere, ut recogitet tres abyssos\apnota{32}{P abyssus.}, scilicet peccatorum suorum, poenarum et iudiciorum dei. Sic b. Bernhardus super Canticis ser. 16. et aliis pluribus locis, Gregorius in moralibus lib. 2. et 5., et item super ca. 29. Job, libro quoque 2. super Ezech., ubi Homilia 19. inquit: Scriptum est: Initium sapientiae timor Domini. Constat procul dubio, quia a timore ad Dominum ascenditur, non autem a sapientia ad timorem reditur. Propheta ergo de caelestibus ad ima loquebatur, coepit magis a sapientia et descendit ad timorem. Sed nos, qui a terrenis ad caelestia tendimus, eosdem [361] gradus ascendendo numeramus, ut a timore ad sapientiam pervenire valeamus. Idem lib. 1. super Ezechiele\apnota{2}{M Ezechielem.} dixit, duas esse pennas, scilicet timorem et poenitentiam, quae tegere corpora dicuntur, id est, peccata.

Ex quibus omnibus, quoad unam particulam de timore, volo habere conclusum, bene praedicasse nostrae tempestatis et superioris praedicatores, quod poenitentia a timore incipiat\apnota{6}{P poenitentiam . . . . incipiunt.}: et quod nituntur sermonibus suis in populo timorem dei seminare, ut sic semen diaboli extirpetur, iuxta illud Origenis lib. 3. super Job: Bonus est timor poenae et iudiciorum, quem nisi daemon expulerit, non poterit seminare semen peccatorum. Ex quo dein\apnota{9}{U deinde.} consectarium est, poenitentiam non incipere ab amore et dilectione iustitiae\apnota{10}{M incipere a dilectione et iustitia, a. R. 'ab amore'.}, quamvis si sic inciperet, faterer\apnota{11}{M inciperer faterer, U inciperet fateor.} eam esse laudabiliorem et perfectiorem, quam quod a timore poenae incipit\apnota{12}{M U inciperet.}. At fragilitas nostra id non patitur, cui Dominus Jesus et praedicatores condescendentes timorem praedicant, tamquam gradum, quo verum iustitiae amorem consequamur.

D. D. Martinus.

Non est iste modus scripturas divinas feliciter intelligendi vel interpretandi, si ex diversis locis diversa decerpantur dicta, nulla habita ratione vel consequentiae vel collationis. Imo iste est canon errandi vulgatissimus in sacris literis. Oportet ergo theologum, si nolit errare, universam scripturam ob oculos ponere et contraria contrariis conferre, et sicut

duo Cherubim adversis vultibus\apnota{22}{U Cherobin statt Cherubim; M fultibus, verbessert vultibus.} utriusque diversitatis consensum in medio propitiatorii invenire, alioquin cuiuslibet Cherubim vultus longe divertet sequacem oculum a propitiatorio, id est vera Christi intelligentia.

Proinde nunquam mihi egregius D. D. remotior visus est a sacris literis, quam hodie, atque eo magis, quod concludendo fatetur, laudabiliorem et perfectiorem esse poenitentiam, si ab amore iustitiae inciperet, sicut ego sapio, quam si a timore poenae inciperet\apnota{27}{fehlt in M.}. Quasi non conandum sit, ut laudabiliter et, ut Johannes ait, dignos\apnota{28}{M dignus.} fructus poenitentiae agamus. Nam hoc prorsus non accipio, quod dicit, fragilitatem\apnota{29}{P fragilitarem.} nostram obstare, quominus ab amore iustitiae poenitentiam incipiamus. Si secundum fragilitatem nostram agendum est et scriptura exponenda, nunquam poenitebimus, sed de die in diem peiores fiemus. Quare antequam respondeam suis obiectis, primum declaro\apnota{33}{M U primo declaro (M declaram).} me ipsum.

Omnis bona vita necesse est ut instituatur per aliquam legem, ideo lex principium est poenitentiae et\apnota{34}{fehlt in M U.} cuiuslibet boni operis. Quare in\apnota{35}{M U et in.} poenitente ante omnia oportet vel revelari vel suggeri legem, contra quam fecerit, et secundum quam facere debet. Lege autem manifestata aut in memoriam revocata mox sequitur augmentum peccati, si desit gratia, quia naturaliter odit voluntas legem, ut sunt expressa testimonia Pauli ad Rom. et ad Galat. Rom.\apnota{39a}{fehlt in U.} 5. [v. 20]: Lex intravit, ut abundaret delictum, Galat. 3. [v. 19]\apnota{39b}{P M U Gal. 5.}: Lex [362] propter transgressiones posita est. Ideo b. Augustinus libro\apnota{1}{fehlt in U.} de spir. et lit. c. 3: Neque enim liberum arbitrium ante gratiam quicquam valet, nisi ad peccandum, non autem ad poenitendum, ut D. D. dicit (hoc enim Pelagianum est). Sequitur Augustinus: Et cum coeperit non latere\apnota{4}{M eum coeperit cognoscere, Mc 'non latere', U cum coeperit cognosci se non latere.}, quod faciendum est, nisi spiritus sanctus diffuderit charitatem in cordibus nostris, non diligitur, non suscipitur, non bene vivitur. Ibi clarissime dicit Augustinus, quod lex dei non potest diligi, nisi accepta gratia spiritus\apnota{7}{U spiritus sancti.}. Si autem non diligitur lex, contrarium eius peccatum non oditur. Ergo impossibile est poenitere ante dilectionem legis. Hoc est, quod Roma. 4. [v. 15] apostolus vult: Lex iram operatur, hoc est, monstrat peccatum, sed non dat gratiam, ut odiatur peccatum. Ideo manet odium legis et dilectio peccati, quantumlibet per increpationes forinsecas\apnota{12a}{U fornisecas, Uv forinsecas.} aut intrinsecas homo

concutiatur timore servili\apnota{12b}{fehlt in M.}. Nam etsi abstinet ab opere peccati, non tamen abstinere potest ab amore peccati. Hoc et Christus Johan. 6. [v. 44] docet, ubi dicit: Nemo venit ad me, nisi pater meus traxerit eum. Concedo ergo, quod lex, recordatio peccatorum, intuitus poenarum possunt terrere peccatorem, sed nunquam facient\apnota{16}{M U faciunt.} poenitentem.

Respondeo ergo ad obiectionem primam, de filio prodigo\apnota{17}{M prodago.} Luc. 15, ubi incepit poenitentiam a magnitudinis praemii recogitatione\apnota{18}{M a 'recogitatione' magnitudinis praemii, U recordatione magn. praem.} dicens: Abundant panibus\apnota{19a}{M U panibus mercenarii.} etc. Dico, quod hic prodigus filius\apnota{19b}{U filius prodigus.} incepit vere ab amore iustitiae, quia conversus in se cognovit primo bonum et ex cognito bono intellexit suum malum. Hanc autem conversionem sui in se ipsum non habuit ex fragilitate sua aut ex timore poenae, sicut ipsemet D. D. hic\apnota{22}{M hic D. D., in U fehlt hic.} dicit, quod timor poenae post praemiorum magnitudinem moverit. Ergo non incepit a poena et timore\apnota{24}{M a poeno timoris.}. Habuit autem eam ex trahente intus patre et dilectionem inspirante paternae domus, dum dixit: Quanti\apnota{25}{M quanto.} mercenarii in domo patris mei. Nam cum prius vixisset\apnota{26}{M dixisset.} in peccatis, bonum nec cognovit nec dilexit, nec peccatum odivit: cum tamen non posset ignorare peccatum. Ergo necessaria fuit alia affectio, id est, amor boni.

Secundo induxit Johannem baptistam, Luc. 3 [v. 7], quod inceperit a timore, dicens: Quis monstravit\apnota{30}{M monstrabit.} vobis fugere etc.? Respondeo: Aliud est praedicare poenitentiam, aliud incipere poenitentiam; aliud est praedicare bonum opus, aliud incipere bonum opus. Praedicator monet, terret\apnota{32}{P movet terret, M monet terrat.}, allicit etc., sed nihil sequitur, nisi gratia moverit voluntatem.

Idem dico ad Psal. 89 [v. 33]\apnota{34a}{U psal. 58.}: Visitabo in virga iniquitates eorum\apnota{34b}{iniquitates eorum fehlt in M.} etc. Percuti potest peccator, sed nisi gratia cooperetur, nihil proficitur\apnota{35}{M proficit.}, ut Jere. 5 [v. 3]: [363] Percussisti eos, et non doluerunt\apnota{1}{P M U percussi eos et non doluerunt (M U voluerunt, Uv = P).}; et Esai. 1 [v. 6]: A planta pedis usque ad verticem non est in eo sanitas: super quo vos percutiam\apnota{2}{M U percutiam vos.} ultra? quia nihil profecit\apnota{3}{M perfecit.} percutiendo.

Et satis admiror egregium D. D., quod oblitus luminis naturae Aristotelis audeat statuere alicuius virtutis initium\apnota{5}{M in ipsum, gestrichen, initium.} in timore urgente, cum ille tot verbis iam receptissimis 3. Ethicorum persuadere conetur: Oportere esse voluntarium opus bonum et libera voluntate fieri. At voluntas certe aut amor est aut soror amoris. Inde etiam dicitur, oportere opus bonum fieri per liberum arbitrium\apnota{9}{U arbitrum.}. Addo et illud, quod Christus nunquam peccatores coegit timore ad poenitentiam\apnota{11}{M poenitentias.}, sed suaviter allexit quoscumque vocavit, ut Zachaem, Magdalenam, apostolos et omnes. Sicut et apud Jere. cap. 31 [v. 3]\apnota{12}{cap 31 fehlt in P M; U cap. 2.} dicit: Perpetua charitate dilexi te, ideo attraxi te miserans\apnota{13}{M U miserans tui.}.

Dico ergo, quod timor Domini quidem necessarius est, sed filialis, quia sine amore impossibile est ferre conversionem sui, in qua terretur, conteritur et humiliatur peccator, iuxta illud I. Reg. 2 [1. Sam. v. 6]: Dominus deducit ad inferos et reducit. Credo autem, etiam ipsius egregii D. D. sententiam et omnium scholasticorum doctorum mecum esse et contra obiecta eius pugnare, cum omnes consentiant, contritionem oportere fieri in charitate, si debet esse bona et meritoria poenitentia. Quod ego sane intelligo contritionem fieri movente et imperante charitate, ut\apnota{20}{M quod.} sic sit voluntaria, hilaris, amorosa poenitentia. Ideo licet\apnota{21}{M licet, Mc a. R. dicit.} Johannes increpaverit Judaeos et terruerit, non ideo sequitur, poenitentes incipere etiam\apnota{22}{M etiam incipere.} a terrore. Aut si incipiunt a terrore, nisi accesserit gratia, hypocritae sunt verius quam poenitentes.

Adiecit D. D. rationem, quod deus res sic\apnota{25}{M sic, gestrichen, deus res sic.} administrat, ut ima per media ducat ad superiora, quod aptavit\apnota{26}{M appavit, Mc aptavit.} ad timorem, volens peccatum esse imum, timorem medium, amorem superius, quod ego transeo et non accipio.

Dicit etiam, esse perfectionis, si homo possit se mero\apnota{29}{M U posset se mero (U mere).} dei obtutu et amore iustitiae erigere ad gratiam capessendam. Sed quis est hic, et laudabimus eum? Respondeo: Nec timore nec amore potest se erigere homo\apnota{31}{M U homo erigere.} ad gratiam capessendam, sed gratia praevenit et movet ad merum dei obtutum et amorem iustitiae.

Ad Augustinum, qui doceat, praecedere timorem ante gratiam et intrante charitate pelli\apnota{35}{M expellit.} timorem, dico: Si recte intelligatur, admitto, hoc est, quod poenitentia nondum est incepta, quando timor praecedit charitatem. Sed intrante charitate incipitur poenitentia, id est, amor iustitiae et odium peccati. Si autem charitas non intraret\apnota{38}{M intret, Mc 'intraret'.}, timor non operaretur nisi maiora peccata. Quare ista similitudo setae\apnota{39}{M cetae.} et lini probatur mihi, modo non intelligatur, quod timor inducat charitatem, quod ex ipsius D. D. verbis capio, quod nisi exeat seta\apnota{40}{M ceta.}, [364] non succedit linum, hoc est, nisi timor, qui impedit poenitentiam veram, pellatur\apnota{2a}{M pellat.} per intrantem charitatem, nunquam vere poenitetur, ut sic dixerim\apnota{2b}{ut sic dixerim fehlt in M.}.

Ambrosium inductum in epistola ad Studium, quod praecedit poenitentia et sequitur gratia, et alia, item Isidorum de summo bono etc., item tres abyssos\apnota{4}{etc. fehlt in U; P abyssus statt abyssos.} recogitationes peccatorum, poenarum, iudiciorum\apnota{5a}{M U recogitandos statt recogitationes; poenarum, iudiciorum fehlt in U.}), Bernhardum super canticis, Gregorium in Moralibus\apnota{5b}{M Greg. in moral Bernh. super cantica (so!)}) et alios patres, qui docent a timore\apnota{6}{M a timore, Mc 'per timorem'.}) ascendendum ad sapientiam et poenitentiam a timore incipiunt, libenter admitto, sed non contra apostolum Paulum de lege et timore legis docentem intelligo. Dico et ego, habita charitate simul moveri hominem ad timorem dei, et sic incipi\apnota{9}{M incipit.}) poenitentiam a timore in charitate, alioquin stat firma sententia, quod timor poenam habet et\apnota{10}{fehlt in U.}) bonum non operator, sed odit legem.

Non ergo egregius D. D. concludit, poenitentiam non incipere ab amore iustitiae, quantumlibet bonus sit timor poenae ex origine\apnota{13}{P M ex Origene.}). Non enim timore poenae, sed timore dei poenitendum est, quod ille sit servus, non mansurus in domo, hic autem filius et heres. Ideo et illud Proverb. 1 [v. 7]: Principium sapientiae timor Domini, non admitto intellectum de timore poenae, qui ante gratiam torquet hominem infructuose, cum expresse dicat: Timor Domini, non timor poenae. Timor poenae potius est principium insipientiae. Videat ergo egreg. D. D., ut non in unum chaos confundat timorem servilem et filialem\apnota{19}{U timorem filialem.}), ne sibi ipsi\apnota{20}{U sibi ipse.}) scripturae et patrum intelligentiam praecludat.

D. D. Eccius.

Reverendus pater fuco verborum conatur se evolvere ex tam expressis sacrae scripturae et s. patrum auctoritatibus. Et ut hoc vobis persuaderet, ausus est dicere, me longe fuisse alienum a sacrae scripturae intelligentia, hoc loco diverticulum quaerens de Cherubim se respicientibus. Sed iudicent, quorum interest, uter rectius\apnota{26}{M recte.}) sentiat de sacra scriptura. Tamen ut suas solutiones impertinentissimas excludam, duo frequenter adducit, quorum in sermone recitando viam communem non meminit\apnota{28}{U secundum viam comm. non memini.}), nec ullus praedicatorum, quod meminerim\apnota{29}{M meminerit.}), negavit, et hunc quoque timorem, in quantum disponit ad veram poenitentiam, praeveniri inspiratione divina. Nam indubitatum est apud Christianos contra Pelagii\apnota{31a}{U Christianum statt Christianos; M Pelagii, Mc Pelagianorum.}) perfidiam, quod salutis nostrae initium\apnota{31b}{M initium nostrae salutis.}) deo inspirante habemus. Quare non fuit necessarium, hoc afferre vel ob hoc modum praedicandi reprehendere.

In uno tamen videtur mihi aequivocatione falli, quod hanc gratiam, qua deus praevenit corda hominum movendo, credit esse charitatem, cum sit aliud gratuitum dei munus. Et per hoc resoluta est b. Augustini de spir. et lit. sententia contra improbos Pelagianos.

Alterum, quod dicit, timorem esse infructuosum, nisi accedente charitate, quis hoc unquam vel scholasticorum vel praedicatorum negavit, qui omnes apostoli Pauli ad Corinth. de charitate sequuntur sententiam?

[365] Porro, quod laudabiliorem dixi poenitentiam, quae ab amore incipiat, assumpsit nos debere iam\apnota{2}{M U illam.}) facere iuxta illud Johannis: Facite dignos fructus poenitentiae [Matth. 3 v. 8]. Nam si secundum fragilitatem nostram esset agendum, nunquam poenitebimus. Dico nos quoque posse facere dignos fructus poenitentiae, etiam si a timore inceperimus et ad charitatem perveniamus.

Et miror, quod reverendus\apnota{5}{fehlt in M U.}) D. pater vult nos facere angelos et fragilitatis nostrae oblitos\apnota{6}{U oblitus.}), cum b. Gregorius lib. 2. super Ezechielem expresse declaraverit\apnota{7}{M declaverit.}), prophetam descendisse de sapientia ad timorem, sed nos eosdem gradus ascendendo a timore ad charitatem pervenimus, sicut multa alia fragilitati nostra condonantur.

Mentem suam declaraturus ostendit, quomodo dilectio legis praecedere debeat poenitentiam: quia per legem abundavit delictum [Röm. 5 v. 20], et ad Gal. 3. [v. 19] etc.\apnota{11}{P M U ad Gal. 5 etc. (etc fehlt in M).}). Dico: Verum est, ut s. Augustinus contra Faustum lib. 22. testatur, peccatum factum vel dictum\apnota{13a}{U dictum, Uv delictum.}) contra legem ideo non esse\apnota{13b}{M esset.}) peccatum, si non esset lex prohibens, sed sicut legis transgressione abundavit\apnota{14}{M U abundat.}) delictum, ita legis observantia augetur meritum. Quare omnino ista diverticula ad propositum non conferunt: et solutiones datas esse contra sacrae scripturae intelligentiam\apnota{26}{(Sin definición en el aparato impreso; la llamada 26) del cuerpo carece de correlato y la nota 16 carece de llamada — probable errata de numeración del editor.)}), post prandium manifestabimus.

Hora secunda pomeridiana continuabitur disputatio.

Hora igitur secunda eiusdem diei

D. D. Eccius prosecutus est replicam suam hodie coeptam\apnota{20}{Die Zeitbestimmung fehlt in M, dafür nur Prosequar replicam; U hora secunda continuavit Eccius semiplenam orationem, a. R. Prosequamur hodie incepta.}).

Reverendus pater ad illud de filio prodigo\apnota{21}{M prodogo.}) respondit, praecessisse amorem iustitiae, iuxta illud: »Et conversus in se«, quia memoriam paternae domus habuerit. Sed hoc non diluit obiectum, quia conversio

illa fuit facta contemplatione poenae, dum neminem haberet, qui eum siliquis satiaret, durante enim pecunia non convertebatur in se, sed famelicus conversus in se etc.\apnota{25}{U est conversus; etc. fehlt in M U.}). Accedit, quod b. Augustinus ait, verba eius fuisse verba poenitentiam meditantis, nondum agentis. Quodsi conversus in se fuisset amore iustitiae, iam poenitentiam incepisset contra Augustinum. Quare solutio sua propositum non infringit et Basilio manifeste reluctatur.

Ad Johannem respondit, aliud esse incipere, aliud praedicare, et alias dixit, praedicatores terrere, sed non facere poenitentem per poenarum cruciatum\apnota{31}{M cruciatu statt per cruciatum.}). Iterum evanida est solutio, quia si Johannes praedicavit, certe intendebat fructum suae praedicationis, et quia insinuabat\apnota{33}{M insinuebat, U insinuavit.}) eis timorem, signum est hoc pacto poenitentiam auspicari. Et nostri praedicatores praedicant, quomodo incipiant, scit deus.

Porro quod poenitentiae initium\apnota{36a}{P nuntium.}) gratiae tribuit, et hi\apnota{36b}{U nostri.}) praedicatores et doctores nunquam negaverunt, quia deus aspirando praeveniat\apnota{37}{M U praevenit, Uv = P.}).

Praeterea Aristotelem non vult acceptare in scholis theologicis, et tamen audet eum mihi opponere. Sed dico, constantem esse ethnicorum sententiam, [366] sistendo in timore poenae non esse perfectam virtutem\apnota{1}{M in timore (so!) . . . . . perfectum retuli.}), iuxta illud: Oderunt peccare mali formidine poenae, oderunt peccare boni virtutis amore\apnota{2}{U amore etc.}). Nec opus virtutis debet esse coactum, sed liberum. Aliud tamen est liberum arbitrium induci, aliud\apnota{4}{M U aliud est.}) liberum arbitrium cogi. Fateor, Christus benevolis verbis vocavit apostolos, Zachaeum, Magdalenam. Sed ipsum aliquando durius vocare, factum Pauli nobis innuit, de quo Augustinus ait 23. quaest. 4. can. Quis nos\apnota{7}{fehlt in U.}): Cui Christus vim intulit? quem coegit? et in evangelio de vocatione ad coenam dicit: Compelle intrare. Sicut pluribus verbis b. Gregorius in homilia prosequitur. Taceat ergo, quaeso, reverendus pater, solum benevolentem afferens\apnota{10}{W. A. verbessert asserens.}) vocationem.

Praeterea solum timorem filialem necessarium arbitratur poenitenti, de quo plurimum miror, cum hodie me ad veram\apnota{12}{M verarum.}) scripturarum intelligentiam monuerit, quod se ipsum primo non emendavit, nam et b. Augustinus in loco hodie allegato super Johannem loquitur de timore servili, quem foras\apnota{14}{U super Johanne; P M foris.}) mittit charitas, et de illo dicit loqui Sapientem: Initium sapientiae timor Domini. [Prov. 1 v. 7] Et ita de timore servili intelligit\apnota{16a}{M U intellexit.}) glossa super Psal. [111 v. 10]: Initium sapientiae etc.\apnota{16b}{fehlt in M.}) Metus iudicii ianua est conversionis ad deum, et est glossa Cassiodori. Sic glossa super illud apostoli ad Roma. 8 [v. 15]: Non enim accepistis

spiritum servitutis in timore, dicit timorem servilem esse bonum et a deo. Sic Paulus testatur illum bonum 1. ad Timoth. 5 [v. 20]\apnota{20}{P M U 1. Tim. 4.}): Peccantem coram omnibus argue, ut et ceteri timorem habeant. Ita Augustinus super Psalmis ait, et Longobardus eum renarrat\apnota{22a}{M enarrat.}) in 3.: Timor servilis est, cum per timorem gehennae\apnota{22b}{M jehennae.}) continet se homo a peccato, quo praesentiam iudicis et poenas metuit. Et infra: Bonus est iste timor et utilis licet insufficiens, per quem paulatim fit consuetudo iustitiae. Beda quoque exponit: Initium sapientiae timor Domini, scilicet servilis.

Quare aliud quaerat sibi patrocinium et aliam det solutionem, quando ista\apnota{27}{M U quoniam; M istam, verbessert ista.}) s. patribus non concordat. Quodsi nolit dictis sanctorum stare, habebit adversum se sacrae scripturae locos. Nam dum Domini timor\apnota{28}{M U timor domini.}) est initium sapientiae, adhuc charitas foras\apnota{29}{P M foris.}) pellit timorem, sed timor filialis iuxta David in aeternum permanet et in saeculum saeculi. Quare alius est timor, quo sapientia initiatur, alius est, qui cum aucta permanet. Unde cum reverendus pater negat timorem inducere charitatem, expresse s. Augustino contradicit, qui in loco allegato hoc totum ait\apnota{33}{M agit.}), timorem scilicet servilem esse initium sapientiae et per charitatem iterum expelli. Subscribit sententiae Augustinianae mellifluus Bernhardus Ser. 52. super Canticis, quem brevitatis gratia non adduco. Sed etiam plus dico, charitatem non sic expellere timorem, sicut [367] unum incompatibile\apnota{1a}{M incompatibile.}) aliud, sed successu ad augmentationem gratiae minui\apnota{1b}{M minuit.}) timorem servilem et aucta gratia timorem omnino auferri.

Nec refert, quod ponderabat hodie D. pater: Timor Domini, non timor poenae est initium sapientiae, quoniam et timor poenae et timor reverentiae uterque est Domini et uterque est a Domino. Sola media timendi sunt diversa, ut Beda et alii per Longobardum inducti testantur praeter Augustinum prius citatum. Et satis liquet ex verbis Christi, si nihil aliud haberemus, Matth. 10 [v. 28]: Nolite timere eos, qui occidunt corpus, animam autem non possunt occidere, sed potius timete eum, qui potest animam et corpus perdere in gehennam\apnota{9}{M jehennam.}). Si enim timor servilis esset damnabilis, sicut dicit reverendus pater, cur Christus nos ad eum invitaret? et Augustinum et alios ita intellexisse, confero in iudicium eorum, quorum intererit.

Insuper dicit, scholasticos omnes concordare, quod contritio nullius sit utilitatis, nisi fiat in charitate. Hoc quidem verum, sed secundum Augustinum ad charitatem non pervenitur, nisi per timorem. Timor medicamentum est, charitas sanitas. Unde non accedente charitate omnes

bene noverunt timorem illum insufficientem, propterea in via illum constituunt\apnota{17}{M illum in via constitutum illum (letzteres gestrichen).}), non in termino.

At\apnota{19}{M U ad.}) rationem ex b. Dionysio sumptam transivit, sed non aperuit D. pater, cum peccatum sit imum et charitas superius, quod sit medium perveniendi\apnota{21a}{M perveniendo.}) a peccato in charitatem, et profecto aliud dare nequit\apnota{21b}{M nequid.}), quam Basilius, Beda, Augustinus, Bernhardus et ceteri\apnota{22}{et ceteri fehlt in M U.}), timorem.

Deinde\apnota{23}{M Dein.}) assumit gratiam praevenire timorem et amorem, quod si loquitur de gratia motionis divinae, qua deus nos aspirando praevenit, fateor quidem\apnota{24}{M fateor (Mc 'accipio') tunc, U fateor quidem tunc.}); si autem loquitur de gratia charitatis, illud non accepto, quia illius\apnota{25}{M tunc illius.}), quae recte\apnota{26}{M recta.}) sapientia in scripturis appellatur, initium est timor Domini, quamvis timor iste absque charitate sit infructuosus.

Porro, quod Ambrosium, Gregorium et alios doctores admittit, sed non contra apostolum de timore legis, nescio quas offundit\apnota{29}{M oppugnat.}) nebulas. Dicat clare, an\apnota{30}{U aut.}) isti s. patres per me allegati repugnent dictis apostoli vel minus. Si repugnant\apnota{31}{M repugnat.}) dicto apostoli, hoc ostendat, quod nunquam fieri posse arbitror, cum in sacris scripturis fuerint exercitatissimi et spiritu sancto repleti, apostolum Paulum aeque bene intelligentes sicut nos. Si non repugnant dictis Pauli, stet eorum sententiae et praedicatores modum contritionis et poenitentiae praefatum docentes non improbet, nec ab hominibus peccatoribus huiusmodi timorem servilem utilem et quasi necessarium medium excludat.

Non relevat eum, quod existimat habita charitate moveri mentem ad timorem, quoniam hoc esset cancrino more retrogradi\apnota{38}{M retrogradi.}) contra b. Gregorii sententiam lib. 2. homilia 19. super Ezechielem\apnota{39}{U super Ezechiele.}). Constat procul dubio, quia a timore ad sapientiam ascenditur, non autem a sapientia ad timorem reditur. Habet ergo timor poenam, quae charitate aucta minuitur et totaliter absorbetur.

[368] Praetereo ob temporis angustiam alia in eodem sermone contenta, quae tamen omnino venirent discutienda, nisi hodiernus dies praefixus esset ad hanc materiam terminandam, sicut est punctus de sigillata confessione peccatorum in specie\apnota{4}{M de sigillato, sonst = P; U de sigillata peccatorum confessione in speeu.}), uti meminit Chrysostomus super caput 12. Matth. de puncto, quod scholastici ponant tres partes poenitentiae, scilicet contritionem, confessionem et satisfactionem, quas prius posuit Chrysostomus sermone 29. de poenitentia, de duplici modo confitendi sacerdoti magna et alia deo cum aliis punctis, quae omnia

prolixe possent\apnota{8}{M possint.}) tractari et excuti. Sed iacto iam fundamento totum simul sermonem et ea, quae adduxi pro parte mea, confero in iudicium iudicum eligendorum.

D. D. Martinus.

Sperabam egregium D. D. confutaturum meas solutiones, praesertim ea, quae pro fundamentis ex Paulo adduxeram\apnota{13}{U induxeram; M ea quae pro fundamento ex Paulo aduxeram (so!) 'quae pro fundamentis' (Mc 'eam quam' am Anfang).}) de lege, quae ante charitatem non operatur nisi iram et auget peccatum, taceo, quod disponat ad gratiam, quemadmodum et\apnota{15}{fehlt in U.}) timor servilis fructus legis iram operatur et auget peccatum. At ipse haec fortiter transiliens, eandem cantilenam recantavit, ac nobis haec non fucum\apnota{17}{U succum.}), sed medullam scripturae contendit persuadere, discurram per singula.

Primum dicit\apnota{19a}{fehlt in M, U proinde dicit.}), non fuisse necessarium afferre, quod initium salutis nobis\apnota{19b}{M U nostrae, Mc 'nobis'.}) sit ex deo inspirante, ubi et timorem disponentem dari intelligit. Respondeo: Haec plane omnia sunt erronea propter verbum Pauli, qui dicit: Impossibile esse legem impleri, imo peccata non augeri, nisi spiritus sanctus diffundat charitatem in cordibus nostris. [Röm. 8 v. 3 ff.] Has auctoritates tam expressas oportuit confutatas, et illam b. Augustini cap. 3. de spir. et lit.: Dum\apnota{24}{M in libro de spir. et lit. c. 3.; M U cum statt dum.}) coeperit non latere, quomodo vivendum est, nisi gratia donetur, non suscipitur, non bene vivitur. Vadat ergo cum suo servili timore, qui non operatur nisi odium legis et dei et cum iniuria vocatur disponens ad gratiam.

Invenit etiam aequivocationem gratiae, aliam esse charitatem, aliam donum, quo primo\apnota{29a}{U primum.}) movemur. Transeat illa\apnota{29b}{M ista, U haec.}) distinctio, nihil ad propositum, est elusio\apnota{30}{M delusio.}) verborum Pauli. Clarissimus est textus Pauli, nisi gratia et charitas faciat nos diligere legem, lex semper iram operatur. Hoc donum autem, quo primum movemur, non diliget\apnota{32}{M primo movemur, non diluet.}) legem, sed charitas spiritus.

Tertio concessit, timorem esse infructuosum sine charitate, et hoc nullum negasse. Respondeo: Cur ergo docent infructuosas poenitentias et resistunt mihi fructuosam docenti? Relinquo, confutationem hanc esse omnino subversam\apnota{35}{M docentem Relinquo hanc omnino esse superstem, U docenti Relinquo hanc esse omnino subversam.}), cum tot opusculis Augustinus\apnota{36a}{W. A. Augustini (unnötige Änderung).}), et si Augustus non faceret tot fulminibus\apnota{36b}{M voluminibus, Mc 'fulminibus'.}), tamen unus Paulus concludit, opera quaecunque ante charitatem esse peccata et damnabilia [369] et

indisponentia ad gratiam. Et per haec dicta ferme ad omnia sua replicata responderi potest. Ipse quidem pro more suo dicta scripturae divinae temperat ad dicta patrum, imo trahit ad intelligentiam suam, quam habet in patribus, cum potius contra scripta patrum debeant conferri et iudicari ad dicta scripturae.

Quarto non accipit, quod assumpsi, nos debere facere laudabiliorem illam poenitentiam, sed dicit, nos posse etiam dignos fructus poenitentiae facere, si a timore inceperimus. Hoc secundum Augustinum\apnota{7}{M U beatum Aug.}) sic expono: Si peccata peccatis addiderimus, cum, ut saepius iam dictum est, quidquid ante sanatam\apnota{9}{M hec gratia (?), gestrichen, sanatam.}) per gratiam voluntatem fit, fructus malae arboris sit, quae non potest fructus bonos\apnota{10}{M U bonos fructus.}) facere. Ergo per timorem servilem nunquam pervenitur ad charitatem.

Quinto, innixus Gregorio super Ezechielem, qui descendat de sapientia ad timorem, sed nos contra ascendere\apnota{13}{U descendere, Uv ascendere.}) a timore ad charitatem docenti, sic respondeo, quod divus Gregorius non debet excludere charitatem nec exclusit a timore incipiente poenitentiam, sicut\apnota{15}{P sic.}) D. D. intelligit; sicut nec a charitate excluditur omnis timor servilis, praesertim in hac vita, cum charitatis officium sit expellere timorem servilem, atque id tota vita, et inducere timorem filialem.

Sexto, ad auctoritates Pauli, quod lex propter transgressiones posita est et auget delictum, dicit, nihil esse ad propositum et me diverticula quaesivisse, adducens Augustinum contra Faustum lib. 22: Peccatum esse dictum vel factum contra legem dei etc.\apnota{21}{fehlt in U.}); nescio quid velit D. D. Non est hic disputatio, quid sit peccatum, ideo superfluis verbis hoc tempus perditum est. Scopus iste versatur, quod auctoritate Pauli timor servilis non possit esse nisi peccatum et augere peccatum per virtutem\apnota{24}{M virtutum.}) legis, si desit gratia. Et sic auctoritates Pauli non sunt diverticula, sed fulmina conterentia sententiam D. doctoris usque ad pulverem. Non ergo ego ex hominibus angelos facio, oblitus fragilitatis, sed doceo, ne deos faciamus ex peccatoribus, dum obliviscuntur fragilitatis suae, qua non possunt nisi malum facere ante gratiam.

Quod conversionem filii prodigi factam dicit contemplatione poenae, nego. Ad probationem, quod neminem haberet, qui eum siliquis satiaret, dico: Nisi intus fuisset tractus et in se conversus, potius fame mortuus esset, quam rediisset\apnota{32a}{M rediisse.}), ut stet sententia: Nemo venit ad me, nisi pater meus traxerit eum [Joh. 6. v. 44]\apnota{32b}{M eum traxerit.}).

Quod Augustinus dicit\apnota{33}{fehlt in M.}), verba eius fuisse verba poenitentiam meditantis, nondum agentis, spero, quod pro me faciat, quando meditare

poenitentiam, praesertim medullitus in corde, sit certe incipere poenitentiam. Ergo illud »nondum agentis« oportet de perfectione ad extra intelligi.

Idem ad Basilium dicetur\apnota{37}{fehlt in M.}).

De verbo Johannis baptistae dicit, evanidam esse solutionem, quod dixerim, aliud esse poenitentiam docere, aliud incipere. Respondeo: Salva [370] reverentia egregii D. D.\apnota{1a}{U egre: d. doct.;, W. A. egregius dom. doctor.}), non videtur Paulum intelligere nec virtutem\apnota{1b}{M virtutis.}) legis cognovisse. Lex enim docet quidem sancta\apnota{2}{P facta.}), iusta et bona, sed sola gratia incipit, facit et perficit ea. Ergo etiam si\apnota{3}{etiam fehlt in M, von Mc hinter si nachgetragen.}) Johannes timorem docuisset esse initium poenitentiae, non ideo sequitur, poenitentiam incipere a timore. Sicut si insinuem quodlibet bonum opus per terrores et minas, non ideo incipit bonum opus\apnota{6}{U opus bonum.}) a terrore et minis, sed a charitate.

Indignatur, quod Aristotelem ei obiecerim in scholis theologicis, quem tamen non acceptarim. Fateor, minor est, quam ut valeat in theologica schola, cum seductor fuerit scholasticorum doctorum. Sed satisfacere volui conclusioni meae, quod contra sacrum suum Aristotelem desipiunt, qui poenitentiam a timore et non libera voluntate incipiunt. Nam distinctionem\apnota{11}{M distinctio.}) de libero arbitrio coacto et 'inducto' transeo. Non enim verum est, quod liberum arbitrium unquam cogatur ad bonum aut cogi etiam possit. Si autem cogitur, in diversa rapitur\apnota{14}{M raptum.}) et coactionem sui odit. Sola autem gratia trahitur, hoc est, vere liberum efficitur, ut Augustus contra Julianum lib. 2\apnota{15}{M contra Pelagianum (lib. 2 fehlt); U contra Jul. lib. 2 et in multis locis.}).

Ubi Christum dixi vocasse apostolos et Magdalenam blando tractu, contra opponit, Paulum vocatum esse durius, et in evangelio iussum esse servum, ut compelleret intrare. Primum satis admiror, quod cum nos exemplum Pauli soleamus adducere pro defendenda gratia, ipsi nobis elabuntur dicentes, hoc esse miraculosum factum et non facere regulam. Hic tamen egregius D. D. quasi non pro miraculo habens, pro regula inducit. Sed hoc dimitto. Dico, non potuisse Paulum intus ex corde converti, nisi trahente gratia, ut b. Augustus contra epistolas Pelagianorum docet, gratiam dei facere ex nolentibus ac resistentibus volentes et sequaces. Ita et illud in evangelio solvitur, quod servus potest compellere verbo, sed nisi Dominus sibilet, nihil sequitur\apnota{26}{U a. R. Ut Esaias loquitur.}).

Ad illud Augustini, intelligentis servilem timorem in verbo Salomonis: Initium sapientiae timor Domini, et illud glossae Psal. 111: Metus iudicii ianua est conversionis ad deum, respondeo: Si adsit\apnota{29}{M assit.})

gratia; alioquin servilis timor sine gratia (neque enim Augustus excludit) non operatur nisi iram. Oportet enim, ut\apnota{31}{M U quod.}) Augustus non pugnet, sicut vere non pugnat, cum Paulo damnante omnia, quae sunt extra gratiam. Glossam super illud\apnota{32}{M id.}) Roma. 8 [v. 15]: Non accepistis spiritum servitutis in timore, dicentem, esse timorem servilem\apnota{33}{M servilem timorem esse, U esse servilem timorem.}) bonum, potius reiicerem, tamquam expresse contra textum loquentem. Apostolus enim dicit: Non accepistis spiritum servitutis in timore, damnans eum. Aut dico, quod glossa\apnota{36}{fehlt in M.}) textum non exponit.

Illud apostoli ad Timoth. [1. 5 v. 20]: Argue coram omnibus, ut et\apnota{37}{U ut de, Uv ut et.}) ceteri timorem habeant, D. D. ad servilem timorem adaptavit, quod relinquo. Ego de filiali timore accipio, donec aliter probet. Quod autem Augustus magistro reci-[371]tante lib. 3. dicit: Timor servilis est, cum per timorem gehennae\apnota{1}{M pro timorem jehennae.}) continet se a peccato, quo praesentiam iudicis et poenas metuit etc.\apnota{2}{M poenae metuat (etc. fehlt); U quo poenitentiam iudicis, quo poenas metuit etc.}): continet, inquam, se a peccato externo tantum, intus tamen auget odium iustitiae gehennam comminantis. Et infra: Bonus est timor et utilis, licet\apnota{4}{M dicit, U licet et.}) insufficiens, per quem paulatim fit consuetudo iustitiae. Hoc est meo iudicio consuetudo desperandi et odiendi deum, si excludatur gratia: verum autem est, si includatur gratia.

Non est ergo\apnota{7}{M ergo est.}) necesse, ut dem aliam solutionem, nisi D. D. egregius primo\apnota{8}{M U primum.}) probaverit, s. patres loqui de servili timore, exclusa gratia, aut charitatem non habere negotium in expellendo timore servili. Quare, quae consequenter induxit de expulsione timoris, de initio sapientiae et quomodo expellat charitas timorem successu per augmentationem gratiae, ex praecedentibus satis intelliguntur\apnota{12}{P intelligitur.}).

Superest robustissima auctoritas Matth. 10 [v. 28.], quam unam sufficere putat: Nolite timere eos, qui occidunt corpus, sed timete eum etc.\apnota{14}{sed timete eum fehlt in M, U etc. statt qui occidunt corpus.}), ideo non esse servilem timorem damnabilem, ad quem Christus nos invitat. Respondeo primum: Si hoc vult D. D., contradicetur praecedentibus, ubi timor servilis insufficiens dicitur, ideo aeque absurdum\apnota{17a}{M absordum.}) est dicere, Christum nos docuisse insufficientia\apnota{17b}{U insufficientiam.}). Dico tamen ego, non esse ibi servilem timorem 'doctum'\apnota{18}{U domini.}), quia et filialis timor timet offendere deum et separari ab eo. Atque etiamsi de mero servili timore loqueretur, nondum intelligitur exclusa gratia, imo inclusa, cum teste apostolo et Augustino omnis lex et doctrina sit litera requirens spiritum.

Consensit etiam scholasticos vere dicere: Contritionem non valere extra charitatem, sed non confutavit. Stat ergo eorum ipsorum auctoritas contra eos: nisi confutet contritionem in charitate esse id, quod a charitate incipit fieri\apnota{24}{P M incipi et fieri.}).

Illud Augustini bene placet: quod timor est medicamentum, charitas est sanitas, timor scilicet in charitate imperfecta, et charitas perfecta.

Rationem illam sumptam ex Dionysio\apnota{27}{M U ex Dionysio sumptam (M sumpta).}), de imo, medio et summo, dicit me transiliisse\apnota{28}{M U transiisse.}) et non aperuisse, quod sit medium inter peccatum imum et charitatem superiorem. Fateor me libenter transiliisse\apnota{29}{M U transiise.}), ut qui crederem ipsi D. D. satis esse perspectum\apnota{30}{M U perspectam.}), hanc rationem nihil esse ad propositum. Dionysius enim loquitur de ordinibus infimis, mediis, summis\apnota{31}{M U infimis (U infinitis) mediis et summis.}). At ego prorsus nullum habeo medium inter peccatum et gratiam, sicut nec Christus, quando dicit: Qui non est mecum, contra me est. Et iterum: Aut facite arborem bonam aut facite arborem malam. Sed et ipsum D. D. credo id ipsum asserere, quod gratia et peccatum apud scholasticos immediate opponuntur.

Non accipit D. D., quod gratiam dixi praevenire amorem et timorem, nisi loquar\apnota{37}{M loquor.}) de gratia primae motionis. Respondeo: Ego cum apostolo et Augustino sentio: quod nisi lex diligatur\apnota{38}{M dirigatur.}) (quod est charitatis et non primae motionis), non bene vivitur. Ergo nec deus timetur seu\apnota{39}{M colitur seu timetur, U timetur nec colitur.}) colitur.

[372] In fine instruit adversus me cornutum syllogismum: aut Ambrosius, Gregorius aliique repugnant dictis apostoli vel non; si sic, ut\apnota{2}{M sic sic ut, U si sic quod.}) ostendam; si non, ut stem eorum sententiae. Respondeo et per medium transeo: non repugnant dictis apostoli, si\apnota{4}{fehlt in M U.}) et sto eorum sententiae, non autem Eccianae intelligentiae, imo errori. Non enim excludunt charitatem a timore sive servili sive filiali.

Dicit D. D. esse\apnota{7}{fehlt in M.}) cancrino more retrogradi, quod dixi: Habita charitate moveri mentem ad timorem. Satis miror sive scorpionem, sive cancrum illum\apnota{8}{M U istum.}), cum et gentilis ille poeta dixerit\apnota{9}{M dixerim.}): Res est solliciti plena timoris amor. Quasi ignoremus, quod amor fons et caput est omnium affectuum. Ideo enim timetur poena et infernus serviliter, quia diligitur vita et voluptas pueriliter. Et ad hunc amorem et timorem\apnota{12}{M timorem et amorem; U pueriliter et serviliter. Ad hunc timorem et amorem.}) extirpandum diffunditur charitas dei, qua amemus aliam vitam et timeamus aliam mortem, id est, separationem a deo.

Transiit D. D. de punctis in sermone meo signatis et de toto

sermone discutiendo\apnota{15}{M et totum sermonem confutandum, U et de toto sermone confutando.}. Dico breviter: Habet calamum et papyrum, aggrediatur eum cum fiducia, videbitur, an sermonem confutaverit, an se ipsum irriserit. Haec\apnota{16}{M U vel se ipsum irriserit. Hoc (M haec).} relinquo iudicio ordinandorum.

D. D. Eccius.

Quia 'ob' temporis angustiam respondere non possum his\apnota{19}{fehlt in M; U iis.}, quae reverendus pater adduxit: refero tamen me\apnota{20}{fehlt in P, in M gestrichen wegen des folgenden me.} ad iudices, me nihil transiliisse neque fuisse de mente praedicatorum aut doctorum, legem adimpleri sine charitate, neque aliquos docuisse infructuosam poenitentiam, sed quomodo per timorem servilem ad fructuosam perveniatur. Et quod magnifico verborum apparatu recessit ab hodierna solutione, [D. D. Martinus dixit: non recessi\apnota{21a}{(En el aparato impreso numerada 24a — errata de numeración del editor.) U a. R. Martinus non recessi respondit. [M a. R. nur Mar.]}] qua\apnota{21b}{(En el aparato impreso numerada 24b — errata de numeración del editor.) U quia.} dixerat Sapientem loqui de timore filiali, iam vero admittit loquentem de timore servili: non tamen excludendo gratiam, quod nec textus, nec s. doctores patiuntur. Nam per sapientiam intelligit charitatem, et sic\apnota{27a}{M U ita.} initium esset ante initium. Totus\apnota{27b}{M initium et Bernhardus et Augustinus ita intelligit sic et totus.} quoque Augustinus non laborat, quomodo timor servilis cum charitate sit initium charitatis, sed quomodo timor servilis primo occupet mentem et primo ingrediatur et sic inducat charitatem, quod etiam est de mente Gregorii improbantis reditum cancrinum, ubi reverendus pater pro Gregorio scorpionem mihi adducit: Res est solliciti etc.\apnota{31}{M solliciti certe, U sollicita etc.}. Omnes ergo s. doctores hodie citati volunt, timorem servilem esse initium charitatis, ad intelligentiam saepe datam, et quam praedicatores solent dare, de quo me refero\apnota{34}{M U refero me.} ad iudices.

Crastina hora secunda propter festum divae Margaritae disputabunt iidem D. doctores in materia absolutionis a poena et a culpa per quemlibet simplicem sacerdotem.

[378] Die Mercurii tertia decima mensis Julii hora secunda pomeridiana\apnota{2}{statt Zwischenbemerkung und Zeitangabe M nur In die Margarethae}.

D. D. Eccius.

Circa materiam conclusionum quartae et quintae, reverende pater, quia tempus nobis praefixum est nimis artum, tamen tangendo fundum

negotii aliqua obiter attingemus, illud scilicet, quod quilibet sacerdos absolvat a poena et culpa, contra communem usum totius ecclesiae, et quod in sermone vulgari et conclusione dicitis, ex nulla scriptura probari, divinam iustitiam aliquam poenam vel satisfactionem expostulare a peccatore. Et illis innitor, quia in absolutione sacramentali culpa remittitur simultanee, sed poena peccato debita aeterna commutatur in temporalem. Hoc liquet ex\apnota{10}{P et.} sacrae scripturae traditione et s. patrum usu. Et ut brevior sim, expresse hoc\apnota{12}{fehlt in M.} voluit Ambrosius super Lucam, Hieronymus lib. 1. contra Jovinianum, Augustinus quaesti. 1. de octo Dulcitii quaestionibus, et pertinenter Ambrosius Luc. 5. testatur satisfactione poenam peccati dissolvi.

Scriptura est in promptu: quoniam peccatum Adae, culpa etiam remissa, punitur in tota posteritate. Quare transeunte culpa remanet poena. Quam rationem tangit Augustinus tractatu 124. super Matth. et glossa 2. Reg. [Sam.] 12. Similiter de David. 2. Reg. 24 [v. 10], ubi David poenituerat et percusserat cor eius: peccavi, dixit, valde in hoc facto, sed peto, transferas\apnota{19}{M U ut transferas.} iniquitatem servi tui. Modo peccatum non transfertur quoad culpam, quia, ut propheta ait, anima ipsa quae peccaverit morietur\apnota{21a}{M U morte morietur.}. Restat ergo, ut peccatum quoad\apnota{21b}{M quantum ad.} poenam transferatur. Sic per prophetam Nathan dicitur David: Transtulit deus peccatum tuum. Ubi glossa: Deus delictum delet, sed inultum non deserit. Aut enim homo in se poenitens punit\apnota{24}{Mc ac enim homo . . . ponit.}, aut deus cum homine vindicans percutit.

Et ut ad s. patres veniamus, a quibus usus et praxis ecclesiae principaliter ab ipsis\apnota{27}{fehlt in M U.} apostolis derivatus est. Nam primo nullum peccatum manet impunitum, et poena est ordinativa culpae. Quas rationes tangit Augustinus et refert Gratianus can.\apnota{29a}{fehlt in U.} Sicut primi de poenitentia, distinctione 1.\apnota{29b}{U dist. 7.}. Nam deus malum fieri non pateretur, nisi illud per iustitiam ordinando melius stare\apnota{30}{M melius dare.} in mundo faceret. Inquit ergo Augustinus lib. de poenitentiae\apnota{31}{M poenitentia.} medicina: Non sufficit mores in melius commutare et a factis malis recedere, nisi etiam de his\apnota{32}{U iis.}, quae facta sunt, satisfiat deo per poenitentiae dolorem, per humilitatis gemitum, per contriti cordis sacrificium, eleemosynis cooperantibus etc. Non enim dictum est, ut tantum abstineatis a peccatis, sed et de praeteritis\apnota{35}{U de peccatis.}, inquit, deprecare Dominum, ut tibi dimittantur etc.

[374] Sed reverendus\apnota{1}{M U Sed si reverendus.} pater, sicut in resolutorio et alias, meo quidem sensu, facere nititur, satisfactionem illam fieri ecclesiae, sed deum

[kalligraphisch geschrieben]; de absolutione sacerdotis (sacerdotali?); a. R. circa conclusionem 4 et 5; U Tertia decima Julii Mcccc. XIX. hora secunda. Circa materiam conclusionum quartae et quintae; a. R. Eccius.

talem poenam non requirere, et quam deus requirat, homo non auferat; contra ego oppono: Primo quod Augustinus in Enchiridio inquit\apnota{4}{fehlt in P.}, cap. 71. (loquitur de oratione dominica): Delet omnino haec oratio minima et quotidiana peccata, delet et illa, a quibus vita fidelium scelerate gesta, sed poenitendo in melius mutata discedit. Ubi constat, orationem dominicam delere venialia et etiam mortalia, non quoad culpam, sed quoad poenam, quia poenitendo mutata, quae mutatio Augustini intelligi nequit, nisi de mutatione poenae aeternae in temporalem.

Eos, qui fuerunt a quadringentis annis, non adduco, ut Wilhelmum Parisiensem, Wilhemum Altisiod.\apnota{11}{M Altiodorensem et alii (so!); U non adduco huc Wilhel: Paris: Altisidio:}, qui id pleno ore affirmant. Accedat Cyprianus in epistola ad Fidum, ubi sic ait\apnota{12}{P Fidium; ubi sic ait fehlt in M.}: Legimus literas tuas, carissime frater, quibus significasti de quodam presbytero Victore, quod ei, antequam plenam poenitentiam egisset et Domino deo, in quem deliquerat, satisfecisset, temere Therapius collega praepropere pacem dederit\apnota{15}{U dederat.}. Hic s. Cyprianus non dicit Victorem ecclesiae non satisfecisse, sed Domino deo.

Idem Cyprianus ad Cornelium papam de haereticis ait: Elaborant\apnota{17}{M elaboram.}, ut opus suum diaboli malitia consumment\apnota{18}{M consumen.}, ne vulneratos divina clementia in ecclesia sua curet. Miserorum poenitentiam mendaciorum fraude corrumpunt, ne deo indignanti satisfiat. Apertissime sanctus\apnota{20}{M U sacer, Mc 'sanctus'.} martyr haereticos increpat, qui poenitentiam ab ecclesia iniunctam existimabant non esse satisfactionem dei.

Et quamvis\apnota{22a}{U quamvis (et fehlt).} reverendus pater citra modestiam saepius mihi improperet\apnota{22b}{M improberet.}, quasi scripturarum sanctarum intelligentiam non idonee tractem, de quo iudicent hodie designati iudices, uter rectius de fide et sacrarum scripturarum\apnota{25}{M literarum,} sensu sentiat: attamen, quantum ego video, universi sacrarum scripturarum tractatores in hanc sententiam descendunt. Ad quid enim ecclesia vellet imponere tale opus\apnota{27}{M U onus.} grave poenitentibus, cum Dominus et caput ecclesiae illud non exigeret? Quam rationem in annotationibus tetigi, non tamen fuit mea, sed Chrysostomi, sicut refert Gratianus 26. quaest., 7. can. Alligant\apnota{29}{M P alligat, U allegat.}: homo cui grave pondus poenitentiae imponis, aut poenitentiam reiiciet, aut dum sufferre nequit\apnota{31}{M nequid.}, scandalizatus amplius peccat. Deinde etsi erramus, poenitentiam modicam imponentes, nonne melius est, propter misericordiam rationem reddere quam propter crudelitatem? Ubi enim paterfamilias largus est, dispensator non debet esse tenax. Si deus benignus, ut quid sacerdos vult esse austerus? Hoc pro confessoribus.

Idem Homilia 31. ait super epistola ad Hebraeos: Peccatum confessione minuitur, sed non minuitur quoad culpam, quia impium est, ut Hieronymus ait, a deo dimidiam sperare veniam, ergo minuitur quoad poenam.

[375] Addo Gregorium, 4. moral. § 49., ubi post multa ita dicit\apnota{1}{post multa ita dicit fehlt in M; U cap. 49 post multa ita dicit.}: Sed quia nullum peccatum deus inultum relaxat, aut enim nos hoc flendo insequimur, aut ipse iudicando, restat, ut ad emendationem suam semper mens solerter invigilet etc.\apnota{4}{fehlt in M U.}.

Hoc sensit et b. Hieronymus, verba sua transcripta sunt can. Mensuram, de poenitent. dis. 1: Mensuram temporis in agenda poenitentia idcirco non satis praefigunt canones pro unoquoque crimine, ut de singulis dicant\apnota{8}{M dicam.}, qualiter unumquodque emendandum sit, sed magis in arbitrio sacerdotis intelligentis relinquendum\apnota{9}{M relinquentis.} statuunt, quia apud deum non tam valet mensura temporis quam doloris.

Augustinum\apnota{11}{W. A. Augustinus.} quoque refert in can. Nullus: Nullus debitae gravioris poenae accipit veniam, nisi qualemcunque etsi longe minorem, quam debeat, poenam solverit, ita enim a deo impertitur\apnota{13}{M U impertitur a deo.} largitas misericordiae, ut non relinquatur iustitiae disciplina. Hic habet egregius D. D. Carolostadius Augustinum in defensione mea adductum, cuius sententiam prosequi non est instituti nostri in praesentia.

Accedat postremo Isidorus lib. 2. de summo bono, cap. 13: Quamvis per poenitentiam propitiatio peccatorum sit, tamen sine metu hoc esse non debet, quia poenitentis satisfactio divino tantum pensatur iudicio, non humano. Quare poenitentia iniuncta non est, quia satisfiat\apnota{20}{U satisfacit.} tantum ecclesiae, sed etiam quia satisfit deo. Nam et apud deum productior est poena quam culpa, ut Augustinus testatur tract. 124\apnota{22}{M tract. 134.} super Iohannem.

Quibus omnibus pro indubitata veritate habitis faciliter patet improbatio conclusionis quintae, ubi dicit reverendus pater, quemlibet sacerdotem debere absolvere\apnota{25}{M U absolvere poenitentem.} a poena et a culpa, et ita quilibet sacerdos villanus esset in sua parochia episcopus, archiepiscopus et papa. Quod esse manifestissime\apnota{27}{M manifestissimum.} falsum praeter usum totius ecclesiae, non solum adulatorum, probatur primo ex his\apnota{28}{U iis.}, quae dicta sunt, quia per absolutionem sacerdotalem culpa deletur, poena manet, licet commutata; deinde\apnota{29}{M dein.}, quia sacramentum poenitentiae est\apnota{30a}{M est quasi.} iudiciale, ut formam huius iudicii describit b. Augustinus ca. forte\apnota{30b}{fehlt in M U.} 2. de poen. medic. Et ad iudicium ferendum pertinet iurisdictio, nisi velit Anaxagoreum chaos et confusionem maximam facere in ecclesia dei.

Et cum iurisdictio se ad duo extendat materialia in hoc sacramento habita\apnota{33}{M habito.}, nam et peccator poenitens et peccata ipsa sunt de materia poenitentiae et nemo negat iurisdictionem in inferioribus praelatis ad confusionem tollendam esse coartatam quoad materiam peccantium, nam hinc proprii sacerdotis habetur determinatio, et quia sententia non lata a suo iudice est nulla, ideo absolvens non subditum nihil facit.

Sic a simili in materia peccatorum potest coar- [376] tari iurisdictio eadem ratione et in criminum detestationem\apnota{1}{U detestatione.}. Quod autem etiam a sibi commissis non possit plenarie absolvere, ex eo patet, quia tanta potestas, nisi auctoritate fundetur, contra totius ecclesiae consensum, cuius praxis debet homini\apnota{4}{M omni.} Christiano esse pro regula, non debet tribui. Sed in superioribus ostensum fuit per maiores clericos, Cyprianum, Chrysostomum et Augustinum, poenae debitum remanere post culpam. Ideo reverendus pater cedat s. patrum auctoritati et consuetudini\apnota{7}{M consuetudine.} totius ecclesiae obtemperet, aut cur id fieri non debeat, auctoritate sacrae scripturae, conciliorum vel sanctorum patrum exponat.

Dixit in fine D. D. Eccius, se nihilominus velle replicare, si necessarium videretur sibi, post responsionem domini doctoris Martini\apnota{9}{Diese Schlussbemerkung Ecks fehlt in M U.}.

D. D. Martinus respondit\apnota{10}{U nur a. R. Martinus.}.

Admitto D. doctori, quod habeat ultimum verbum, quia sic vult\apnota{11}{U a. R. Adduxit tamen se velle replicare, si dominus d. Eccius plus iusto responderet.}. Arguit autem egregius D. D. contra conclusionem meam, volens probare poenam requiri a deo, et pro primo inducit scripturam Gen. 3, ubi peccatum Adae etiam remissa culpa punitur in posteritate, ergo poena remanet, transeunte culpa. Volo D. doctorem obstrictum hoc exemplo, et non replicet, nisi confirmet, aut sententia\apnota{16}{P M U sententiam.} cedat. Si ista poena requiritur pro peccato et similiter sentiendum est de qualibet poena, ut infert, habeo propositum, quod poena non possit remitti per papam aut ullum sacerdotem, quia nullus adhuc remisit mortem, infinita genera poenarum, morborum et similium miseriarum, quae omnes pro peccato primo sunt illatae. Et sic probatio D. doctoris probabit\apnota{20}{M U probavit.} conclusionem meam contra se ipsum\apnota{21}{U se ipsam, Uv se ipsum.}.

Secundo, quod a David translato peccato non est translata tamen\apnota{22}{M tamen translata.} poena, etiam pro me facit, sicut et in sermone eodem vulgari sequenter\apnota{23}{M in eodem sermone vulgari frequenter.} dixi: Quia erat poena, quam deus requisivit, ideo nemo

auferre potest\apnota{24}{M poterit, U potuit}. Si D. D. non intellexit me, legat diligentius, ego enim dixi et dico, quod hanc poenam deus non requirat\apnota{26}{M deus hanc poenam non requirit.}, quam papa vel homo possit solvere; quam autem ipse requirat, scilicet immediate\apnota{27a}{scil. immediate fehlt in M; U quam aut ipse requirit (scil. immediate loquitur).}, non possit\apnota{27b}{P posset.} homo solvere.

Tertio, induxit glossam: Deus delictum delet, sed inultum non dimittit; nihil pugnat contra me. Ulciscitur enim, sive\apnota{29}{U sine, Uv sive.} per ipsummet hominem conterendo, sive per ecclesiam emendando, sive per se ipsum iudicando. Et hoc ultimum et primum genus poenae in nullius hominis arbitrio est, sicut apostolus 1. Cor. 11 [v. 31 f.] dicit: Si nos ipsos iudicaremus, non utique iudicaremur a Domino, cum autem iudicamur a Domino, corripimur, ut non cum hoc mundo damnemur. Ibi clarum est, quomodo deus poenam requirat et non requirat.

Quarto, Augustinum per Gratianum relatum, quod deus non permitteret\apnota{35}{M pateretur.} malum\apnota{36}{P mala.} fieri, nisi per iustitiam illud ordinando melius stare faceret: miror egregium\apnota{37a}{fehlt in M U.} D. D.: quod has et similes auctoritates ita consarcinat, cum nemo eas neget\apnota{37b}{M dum . . . . . neget; U negat.} aut contrarium sapiat. Ego enim hoc solum impugnavi semper, quod virtute [377] clavium iactant solvi poenas per iustitiam divinam requisitas. Hoc enim non credo esse verum, nec probabitur. Non enim dixit ad Petrum: Quod ego ligo, tu solves, sed: Quodcunque tu solveris, solutum erit [Matth. 16 v. 19].

Quinto, illud de poenitentiae medicina: Non sufficit mores in melius commutare\apnota{5a}{M U commutari.}, nisi etiam satisfiat deo de his\apnota{5b}{U iis.}, quae facta sunt, per poenitentiae dolorem, per humilitatis gemitum, per contriti cordis sacrificium. Haec omnia ego quam maxime volui semper, ut quid ergo per indulgentias ista iactantur remitti, si non sufficit mores in melius commutare et divina iustitia has requirit\apnota{8}{M U requirat.} satisfactiones? Frustra ergo gloriatur D. D., quod non dictum sit, ut\apnota{9}{U quod.} tantum abstineatis a peccatis, sed de praeteritis deprecare Dominum\apnota{10}{M Dominum etc.}. Et sic patet, quod adhuc nulla scriptura contra me adducta est, quanquam ego possem multo fortiores inducere pro me\apnota{12}{P U contra me; M multa fortiora introducere contra me.}.

Post haec adducit Augustinum in Enchiridio: Delet omnino haec oratio minima et quotidiana peccata, delet et illa, a quibus vita fidelium scelerate gesta, sed poenitendo in melius mutata discedit\apnota{15}{M commutata discedit, U mutata etc.; discedit fehlt in P.}. Hoc est, quod

dixi, peccatorem post conversionem teneri ad crucem et passiones vitae, secundum quod deus intulerit, quas non possit homo solvere, et iterum haec auctoritas pro me facit\apnota{17}{fehlt in M.}.

Eos, qui fuerunt a quadringentis annis\apnota{18}{U qui a ecce (M quadragentis) annis fuerunt.}, non adducit, et placet.

Ad Cyprianum, in epistola ad Fidum, ubi damnat Therapium\apnota{19}{P Fidium; M in (fehlt in U) epistola ad filum, ubi damnat Herapium.}, quod praepropere pacem dederit\apnota{20a}{U dederat.} Victori, antequam plenam egisset\apnota{20b}{M fecisset.} poenitentiam et Domino satisfecisset: Ecce, inquit, non ecclesiae, sed Domino 'non' satisfecisse dicit Victorem\apnota{22}{fehlt in M.}. Respondeo: Legat et conferat Cyprianum bene D. D. et inveniet eos etiam, quos paci dederant\apnota{23}{M dederam.}, ideo datos esse paci etiam praepropere, ut cruces et martyria expeditius sustinerent. Quae ipse per multas epistolas exponit esse poenas et flagella peccatis\apnota{25}{M flagellas, M U pro peccatis.} a deo inflicta. Quare Victor hoc modo nondum satisfecit deo; et tamen satisfecit, quia ecclesiae, quam nos vult audire deus, satisfecit. Hoc enim habent verba Christi, quod ecclesia debet imponere poenas, quando dicit: Quodcunque ligaveris. Et hoc modo possem admittere, deum requirere poenas, quas ecclesia possit solvere, quia pactum fecit cum illa.

Altera auctoritas Cypriani ad Cornelium, qua iterum scribit, haereticos impedire peccantes, ne indignanti deo satisfaciant: Quanquam Cyprianus ibi de poenitentibus non loquatur\apnota{32}{U loquitur.}, sed de haereticis, qui tanquam iusti et quasi bene fecissent excusaverunt et defenderunt se in peccatis suis, tamen respondeo sicut ad priora\apnota{34}{M U priorem.}.

Post haec addidit rationem: Ad quid ecclesia vellet onus imponere poenitentibus\apnota{36}{M poenitenti.}, si illud Dominus non exigeret, et longam post auctoritatem ca. Alligant inducit, quae omnia transeo\apnota{37a}{M U non exigit et longam post (fehlt in U.) auctoritatem capituli allegat (U allegat capitulum) inducit quae (U quae inducit) omnia transeo.}, et dico, quod ecclesiae mandatum est, quod\apnota{37b}{U ut.} castiget et iudicet peccatores, quod si non fecerit, deus irremissibiliter faciet iuxta sententiam Pauli superius inductam 1. Cor. 11 (v. 31 f.\apnota{39}{M U adductam 1. Cor. 11 (U 1. Cor. 13, Uv 1. Cor. 11).}) et sic non potest solvi.

[378] Item idem de Homilia 31. super Hebraeos\apnota{1}{de fehlt in M U; M homilia Chrysostomi, U = P.}: Peccatum confessione minuitur. Ex quo verbo colligit D. D. in hunc modum: Non minuitur quoad culpam, quod impium sit a deo sperare\apnota{3}{M sperare a deo.} dimidiam veniam: ergo quoad poenam. Admitto totum, secundum praedicta.

Iam illa auctoritas Gregorii\apnota{5}{fehlt in M.} 4. Mor.: Aut enim nos hoc flendo insequimur, aut ipse iudicando; vix aliud aptius pro me adduci potuit. Similiter et illud Hiero. sequens c. Mensuram, quod tamen Hieronymi esse dubito, admitto totum\apnota{8a}{M sequens Hieronymi c. mensura admitto totum, U Hieronymi c. mensuram admitto, quod . . . . . dubito. [Absatz.] Admitto totum.}, quod canones non satis praefigunt, ideo relinqui\apnota{8b}{U relinquo.} arbitrio sacerdotis. Addo et ego\apnota{9}{M at nego statt addo et ego.}: multo magis arbitrio dei, qui solus est ponderator spirituum et non ignarus neque iniquus iudicat\apnota{10}{M inique indicat, U iniquus indicat.}.

Placet et illa auctoritas Augustini in\apnota{11}{fehlt in U.} cap. Nullus: Ita impartitur a deo largitas misericordiae, ut non\apnota{12}{M quod nulli (letzteres gestrichen, Mc 'non').} relinquatur iustitiae disciplina. Potest pro me et pro D. D. valere\apnota{13}{M pro me et domino doct. esse, ideo transeo, U pro me et pro d. d. valere transeo.}.

Et Isidorus quamquam\apnota{14}{M U quamvis.} non satis gravis auctor in his rebus, placet tamen, quod poenitentis satisfactionem tantum divino pensari iudicio dictat, non humano. Ergo multo minus remitti potest humano iudicio, cum clavis potestatis non debeat\apnota{17}{M debet.} operari, nisi prior sit clavis scientiae, quae sciat, quid et quantum solvat\apnota{18}{M solveat.}.

Apud deum productiorem esse poenam quam culpam, ex Augustino concedo salva tamen auctoritate apostoli\apnota{20}{M U Pauli.} Roman. 7. [v. 18] dicentis: Non invenio in me, hoc est\apnota{21}{M id est.}, in carne mea, bonum. Cuius sententia est, poenam et peccatum simul finiri.

Haec de prima conclusione. — Contra\apnota{24}{U Sed contra.} aliam dicit, esse manifestissime falsum et praeter usum totius ecclesiae, quod quilibet sacerdos debeat\apnota{25}{M U debet.} absolvere a poena et culpa poenitentem. Et hoc probat primo ex dictis, id est ex nihilo. Secundo per rationem, quia sacramentum poenitentiae sit quoddam iudiciale, et ad iudicium ferendum pertinet iurisdictio. Iurisdictionem autem esse coartatam in inferioribus praelatis ad tollendam confusionem, tam in peccantibus quam in peccatis, secundo in detestationem criminum, alioquin quilibet villanus sacerdos esset episcopus, archiepiscopus et papa.

Respondeo et dico duo, primo\apnota{31}{U primum.} me nescire usque in hodiernum diem, an coartatio istius\apnota{32}{M illius.} iurisdictionis fecerit hoc, quod praetenditur, scilicet criminum\apnota{33}{M criminis.} detestationem et confusionis ablationem. Hoc certum est, longe secius evenisse, nam peccata crassissima etiam ridentur in maioribus curiis, quae in propriis parochiis possent pulcherrime puniri, si servatus esset modus, quem instituerunt apostoli et servaverunt s. patres usque post Nicenum concilium, ubi definitum est et longe post servatum, ut dioeceses non permiscerentur et unusquisque poeniteret in sua dioecesi. De [379] quo exstat cum aliis praeclarissima epistola Cypriani, ad Cornelium Romanum pontificem epistola III\apnota{2}{U pontificem 3.}: Nam cum statutum sit ab omnibus nobis et aequum sit pariter et iustum, ut uniuscuiusque causa illic\apnota{3}{M et illic.} audiatur, ubi est crimen admissum, et singulis pastoribus portio gregis sit adscripta, quam regat unusquisque et gubernet, rationem sui\apnota{5}{M suam, verbessert sui.} actus Domino redditurus, oportet utique eos, quibus praesumus, non circumcursare nec episcoporum concordiam cohaerentem sua subdola et fallaci temeritate collidere\apnota{7}{U confidere, Uv collidere.}, sed agere illic causam suam, ubi et accusatores haberi et testes sui criminis possint\apnota{8a}{M U habere et testes . . . . . possint etc.}. Loquitur enim de his\apnota{8b}{U iis.}, qui in Africa peccaverant et ad\apnota{9}{U peccaverunt; ad fehlt in M.} Romanum pontificem Cornelium cucurrerant. Et sic patet usus primitivae ecclesiae de ligandis et solvendis peccatoribus.

Cum autem, ut apostolus Act. 20. [v. 17 u. 28] ostendit, idem sit quod\apnota{11}{fehlt in U.} episcopus et presbyter, et ad Tit. 1. [v. 5] quaelibet civitas suum episcopum habere debeat iure divino, longe utilior esset\apnota{13}{fehlt in M.} ad corripienda peccata modus, si quilibet sacerdos in sua parochia ligaret et solveret poenitentem. Quod exemplum monstravit apostolus 1. Corinth. 5, ubi cum Corinthiis praesens spiritu tradidit fornicarium satanae, obiurgans quod ipsi non fecissent.

Quo vero iure aut qua ecclesiae felicitate hic modus divino iure praescriptus et in tantum tempus roboratus sit sublatus, viderint alii. Ego quidem negare\apnota{19a}{M negavi.} non possum, quia ad oculum videmus ita fieri, ita statui, ut\apnota{19b}{U quod.} tam peccantes quam peccata reserventur et uni animae sex vel septem pastores secundum, sub et supra imponantur. Sed an ita debeat fieri aut expediat, non definio\apnota{22}{M U diffinio.}. Scio, quod inferior tenetur obedire restringenti et vexanti, quanquam nullo iure divino. Superior tamen, ut dicit conclusio mea, gravissime peccat, si reservet occulta peccata sine rationabilissima causa, imo adhuc dubito, et quantum capio, credo sine temeritate, nullum occultum peccatum debere\apnota{26}{M deberi.} reservari aut posse, optans audire probationem contrariam.

Secundo dico, ecclesia\apnota{27}{U ecclesiam, Uv ecclesia.} non esset ruitura, si idem plebanus episcopus, archiepiscopus et papa esset ac sola concordia cohaerente\apnota{28}{P hac sola statt ac sola; M adhaerente, Mc 'cohaerente'}, ut Cyprianus ait, et sicut usus prioris ecclesiae fuit, iungerentur.

Proinde, quod et damnatus fertur articulus\apnota{29}{M Proinde volo damnatos ostendere articulos.} in Constantiensi concilio de istis\apnota{30}{M his.} reservationibus, non satis curo. Hoc scio, quod probatus fuit et reprobata ista reservatio\apnota{31}{M reservatio huiusmodi.} in primitiva ecclesia et institutione apostolorum, et nunc quoque, ut miserrima ecclesiae experientia\apnota{32}{M experientia ecclesiae.} docet, esset utilissimus et saluberrimus ad coërcenda peccata et tollendam detestabilem confusionem omnium episcopatuum, quam hodie videmus. Qua causa autem sit suo tempore mutatus, transeo. Mutabilem invenio, mutabilem relinquo. Haec\apnota{35}{W. A. verbessert Mutabilem invenio. Relinquo haec.} in arbitrium iudicum\apnota{36}{U iudicum. Hora transit.}.

[380] D. D. Eccius\apnota{3}{U hat hier die Zeitangabe Quarta decima die Julii mane hora septima continuavit Eccius praeter pactum; vgl. Z. 8.}.

Reverendus pater ab initio satis glorianter parvi pendit per me adducta, quasi sibi minime obstent, quae maxime adversantur. Nam ipse voluit in sermone vulgari, non fieri commutationem poenae aeternae in\apnota{6}{U in poenam; M aeternae mutatae in poenam.} temporalem, et opinioni communi contradicit, quasi non possit probari\apnota{7}{M mutari.}, deum exigere aliquam satisfactionem, praeter portationem crucis. Et in sermone latino collaudat

Haec coepta fuerunt tantum per D. D. Eccium circa horam quartam vespertinam, sed intercepta et reassumpta in exordio disputationis postridianae, videlicet

hora septima antemeridiana die Jovis XIV. mensis Julii\apnota{8}{Die Zwischenbemerkung fehlt in M; an der Schrift ist der neue Ansatz zu merken, allerdings schon hinter sermone vulgari. U setzt die Zeitangabe an den Anfang der Rede, vgl Z. 3; continuavit wird dadurch sinnlos.}.

dictum vulgarium ultra omnem doctrinam doctorum scholasticorum, de poenitentia datam: NVMMER THVN DIE HOCHSTE BVES\apnota{10a}{P CHVN statt THVN; M nummer thun die hochteste busse; U nymmer thuen die hochste puess.}. Optima poenitentia nova vita\apnota{10b}{M secundum vitam, Mc 'nova vita'.}, secundum glossam eius. Addiderunt talem persuasionem ipse et propugnator: si averterit se impius etc.\apnota{12}{U quia siaverterit se impius a iustitia (so!) sua etc.}. Si ad novam poenam imputat, quomodo dicitur\apnota{13}{M dicit.} non recordari? et patet conclusio sua manifeste dicens, peccare sacerdotem, qui non absolvat a culpa et poena. Huic errori ego contradixi s. patrum auctoritatibus, quibus ipse voluit illudere et auditores fascinare, quasi apud eum solum esset potestas interpretandi sacram scripturam.

Unde bene adduxi contra eum Augustinum: Non sufficit mores in melius commutare et a malis factis recedere\apnota{19}{M U recedere etc.}. Ubi est liquidissimum, doctrinam ex sermone\apnota{20}{M ex 'eius' sermone.} allegatam esse falsam, quia nova vita non est optima poenitentia, cum non sufficiat secundum Augustinum. Et bene etiam illa doctrina\apnota{22}{M U doctrina illa.} improbatur per b. Ambrosium, lib. 2. de poeniten. cap. 5: Apostoli secundum Christi magisterium docuerunt poenitentiam. Et infra: Qui enim agit\apnota{23}{M ait.} poenitentiam, non solum diluere debet lacrimis peccatum suum, sed etiam emendatioribus factis operire et tegere peccata sua. Clarissima s. patris verba, quod emendatioribus factis agendo poenitentiam tegere debemus peccata.

At cum auctoritates forent expressissimae, deum non remittere peccatum impunitum, confugit ad mirabilem poenarum distinctionem, qui tamen solitus est scholasticos ob distinctionum usum improbare et dixit poenam, qua deus vult peccatum punire, non posse\apnota{30}{M possit, U posset.} auferri per hominem vel per papam, quod est omnium falsissimum et potestatis clavium annullativum. Nam in poenitentia facta commutatione poenae certe homo potest solvere illam poenam, per ipsum apostolum, per reverendum patrem inductum, 1. Cor. 11: Si nos ipsos iudicaremus, non utique iudicaremur a Domino. Quare si nos pro hac poena satisfacimus, deus pro peccato a nobis aliam non exigit\apnota{35}{M pro hoc peccato aliam non exigat.}; alioquin contra prophetam, illud quod nos puniremus, si deus vellet punire, iam bis puniret id ipsum.

[381] Praeterea clara fuerunt Cypriani verba, Chrysostomi, Gregorii et\apnota{1}{U Grego: 4 statt Greg. et.} Hieronymi, quod poena iniuncta a sacerdote in satisfactione\apnota{2}{M iuncta a sacerdoti; M U in satisfactionem.} est poena deo debita, et ita Victor, quia non impleverat poenitentiam iniunctam, a Cypriano dicitur nondum satisfecisse deo. Quod et Theodorus in suo poenitentiali, quem\apnota{4}{M quam.} sequitur quasi ad verbum Beda, id ipsum testatur, quamvis forte nec ii\apnota{5}{M quamvis forte hi.} auctores sicut nec Isidorus ei satisfaciunt\apnota{6}{M U satisfacit.}.

Unde quidquid dicat ore, auctoritates omnes facere pro eo, tamen necessarium est, quod corde dissentiat: cum hae\apnota{8}{M haec.} auctoritates doctrinam eius et conclusionem subvertant. Si enim sacerdos peccat non absolvendo a poena et culpa, tunc episcopi peccassent non absolventes Victorem a poena et culpa, et omnes sacerdotes peccarent per orbem Christianum non absolventes\apnota{11}{non fehlt in P M U; M absolvendo.} extra indulgentias.

Neque relevatur in eo, quod debeat crucem portare et hanc poenam exigat deus, quia ista crucis portatio non est aliud, quam vita Christiniana, sicut ipse reverendus pater docte hoc exposuit; sed cum hoc oportet de praeteritis satisfacere et Dominum deprecari de praeteritis.

De reservatione casuum conclusio mea nihil habet, tamen reservationes moderatas credo esse utiles, quod etiam praelati in monasteriis experiuntur. Fateor ingenue, me esse in voto Gersonis, in concilio\apnota{19}{U consilio.} Constantiensi interpellantis pro delenda immodica reservatione casuum, et potissimum displicet mihi illa reservatio, sicut et domino patri, quando habet avaritiam comitem, id est, pecuniariam\apnota{22}{U pecuniarum.} poenam annexam.

Ecclesiam non dicit ruituram, si sacerdos in sua parochia esset episcopus et papa; certe apparet mihi, si esset illius pulcherrimi ordinis hierarchici corruptio, quod etiam esset\apnota{24}{M esset etiam.} ecclesiae subversio.

Sed maneam in principali, quia peccatum non manet impunitum, secundum Augustinum, Gregorium etc.\apnota{26}{fehlt in U.}. Quare satisfactio merito tertia pars poenitentiae asseritur. Et plene probavit Augustinus in Ench., per orationem dominicam nos satisfacere pro scelerate gestis. Et ut vere dixit heri reverendus pater, deus fecit pactum cum ecclesia: Si ecclesia non facit, deus facit. Ergo habeo propositum scholasticorum et praedicantium, quod oportet vel nos satisfacere, aut deus exiget. Quodsi per orationes nostras vel bona opera satisfacimus, cur illud non posset fieri accedente virtute clavium, quas deus non frustra suae sponsae ecclesiae contulit?\apnota{33}{M tradidit.} Et cum iuxta Gregorii sententiam cap. Decreto 2, q. 6,\apnota{34}{P M U de cetero 2 q. 7.} alii episcopi sunt vocati in partem sollicitudinis, papa habet plenitudinem potestatis, per indulgentias ab eo datas satisfit poenae pro peccatis debitae solutione ex thesauro ecclesiae facta\apnota{36}{M factae.}, ut post Sixtum declaravit papa modernus, praecipiendo sub poena excommunicationis ita doceri\apnota{37}{fehlt in M.}, [382] teneri et praedicari. Quare si reverendus pater contra doceat, praedicet vel disputet, iam est anathemate percussus. Placuit tamen mihi, quod ultra clavem potestatis heri posuit clavem discretionis, cum tamen plures claves contra scholasticos doctores neget conclusione 7. Resolutorii.

Haec volui omnino adiicere, ut futuri iudices apertius intelligerent, quae nostra esset controversia in hoc puncto, quoniam, si reverendus pater stat sententiae doctorum per me allegatorum, non adversabitur nec scholasticis nec praedicatoribus nec mihi. Quare si voluerit, etiam

sententiam suam pro informatione dominorum iudicum poterit clarius exprimere.

D. D. Martinus
non ascensa cathedra, sed ex loco sedendi solito
assurgens respondit\apnota{10}{M U Martinus (die Bemerkung fehlt).}.

Ad istas naenias et ineptias D. doctoris heri satis respondi, repetit enim eadem et velut ridiculus citharoedus chorda semper oberrat eadem.

Secundo scopum controversiae non attigit Non enim quaestio est, an deus peccatum inultum dimittat, quod satis copiose probavit\apnota{14}{U probavi, M probatum est.}, sed an papa vel ecclesia remittat, quas poenas deus exigat. De hoc nihil probavit, quod relinquo iudicio iudicum et omnium auditorum.

Tertio, obticuit hodie scripturas sanctas, ideo sto in eius probatione\apnota{17}{M probatione eius.} hesterna prima, Gene. 3. ca.\apnota{18}{M U ex Genes. 3 (U 30) ca.}, ubi probavit, poenas requiri a deo, quas ibi scriptura ostendit irremissibiles. Doleo, quod D. D. ita profunde penetrat scripturas sicut tipula aquas, imo videtur fugere a facie earum sicut diabolus crucem. Quare salvis reverentiis patrum, praefero ego auctoritatem\apnota{21}{M auctoritates} scripturae, quod commendo iudicibus futuris\apnota{22}{'M iudicio 'iudicum'.}.

D. D. Eccius.

Quia impatiens monachus scurrilia quaedam addidit\apnota{24}{M scurrilia quaedam addit.} praeter gravitatem theologicam, de quo integri viri iudicent\apnota{25}{M in quae (Mc 'de quo') 'interim' iudicem (verbessert iudicent, folgt ein gestrichenes Wort).}, an recte induxerim contra eum, iudices iudicabunt. Sed hanc fuisse materiam nostram, patet ex conclusione 4.: dicere, deum remittendo culpam remittere poenam etc. Hoc erat saxum volvendum\apnota{28}{volvendum fehlt in M, Lücke ausgefüllt Mc per nos solvendam; U per nos volvendum.}, et quia praefert auctoritatem sacrae scripturae patribus, quasi ipse velut alterum oraculum Apollinis solus habeat scripturarum\apnota{29}{M solus ipse habeat scripturam.} intelligentiam ultra s. patres, et apparenter inducit auctoritatem heri per me citatam, dico duo.

Primo me adduxisse in eum finem, ut ostenderetur, cum impius avertisset se ab iniustitia sua, deum adhuc recordari, non quoad culpam, sed quoad poenam, quod facit pro scholasticis et praedicantibus.

Secundo, cum reverendus pater reflectit auctoritatem illam contra me, quia poenae istae a deo peccato Adae impositae sint irremissibiles a papa et homine, verum est, et fateor, quia istae sunt poenae, quae

non solum consequuntur per-[383]sonam, sed etiam consequuntur naturam, et ergo non est mirum, quod istae poenae non sunt remissibiles\apnota{2}{U sunt irremissibiles.} ab homine, sed per hoc non probatur, poenas personales pro peccato debitas non esse a sacerdote vel papa\apnota{3}{U a papa vel sacerdote.} remissibiles. At in his refero\apnota{4}{U remitto.} me ad iudices, paratus mutare sententiam, si me aliter et melius\apnota{5}{M aliter in melius, U aliter ad melius.} docuerint\apnota{*}{In P folgt hier die Bemerkung: Et sic imposita fuit extrema manus disputationi inter E. D. Doctores praenominatos, inchoato deinceps evestigio, imo instaurata rursus disceptatione inter praestantissimos D. Doctores, videlicet eundem Johannem Eccium et Andream Carolostadium in prima fronte supra memoratos etc. Den dann folgenden actenmässigen Schluss des Protokolls s. unten. — U hat Finita sunt haec 14. die Julii hora octava praesente frequenti concione auditorum. Soli Deo honor et gloria. Anno M. D. xix.}.

\noindent\textit{Notas del aparato sin llamada en el cuerpo:} [16] M sententias.; [27] M quodam.

\end{document}
